\documentclass[11pt]{amsart}

\usepackage[T1]{fontenc}
\usepackage[utf8]{inputenc}
\usepackage{lmodern}
\usepackage{amsmath,amssymb,amsthm}
\providecommand{\mathscr}{\mathcal}
\usepackage{booktabs,array,longtable}
\usepackage{xcolor}
\usepackage[colorlinks=true,linkcolor=blue!55!black,
  citecolor=blue!55!black,urlcolor=blue!55!black]{hyperref}

\setlength{\emergencystretch}{5em}
\hbadness=10000
\raggedbottom
\allowdisplaybreaks
\numberwithin{equation}{section}

\newtheorem{theorem}{Theorem}[section]
\newtheorem{proposition}[theorem]{Proposition}
\newtheorem{lemma}[theorem]{Lemma}
\newtheorem{corollary}[theorem]{Corollary}
\theoremstyle{definition}
\newtheorem{definition}[theorem]{Definition}
\newtheorem{contract}[theorem]{Proof contract}
\newtheorem{obligation}[theorem]{Intermediate obligation}
\newtheorem{checkpoint}[theorem]{Closure checkpoint}
\theoremstyle{remark}
\newtheorem{remark}[theorem]{Remark}
\newtheorem{warning}[theorem]{Audit warning}

\newcommand{\T}{\mathbb T}
\newcommand{\Sone}{\mathbb S^1}
\newcommand{\R}{\mathbb R}
\newcommand{\PiNG}{\Pi_{\rm ng}}
\newcommand{\PiZero}{\Pi_0}
\newcommand{\Dop}{|D|}
\newcommand{\Hil}{\mathcal H}
\newcommand{\HH}{\mathcal H}
\newcommand{\Jfour}{\mathfrak J_4}
\newcommand{\Dvar}{\mathfrak D}
\newcommand{\Def}{\mathfrak D}
\newcommand{\G}{\mathfrak G}
\newcommand{\Gtilde}{\widetilde{\mathfrak G}}
\newcommand{\Cnull}{\mathcal C}
\newcommand{\qD}{q_{\mathbb D}}
\newcommand{\qDzero}{q_{\mathbb D}^{0}}
\newcommand{\qphys}{\mathbf q_{\mathbb D}}
\newcommand{\qplus}{q_+}
\newcommand{\qdisk}{q_D}
\newcommand{\qdot}{q_{\dot H}}
\newcommand{\qzero}{q_0}
\newcommand{\kform}{\mathcal K}
\newcommand{\Pplus}{\mathcal P_+^0}
\newcommand{\PGdot}{\mathcal P_{G,|D|}^0}
\newcommand{\PGplus}{\mathcal P_{G,+}^0}
\newcommand{\PD}{\mathcal P_D^{\rm ex}}
\newcommand{\Eaug}{\mathcal E_\#}
\newcommand{\Rone}{\mathcal R_1}
\newcommand{\Rfull}{\mathcal R_1^{\rm full}}
\newcommand{\Rgeo}{\mathcal R_{\rm geo}}
\newcommand{\Rphys}{\mathcal R_1^{\rm phys}}
\newcommand{\Rfour}{\mathcal R_4}
\newcommand{\Pthree}{\mathcal P_3}
\newcommand{\Gjp}{\mathfrak G^{\rm JP}}
\newcommand{\Dred}{\mathfrak D^{\rm red}}
\newcommand{\Xjp}{\mathcal X^{\rm JP}}
\providecommand{\mathscr}[1]{\mathcal{#1}}
\newcommand{\X}{\mathcal X}
\newcommand{\Y}{\mathcal Y}
\newcommand{\N}{\mathcal N}
\newcommand{\Kclass}{\mathcal K}
\newcommand{\Pclass}{\mathcal P}
\newcommand{\LambdaF}{\Lambda}


\title[Audited convex large-$N$ polygonal assembly]
{An audited assembly for the convex large-$N$ polygonal
P\'olya--Szeg\H{o} theorem with a pathwise LCA frontier}
\author{}
\date{July 20, 2026}

\begin{document}

\begin{abstract}
We assemble the argument in logical order and separate proved interfaces
from conditional ones.  The disk endpoint, physical trace transfer,
complete cubic tensor, structured collapsed-bubble commutator, material
transport, exact fixed-normal support Hessian, fan-normal metric, strict
side activation, and compact global implication are proved.  Revision 24
separates the reduced disk form $q_D^0$ from the physical half-Hessian
$\mathbf q_D=2j_{0,1}^2q_D^0$, corrects the physical corner exponent in
the support-time specialization, and retracts the former near-regular
three-periodic material transfer for the flux-matching clause JE--FM.
The unproved separated JP--SC/JP--JE/JP--FS route is removed from the
active DAG.  The closed $D0_{\rm ns}$, JE--P, and JE--V
rows form the disk endpoint JP--DE.  A new honest fixed-fan scalar residual
keeps the exact forcing and integrated support Hessian in the same Schur
comparison.  The exact completed-action row JC--ALG and the support-time
logarithm cancellation JC--TL are proved.  Revision 24 retains the complete
regular block JC--AR for every character: it combines the genuine
Schwarz--Christoffel mixed finite part, a rational completed-covector ratio
certificate, and a uniform finite-$N$ true-flux/reduced-resolvent bridge.
The exact open cut inside the remaining pathwise payment JC--PL is now
JC--SP$+$JC--CL, namely the completed-action Hessian modulus and coarse/far
action localization.  The jointly renormalized pure-support row, canonical
super-near finite energy bridge, and exact completed-path activity are
proved subrows.  The former fixed-reserve RF--IM target is rigorously
disproved by an active period-three family and replaced by the open
relative-root interface RF--REL.  The terminal-radius pure-power row is
resummed exactly; the noninvariant DAR row is retracted, while the exact
completed third-jet/QRCI estimate remains open.  The regular coarse leading
row is reduced to an unproved monotone-transport gap-simplex inequality.
JC--PL implies the former one-sided comparison
JP--JC and hence J0 by exact disk completion.  JC--PL is the sole open
mathematical input.  Hence this file is a complete conditional assembly,
not an unconditional certificate.
\end{abstract}

\maketitle
\tableofcontents

\section{Main theorem and assembly order}
\label{sec:assembled-main}

For $A>0$, let $R_N(A)$ be the regular $N$-gon of area $A$.

\begin{theorem}[Conditional convex large-$N$ P\'olya--Szeg\H{o} theorem]
\label{thm:assembled-main}
Assume the pathwise LCA payment JC--PL,
Contract~\ref{JC:con:JC-PL}.  Then there exists $N_0$ such that, for every
$N\ge N_0$ and every bounded convex planar polygon $P$ of area $A$ with
exactly $N$ genuine sides,
\[
 \lambda_1(P)\ge \lambda_1(R_N(A)).
\]
Equality holds if and only if $P$ is congruent to $R_N(A)$.
\end{theorem}

The proof is organized by the following acyclic implication chain:
\begin{scriptsize}
\[
\begin{gathered}
 \mathrm{F51+H0+QI}\longrightarrow\mathrm{F65},\qquad
 \mathrm{F51}\longrightarrow\mathrm{F129},\qquad
 \mathrm{F51+F65+F129}\longrightarrow\mathrm{TR},\\
 \mathrm{F51}\longrightarrow\mathrm{CT_{core}},\qquad
 \mathrm{F51+TR+CT_{core}}\longrightarrow\mathrm{QB},\\
 \mathrm{F144}\longrightarrow\mathrm{F148}\longrightarrow\mathrm{DF0B},
 \qquad \mathrm{F144}\longrightarrow\mathrm{DF0A},\\
 \mathrm{F144+DF0A+DF0B}\longrightarrow\mathrm{ZG}
 \longrightarrow\mathrm{DF},\\
 \mathrm{F51+CT_{core}}\longrightarrow\mathrm{PT}
 \longrightarrow\mathrm{DM},\\
 \mathrm{CT_{core}}\longrightarrow
 \{\mathrm{JP\text{-}EH},\mathrm{JP\text{-}NS}\},\\
 \mathrm{JP\text{-}NS+TR+CT_{core}+QB+DF+DM}
 \longrightarrow\mathrm{JP\text{-}DE}
 =\{D0_{\rm ns},\mathrm{JE\text{-}P},\mathrm{JE\text{-}V}\},\\
 \mathrm{JP\text{-}EH+JP\text{-}NS+JP\text{-}DE}
 \longrightarrow\{\mathrm{JC\text{-}ALG},\mathrm{JC\text{-}TL}\},\\
 \textcolor{red}{\mathrm{JC\text{-}ALG+JC\text{-}TL}
 \longrightarrow\mathrm{JC\text{-}PL}}
 \longrightarrow\mathrm{JP\text{-}JC},\\
 \mathrm{JP\text{-}DE+JP\text{-}JC}\longrightarrow\mathrm{NR}
 \longrightarrow\mathrm{J0},\\
 \mathrm{F57}\longrightarrow\mathrm{A0},\qquad
 \mathrm{C0+A0+J0}\longrightarrow\mathrm{G0}
 \longrightarrow\text{Theorem~\ref{thm:assembled-main}}.
\end{gathered}
\]
\end{scriptsize}
Every imported label is namespaced by its module identifier.  The final
ledger below distinguishes closed, retired, open, and dependent rows.

\part{Disk endpoint foundations and physical trace}

\begingroup
\renewcommand{\qD}{q_{\mathbb D}^{0}}
\section{Foundation F51: exact disk Hessian and Bessel--Schur normalization}
\label{module:F51}
\noindent\textit{Audited source: \nolinkurl{convex_largeN_polya_szego_stage51_exact_disk_normalization.tex}.}
\subsection{Fourier and energy conventions}

For a real $2\pi$-periodic function $f$, use
\begin{equation}
 \widehat f_n=\frac1{2\pi}\int_0^{2\pi}f(\theta)e^{-in\theta}\,d\theta,
 \qquad
 f(\theta)=\sum_{n\in\mathbb Z}\widehat f_ne^{in\theta}.
\end{equation}
Our normalized circle inner product is
\begin{equation}\label{F51:eq:L2-star}
 \langle f,g\rangle_{L^2_*}
 :=\frac1{2\pi}\int_0^{2\pi}f(\theta)\overline{g(\theta)}\,d\theta
 =\sum_{n\in\mathbb Z}\widehat f_n\overline{\widehat g_n}.
\end{equation}
Let
\begin{equation}
 H_{\rm ng}^{s}
 :=\left\{f\in H^s(\Sone):
 \widehat f_0=\widehat f_1=\widehat f_{-1}=0\right\},
 \qquad
 \PiNG f=\sum_{|n|\ge2}\widehat f_ne^{in\theta}.
\end{equation}
The normalized homogeneous and inhomogeneous half-order forms are
\begin{align}
 \|f\|_{\dot H^{1/2}_*}^2
 &:=\sum_{|n|\ge2}|n|\,|\widehat f_n|^2,\label{F51:eq:hhalf}\\
 \qplus(f,g)
 &:=\sum_{|n|\ge2}(|n|+1)\widehat f_n\overline{\widehat g_n}.
 \label{F51:eq:qplus}
\end{align}
If $U$ is the decaying harmonic extension of $f$ to the half-cylinder
$\Sone\times(0,\infty)$, Parseval's identity gives
\begin{equation}\label{F51:eq:cylinder}
 \int_{\Sone\times(0,\infty)}|\nabla U|^2\,d\theta\,dy
 =2\pi\|f\|_{\dot H^{1/2}_*}^2.
\end{equation}
Thus a strip computation made in cylinder energy must be divided by $2\pi$
before it is inserted into \eqref{F51:eq:hhalf}.  This is the source of the two
constants
\begin{equation}\label{F51:eq:twoK}
 K_{\rm cyl}=\frac{\zeta(3)}{\pi^3},
 \qquad
 K_*:=\frac{K_{\rm cyl}}{2\pi}
 =\frac{\zeta(3)}{2\pi^4}.
\end{equation}
They are not competing normalizations of the same displayed form.

\subsection{The exact fan/support affine class}

Let
\begin{equation}
 \alpha_i>0,\qquad \sum_{i=1}^N\alpha_i=2\pi,
 \qquad \theta_{i+1}-\theta_i=\alpha_i
\end{equation}
cyclically, and put $\nu_i=(\cos\theta_i,\sin\theta_i)$.  On the edge cell
\begin{equation}
 K_i=[-\alpha_{i-1}/2,\alpha_i/2]
\end{equation}
write
\begin{equation}
 a=\alpha_{i-1},\quad b=\alpha_i,\quad
 s=\frac{x+a/2}{(a+b)/2},\quad
 r_- =\frac{z_i-z_{i-1}}a,\quad
 r_+=\frac{z_{i+1}-z_i}b.
\end{equation}

\begin{definition}[Exact trigonometric support trace]\label{F51:def:T}
For $z\in\R^N$, define $T_\alpha z$ on $K_i$ by
\begin{equation}\label{F51:eq:T}
 (T_\alpha z)_i(x)
 =z_i\cos x+\bigl((1-s)a_i^-+sa_i^+\bigr)\sin x,
\end{equation}
where
\begin{equation}\label{F51:eq:apm}
 a_i^-=r_-\frac a{\sin a}-z_i\tan\frac a2,
 \qquad
 a_i^+=r_+\frac b{\sin b}+z_i\tan\frac b2.
\end{equation}
The formulas on adjacent cells agree at their common endpoint, so
$T_\alpha z$ is continuous and belongs to $H^1(\Sone)$.
\end{definition}

Indeed, at the right endpoint of $K_i$, direct substitution gives
\begin{equation}
 (T_\alpha z)_i(b/2)
 =\frac{z_i+z_{i+1}}{2\cos(b/2)}
 =(T_\alpha z)_{i+1}(-b/2).
\end{equation}
The exact area row is
\begin{equation}\label{F51:eq:area-row}
 L_\alpha z=\sum_{i=1}^N\ell_i z_i,
 \qquad
 \ell_i=\tan\frac{\alpha_{i-1}}2+\tan\frac{\alpha_i}2,
\end{equation}
and the admissible coefficient space is $Z_\alpha=\ker L_\alpha$.  Let
\begin{equation}\label{F51:eq:V}
 V_\alpha:=\PiNG T_\alpha Z_\alpha\subset H_{\rm ng}^{1/2}.
\end{equation}

Let $B_\alpha^{\rm ex}\in H^{1/2}(\Sone)$ be the exact canonical fan
profile in this fixed material convention.  The normalization of that
profile is part of the definition; no factor is silently absorbed into it.
Set
\begin{equation}\label{F51:eq:b}
 b_\alpha:=\PiNG B_\alpha^{\rm ex},
 \qquad
 \mathcal A_\alpha=b_\alpha+V_\alpha.
\end{equation}
All endpoint values below depend only on the affine subspace
$\mathcal A_\alpha$, so replacing $B_\alpha^{\rm ex}$ by
$B_\alpha^{\rm ex}+T_\alpha z_0$ with $z_0\in Z_\alpha$ changes no value.

\subsubsection{Translations}

\begin{lemma}[Exact translation identity]\label{F51:lem:translations}
For $c\in\R^2$, put $z_i=c\cdot\nu_i$.  Then $L_\alpha z=0$ and
\begin{equation}\label{F51:eq:translation}
 T_\alpha z(\theta)=c\cdot(\cos\theta,\sin\theta).
\end{equation}
Consequently $\PiNG T_\alpha z=0$.
\end{lemma}

\begin{proof}
Let $\tau_i=(-\sin\theta_i,\cos\theta_i)$.  From
\begin{equation}
 \nu_{i-1}=\nu_i\cos a-\tau_i\sin a,
 \qquad
 \nu_{i+1}=\nu_i\cos b+\tau_i\sin b
\end{equation}
one obtains from \eqref{F51:eq:apm}
\begin{equation}
 a_i^-=a_i^+=c\cdot\tau_i.
\end{equation}
Equation \eqref{F51:eq:T} is therefore exactly
\begin{equation}
 (c\cdot\nu_i)\cos x+(c\cdot\tau_i)\sin x
 =c\cdot\nu(\theta_i+x),
\end{equation}
which proves \eqref{F51:eq:translation}.  The area row vanishes because
translations preserve area.  Algebraically, it is the closure identity
$\sum_i\ell_i\nu_i=0$ for the tangential polygon with normal fan $\alpha$.
The last assertion follows from the Fourier support of
$c\cdot\nu(\theta)$.
\end{proof}

Thus centering only chooses a representative of the translation orbit.  It
does not add a finite-rank loss to the nongeometric endpoint Schur value.

\subsection{The disk multiplier}

Let $j=j_{0,1}$ be the first positive zero of $J_0$.  For $n\ge2$, set
\begin{equation}\label{F51:eq:kn}
 k_n=j\frac{J_{n+1}(j)}{J_n(j)}.
\end{equation}

\begin{lemma}[A rational enclosure for $j$]\label{F51:lem:j-bound}
One has
\begin{equation}\label{F51:eq:j-bound}
 2.4<j<2.405,
 \qquad \frac{j^2}{4}<1.447.
\end{equation}
\end{lemma}

\begin{proof}
The series
\begin{equation}
 J_0(x)=\sum_{m\ge0}\frac{(-1)^m(x^2/4)^m}{(m!)^2}
\end{equation}
is alternating with decreasing terms on the indicated interval after its
first term.  Rational partial sums through orders eight and nine give
\begin{align*}
 0.0025076832943&<J_0(12/5)<0.0025076834966,\\
 -0.000090558004&<J_0(481/200)<-0.000090557793.
\end{align*}
Moreover $J_0'=-J_1$, and
\begin{equation}
 J_1(x)=\frac x2\sum_{m\ge0}
 \frac{(-1)^m(x^2/4)^m}{m!(m+1)!}>0
 \qquad(0<x\le2.405),
\end{equation}
because the alternating sum is at least $1-(2.405)^2/8>0$.
Hence $J_0$ is strictly decreasing there and its first zero lies in the
displayed interval.  The final bound follows by squaring $2.405$.
\end{proof}

\begin{theorem}[Positive order-$(-1)$ splitting]\label{F51:thm:bessel}
For every $n\ge2$,
\begin{equation}\label{F51:eq:recurrence}
 1+j\frac{J_n'(j)}{J_n(j)}=n+1-k_n,
\end{equation}
and
\begin{equation}\label{F51:eq:k-bounds}
 0<k_n\le \frac{2x}{n+1-x}<2,
 \qquad k_n<\frac4n,
 \qquad x:=\frac{j^2}{4}<1.447.
\end{equation}
In particular
\begin{equation}\label{F51:eq:weight-equivalence}
 \frac13(n+1)\le n+1-k_n\le n+1
 \qquad(n\ge2).
\end{equation}
\end{theorem}

\begin{proof}
The Bessel recurrence
\begin{equation}
 jJ_n'(j)=nJ_n(j)-jJ_{n+1}(j)
\end{equation}
gives \eqref{F51:eq:recurrence}.  Write
\begin{equation}
 J_n(j)=\frac{(j/2)^n}{n!}F_n,
 \qquad
 F_n=\sum_{m\ge0}\frac{(-1)^mx^m}{m!(n+1)_m}.
\end{equation}
For $n\ge2$, the ratio of consecutive absolute terms is at most
$x/(n+1)<1/2$.  Alternating-series bounds give
\begin{equation}
 0<1-\frac{x}{n+1}\le F_n\le1.
\end{equation}
Consequently
\begin{equation}
 0<k_n=\frac{2x}{n+1}\frac{F_{n+1}}{F_n}
 \le\frac{2x}{n+1-x}.
\end{equation}
The right-hand side is decreasing in $n$ and is less than $1.864$ at
$n=2$.  Direct cross multiplication, using $x<1.447$, also gives
$2x/(n+1-x)<4/n$ for every $n\ge2$.  Hence
$n+1-k_n\ge n-1\ge(n+1)/3$, proving
\eqref{F51:eq:weight-equivalence}.
\end{proof}

Define the positive self-adjoint multiplier
\begin{equation}
 K_{-1}e^{in\theta}=k_{|n|}e^{in\theta},
 \qquad |n|\ge2,
\end{equation}
and the reduced disk form
\begin{equation}\label{F51:eq:qD}
 \qD(f,g):=\qplus(f,g)-k(f,g),
 \qquad
 k(f,g):=\langle f,K_{-1}g\rangle_{L^2_*}.
\end{equation}
The physical disk quadratic Taylor form (equivalently, one half of the
second variation of the scale-invariant deficit at the disk) is
\begin{equation}\label{F51:eq:physical-Q}
 \mathsf Q_D(f,g):=2j^2\qD(f,g).
\end{equation}
Thus the full Hessian is $2\mathsf Q_D=4j^2\qD$.  The factor $2j^2$ is
not included in either $\qplus$ or $\qD$.

\begin{corollary}[Coercivity and negative order]\label{F51:cor:coercivity}
For $f\in H_{\rm ng}^{1/2}$,
\begin{equation}\label{F51:eq:coercivity}
 \frac13\qplus(f,f)\le\qD(f,f)\le\qplus(f,f).
\end{equation}
Moreover
\begin{equation}\label{F51:eq:minus}
 |k(f,g)|\le C_j\|f\|_{H^{-1/2}}\|g\|_{H^{-1/2}}.
\end{equation}
\end{corollary}

\begin{proof}
Equation \eqref{F51:eq:coercivity} is \eqref{F51:eq:weight-equivalence} summed over
Fourier modes.  By \eqref{F51:eq:k-bounds}, $k_n\le C_j/n$; Cauchy--Schwarz with
weight $1/n$ gives \eqref{F51:eq:minus}.
\end{proof}

\subsection{Exact terminal Schur identity}

For the affine class \eqref{F51:eq:b}, define
\begin{align}
 \Pplus(\alpha)
 &:=\inf_{v\in V_\alpha}\qplus(b_\alpha+v,b_\alpha+v),\label{F51:eq:Pplus}\\
 \widetilde{\PD}(\alpha)
 &:=\inf_{v\in V_\alpha}\qD(b_\alpha+v,b_\alpha+v),\label{F51:eq:PDtilde}\\
 \PD(\alpha)&:=2j^2\widetilde{\PD}(\alpha).
 \label{F51:eq:PD}
\end{align}
Let $f_\alpha\in\mathcal A_\alpha$ be the unique $\qplus$-minimal
residual.  Thus
\begin{equation}\label{F51:eq:galerkin}
 \qplus(f_\alpha,v)=0\qquad(v\in V_\alpha).
\end{equation}
On $V_\alpha$, define
\begin{equation}
 A_{D,\alpha}[v,w]=\qD(v,w),
 \qquad g_\alpha(v)=k(f_\alpha,v).
\end{equation}

\begin{theorem}[Terminal Bessel--Schur formula]\label{F51:thm:schur}
For every positive fan,
\begin{equation}\label{F51:eq:schur}
 \widetilde{\PD}(\alpha)
 =\Pplus(\alpha)-k(f_\alpha,f_\alpha)
 -\langle g_\alpha,A_{D,\alpha}^{-1}g_\alpha\rangle.
\end{equation}
The inverse and dual pairing are taken on the quotient of $V_\alpha$ by its
zero trace; by \eqref{F51:eq:coercivity} they are well defined.
\end{theorem}

\begin{proof}
Every element of $\mathcal A_\alpha$ has a unique representation
$f_\alpha+w$ modulo the zero trace, with $w\in V_\alpha$.  Using
\eqref{F51:eq:galerkin},
\begin{align*}
 \qD(f_\alpha+w,f_\alpha+w)
 ={}&\qplus(f_\alpha,f_\alpha)-k(f_\alpha,f_\alpha)\\
 &+A_{D,\alpha}[w,w]-2g_\alpha(w).
\end{align*}
Completing the square in the coercive form $A_{D,\alpha}$ gives
\eqref{F51:eq:schur}.
\end{proof}

\subsection{The uniform fan}

Let $\alpha_*=a\mathbf1$, $a=2\pi/N$.  Rotation by $a$ acts on coefficient
vectors by the cyclic shift and on traces by the unitary rotation
$R_af(\theta)=f(\theta-a)$.

\begin{lemma}[Uniform Bessel forcing vanishes]\label{F51:lem:gzero}
Assume the canonical choice of $B_{\alpha_*}^{\rm ex}$ is cyclically
equivariant.  Then
\begin{equation}\label{F51:eq:gzero}
 g_{\alpha_*}=0.
\end{equation}
\end{lemma}

\begin{proof}
The affine class, $\qplus$, and the canonical datum are invariant under
$R_a$, so uniqueness of the $\qplus$ projection gives
$R_af_{\alpha_*}=f_{\alpha_*}$.  The multiplier $K_{-1}$ commutes with
$R_a$.  Hence for $z\in Z_{\alpha_*}$,
\begin{equation}
 k(f_{\alpha_*},\PiNG T_{\alpha_*}z)
 =k\left(f_{\alpha_*},\PiNG T_{\alpha_*}
 \frac1N\sum_{m=0}^{N-1}\sigma^m z\right),
\end{equation}
where $\sigma$ is the cyclic shift.  At the uniform fan,
$L_{\alpha_*}z=0$ is equivalent to $\sum_i z_i=0$, so the averaged
coefficient vector is zero.  Therefore the displayed pairing vanishes for
every element of $V_{\alpha_*}$.
\end{proof}

Combining Theorem~\ref{F51:thm:schur} and Lemma~\ref{F51:lem:gzero} gives the exact
endpoint difference
\begin{align}\label{F51:eq:endpoint-difference}
 \widetilde{\PD}(\alpha)-\widetilde{\PD}(\alpha_*)
 ={}&\Pplus(\alpha)-\Pplus(\alpha_*)\\
 &-\bigl[k(f_\alpha,f_\alpha)
          -k(f_{\alpha_*},f_{\alpha_*})\bigr]\notag\\
 &-\langle g_\alpha,A_{D,\alpha}^{-1}g_\alpha\rangle.\notag
\end{align}
\endgroup

\begingroup
\section{Node H0: the correct-source homogeneous gap}
\label{module:H0}
\noindent\textit{Audited source: \nolinkurl{convex_largeN_polya_szego_stage64_correct_polygonal_source.tex}.}
\subsection{Exact tangential-polygon geometry}

Let
\begin{equation}\label{H0:eq:fan}
 \alpha_i>0,\qquad \sum_{i=1}^N\alpha_i=2\pi,
 \qquad \theta_{i+1}-\theta_i=\alpha_i
\end{equation}
cyclically, and let
$\nu_i=(\cos\theta_i,\sin\theta_i)$.  Put
\begin{equation}\label{H0:eq:vertices}
 \phi_i=\theta_i+\frac{\alpha_i}{2},
 \qquad K_i=[\phi_{i-1},\phi_i].
\end{equation}
The polygon with all support numbers equal to one is
\begin{equation}
 P(\alpha,\mathbf 1)=\bigcap_i\{x:x\cdot\nu_i<1\}.
\end{equation}
Its area and its area-$\pi$ scaling factor are
\begin{equation}\label{H0:eq:scale}
 A_\alpha=\sum_i\tan\frac{\alpha_i}{2},
 \qquad h_\alpha=\left(\frac{\pi}{A_\alpha}\right)^{1/2}.
\end{equation}
On $K_i$ use
\begin{equation}\label{H0:eq:edge-coordinate}
 x=\theta-\theta_i,
 \qquad -\frac{\alpha_{i-1}}2\le x\le\frac{\alpha_i}2.
\end{equation}
The exact radial defect of the translated, area-normalized tangential
polygon is
\begin{equation}\label{H0:eq:exact-source}
 \rho_\alpha^{\rm tan}(\theta)=h_\alpha\sec x-1,
 \qquad \theta\in K_i.
\end{equation}
Translation of the polygon does not change its eigenvalue, so the auxiliary
discrete centering translation is absent from \eqref{H0:eq:exact-source}.

\begin{definition}[Correct principal source]\label{H0:def:G}
Define the continuous edge source
\begin{equation}\label{H0:eq:G}
 G_\alpha(\theta)=\frac{x^2}{2},
 \qquad \theta\in K_i,
\end{equation}
with $x$ as in \eqref{H0:eq:edge-coordinate}.
\end{definition}

\begin{lemma}[Exact source expansion]\label{H0:lem:source-expansion}
Let $S=\alpha_{\max}$.  There is a constant $c_\alpha$ such that
\begin{equation}\label{H0:eq:source-expansion}
 \rho_\alpha^{\rm tan}=c_\alpha+G_\alpha+R_\alpha,
\end{equation}
where, on $K_i$,
\begin{equation}\label{H0:eq:cell-remainder}
 |R_\alpha(\theta)|\le CS^2x^2,
 \qquad |R_\alpha'(\theta)|\le CS^2|x|.
\end{equation}
Consequently
\begin{equation}\label{H0:eq:global-remainder}
 \|R_\alpha\|_{H^{1/2}(\Sone)}
 \le CS\|G_\alpha\|_{H^{1/2}(\Sone)}.
\end{equation}
The same estimate holds after removal of the modes $0,\pm1$.
\end{lemma}

\begin{proof}
Taylor's formula and \eqref{H0:eq:fan} give
\begin{equation}
 A_\alpha
 =\pi+\frac1{24}\sum_i\alpha_i^3+O(S^4),
 \qquad h_\alpha=1+O(S^2).
\end{equation}
Take $c_\alpha=h_\alpha-1$.  Since
\begin{equation}
 \sec x-1=\frac{x^2}{2}+O(x^4),
 \qquad (\sec x-1)'=x+O(x^3),
\end{equation}
uniformly for $|x|\le S$, equations
\eqref{H0:eq:source-expansion}--\eqref{H0:eq:cell-remainder} follow.  The
remainder is continuous at each $\phi_i$, because the two adjacent endpoint
values use the same gap $\alpha_i$.

More intrinsically, write $R_\alpha=m_\alpha G_\alpha$, extending the
quotient continuously at $x=0$.  The even analytic quotient
\begin{equation}
 m_\alpha(x)=2h_\alpha\frac{\sec x-1}{x^2}-1
\end{equation}
satisfies
\begin{equation}
 \|m_\alpha\|_{L^\infty}\le CS^2,
 \qquad \|m_\alpha'\|_{L^\infty(K_i)}\le CS.
\end{equation}
The endpoint values from the two incident edges agree, so $m_\alpha$ is a
global Lipschitz multiplier.  The standard Lipschitz multiplier estimate on
$H^{1/2}(\Sone)$ now gives \eqref{H0:eq:global-remainder}.  A fixed Fourier
projection is bounded on $H^{1/2}$.
\end{proof}

\subsection{The old normal-cell source is not the polygonal source}

For comparison, define the normal-cell bubble
\begin{equation}\label{H0:eq:B}
 B_\alpha(\theta_i+y)=y(\alpha_i-y),
 \qquad 0\le y\le\alpha_i.
\end{equation}
Let $T_\alpha z$ be any vertex-compatible support trace obtained from a
vector field affine along each straight edge.  In particular, on $K_i$ it
has the exact trigonometric form
\begin{equation}\label{H0:eq:exact-trace-form}
 (T_\alpha z)(\theta)
 =z_i\cos x+(A_i+B_i x)\sin x.
\end{equation}
The coefficients are constrained only by matching the common vertex
velocities.  Thus $T_\alpha z$ is smooth in the interior of every $K_i$.

\begin{proposition}[Principal source obstruction]\label{H0:prop:obstruction}
For every positive fan and every $c\ne0$,
\begin{equation}\label{H0:eq:not-affine}
 G_\alpha\notin cB_\alpha+
 \{T_\alpha z:z\in\R^N\}+\operatorname {span}\{1,\cos\theta,\sin\theta\}.
\end{equation}
The same assertion holds with the homogeneous support trace
\begin{equation}\label{H0:eq:T0-pre}
 T_\alpha^0z
 =z_i+x\bigl((1-s)r_{i-1}+sr_i\bigr)
 \quad\hbox{on }K_i.
\end{equation}
Hence an effective polygonal source of the form
$cB_\alpha+O_{\rm scaled}(S^3)$ cannot follow from exact polygon geometry
and vertex-compatible support completion.
\end{proposition}

\begin{proof}
In the sense of periodic distributions,
\begin{align}
 G_\alpha''
 &=1-\sum_i\alpha_i\delta_{\phi_i},\label{H0:eq:Gdd}\\
 B_\alpha''
 &=-2+\sum_i(\alpha_{i-1}+\alpha_i)\delta_{\theta_i}.
 \label{H0:eq:Bdd}
\end{align}
Indeed, $G_\alpha'=x$ on $K_i$ and its jump at $\phi_i$ is
$-\alpha_i$.  Likewise, $B_\alpha''=-2$ between consecutive normals and
the derivative jump at $\theta_i$ is $\alpha_{i-1}+\alpha_i$.

Every normal $\theta_i$ lies strictly inside $K_i$, whereas every
$\phi_i$ is a polygonal vertex direction.  By
\eqref{H0:eq:exact-trace-form}, the singular part of
$(T_\alpha z)''$ is supported only on the set $\{\phi_i\}$.  The same is
true for \eqref{H0:eq:T0-pre}; the three geometric modes are smooth.  If
\eqref{H0:eq:not-affine} failed, comparison of the atoms at $\theta_i$ in the
second distributional derivatives would give
$c(\alpha_{i-1}+\alpha_i)=0$ for every $i$, contrary to $c\ne0$.

The two atom measures in \eqref{H0:eq:Gdd}--\eqref{H0:eq:Bdd} have coefficients
of order $S$ and their twice-integrated profiles have size $S^2$ on a
normalized cell.  Thus the discrepancy occurs at principal order, not in
an endpoint-weighted $O(S^3)$ remainder.
\end{proof}

\begin{remark}
The bulk curvatures explain the formerly proposed constant: choosing
$c=-1/2$ makes the absolutely continuous parts of
$G_\alpha''$ and $(cB_\alpha)''$ agree.  It does not move the atomic part
from the normals to the vertices.  Keeping only the bulk Taylor coefficient
therefore loses exactly the geometric datum that the critical
$H^{1/2}$ energy detects.
\end{remark}

\subsection{The corrected homogeneous Schur problem}

Put
\begin{equation}\label{H0:eq:r}
 r_i=\frac{z_{i+1}-z_i}{\alpha_i}.
\end{equation}
On $K_i$, let
\begin{equation}\label{H0:eq:T0}
 T_\alpha^0z
 =z_i+x\bigl((1-s)r_{i-1}+sr_i\bigr),
 \qquad
 s=\frac{x+\alpha_{i-1}/2}{(\alpha_{i-1}+\alpha_i)/2}.
\end{equation}
For
\begin{equation}
 \widehat f_n=\frac1{2\pi}\int_0^{2\pi}
 f(\theta)e^{-in\theta}\,d\theta,
 \qquad
 \|f\|_{\dot H_*^{1/2}}^2
 =\sum_{n\ne0}|n|\,|\widehat f_n|^2,
\end{equation}
define
\begin{align}
 \PGdot(\alpha)
 :=\inf_{z\in Z_\alpha^0}
 \|G_\alpha+T_\alpha^0z\|_{\dot H_*^{1/2}}^2,
 \label{H0:eq:PGdot}\\
 \PGplus(\alpha)
 :=\inf_{z\in Z_\alpha^0}\left(
 \|G_\alpha+T_\alpha^0z\|_{\dot H_*^{1/2}}^2
 +\|G_\alpha+T_\alpha^0z\|_{L_*^2}^2\right),
 \label{H0:eq:PGplus}
\end{align}
Here $Z_\alpha^0$ may be the full coefficient space or any homogeneous
area/centering subspace for which $z=0$ is admissible at the uniform fan.
The lower bounds below are pointwise in $z$ and are therefore unaffected by
these rows.

Let $\mathcal N_i(f)$ denote the least Dirichlet energy in the half-strip
$K_i\times(0,\infty)$ with bottom trace $f|_{K_i}$ and homogeneous Neumann
data on the two vertical walls.  Restriction of the periodic harmonic
extension gives
\begin{equation}\label{H0:eq:strip-localization}
 2\pi\|f\|_{\dot H_*^{1/2}}^2
 \ge\sum_i\mathcal N_i(f).
\end{equation}

\subsection{Exact one-edge energy and the rational Bellman inequality}

On one edge put
\begin{equation}\label{H0:eq:normalized-variables}
 a=\alpha_{i-1},\quad b=\alpha_i,\quad w=a+b,
 \quad t=\frac aw,\quad X=1-2t,
\end{equation}
and
\begin{equation}
 \rho_-=\frac{r_{i-1}}w,\qquad
 \rho_+=\frac{r_i}w,\qquad
 \delta=\rho_+-\rho_-,\qquad
 q=t\rho_-+(1-t)\rho_+.
\end{equation}
After subtracting the harmless constant $z_i$ and using $y=s$, the bottom
trace is $w^2h(y)$, where
\begin{equation}\label{H0:eq:h}
 h(y)=\frac18(y-t)^2
 +\frac12(y-t)\bigl((1-y)\rho_-+y\rho_+\bigr).
\end{equation}

\begin{lemma}[Exact edge energy]\label{H0:lem:edge-energy}
With
\begin{equation}
 K_{\rm cyl}=\frac{\zeta(3)}{\pi^3},
\end{equation}
the Neumann half-strip energy of \eqref{H0:eq:h} is
\begin{equation}\label{H0:eq:edge-energy}
 \mathcal N_i(G_\alpha+T_\alpha^0z)
 =\frac{K_{\rm cyl}}4w^4
 \left[\left(\delta+\frac14\right)^2
 +7\left(q+\frac X4\right)^2\right].
\end{equation}
\end{lemma}

\begin{proof}
Let $u=h'(0)$ and $v=h'(1)$.  Twice integrating the cosine coefficients by
parts gives, for $k\ge1$,
\begin{equation}
 2\int_0^1h(y)\cos(k\pi y)\,dy
 =\frac{2\bigl((-1)^kv-u\bigr)}{(k\pi)^2}.
\end{equation}
The energy of this cosine mode in the unit half-strip is $k\pi/2$ times
the square of its coefficient.  Since
\begin{equation}
 \sum_{k\ {\rm even}}\frac1{k^3}=\frac18\zeta(3),
\qquad
 \sum_{k\ {\rm odd}}\frac1{k^3}=\frac78\zeta(3),
\end{equation}
the total energy is
\begin{equation}
 \frac{K_{\rm cyl}}4w^4
 \bigl((v-u)^2+7(v+u)^2\bigr).
\end{equation}
Direct differentiation of \eqref{H0:eq:h} gives
\begin{equation}
 v-u=\delta+\frac14,
 \qquad v+u=q+\frac X4,
\end{equation}
which proves \eqref{H0:eq:edge-energy}.
\end{proof}

\begin{lemma}[Rational edge Bellman inequality]\label{H0:lem:bellman}
For every $0\le t\le1$ and every $\rho_-,\rho_+\in\R$,
\begin{align}
 &\frac14\left[\left(\delta+\frac14\right)^2
 +7\left(q+\frac X4\right)^2\right]\notag\\
 &\quad\ge
 \frac1{10}\bigl(t^4+(1-t)^4\bigr)+\frac1{320}
 +(1-t)^3\rho_+-t^3\rho_- .
 \label{H0:eq:bellman}
\end{align}
Equality holds at $t=1/2$ and $\rho_-=\rho_+=0$.
\end{lemma}

\begin{proof}
Put
\begin{equation}
 s=t(1-t),\qquad
 d=s(1-2s),\qquad e=X(1-s).
\end{equation}
The last two terms on the right of \eqref{H0:eq:bellman} that depend on the
slopes satisfy
\begin{equation}
 (1-t)^3\rho_+-t^3\rho_-=d\delta+eq.
\end{equation}
Completing the two squares shows that the left side of
\eqref{H0:eq:bellman} minus its right side is
\begin{align}
 &\frac14\left(\delta+\frac14-2d\right)^2
 +\frac74\left(q+\frac X4-\frac{2e}{7}\right)^2\notag\\
 &\quad+X^2\left(
 \frac9{280}-\frac{X^2}{280}
 -\frac{X^4}{112}-\frac{X^6}{64}\right).
 \label{H0:eq:bellman-squares}
\end{align}
The bracket is decreasing as a function of $X^2\in[0,1]$ and its value at
$X^2=1$ is $9/2240$.  Hence every term in
\eqref{H0:eq:bellman-squares} is nonnegative.  The stated equality is immediate.
\end{proof}

Restoring $a,b,w$ in Lemma~\ref{H0:lem:bellman} and using
Lemma~\ref{H0:lem:edge-energy} gives the pointwise inequality
\begin{equation}\label{H0:eq:physical-bellman}
 \mathcal N_i(G_\alpha+T_\alpha^0z)
 \ge K_{\rm cyl}\left\{\frac1{10}(a^4+b^4)
 +\frac1{320}(a+b)^4+b^3r_i-a^3r_{i-1}\right\}.
\end{equation}

\subsection{The replacement no-ratio gap}

\begin{theorem}[Corrected homogeneous $|D|$ gap]\label{H0:thm:dot-gap}
Let
\begin{equation}
 \bar\alpha=\frac{2\pi}{N},
 \qquad
 \Jfour(\alpha)=\sum_i\alpha_i^4-N\bar\alpha^4.
\end{equation}
For every positive cyclic fan,
\begin{equation}\label{H0:eq:gap}
 \boxed{
 \PGdot(\alpha)-\PGdot(\bar\alpha\mathbf1)
 \ge\frac{\zeta(3)}{10\pi^4}\Jfour(\alpha).}
\end{equation}
The conclusion remains valid after imposing the homogeneous area and
centering rows described after \eqref{H0:eq:PGplus}.
\end{theorem}

\begin{proof}
Sum \eqref{H0:eq:physical-bellman}.  The slope potential telescopes exactly:
\begin{equation}
 \sum_i(\alpha_i^3r_i-\alpha_{i-1}^3r_{i-1})=0.
\end{equation}
Together with \eqref{H0:eq:strip-localization}, this gives, pointwise in $z$,
\begin{equation}\label{H0:eq:global-lower}
 2\pi\|G_\alpha+T_\alpha^0z\|_{\dot H_*^{1/2}}^2
 \ge K_{\rm cyl}\left[
 \frac15\sum_i\alpha_i^4
 +\frac1{320}\sum_i(\alpha_{i-1}+\alpha_i)^4\right].
\end{equation}
For the uniform fan, $z=0$ makes the periodic harmonic extension even
across every vertical wall.  Thus strip localization is an equality.
Lemmas~\ref{H0:lem:edge-energy} and \ref{H0:lem:bellman} are also equalities, and
\begin{equation}\label{H0:eq:uniform}
 2\pi\PGdot(\bar\alpha\mathbf1)
 =\frac{K_{\rm cyl}}4N\bar\alpha^4.
\end{equation}
Finally,
\begin{equation}
 \sum_i(\alpha_{i-1}+\alpha_i)^4
 \ge16N\bar\alpha^4
\end{equation}
by Jensen's inequality.  Subtract \eqref{H0:eq:uniform} from
\eqref{H0:eq:global-lower}, take the infimum over $z$, and use
$K_{\rm cyl}/(2\pi)=\zeta(3)/(2\pi^4)$.  This proves
\eqref{H0:eq:gap}.
\end{proof}

\subsection{The identity multiplier has a nonnegative relative defect}

For a function $f$ on the circle, set
\begin{equation}
 \|f\|_{L_*^2}^2
 =\frac1{2\pi}\int_{\Sone}|f-\mathop{\rm avg}f|^2.
\end{equation}
On one edge the trace, after a cell constant is removed, is $w^2h(y)$ with
$h$ given by \eqref{H0:eq:h} and $d\theta=(w/2)dy$.  Write
\begin{equation}
 \operatorname {Var}(h)=\int_0^1h^2-\left(\int_0^1h\right)^2.
\end{equation}

\begin{lemma}[Exact rational variance inequality]\label{H0:lem:variance}
For all $0\le t\le1$ and all $\rho_-,\rho_+\in\R$,
\begin{equation}\label{H0:eq:variance}
 \frac12\operatorname {Var}(h)
 \ge\frac1{23040}
 +\frac1{180}\left((1-t)^4\rho_+-t^4\rho_-\right).
\end{equation}
Equality holds at $t=1/2$ and $\rho_-=\rho_+=0$.
\end{lemma}

\begin{proof}
Up to an irrelevant constant, write $h(y)=Ay^2+By$, where
\begin{equation}
 A=\frac18+\frac\delta2,
 \qquad B=\frac q2-\frac\delta2-\frac t4.
\end{equation}
The elementary moments of $y$ and $y^2$ give
\begin{equation}\label{H0:eq:variance-polynomial}
 \frac12\operatorname {Var}(h)
 =\frac{2A^2}{45}+\frac{B^2}{24}+\frac{AB}{12}.
\end{equation}
Put
\begin{equation}
 s=t(1-t),\quad X=1-2t,\quad
 d_4=s(1-3s),\quad e_4=X(1-2s).
\end{equation}
Then
\begin{equation}
 (1-t)^4\rho_+-t^4\rho_-=d_4\delta+e_4q.
\end{equation}
Define
\begin{equation}
 \delta_*=4d_4-\frac14,
 \qquad q_*=\frac4{15}e_4-\frac X4.
\end{equation}
Substitution in \eqref{H0:eq:variance-polynomial} and completion of the two
independent squares gives the exact identity
\begin{align}
 &\frac12\operatorname {Var}(h)-\frac1{23040}
 -\frac1{180}(d_4\delta+e_4q)\notag\\
 &\quad=\frac{(\delta-\delta_*)^2}{1440}
 +\frac{(q-q_*)^2}{96}\notag\\
 &\qquad+\frac{X^2}{345600}
 \left(176+52X^2+116X^4-135X^6\right).
 \label{H0:eq:variance-squares}
\end{align}
For $|X|\le1$, the last polynomial is at least
$176-135=41$.  This proves \eqref{H0:eq:variance} and its equality statement.
\end{proof}

Restoring the physical widths in \eqref{H0:eq:variance} gives
\begin{equation}\label{H0:eq:physical-variance}
 \min_{C\in\R}\int_{K_i}
 |G_\alpha+T_\alpha^0z+C|^2
 \ge\frac{(a+b)^5}{23040}
 +\frac1{180}(b^4r_i-a^4r_{i-1}).
\end{equation}
Summing removes the last potential.  Since a single global mean is no better
than independent cell constants,
\begin{equation}\label{H0:eq:global-L2}
 2\pi\|G_\alpha+T_\alpha^0z\|_{L_*^2}^2
 \ge\frac1{23040}\sum_i(\alpha_{i-1}+\alpha_i)^5.
\end{equation}
At the uniform fan, $z=0$ gives equality in \eqref{H0:eq:global-L2}.  Jensen's
inequality therefore shows that the identity-multiplier relative defect is
nonnegative.

\begin{corollary}[Corrected homogeneous $|D|+I$ gap]\label{H0:cor:plus-gap}
For every positive cyclic fan,
\begin{equation}\label{H0:eq:plus-gap}
 \boxed{
 \PGplus(\alpha)-\PGplus(\bar\alpha\mathbf1)
 \ge\frac{\zeta(3)}{10\pi^4}\Jfour(\alpha).}
\end{equation}
The conclusion remains valid with the homogeneous area and centering rows.
\end{corollary}

\begin{proof}
Add \eqref{H0:eq:global-L2} to the pointwise order-one estimate
\eqref{H0:eq:global-lower}, subtract the two equalities at the uniform fan,
and use
\begin{equation}
 \sum_i(\alpha_{i-1}+\alpha_i)^5
 \ge N(2\bar\alpha)^5.
\end{equation}
Taking the infimum over the same coefficient space proves
\eqref{H0:eq:plus-gap}.
\end{proof}
\endgroup

\begingroup
\section{Foundation F65: the exact angular disk endpoint}
\label{module:F65}
\noindent\textit{Audited source: \nolinkurl{convex_largeN_polya_szego_stage65_correct_source_disk_endpoint.tex}.}
\subsection{Fan, correct source, and the two support spaces}

Let
\begin{equation}\label{F65:eq:fan}
 \alpha_i>0,\qquad \sum_i\alpha_i=2\pi,
 \qquad \theta_{i+1}-\theta_i=\alpha_i,
 \qquad \bar\alpha=\frac{2\pi}{N},
\end{equation}
and put
\begin{equation}\label{F65:eq:SJ}
 S=\alpha_{\max}+\bar\alpha,
 \qquad
 \Jfour(\alpha)=\sum_i\alpha_i^4-N\bar\alpha^4.
\end{equation}
Set $\phi_i=\theta_i+\alpha_i/2$ and
$K_i=[\phi_{i-1},\phi_i]$.  On $K_i$ write
\begin{equation}
 a=\alpha_{i-1},\quad b=\alpha_i,\quad
 x=\theta-\theta_i,\quad
 s=\frac{x+a/2}{(a+b)/2},
\end{equation}
and define the correct principal source
\begin{equation}\label{F65:eq:G}
 G_\alpha(\theta)=\frac{x^2}{2}.
\end{equation}

For a coefficient vector $z$, put
\begin{equation}
 r_- =\frac{z_i-z_{i-1}}a,
 \qquad r_+=\frac{z_{i+1}-z_i}b.
\end{equation}
The homogeneous and exact vertex-compatible traces are
\begin{align}
 (T_\alpha^0z)_i(x)
 &=z_i+x\bigl((1-s)r_-+sr_+\bigr),\label{F65:eq:T0}\\
 (T_\alpha z)_i(x)
 &=z_i\cos x+\bigl((1-s)a_i^-+sa_i^+\bigr)\sin x,
 \label{F65:eq:T}
\end{align}
where
\begin{equation}\label{F65:eq:apm}
 a_i^-=r_-\frac a{\sin a}-z_i\tan\frac a2,
 \qquad
 a_i^+=r_+\frac b{\sin b}+z_i\tan\frac b2.
\end{equation}
Both traces are continuous at the vertices $\phi_i$.

The homogeneous and exact area rows are
\begin{align}
 L_\alpha^0z
 &=\frac12\sum_i(\alpha_{i-1}+\alpha_i)z_i,\label{F65:eq:L0}\\
 L_\alpha z
 &=\sum_i\left(\tan\frac{\alpha_{i-1}}2
 +\tan\frac{\alpha_i}2\right)z_i.\label{F65:eq:L}
\end{align}
Let $Z_\alpha^0=\ker L_\alpha^0$ and $Z_\alpha=\ker L_\alpha$.
Centering rows may be added throughout; translations will be removed exactly
in Section~\ref{F65:sec:translations}.

Work in $H=H^{1/2}(\Sone)/\R$ with
\begin{equation}\label{F65:eq:qplus}
 \qplus(f)=\sum_{n\ne0}(|n|+1)|\widehat f_n|^2,
 \qquad
 \widehat f_n=\frac1{2\pi}\int_0^{2\pi}f(\theta)e^{-in\theta}\,d\theta.
\end{equation}
Set
\begin{equation}
 V_\alpha^0=T_\alpha^0Z_\alpha^0,
 \qquad V_\alpha=T_\alpha Z_\alpha,
\end{equation}
and
\begin{align}
 P_0^G(\alpha)&=\inf_{v\in V_\alpha^0}\qplus(G_\alpha+v),
 \label{F65:eq:P0}\\
 P_{\rm tr}^G(\alpha)&=\inf_{v\in V_\alpha}\qplus(G_\alpha+v).
 \label{F65:eq:Ptr}
\end{align}
The corrected homogeneous theorem gives
\begin{equation}\label{F65:eq:input-gap}
 P_0^G(\alpha)-P_0^G(\bar\alpha\mathbf1)
 \ge\frac{\zeta(3)}{10\pi^4}\Jfour(\alpha).
\end{equation}

\subsection{No-ratio exact-trace graph}

\begin{lemma}[Local jet coercivity]\label{F65:lem:jet}
There are absolute constants $0<c<C$ such that, with $w=a+b$,
\begin{equation}\label{F65:eq:jet}
 cw\bigl(z_i^2+w^2r_-^2+w^2r_+^2\bigr)
 \le\|T_\alpha^0z\|_{L^2(K_i)}^2
 \le Cw\bigl(z_i^2+w^2r_-^2+w^2r_+^2\bigr).
\end{equation}
The constants are independent of $a/b$.
\end{lemma}

\begin{proof}
After $x=(w/2)(y-t)$ with $t=a/w$, the three normalized polynomials are
\begin{equation}
 1,\qquad (y-t)(1-y),\qquad (y-t)y,
 \qquad 0\le y\le1.
\end{equation}
The matrix carrying the two slope variables to the linear and quadratic
coefficients has determinant one and it and its inverse are uniformly
bounded for $0\le t\le1$.  Finite-dimensional norm equivalence proves
\eqref{F65:eq:jet}.
\end{proof}

\begin{lemma}[Exact-minus-homogeneous graph]\label{F65:lem:graph}
For $S$ sufficiently small there is a bijection
$C_\alpha:Z_\alpha^0\to Z_\alpha$ for which
\begin{equation}
 J_\alpha(T_\alpha^0z)=T_\alpha C_\alpha z
\end{equation}
satisfies
\begin{equation}\label{F65:eq:Jclose}
 \|(J_\alpha-I)v\|_H\le CS^{3/2}\|v\|_H,
 \qquad v\in V_\alpha^0.
\end{equation}
\end{lemma}

\begin{proof}
Taylor's formula in \eqref{F65:eq:T}--\eqref{F65:eq:apm} gives on $K_i$
\begin{align}
 \|T_\alpha z-T_\alpha^0z\|_{L^\infty(K_i)}
 &\le C\bigl(w^2|z_i|+w^3(|r_-|+|r_+|)\bigr),\notag\\
 \|\partial_\theta(T_\alpha z-T_\alpha^0z)\|_{L^\infty(K_i)}
 &\le C\bigl(w|z_i|+w^2(|r_-|+|r_+|)\bigr).
 \label{F65:eq:trace-jets}
\end{align}
The error is continuous at every wall.  Summing \eqref{F65:eq:trace-jets},
using Lemma~\ref{F65:lem:jet}, cyclic Poincar\'e on the area kernel, and
$\|f\|_{H^{1/2}}^2\le C\|f\|_{L^2}\|f\|_{H^1}$ gives
\begin{equation}\label{F65:eq:trace-close}
 \|T_\alpha z-T_\alpha^0z\|_H
 \le CS^{3/2}\|T_\alpha^0z\|_H.
\end{equation}

For $z\in Z_\alpha^0$, define
\begin{equation}
 C_\alpha z=z-\frac{L_\alpha z}{L_\alpha\mathbf1}\mathbf1.
\end{equation}
Since
\begin{equation}
 \left|\tan\frac u2-\frac u2\right|\le Cu^3,
\end{equation}
the coefficient correction is $O(S^2)$ in the trace metric.  The analogous
formula with $L_\alpha^0$ is the inverse map.  Combining this row correction
with \eqref{F65:eq:trace-close} proves \eqref{F65:eq:Jclose}.
\end{proof}

\begin{lemma}[Hilbert Schur perturbation]\label{F65:lem:schur}
If $J:V_0\to V_1$ is a bijection with
$\|(J-I)v\|\le\eta\|v\|$, $\eta<1/2$, then
\begin{equation}\label{F65:eq:schur}
 \left|\operatorname {dist}(b,V_1)^2
 -\operatorname {dist}(b,V_0)^2\right|
 \le C\eta\|b\|^2.
\end{equation}
\end{lemma}

\begin{proof}
Map the orthogonal projection of $-b$ onto $V_0$ through $J$ and use it as
a competitor in $V_1$.  Apply the same argument to $J^{-1}$.
\end{proof}

The continuous edge quadratic satisfies the no-ratio spline bound
\begin{equation}\label{F65:eq:Gbound}
 \qplus(G_\alpha)\le C\sum_i\alpha_i^4\le CS^3.
\end{equation}
Indeed, affine normalization on each $K_i$ controls the same-cell and
adjacent-cell terms, while dyadic annuli control the separated-cell part.
Lemma~\ref{F65:lem:schur} therefore yields, for
\begin{equation}
 R_{\rm tr}^G=P_{\rm tr}^G-P_0^G,
\end{equation}
the absolute estimate
\begin{equation}\label{F65:eq:trace-absolute}
 |R_{\rm tr}^G(\alpha)|\le CS^{9/2}.
\end{equation}

\subsection{Relative trace comparison}

\begin{proposition}[Quasiuniform trace Hessian]\label{F65:prop:trace-hessian}
Suppose
\begin{equation}\label{F65:eq:quasiuniform}
 (1-\eta_0)a\le\gamma_i\le(1+\eta_0)a,
 \qquad \sum_i\gamma_i=2\pi,
\end{equation}
where $0<\eta_0<1/4$.  If $\sum_i\delta_i=0$, then
\begin{equation}\label{F65:eq:trace-hessian}
 |D^2R_{\rm tr}^G(\gamma)[\delta,\delta]|
 \le Ca^{1/2}\sum_i\gamma_i^2\delta_i^2.
\end{equation}
\end{proposition}

\begin{proof}
If $\delta\ne0$, complexify
$\gamma(z)=\gamma+z\delta$ in the disk
\begin{equation}
 |z|<r,\qquad r=\frac{c_0a}{\|\delta\|_{\ell^\infty}}.
\end{equation}
For every cyclic arc $I$,
\begin{equation}
 \left|z\sum_{i\in I}\delta_i\right|
 \le c_0a\#I
 \le Cc_0\sum_{i\in I}\gamma_i.
\end{equation}
Thus the cellwise affine material map changes every arc length by a fixed
small relative amount.  The pulled-back Gagliardo kernel, the $L^2$ row,
the two traces, the area graph, and the polynomial source $G_{\gamma(z)}$
are holomorphic.  Their Gram inverses stay uniformly bounded by
Lemma~\ref{F65:lem:jet}.  The proof of \eqref{F65:eq:trace-absolute} applies to the
complex bilinear Schur formula and gives
\begin{equation}
 |R_{\rm tr}^G(\gamma(z))|\le Ca^{9/2}.
\end{equation}
Cauchy's estimate gives
\begin{equation}
 |D^2R_{\rm tr}^G(\gamma)[\delta,\delta]|
 \le Ca^{5/2}\|\delta\|_{\ell^\infty}^2.
\end{equation}
Use
$\sum_i\gamma_i^2\delta_i^2\ge c a^2
\|\delta\|_{\ell^\infty}^2$ to obtain \eqref{F65:eq:trace-hessian}.
\end{proof}

\begin{lemma}[Relative trace remainder]\label{F65:lem:trace-relative}
For every sufficiently fine positive fan,
\begin{equation}\label{F65:eq:trace-relative}
 |R_{\rm tr}^G(\alpha)-R_{\rm tr}^G(\bar\alpha\mathbf1)|
 \le C(S^{1/4}+S^{1/2})\Jfour(\alpha).
\end{equation}
\end{lemma}

\begin{proof}
Put $a=\bar\alpha$ and $d_i=\alpha_i-a$.  The exact identity
\begin{equation}\label{F65:eq:Jidentity}
 \Jfour(\alpha)
 =\sum_i d_i^2(\alpha_i^2+2a\alpha_i+3a^2)
\end{equation}
implies $\max_i|d_i|^4\le C\Jfour(\alpha)$.
If $\Jfour\ge S^{17/4}$, use \eqref{F65:eq:trace-absolute}.  Otherwise
$\max|d_i|/S\le CS^{1/16}$, so the segment
$a\mathbf1+td$ is quasiuniform.  Cyclic invariance makes the first
derivative of $R_{\rm tr}^G$ vanish at $a\mathbf1$ on the zero-sum
direction.  Taylor's formula, Proposition~\ref{F65:prop:trace-hessian}, and
\eqref{F65:eq:Jidentity} give the $CS^{1/2}\Jfour$ bound.  The two cases prove
\eqref{F65:eq:trace-relative}.
\end{proof}

\begin{theorem}[Correct-source exact-trace gap]\label{F65:thm:trace-gap}
There is $S_{\rm tr}>0$ such that, for $S\le S_{\rm tr}$,
\begin{equation}\label{F65:eq:trace-gap}
 \boxed{
 P_{\rm tr}^G(\alpha)-P_{\rm tr}^G(\bar\alpha\mathbf1)
 \ge\frac{\zeta(3)}{20\pi^4}\Jfour(\alpha).}
\end{equation}
\end{theorem}

\begin{proof}
Combine \eqref{F65:eq:input-gap} and \eqref{F65:eq:trace-relative}, decreasing
$S_{\rm tr}$ until the latter coefficient is at most
$\zeta(3)/(20\pi^4)$.
\end{proof}

\subsection{Exact translations}\label{F65:sec:translations}

\begin{lemma}[Translation quotient]\label{F65:lem:translations}
For $c\in\R^2$, let $z_i=c\cdot(\cos\theta_i,\sin\theta_i)$.  Then
\begin{equation}\label{F65:eq:translation}
 L_\alpha z=0,
 \qquad T_\alpha z(\theta)=c\cdot(\cos\theta,\sin\theta).
\end{equation}
Consequently
\begin{equation}\label{F65:eq:translation-quotient}
 P_{\rm tr}^G(\alpha)
 =\inf_{v\in\PiNG V_\alpha}
 \qplus(\PiNG G_\alpha+v),
\end{equation}
where $\PiNG$ removes the modes $0,\pm1$.
\end{lemma}

\begin{proof}
With $\tau_i=(-\sin\theta_i,\cos\theta_i)$, the identities
\begin{equation}
 \nu_{i-1}=\nu_i\cos a-\tau_i\sin a,
 \qquad
 \nu_{i+1}=\nu_i\cos b+\tau_i\sin b
\end{equation}
give $a_i^-=a_i^+=c\cdot\tau_i$ in \eqref{F65:eq:apm}.  Hence
\eqref{F65:eq:T} equals $c\cdot\nu(\theta)$.  Translation invariance of area
gives $L_\alpha z=0$.  Orthogonality of the Fourier modes proves
\eqref{F65:eq:translation-quotient}.
\end{proof}

\subsection{Bessel order-minus-one reduction}

Let $j=j_{0,1}$ and, for $n\ge2$, set
\begin{equation}
 k_n=j\frac{J_{n+1}(j)}{J_n(j)}.
\end{equation}
On the range of $\PiNG$, define
\begin{align}
 k(f,g)&=\sum_{|n|\ge2}k_{|n|}\widehat f_n\widehat g_{-n},\notag\\
 \qdisk(f,g)&=\qplus(f,g)-k(f,g).
 \label{F65:eq:qdisk}
\end{align}
The Bessel recurrence and its alternating fixed-argument series give
\begin{equation}\label{F65:eq:bessel-bounds}
 0<k_n<\frac4n,
 \qquad \frac13\qplus(f)\le\qdisk(f)\le\qplus(f).
\end{equation}
Put
\begin{equation}
 b_\alpha^G=\PiNG G_\alpha,
 \qquad W_\alpha=\PiNG V_\alpha,
\end{equation}
and
\begin{equation}\label{F65:eq:PDG}
 P_D^G(\alpha)=\inf_{v\in W_\alpha}\qdisk(b_\alpha^G+v).
\end{equation}

\begin{lemma}[No-ratio exact-trace quasi-interpolant]
\label{F65:lem:QI}
There is a linear operator
\begin{equation}
 I_\alpha:H_{\rm ng}^{3/2}(\Sone)\longrightarrow W_\alpha
\end{equation}
such that, for every sufficiently fine positive fan,
\begin{equation}\label{F65:eq:interpolant}
 \|\psi-I_\alpha\psi\|_{H^{1/2}}
 \le CS\|\psi\|_{H^{3/2}}.
\end{equation}
The constant is independent of the smallest cell and every adjacent
mesh ratio.
\end{lemma}

\begin{proof}
It is enough to work first in the homogeneous trace space.  Put
\(M=\lfloor S^{-1}\rfloor\) and let \(\phi=P_{\le M}\psi\) be the
Fourier projection onto \(|n|\le M\).  Directly from the Fourier norms,
\begin{equation}\label{F65:eq:QI-tail}
 \|\psi-\phi\|_{H^{1/2}}
 \le CS\|\psi\|_{H^{3/2}},
 \qquad
 \|\phi''\|_{L^2}
 \le CS^{-1/2}\|\psi\|_{H^{3/2}}.
\end{equation}

Take the nodal values \(z_i=\phi(\theta_i)\) and set
\(Q_\alpha\phi=T_\alpha^0z\).  We give the global estimate rather than
invoke a mesh-regular interpolation theorem.  On
\(K_i=[\phi_{i-1},\phi_i]\), write
\(a=\alpha_{i-1}\), \(b=\alpha_i\), \(w=a+b\), and
\(\omega_i=[\theta_{i-1},\theta_{i+1}]\).  Since
\begin{equation*}
 r_-=\frac1a\int_{\theta_{i-1}}^{\theta_i}\phi'(u)\,du,
 \qquad
 r_+=\frac1b\int_{\theta_i}^{\theta_{i+1}}\phi'(u)\,du,
\end{equation*}
Cauchy--Schwarz and the fundamental theorem of calculus give
\begin{equation}\label{F65:eq:QI-slope}
 |r_--\phi'(\theta_i)|+|r_+-\phi'(\theta_i)|
 \le Cw^{1/2}\|\phi''\|_{L^2(\omega_i)}.
\end{equation}
The trace formula \eqref{F65:eq:T0} reproduces every affine function.  Using
\eqref{F65:eq:QI-slope} in that formula, and using Taylor's integral formula
for \(\phi\), yields for \(x\in K_i\)
\begin{align}
 |\phi(x)-Q_\alpha\phi(x)|
 &\le Cw^{3/2}\|\phi''\|_{L^2(\omega_i)},\label{F65:eq:QI-point}\\
 |\phi'(x)-(Q_\alpha\phi)'(x)|
 &\le Cw^{1/2}\|\phi''\|_{L^2(\omega_i)}.\label{F65:eq:QI-derivative}
\end{align}
Indeed, after subtracting the affine Taylor polynomial at \(\theta_i\),
the first row consists of \(x\) times a convex combination of the two
slope errors, and the differentiated row contains in addition
\(2x(r_+-r_-)/w\); all its coefficients are bounded for
\(-a/2\le x\le b/2\), including when \(a/b\) tends to zero or infinity.

The intervals \(\omega_i\) have overlap at most two when their
\(L^2\)-mass is assigned back to the two constituent normal cells.
Since \(w\le2S\), integration of
\eqref{F65:eq:QI-point}--\eqref{F65:eq:QI-derivative} and summation give
\begin{equation}\label{F65:eq:QI-L2H1}
 \|\phi-Q_\alpha\phi\|_{L^2}
 \le CS^2\|\phi''\|_{L^2},
 \qquad
 \|\phi-Q_\alpha\phi\|_{H^1}
 \le CS\|\phi''\|_{L^2}.
\end{equation}
The Fourier interpolation inequality
\(\|e\|_{H^{1/2}}^2\le\|e\|_{L^2}\|e\|_{H^1}\),
\eqref{F65:eq:QI-tail}, and \eqref{F65:eq:QI-L2H1} imply
\begin{equation}\label{F65:eq:QI-homogeneous}
 \|\phi-Q_\alpha\phi\|_{H^{1/2}}
 \le CS^{3/2}\|\phi''\|_{L^2}
 \le CS\|\psi\|_{H^{3/2}}.
\end{equation}

It remains only to impose the exact finite rows.  Replace \(z\) by
\begin{equation*}
 z^0=z-\frac{L_\alpha^0z}{L_\alpha^0\mathbf1}\mathbf1.
\end{equation*}
Because \(T_\alpha^0\mathbf1=1\), this changes the homogeneous trace only
by a constant.  Put \(v^0=T_\alpha^0z^0\in V_\alpha^0\) and define
\begin{equation}
 I_\alpha\psi=\PiNG J_\alpha v^0\in W_\alpha,
\end{equation}
where \(J_\alpha\) is the exact area graph of
Lemma~\ref{F65:lem:graph}.  The projection \(\PiNG\) is an orthogonal
contraction in every Fourier Sobolev norm, \(\psi=\PiNG\psi\), and
\eqref{F65:eq:Jclose} gives
\begin{equation*}
 \|\PiNG(J_\alpha-I)v^0\|_{H^{1/2}}
 \le CS^{3/2}\|v^0\|_{H^{1/2}}
 \le CS\|\psi\|_{H^{3/2}}.
\end{equation*}
Together with \eqref{F65:eq:QI-tail} and
\eqref{F65:eq:QI-homogeneous}, this is \eqref{F65:eq:interpolant}.
\end{proof}

\begin{lemma}[Galerkin negative-norm gain]\label{F65:lem:negative}
Let $f_\alpha=b_\alpha^G+v_\alpha$ be the $\qplus$-minimizing residual.
Then
\begin{equation}\label{F65:eq:negative}
 \|f_\alpha\|_{H^{-1/2}}
 \le CS\|f_\alpha\|_{H^{1/2}}.
\end{equation}
\end{lemma}

\begin{proof}
For $g\in H_{\rm ng}^{1/2}$ solve $(|D|+I)\psi=g$.  Galerkin
orthogonality and Lemma~\ref{F65:lem:QI} give
\begin{equation}
 |\langle f_\alpha,g\rangle|
 =|\qplus(f_\alpha,\psi-I_\alpha\psi)|
 \le CS\|f_\alpha\|_{H^{1/2}}\|g\|_{H^{1/2}}.
\end{equation}
Taking the dual supremum proves \eqref{F65:eq:negative}.
\end{proof}

Let
\begin{equation}
 R_B^G(\alpha)=P_{\rm tr}^G(\alpha)-P_D^G(\alpha).
\end{equation}
Completing the disk square gives the exact identity
\begin{equation}\label{F65:eq:bessel-schur}
 R_B^G
 =k(f_\alpha,f_\alpha)
 +\langle g_\alpha,A_{D,\alpha}^{-1}g_\alpha\rangle,
 \qquad g_\alpha(v)=k(f_\alpha,v),
\end{equation}
where $A_{D,\alpha}=\qdisk|_{W_\alpha}$.
Equations \eqref{F65:eq:bessel-bounds}, \eqref{F65:eq:negative},
\eqref{F65:eq:Gbound}, and \eqref{F65:eq:bessel-schur} imply
\begin{equation}\label{F65:eq:bessel-absolute}
 0\le R_B^G(\alpha)\le CS^5.
\end{equation}

\begin{proposition}[Relative Bessel correction]\label{F65:prop:bessel-relative}
For every sufficiently fine positive fan,
\begin{equation}\label{F65:eq:bessel-relative}
 |R_B^G(\alpha)-R_B^G(\bar\alpha\mathbf1)|
 \le C(S^{1/2}+S)\Jfour(\alpha).
\end{equation}
\end{proposition}

\begin{proof}
In a quasiuniform band, pull all functions to a fixed real fan.  If
$h=\|\delta\|_{\ell^\infty}/a$, the first two material derivatives of
$\qplus$ are bounded by $Ch^r$ in $H^{1/2}$, while those of $k$ are bounded
by $Ch^r$ in $H^{-1/2}$.  The latter follows from
\begin{equation}
 |\Delta^rk_n|\le C_r(1+n)^{-1-r},\qquad 0\le r\le3,
\end{equation}
and three discrete summations by parts in the pulled-back kernel.
Differentiating the Galerkin equation and \eqref{F65:eq:negative} twice gives
\begin{equation}
 \|\partial_t^rf_t\|_{H^{1/2}}\le Ch^ra^{3/2},
 \qquad
 \|\partial_t^rf_t\|_{H^{-1/2}}\le Ch^ra^{5/2},
 \quad r=0,1,2.
\end{equation}
Twice differentiating \eqref{F65:eq:bessel-schur} therefore yields
\begin{equation}\label{F65:eq:bessel-hessian}
 |D^2R_B^G(\gamma)[\delta,\delta]|
 \le Ca\sum_i\gamma_i^2\delta_i^2.
\end{equation}

Use \eqref{F65:eq:bessel-absolute} when
$\Jfour\ge S^{9/2}$.  In the complementary regime,
\eqref{F65:eq:Jidentity} makes the radial segment from the uniform fan
quasiuniform; cyclic invariance kills the first derivative there, and
\eqref{F65:eq:bessel-hessian} with Taylor's formula gives
$CS\Jfour$.  This proves \eqref{F65:eq:bessel-relative}.
\end{proof}

\begin{corollary}[Correct-source Bessel gap]\label{F65:cor:bessel-gap}
There is $S_D>0$ such that, for $S\le S_D$,
\begin{equation}\label{F65:eq:bessel-gap}
 \boxed{
 P_D^G(\alpha)-P_D^G(\bar\alpha\mathbf1)
 \ge\frac{\zeta(3)}{40\pi^4}\Jfour(\alpha).}
\end{equation}
\end{corollary}

\begin{proof}
Subtract \eqref{F65:eq:bessel-relative} from \eqref{F65:eq:trace-gap} and decrease
$S_D$ until the relative error is at most $\zeta(3)/(40\pi^4)$.
\end{proof}

\subsection{The exact secant source}

Let
\begin{equation}\label{F65:eq:secant}
 A_\alpha=\sum_i\tan\frac{\alpha_i}{2},
 \qquad h_\alpha=\left(\frac\pi{A_\alpha}\right)^{1/2},
 \qquad
 \rho_\alpha^{\rm tan}=h_\alpha\sec x-1
 \quad\hbox{on }K_i.
\end{equation}
Writing $c_\alpha=h_\alpha-1$, one has
\begin{equation}\label{F65:eq:secant-split}
 \rho_\alpha^{\rm tan}=c_\alpha+G_\alpha+E_\alpha,
 \qquad
 E_\alpha=h_\alpha(\sec x-1)-\frac{x^2}{2}.
\end{equation}

\begin{lemma}[Scaled-spline secant remainder]\label{F65:lem:secant-remainder}
For every sufficiently fine positive fan,
\begin{equation}\label{F65:eq:secant-remainder}
 \|E_\alpha\|_{H^{1/2}}
 \le CS^2\|G_\alpha\|_{H^{1/2}}.
\end{equation}
The estimate is independent of the smallest cell.
\end{lemma}

\begin{proof}
On $K_i$, normalized by $x=(a+b)(y-t)/2$, both $G_\alpha$ and
$E_\alpha$ belong to a fixed compact family of even analytic profiles.
Taylor's formula gives, through one normalized derivative,
\begin{equation}\label{F65:eq:normalized-remainder}
 |E_\alpha|+(a+b)|E_\alpha'|
 \le CS^2(a+b)^2.
\end{equation}
The values from the two incident cells agree at each vertex.

For completeness, the mesh-independent assembly is as follows.  Split the
Slobodeckij integral into same-cell, adjacent-cell, and separated-cell
pieces.  Affine normalization and finite-dimensional norm equivalence apply
to the first two pieces.  For a separated pair $K_i,K_j$, subtract the cell
means and use
\begin{equation}
 \iint_{K_i\times K_j}\frac{d\theta\,d\varphi}
 {\operatorname {dist}_{\Sone}(\theta,\varphi)^2}
 \le C\frac{|K_i||K_j|}{d_{ij}^2}.
\end{equation}
Summing in dyadic annuli bounds the mean-difference energy by the same-cell
and adjacent-cell charges.  Equation \eqref{F65:eq:normalized-remainder} makes
every charge of $E_\alpha$ at most $CS^2$ times the corresponding charge of
$G_\alpha$.  This proves \eqref{F65:eq:secant-remainder} without a mesh-ratio
assumption.
\end{proof}

Define the exact tangential-source disk value
\begin{equation}\label{F65:eq:PDtan}
 P_D^{\rm tan}(\alpha)
 =\inf_{v\in W_\alpha}
 \qdisk\bigl(\PiNG(\rho_\alpha^{\rm tan}-c_\alpha)+v\bigr)
\end{equation}
and
\begin{equation}
 R_{\rm sec}(\alpha)=P_D^{\rm tan}(\alpha)-P_D^G(\alpha).
\end{equation}
Coercivity of \eqref{F65:eq:bessel-bounds}, the Hilbert distance inequality,
Lemma~\ref{F65:lem:secant-remainder}, and \eqref{F65:eq:Gbound} give
\begin{equation}\label{F65:eq:secant-absolute}
 |R_{\rm sec}(\alpha)|\le CS^5.
\end{equation}

\begin{lemma}[Relative secant correction]\label{F65:lem:secant-relative}
For every sufficiently fine positive fan,
\begin{equation}\label{F65:eq:secant-relative}
 |R_{\rm sec}(\alpha)-R_{\rm sec}(\bar\alpha\mathbf1)|
 \le C(S^{1/2}+S)\Jfour(\alpha).
\end{equation}
\end{lemma}

\begin{proof}
Work in a real quasiuniform material segment
$\gamma(t)=\gamma+t\delta$ and put
$h=\|\delta\|_{\ell^\infty}/a$.  Affine normalization of
\eqref{F65:eq:secant} and two differentiations of its elementary trigonometric
coefficients give, for $r=0,1,2$,
\begin{equation}
 \|\partial_t^rG_{\gamma(t)}\|_{H^{1/2}}
 \le Ch^ra^{3/2},
 \qquad
 \|\partial_t^rE_{\gamma(t)}\|_{H^{1/2}}
 \le Ch^ra^{7/2}.
\end{equation}
The real-material $q_D$ form and the exact trace/area graph have first and
second derivatives bounded by $Ch^r$ in their natural energy metrics; for
the Bessel part this is the order-minus-one kernel estimate used in
Proposition~\ref{F65:prop:bessel-relative}.  Differentiating the two coercive
Schur equations defining $P_D^{\rm tan}$ and $P_D^G$ therefore shows that
every second-derivative term in their difference is bounded by
\begin{equation}
 Ch^2(a^{3/2})(a^{7/2})=Ch^2a^5.
\end{equation}
Thus
\begin{equation}
 |D^2R_{\rm sec}(\gamma)[\delta,\delta]|
 \le Ca\sum_i\gamma_i^2\delta_i^2.
\end{equation}
The two-regime argument used in Proposition~\ref{F65:prop:bessel-relative},
with threshold $\Jfour=S^{9/2}$, proves \eqref{F65:eq:secant-relative}.
\end{proof}

\begin{theorem}[Exact tangential-polygon disk endpoint]\label{F65:thm:exact-disk}
There is $S_*>0$ such that, whenever $S\le S_*$,
\begin{equation}\label{F65:eq:exact-disk-gap}
 \boxed{
 P_D^{\rm tan}(\alpha)-P_D^{\rm tan}(\bar\alpha\mathbf1)
 \ge\frac{\zeta(3)}{80\pi^4}\Jfour(\alpha).}
\end{equation}
In particular, equality
$P_D^{\rm tan}(\alpha)=P_D^{\rm tan}(\bar\alpha\mathbf1)$ forces
$\alpha_i=\bar\alpha$ for every $i$.
\end{theorem}

\begin{proof}
Combine \eqref{F65:eq:bessel-gap} and \eqref{F65:eq:secant-relative}, decreasing
$S_*$ until the latter coefficient is at most
$\zeta(3)/(80\pi^4)$.  Strict convexity of $x^4$ gives
$\Jfour(\alpha)=0$ only at the uniform fan.
\end{proof}
\endgroup

\begingroup
\renewcommand{\qD}{q_D}
\section{Foundation F129: the exact disk-source envelope}
\label{module:F129}
\noindent\textit{Audited source: \nolinkurl{convex_largeN_polya_szego_stage129_EF_sharp_strip_gluing_closure.tex}.}
\subsection{Exact source and the three disk rows}

Let
\begin{equation}\label{F129:eq:fan}
 \alpha_i>0,\qquad \sum_i\alpha_i=2\pi,\qquad
 a=\frac{2\pi}{N},\qquad
 \Jfour(\alpha)=\sum_i\alpha_i^4-Na^4.
\end{equation}
Choose consecutive side normals \(\theta_i\), put
\(I_i=[\theta_i,\theta_{i+1}]\), and use \(x=\theta-\theta_i\) on
\(I_i\).  After removal of a constant, the unnormalized exact secant source
is
\begin{equation}\label{F129:eq:B}
 B_\alpha(\theta_i+x)
 =F_{\alpha_i}(x)
 :=\sec\bigl(\min\{x,\alpha_i-x\}\bigr)-1,
 \qquad 0\le x\le\alpha_i.
\end{equation}
Both \(F_h\) and \(F_h'\) vanish at \(0,h\), so the pieces form one
continuous periodic source.  The exact area-normalized source is
\begin{equation}\label{F129:eq:b}
 b_\alpha=h_\alpha\PiNG B_\alpha,\qquad
 h_\alpha^2=\frac{\pi}{A_\alpha},\qquad
 A_\alpha=\sum_i\tan\frac{\alpha_i}{2}.
\end{equation}

For \(n\ge2\), set
\begin{equation}\label{F129:eq:symbols}
 k_n=j_{0,1}\frac{J_{n+1}(j_{0,1})}{J_n(j_{0,1})},
 \qquad m_n=n+1-k_n.
\end{equation}
On nongeometric modes,
\begin{equation}\label{F129:eq:split}
 \qD(f)=\qdot(f)+\qzero(f)-\kform(f),
\end{equation}
where
\begin{align}
 \qdot(f)&=\sum_{|n|\ge2}|n|\,|\widehat f_n|^2,\label{F129:eq:qdot}\\
 \qzero(f)&=\sum_{|n|\ge2}|\widehat f_n|^2,\label{F129:eq:qzero}\\
 \kform(f)&=\sum_{|n|\ge2}k_{|n|}|\widehat f_n|^2.\label{F129:eq:kform}
\end{align}
Stages 51 and 65 give, for \(0\le r\le3\),
\begin{equation}\label{F129:eq:k-symbol}
 0<k_n\le\frac Cn,\qquad
 |\Delta^rk_n|\le C_r(1+n)^{-1-r}.
\end{equation}

\subsection{The sharp broken-strip gluing lemma}

For \(0<h\le h_0<\pi\), let \(V_h\) be the Neumann harmonic extension of
\(F_h\) to the half-strip
\(\Sigma_h=(0,h)\times(0,\infty)\):
\begin{equation}\label{F129:eq:strip-problem}
 \Delta V_h=0,\qquad V_h(\cdot,0)=F_h,\qquad
 \partial_xV_h(0,\cdot)=\partial_xV_h(h,\cdot)=0.
\end{equation}
Its limit at infinity is the cell mean
\begin{equation}\label{F129:eq:mean}
 \mathfrak m(h)=\frac1h\int_0^hF_h(x)\,dx,
\end{equation}
and put
\begin{equation}\label{F129:eq:N}
 \mathcal N(h)=\int_{\Sigma_h}|\nabla V_h|^2\,dx\,dy.
\end{equation}

\begin{lemma}[Scaled strip data]\label{F129:lem:scaled}
For \(0<h\le h_0\),
\begin{equation}\label{F129:eq:scaled-bounds}
 |\mathfrak m(h)|\le Ch^2,\qquad
 0\le\mathcal N(h)\le Ch^4,\qquad
 |\mathcal N''(h)|\le Ch^2.
\end{equation}
If \(w_h(y)=V_h(0,y)=V_h(h,y)\), then
\begin{equation}\label{F129:eq:wall-scale}
 w_h(y)=h^2W_h(y/h),
\end{equation}
where \(h\mapsto W_h\) is bounded with two derivatives in
\(H^{1/2}(\mathbb R_+)\cap L^2(\mathbb R_+)\) after its constant limit is
removed.
\end{lemma}

\begin{proof}
Write
\begin{equation}\label{F129:eq:profile}
 F_h(hs)=h^2\Phi_h(s),\qquad
 \Phi_h(s)=\frac{\sec(h\min\{s,1-s\})-1}{h^2}.
\end{equation}
The family \(\Phi_h\), including its analytic value at \(h=0\), is bounded
with two \(h\)-derivatives and four one-sided \(s\)-derivatives.  Its value
and first derivative vanish at \(s=0,1\).  Expand it in the Neumann cosine
basis:
\begin{equation}
 \Phi_h(s)=\overline\Phi_h+\sum_{\ell\ge1}c_\ell(h)\cos(\ell\pi s).
\end{equation}
The midpoint derivative jump gives
\(|\partial_h^r c_\ell(h)|\le C_r(1+\ell)^{-2}\), \(0\le r\le2\).
Consequently
\begin{align}
 V_h(hs,hy)
 &=h^2\left\{\overline\Phi_h+
   \sum_{\ell\ge1}c_\ell(h)
   \cos(\ell\pi s)e^{-\ell\pi y}\right\},\label{F129:eq:scaled-V}\\
 \mathcal N(h)
 &=\frac{\pi h^4}{2}\sum_{\ell\ge1}\ell\,|c_\ell(h)|^2.\label{F129:eq:N-series}
\end{align}
Equations \eqref{F129:eq:scaled-bounds} and \eqref{F129:eq:wall-scale} now follow by
termwise differentiation; the series are dominated by
\(\sum\ell^{-3}\).
\end{proof}

For \(H>0\), use the wall-trace norm
\begin{equation}\label{F129:eq:XH}
 \|g\|_{X_H}^2
 :=\|g\|_{\dot H^{1/2}(\mathbb R_+)}^2
   +H^{-1}\|g\|_{L^2(\mathbb R_+)}^2.
\end{equation}
Functions below vanish at \(y=0\) and at infinity, so the homogeneous
half-line norm has no constant ambiguity.

\begin{lemma}[Wall equalization and lifting]\label{F129:lem:wall}
Fix \(a\le h_0\).  There are modified strip fields
\(\widetilde V_h\) with the same bottom value \(F_h\), common limit
\(\mathfrak m(a)\) at infinity, and
\begin{equation}\label{F129:eq:modified-energy}
 \int_{\Sigma_h}|\nabla\widetilde V_h|^2
 \le\mathcal N(h)+C\{\mathfrak m(h)-\mathfrak m(a)\}^2.
\end{equation}
For two adjacent widths \(h,k\), \(H=\max\{h,k\}\), their wall jump
\(j_{h,k}\) satisfies
\begin{equation}\label{F129:eq:jump}
 \|j_{h,k}\|_{X_H}^2
 \le C\{(h^2-a^2)^2+(k^2-a^2)^2\}.
\end{equation}
Moreover the jump can be lifted into the wider of the two adjacent strips,
with zero bottom value, zero limit at infinity, and Dirichlet energy at
most \(C\|j_{h,k}\|_{X_H}^2\).
\end{lemma}

\begin{proof}
Choose a fixed smooth \(\chi\) on \([0,\infty)\) with
\(\chi(0)=0\) and \(\chi=1\) on \([1,\infty)\), and set
\begin{equation}\label{F129:eq:tail-correction}
 C_h(x,y)=-\{\mathfrak m(h)-\mathfrak m(a)\}\chi(y/h),
 \qquad \widetilde V_h=V_h+C_h.
\end{equation}
The correction is independent of \(x\).  Since the horizontal mean of
\(\partial_yV_h\) is zero, its Dirichlet cross term with \(C_h\) vanishes.
Direct scaling gives
\begin{equation}
 \int_{\Sigma_h}|\nabla C_h|^2
 =|\mathfrak m(h)-\mathfrak m(a)|^2
   \int_0^\infty|\chi'(s)|^2\,ds,
\end{equation}
which proves \eqref{F129:eq:modified-energy}.

The modified wall trace is zero at \(y=0\) and tends to the common value
\(\mathfrak m(a)\).  Subtract that common value.  Equations
\eqref{F129:eq:wall-scale} and \eqref{F129:eq:tail-correction}, followed by scaling
in \eqref{F129:eq:XH}, prove \eqref{F129:eq:jump}.  For completeness, first suppose
\(h,k\in[a/2,2a]\).  The fundamental theorem of calculus in the scale
parameter gives
\[
 \|j_{h,k}\|_{X_H}^2\le C(h-k)^2(h+k)^2.
\]
If at least one width lies outside \([a/2,2a]\), the triangle inequality
and scaling of the two tail-equalized traces give
\[
 \|j_{h,k}\|_{X_H}^2
 \le C(h^4+k^4+a^4)
 \le C\{(h^2-a^2)^2+(k^2-a^2)^2\}.
\]
When one width remains in \([a/2,2a]\), insert the nearer endpoint
\(a/2\) or \(2a\) as an intermediate scale, use the first estimate on the
short piece, and the coarse estimate on the other.  This proves
\eqref{F129:eq:jump} in every ratio.  The \(a^4\) term is essential when both
\(h\) and \(k\) are much smaller than \(a\): it pays for changing their
common strip tail to the regular-cell mean.

Finally, the Poisson extension in a quadrant is a right inverse of the
vertical trace with energy
\(\|j\|_{\dot H^{1/2}(\mathbb R_+)}^2\).  Multiply it by a horizontal cutoff
supported before distance \(H/3\).  The cutoff row costs
\(CH^{-1}\|j\|_2^2\), proving the \(X_H\) estimate.  Put the lifting in the
wider adjacent strip.  The two possible liftings assigned to one strip
occupy its left and right thirds, so their overlap is uniformly bounded.
\end{proof}

\begin{proposition}[Sharp no-ratio strip gluing]\label{F129:prop:gluing}
For every sufficiently fine positive fan,
\begin{equation}\label{F129:eq:gluing}
 \sum_{n\ne0}|n|\,|\widehat B_{\alpha,n}|^2
 \le\frac1{2\pi}\sum_i\mathcal N(\alpha_i)
   +C\sum_i(\alpha_i^2-a^2)^2.
\end{equation}
At the regular fan equality holds without the last term:
\begin{equation}\label{F129:eq:uniform-strip}
 \sum_{n\ne0}|n|\,|\widehat B_{a\mathbf1,n}|^2
 =\frac N{2\pi}\mathcal N(a).
\end{equation}
\end{proposition}

\begin{proof}
Place \(\widetilde V_{\alpha_i}\) over \(I_i\).  Lemma~\ref{F129:lem:wall}
gives wall jumps which vanish at the bottom and at infinity.  Lift all
jumps as in that lemma, assigning each lifting to the wider adjacent strip.
The resulting field \(U\) belongs to
\(H^1(\Sone\times\mathbb R_+)\), has bottom trace \(B_\alpha\), and tends
to the single constant \(\mathfrak m(a)\).

The original \(V_{\alpha_i}\) is harmonic and has zero normal derivative
on its vertical walls.  Hence it is Dirichlet-orthogonal to every jump
lifting, which vanishes on the bottom and at infinity.  It is also
orthogonal to the constant-in-\(x\) tail correction.  The only remaining
cross term is between tail corrections and jump liftings; Cauchy--Schwarz
absorbs it into their two energies.  Lemma~\ref{F129:lem:wall} therefore gives
\begin{equation}
 \int|\nabla U|^2
 \le\sum_i\mathcal N(\alpha_i)
 +C\sum_i(\alpha_i^2-a^2)^2.
\end{equation}
The harmonic-extension characterization of the periodic
\(\dot H^{1/2}\) norm proves \eqref{F129:eq:gluing}.

For \(\alpha_i=a\), all strip fields and their wall traces agree.  No tail
correction or lifting is present, and reflection of \(V_a\) across every
vertical wall is exactly the periodic harmonic extension.  This proves
\eqref{F129:eq:uniform-strip}.
\end{proof}

\subsection{The order-one and identity source envelopes}

Recall the exact identity
\begin{equation}\label{F129:eq:Jidentity}
 \Jfour(\alpha)
 =\sum_i(\alpha_i-a)^2
   (\alpha_i^2+2a\alpha_i+3a^2).
\end{equation}
In particular,
\begin{equation}\label{F129:eq:square-charge}
 \sum_i(\alpha_i^2-a^2)^2
 \le\Jfour(\alpha).
\end{equation}

\begin{theorem}[Sharp order-one source envelope]\label{F129:thm:qdot}
For every sufficiently fine positive fan,
\begin{equation}\label{F129:eq:qdot-envelope}
 \qdot(\PiNG B_\alpha)-\qdot(\PiNG B_{a\mathbf1})
 \le C\Jfour(\alpha).
\end{equation}
\end{theorem}

\begin{proof}
Removing the modes \(\pm1\) only decreases the left energy at \(\alpha\);
they vanish at the regular fan by cyclic symmetry.  Apply
Proposition~\ref{F129:prop:gluing} and subtract \eqref{F129:eq:uniform-strip}.
By Lemma~\ref{F129:lem:scaled} and Taylor's formula,
\begin{align}
 \sum_i\{\mathcal N(\alpha_i)-\mathcal N(a)\}
 &=\sum_i\{\mathcal N(\alpha_i)-\mathcal N(a)
           -\mathcal N'(a)(\alpha_i-a)\}\notag\\
 &\le C\sum_i(\alpha_i-a)^2(\alpha_i+a)^2
 \le C\Jfour(\alpha).\label{F129:eq:N-Bregman}
\end{align}
Use \eqref{F129:eq:square-charge} for the gluing row.
\end{proof}

\begin{lemma}[Identity source envelope]\label{F129:lem:qzero}
For every sufficiently fine positive fan,
\begin{equation}\label{F129:eq:qzero-envelope}
 \qzero(\PiNG B_\alpha)-\qzero(\PiNG B_{a\mathbf1})
 \le C\Jfour(\alpha).
\end{equation}
\end{lemma}

\begin{proof}
Put
\begin{equation}
 R(h)=\int_0^hF_h(x)^2\,dx,\qquad
 M(h)=\int_0^hF_h(x)\,dx
     =2\int_0^{h/2}(\sec x-1)\,dx.
\end{equation}
Scaling gives \(|R''(h)|\le Ch^3\).  The linear terms cancel under
\(\sum_i(\alpha_i-a)=0\), so
\begin{equation}
 \sum_i\{R(\alpha_i)-R(a)\}
 \le C h_0\sum_i(\alpha_i-a)^2(\alpha_i+a)^2
 \le C\Jfour(\alpha).
\end{equation}
Moreover
\begin{equation}
 M''(h)=\frac12\sec\frac h2\tan\frac h2>0.
\end{equation}
Thus Jensen's inequality gives
\(\sum_iM(\alpha_i)\ge NM(a)>0\); subtracting the squared global mean can
only improve the relative upper bound.  Finally, the regular source has no
modes \(\pm1\), while deleting those modes at \(\alpha\) again decreases
the energy.  Parseval proves \eqref{F129:eq:qzero-envelope}.
\end{proof}

\subsection{The smoother Bessel row}

\begin{lemma}[Regular Bessel baseline]\label{F129:lem:kbaseline}
For every sufficiently small \(a\),
\begin{equation}\label{F129:eq:kbaseline}
 0\le\kform(\PiNG B_{a\mathbf1})\le Ca^5.
\end{equation}
\end{lemma}

\begin{proof}
Cyclic symmetry restricts the Fourier support to \(N\mathbb Z\).
The scaled cosine expansion in Lemma~\ref{F129:lem:scaled} gives
\(|\widehat B_{\ell N}|\le Ca^2(1+|\ell|)^{-2}\).  Since
\(k_{\ell N}\le C/(|\ell|N)\le Ca/|\ell|\), summation proves
\eqref{F129:eq:kbaseline}.
\end{proof}

\begin{lemma}[Quasiuniform Bessel material Hessian]\label{F129:lem:kHessian}
Suppose
\begin{equation}
 \frac a2\le\gamma_i\le\frac{3a}{2},\qquad
 \sum_i\gamma_i=2\pi,\qquad \sum_i d_i=0.
\end{equation}
Then
\begin{equation}\label{F129:eq:kHessian}
 \left|D^2\{\kform(\PiNG B_\gamma)\}[d,d]\right|
 \le Ca\sum_i\gamma_i^2d_i^2.
\end{equation}
\end{lemma}

\begin{proof}
Pull every normal interval affinely to the regular material interval.
The profile \eqref{F129:eq:profile} and its first two material derivatives are
uniformly bounded after the factors \(\gamma_i^2\), \(\gamma_id_i\), and
\(d_i^2\) are removed.  The multiplier bounds \eqref{F129:eq:k-symbol} allow
three cyclic Abel summations in the pulled-back kernel.  The first removes
the arbitrary root velocity, because consecutive midpoint velocities
differ by \((d_i+d_{i+1})/2\); the next two give a summable
\((1+|r|)^{-2}\) interaction between material cells at distance \(r\).
Cauchy--Schwarz after the full cyclic sum yields
\begin{equation}
 |D^2\kform|
 \le Ca\sum_i\gamma_i^2d_i^2.
\end{equation}
The factor \(a\) is the order-minus-one gain.  This is the source-only
specialization of the real-material Bessel Hessian estimate proved in
stage 65; no Galerkin or support-space term is present here.
\end{proof}

\begin{proposition}[Relative Bessel source row]\label{F129:prop:krelative}
For every sufficiently fine positive fan,
\begin{equation}\label{F129:eq:krelative}
 -\{\kform(\PiNG B_\alpha)-\kform(\PiNG B_{a\mathbf1})\}
 \le C\Jfour(\alpha).
\end{equation}
\end{proposition}

\begin{proof}
If \(\Jfour(\alpha)\ge a^5\), positivity and
Lemma~\ref{F129:lem:kbaseline} give
\begin{equation}
 -\{\kform(B_\alpha)-\kform(B_{a\mathbf1})\}
 \le\kform(B_{a\mathbf1})\le C\Jfour(\alpha).
\end{equation}
If \(\Jfour(\alpha)<a^5\), \eqref{F129:eq:Jidentity} implies
\(\max_i|\alpha_i-a|^4\le\Jfour<a^5\).  For small \(a\), the whole affine
segment \(\gamma_i(t)=a+t(\alpha_i-a)\) lies in
\([a/2,3a/2]\).  Cyclic invariance kills the first derivative at
\(a\mathbf1\).  Lemma~\ref{F129:lem:kHessian}, Taylor's formula, and
\begin{equation}
 \Jfour(\alpha)
 =12\int_0^1(1-t)\sum_i\gamma_i(t)^2(\alpha_i-a)^2\,dt
\end{equation}
prove the stronger bound \(Ca\Jfour\).
\end{proof}

\subsection{Exact source envelope and endpoint forcing}

\begin{theorem}[Unnormalized exact disk-source envelope]\label{F129:thm:B}
For every sufficiently fine positive fan,
\begin{equation}\label{F129:eq:B-envelope}
 \qD(\PiNG B_\alpha)-\qD(\PiNG B_{a\mathbf1})
 \le C\Jfour(\alpha).
\end{equation}
\end{theorem}

\begin{proof}
Add Theorem~\ref{F129:thm:qdot}, Lemma~\ref{F129:lem:qzero}, and
Proposition~\ref{F129:prop:krelative} according to the exact multiplier split
\eqref{F129:eq:split}.
\end{proof}

\begin{theorem}[Area-normalized exact source envelope]\label{F129:thm:b}
For every sufficiently fine positive fan,
\begin{equation}\label{F129:eq:b-envelope}
 \boxed{\qD(b_\alpha)-\qD(b_{a\mathbf1})
 \le C\Jfour(\alpha).}
\end{equation}
\end{theorem}

\begin{proof}
Convexity of \(\tan(x/2)\) and \(\sum_i\alpha_i=Na\) give
\(A_\alpha\ge A_{a\mathbf1}\), hence
\(h_\alpha^2\le h_{a\mathbf1}^2\).  From \eqref{F129:eq:b},
\begin{align}
 \qD(b_\alpha)-\qD(b_{a\mathbf1})
 &=h_\alpha^2\{\qD(B_\alpha)-\qD(B_{a\mathbf1})\}\notag\\
 &\quad +(h_\alpha^2-h_{a\mathbf1}^2)\qD(B_{a\mathbf1}).\label{F129:eq:area-sign}
\end{align}
The second term is nonpositive.  The factors \(h_\alpha\) are uniformly
bounded for fine fans, so Theorem~\ref{F129:thm:B} proves the claim.
\end{proof}
\endgroup

\begingroup
\section{Node TR: physical trace and relative Schur transfer}
\label{module:TR}
\noindent\textit{Audited source: \nolinkurl{TR_physical_trace.tex}.}
\subsection{The true straight-edge material trace}

Let $\alpha_i>0$, $\sum_i\alpha_i=2\pi$, and
$\theta_{i+1}-\theta_i=\alpha_i$.  Put
\begin{equation}\label{TR:eq:fan}
 a=\alpha_{i-1},\qquad b=\alpha_i,\qquad
 K_i=[\theta_i-a/2,\theta_i+b/2],\qquad x=\theta-\theta_i.
\end{equation}
The area-normalized tangential polygon has common support
\begin{equation}\label{TR:eq:h}
 h_\alpha=\left(\frac\pi{\sum_j\tan(\alpha_j/2)}\right)^{1/2}
\end{equation}
and radial boundary $h_\alpha\sec x$ on $K_i$.

For support increments $z=(z_i)$, the two endpoint displacement vectors
of side $i$ have components in the frame $(\nu_i,\tau_i)$
\begin{equation}\label{TR:eq:endpoint-components}
 (z_i,a_i^-),\qquad (z_i,a_i^+),
\end{equation}
where
\begin{equation}\label{TR:eq:apm}
 a_i^-=\frac{z_i-z_{i-1}}{\sin a}-z_i\tan\frac a2,
 \qquad
 a_i^+=\frac{z_{i+1}-z_i}{\sin b}+z_i\tan\frac b2.
\end{equation}
Indeed these are obtained by differentiating the two adjacent support-line
intersection equations.

The physical affine coordinate of the radial point
$h_\alpha\nu_i+h_\alpha\tan x\,\tau_i$ is
\begin{equation}\label{TR:eq:t}
 t_i(x)=\frac{\tan x+\tan(a/2)}
 {\tan(a/2)+\tan(b/2)}.
\end{equation}
Hence the genuine radial component of the straight-edge affine material
velocity is
\begin{equation}\label{TR:eq:Tphysical}
 (\widetilde{T}_\alpha z)_i(x)
 =z_i\cos x+
 \bigl((1-t_i(x))a_i^-+t_i(x)a_i^+\bigr)\sin x.
\end{equation}
For comparison, the stage-65 angular trace is
\begin{equation}\label{TR:eq:Tangular}
 (T_\alpha z)_i(x)
 =z_i\cos x+
 \bigl((1-s_i(x))a_i^-+s_i(x)a_i^+\bigr)\sin x,
 \qquad
 s_i(x)=\frac{x+a/2}{(a+b)/2}.
\end{equation}

We use the Fourier norm
\begin{equation}\label{TR:eq:Hhalf}
 \|f\|_{H^{1/2}}^2
 =\sum_{n\in\mathbb Z}(1+|n|)|\widehat f(n)|^2.
\end{equation}
All quotient estimates below are unchanged if the equivalent disk form
$q_{\mathbb D}$ is used.

\begin{lemma}[Endpoint-uniform interpolation jets]
\label{TR:lem:weight-jets}
Put $w=a+b$, $p=a/w$, $y=s_i(x)$, and assume $0<w\le w_0<\pi$.
Then, uniformly for $0\le p\le1$,
\begin{align}
 |t_i-s_i|&\le Cw^2s_i(1-s_i),
 \label{TR:eq:weight-value}\\
 w|t_i'-s_i'|&\le Cw^2.
 \label{TR:eq:weight-derivative}
\end{align}
In particular, no endpoint factor occurs in
\eqref{TR:eq:weight-derivative}.
\end{lemma}

\begin{proof}
In normalized variables,
\begin{equation}\label{TR:eq:normalized-t}
 \vartheta_{w,p}(y)=
 \frac{\tan(\frac w2(y-p))+\tan(\frac{wp}2)}
 {\tan(\frac{wp}2)+\tan(\frac{w(1-p)}2)},
 \qquad t_i(x)=\vartheta_{w,p}(y).
\end{equation}
The denominator is between $w/2$ and $C_{w_0}w$.  Direct
differentiation gives
\begin{equation}\label{TR:eq:normalized-second}
 \partial_y^2\vartheta_{w,p}(y)
 =\frac{w^2}{2}
 \frac{\sec^2(\frac w2(y-p))\tan(\frac w2(y-p))}
 {\tan(\frac{wp}2)+\tan(\frac{w(1-p)}2)},
\end{equation}
and hence $|\partial_y^2\vartheta_{w,p}|\le Cw^2$.
The function $f=\vartheta_{w,p}-y$ vanishes at $0$ and $1$.
The Green function of $-\partial_y^2$ on $[0,1]$ therefore gives
$|f(y)|\le Cw^2y(1-y)$.  Moreover
$\int_0^1f'(y)\,dy=0$, so
$|f'(y)|\le\|f''\|_\infty\le Cw^2$.  Since
$\partial_x=(2/w)\partial_y$, these are exactly
\eqref{TR:eq:weight-value}--\eqref{TR:eq:weight-derivative}.
\end{proof}

For the next lemma, write $w_i=\alpha_{i-1}+\alpha_i$ and
$d_i=a_i^+-a_i^-$.  The length of $K_i$ is $w_i/2$.

\begin{lemma}[Stable extraction of the internal quadratic coefficient]
\label{TR:lem:coefficient-stability}
For $S$ sufficiently small,
\begin{equation}\label{TR:eq:local-coefficient}
 w_i^2|d_i|^2
 \le C\left(
 [T_\alpha z]_{H^{1/2}(K_i)}^2+w_i^4|z_i|^2\right).
\end{equation}
On the exact area kernel, modulo exact support translations,
\begin{equation}\label{TR:eq:nodal-sampling}
 \sum_iw_i|z_i|^2
 \le C\|\PiNG T_\alpha z\|_{H^{1/2}(\T)}^2.
\end{equation}
Both constants are independent of adjacent ratios.
\end{lemma}

\begin{proof}
Pull $K_i$ to $0\le y\le1$ and put
\begin{equation*}
 Z=z_i,\qquad A=w_i a_i^-,\qquad B=w_i a_i^+,
 \qquad \xi=\frac{w_i}{2}(y-p).
\end{equation*}
The normalized trace is
\begin{equation}\label{TR:eq:normalized-trace}
 F_{w_i,p}(y)=Z\cos\xi+\{(1-y)A+yB\}\frac{\sin\xi}{w_i}.
\end{equation}
At $w_i=0$ its continuous limit is
\begin{equation}\label{TR:eq:homogeneous-normalized-trace}
 F_{0,p}(y)=Z+\frac12(y-p)\{(1-y)A+yB\}.
\end{equation}
Modulo constants, the last two normalized polynomials have coefficient
matrix
\begin{equation*}
 \begin{pmatrix}1+p&-p\\-1&1\end{pmatrix},
\end{equation*}
whose determinant is one and whose inverse is uniformly bounded for
$0\le p\le1$.  Finite-dimensional norm equivalence therefore gives
\begin{equation}\label{TR:eq:homogeneous-coercivity}
 |A|^2+|B|^2\le C[F_{0,p}]_{H^{1/2}(0,1)}^2.
\end{equation}
Taylor's formula in \eqref{TR:eq:normalized-trace}, in the same
finite-dimensional space, gives
\begin{equation}\label{TR:eq:normalized-perturbation}
 [F_{w_i,p}-F_{0,p}]_{H^{1/2}(0,1)}
 \le Cw_i^2(|Z|+|A|+|B|).
\end{equation}
After decreasing the mesh threshold, the $A,B$ terms are absorbed into
\eqref{TR:eq:homogeneous-coercivity}.  Since
$B-A=w_i d_i$ and the $H^{1/2}$ seminorm is invariant under affine
rescaling in one dimension, this proves \eqref{TR:eq:local-coefficient}.

We also need the constant and translation gauges.  The full normalized
$L^2$ three-jet matrix associated with \eqref{TR:eq:normalized-trace} is
uniformly invertible.  Hence, for every constant $c$,
\begin{equation}\label{TR:eq:local-sampling}
 \sum_iw_i|z_i-c|^2
 \le C\|T_\alpha z-c\|_{L^2(\T)}^2.
\end{equation}
Every first harmonic is exactly the angular trace of a support
translation $z_i^{\rm tr}=q\cdot\nu_i$, and $L_\alpha z^{\rm tr}=0$.
Subtract the unique translation producing the $n=\pm1$ part of
$T_\alpha z$.  For the resulting representative,
\begin{equation*}
 T_\alpha z=g+c,\qquad g=\PiNG T_\alpha z.
\end{equation*}
The exact area row is
\begin{equation}\label{TR:eq:exact-area-row}
 L_\alpha z=\sum_i\ell_i z_i=0,\qquad
 \ell_i=\tan\frac{\alpha_{i-1}}2+
        \tan\frac{\alpha_i}2\asymp w_i.
\end{equation}
By Cauchy--Schwarz and \eqref{TR:eq:local-sampling},
\begin{equation*}
 |c|L_\alpha\mathbf1=|L_\alpha(z-c\mathbf1)|
 \le C\Big(\sum_iw_i|z_i-c|^2\Big)^{1/2},
\end{equation*}
where $L_\alpha\mathbf1\asymp1$.  Thus
$|c|\le C\|g\|_{L^2}$.  Applying
\eqref{TR:eq:local-sampling} once more proves
\eqref{TR:eq:nodal-sampling}.  Exact translation subtraction does not
change $d_i$, because $a_i^-=a_i^+=q\cdot\tau_i$ for a translation.
\end{proof}

\begin{lemma}[Stable zero-wall bubble decomposition]
\label{TR:lem:zero-wall}
Let $E$ equal $E_i$ on $K_i$, where
\begin{equation}\label{TR:eq:trace-error-local}
 E_i=(t_i-s_i)d_i\sin x.
\end{equation}
Then
\begin{equation}\label{TR:eq:zero-wall-bound}
 \|E\|_{H^{1/2}(\T)}^2\le C\sum_iw_i^6|d_i|^2.
\end{equation}
\end{lemma}

\begin{proof}
Lemma~\ref{TR:lem:weight-jets} gives, in normalized coordinates,
\begin{equation}\label{TR:eq:E-normalized-jets}
 |E_i(y)|\le Cw_i^3|d_i|y(1-y),
 \qquad
 |\partial_yE_i(y)|\le Cw_i^3|d_i|.
\end{equation}
In particular, $E_i$ vanishes at both endpoints.  The same-cell
Gagliardo integral is at most $Cw_i^6|d_i|^2$.  For the cross-cell
integral use
\begin{equation*}
 |E_i(\theta)-E_j(\varphi)|^2
 \le2|E_i(\theta)|^2+2|E_j(\varphi)|^2.
\end{equation*}
For $\theta\in K_i$, integration of the circular kernel over
$\T\setminus K_i$ is bounded by
\begin{equation*}
 C\left(1+\operatorname {dist}(\theta,\partial K_i)^{-1}\right).
\end{equation*}
The endpoint factor in \eqref{TR:eq:E-normalized-jets} makes the resulting
Hardy integral at most $Cw_i^6|d_i|^2$.  Summing over $i$, and adding the
$L^2$ row, proves \eqref{TR:eq:zero-wall-bound}.  This argument sums the
whole complement of each cell at once, so it has no mesh-depth or
adjacent-ratio constant.
\end{proof}

\begin{lemma}[Physical versus angular interpolation]
\label{TR:lem:trace-close}
Let $S=\max_i\alpha_i+2\pi/N$.  On the exact area and translation
quotient,
\begin{equation}\label{TR:eq:trace-close}
 \|\PiNG(\widetilde T_\alpha-T_\alpha)z\|_{H^{1/2}(\T)}
 \le CS^2\|\PiNG T_\alpha z\|_{H^{1/2}(\T)}.
\end{equation}
The constant is independent of every adjacent-cell ratio.
\end{lemma}

\begin{proof}
Projection is contractive in \eqref{TR:eq:Hhalf}.  By
Lemmas~\ref{TR:lem:coefficient-stability} and \ref{TR:lem:zero-wall},
\begin{align*}
 \|\PiNG E\|_{H^{1/2}}^2
 &\le C\sum_iw_i^6|d_i|^2\\
 &\le CS^4\sum_i[T_\alpha z]_{H^{1/2}(K_i)}^2
   +C\sum_iw_i^8|z_i|^2\\
 &\le CS^4\|\PiNG T_\alpha z\|_{H^{1/2}}^2
   +CS^7\sum_iw_i|z_i|^2\\
 &\le CS^4\|\PiNG T_\alpha z\|_{H^{1/2}}^2.
\end{align*}
Here we use the representative constructed in the proof of
Lemma~\ref{TR:lem:coefficient-stability}, for which
$T_\alpha z=\PiNG T_\alpha z+c$; consequently its local seminorms are
those of $\PiNG T_\alpha z$.  Also $w_i\le2S$, and the same-cell
Gagliardo integrals are positive parts of the global seminorm.  Taking square roots proves
\eqref{TR:eq:trace-close}.
\end{proof}

\begin{remark}[The repaired terminology]
$T_\alpha z$ remains a useful trigonometric spline and the stage-65 disk
gap for it is valid.  Equation \eqref{TR:eq:Tphysical}, not
\eqref{TR:eq:Tangular}, is the trace of an actual affine map of a straight
edge.  The distinction is a small relative error, but it is not an exact
identity and may not be suppressed in a value comparison.
\end{remark}

\subsection{Transfer of D0 to the physical trace}

We first fix a normalization which was ambiguous in the historical
notation.  The reduced Fourier--Bessel form of F51 is denoted by
$\qDzero$; the physical half-Hessian form is
\begin{equation}\label{TR:eq:physical-disk-normalization}
 \boxed{\qphys(f,g):=2j_{0,1}^{\,2}\qDzero(f,g)
 =\frac12D^2\mathcal L(0)[fe_r,ge_r].}
\end{equation}
Every disk endpoint, metric, pairing, and distance below uses $\qphys$.
Thus the angular endpoint imported from F65 is first multiplied by
$2j_{0,1}^{\,2}$; this changes only its positive universal gap constant,
and changes neither its minimizer nor any relative trace comparison.

Let
\begin{equation}\label{TR:eq:source}
 b_\alpha=\PiNG\{h_\alpha\sec x-1\}
\end{equation}
on the cells $K_i$.  Define the angular and physical disk values
\begin{align}
 P_D(\alpha)&=\inf_{z\in Z_\alpha}
 \qphys(b_\alpha+\PiNG T_\alpha z),\label{TR:eq:PD}\\
 \widetilde P_D(\alpha)&=\inf_{z\in Z_\alpha}
 \qphys(b_\alpha+\PiNG\widetilde T_\alpha z).
 \label{TR:eq:PDphysical}
\end{align}
Here $Z_\alpha$ includes the exact area row and is quotiented by
translations.  Let
\begin{equation}\label{TR:eq:Gphysical}
 \widetilde\G_\alpha(z)
 =\qphys(\PiNG\widetilde T_\alpha z),
\end{equation}
and let $\widetilde z_\alpha^\sharp$ be the corresponding Schur corrector.

Write $H_{\rm ng}=\PiNG H^{1/2}(\T)$ and let
$\langle\cdot,\cdot\rangle_D$ be the bilinear form associated with
$\qphys$.  The Fourier--Bessel bounds from the disk Hessian give
\begin{equation}\label{TR:eq:Bessel-coercivity}
 c\|f\|_{H^{1/2}}^2\le \qphys(f)\le
 C\|f\|_{H^{1/2}}^2,\qquad f\in H_{\rm ng}.
\end{equation}
Set
\begin{equation}\label{TR:eq:trace-spaces}
 V_\alpha=\PiNG T_\alpha Z_\alpha,\qquad
 \widetilde V_\alpha=\PiNG\widetilde T_\alpha Z_\alpha.
\end{equation}
Thus $P_D(\alpha)$ and $\widetilde P_D(\alpha)$ are the squared
$\qphys$-distances from $-b_\alpha$ to $V_\alpha$ and
$\widetilde V_\alpha$, respectively.

\begin{lemma}[Differentiated Hilbert projection calculus]
\label{TR:lem:projection-calculus}
Let $q_t$ be uniformly coercive $C^2$ Hilbert forms on a fixed Hilbert
space, let $P_t,\widetilde P_t$ be the corresponding orthogonal
projections onto two $C^2$ families of closed subspaces, and put
$\Delta P_t=\widetilde P_t-P_t$.  Suppose, for $0\le r\le2$,
\begin{align}
 \|\partial_t^rq_t\|&\le C_rh^r,
 \label{TR:eq:q-material-bound}\\
 \|\partial_t^rP_t\|+\|\partial_t^r\widetilde P_t\|
 &\le C_rh^r,
 \qquad
 \|\partial_t^r\Delta P_t\|\le C_rh^r\varepsilon,
 \label{TR:eq:P-material-bound}\\
 \|\partial_t^rb_t\|&\le C_rh^rB.
 \label{TR:eq:b-material-bound}
\end{align}
Then
\begin{equation}\label{TR:eq:projection-calculus-output}
 \left|\partial_t^2\left\{
 \operatorname {dist}_{q_t}(-b_t,\widetilde V_t)^2-
 \operatorname {dist}_{q_t}(-b_t,V_t)^2\right\}\right|
 \le Ch^2\varepsilon B^2.
\end{equation}
\end{lemma}

\begin{proof}
Orthogonal completion of the square gives
\begin{equation}\label{TR:eq:projection-difference-identity}
 \operatorname {dist}_{q_t}(-b_t,\widetilde V_t)^2-
 \operatorname {dist}_{q_t}(-b_t,V_t)^2
 =-q_t(\Delta P_tb_t,(\widetilde P_t+P_t)b_t).
\end{equation}
Differentiate this identity twice.  Every resulting term contains one of
$\Delta P_t$, $\partial_t\Delta P_t$, or
$\partial_t^2\Delta P_t$; the remaining factors are bounded by
\eqref{TR:eq:q-material-bound}--\eqref{TR:eq:b-material-bound}.  Each term is
therefore at most $Ch^2\varepsilon B^2$.
\end{proof}

\begin{lemma}[Uniform material calculus for the paired trace spaces]
\label{TR:lem:material-calculus}
Let
\begin{equation}\label{TR:eq:quasiuniform-band}
 (1-\eta_0)a\le\gamma_i\le(1+\eta_0)a,\qquad
 \sum_i\gamma_i=2\pi,\qquad 0<\eta_0<\frac14,
\end{equation}
and let $\sum_i\delta_i=0$.  If
\begin{equation}\label{TR:eq:Rtr}
 R_{\rm tr}(\alpha)=\widetilde P_D(\alpha)-P_D(\alpha),
\end{equation}
then
\begin{equation}\label{TR:eq:transfer-hessian}
 |D^2R_{\rm tr}(\gamma)[\delta,\delta]|
 \le Ca\sum_i\gamma_i^2\delta_i^2.
\end{equation}
\end{lemma}

\begin{proof}
The assertion is trivial for $\delta=0$.  Otherwise put
\begin{equation}\label{TR:eq:material-h}
 h=\frac{\|\delta\|_{\ell^\infty}}a,
 \qquad \gamma(t)=\gamma+t\delta,
 \qquad |t|<c_0h^{-1}.
\end{equation}
Let $\Phi_t$ be the piecewise affine circle map carrying the normal nodes
of $\gamma$ to those of $\gamma(t)$, and pull all functions back by
$\Phi_t$.  For every oriented cyclic arc $I$ of length at most $\pi$,
\begin{equation}\label{TR:eq:arc-material}
 \left|t\sum_{i\in I}\delta_i\right|
 \le c_0a\# I\le Cc_0\sum_{i\in I}\gamma_i.
\end{equation}
Consequently $\Phi_t'$ and every chord ratio are a fixed small relative
perturbation of one.  Differentiating the pulled Gagliardo kernel
\begin{equation}\label{TR:eq:pulled-Gagliardo}
 \frac{\Phi_t'(\xi)\Phi_t'(\eta)}
 {\sin^2((\Phi_t(\xi)-\Phi_t(\eta))/2)}
\end{equation}
zero, one, or two times proves
\begin{equation}\label{TR:eq:form-material}
 \|\partial_t^rq_t\|\le C_rh^r,\qquad 0\le r\le2.
\end{equation}
Here the order-one part is exactly \eqref{TR:eq:pulled-Gagliardo}; the
$L^2$ row is elementary.  For the Bessel remainder use its established
order-minus-one coefficients
\begin{equation}\label{TR:eq:Bessel-differences}
 |\Delta^rk_n|\le C_r(1+|n|)^{-1-r},\qquad 0\le r\le3,
\end{equation}
and discrete summation by parts.  It satisfies the same estimate, with
one extra order of decay.

The normalized secant formula and two material differentiations give
\begin{equation}\label{TR:eq:source-material}
 \|\partial_t^rb_{\gamma(t)}\|_{H^{1/2}}
 \le C_rh^ra^{3/2},\qquad 0\le r\le2.
\end{equation}
Indeed a cell source has amplitude $O(a^2)$ and scale-invariant
$H^{1/2}$ charge $O(a^4)$; summing $O(a^{-1})$ cells gives $O(a^3)$.
Every normalized material derivative costs at most $h$.  The global
factor $h_{\gamma(t)}$ obeys the same bounds after differentiating its
explicit tangent sum.

We next fix the coefficient space.  Let
\[
 \tau_t^1=(\cos\theta_i(t))_i,\qquad
 \tau_t^2=(\sin\theta_i(t))_i
\]
be the two exact support-translation vectors.  They satisfy
$L_{\gamma(t)}\tau_t^\mu=0$.  With the area weights
$\ell_i(t)$ from \eqref{TR:eq:exact-area-row}, the $2\times2$ Gram matrix
\[
 M_t^{\mu\nu}=\sum_i\ell_i(t)\tau_{t,i}^\mu\tau_{t,i}^\nu
\]
is uniformly comparable with the identity in the quasiuniform band.
Define the explicit translation gauge
\begin{equation}\label{TR:eq:translation-gauge}
 \mathcal Q_tz=z-\sum_{\mu,\nu=1}^2
 (M_t^{-1})_{\mu\nu}
 \left(\sum_i\ell_i(t)z_i\tau_{t,i}^\nu\right)\tau_t^\mu
\end{equation}
and transport the area kernel by
\begin{equation}\label{TR:eq:area-transport}
 C_tz=\mathcal Q_t\left(
 z-\frac{L_{\gamma(t)}z}{L_{\gamma(t)}\mathbf1}\mathbf1\right).
\end{equation}
The coefficients of $L_{\gamma(t)}$ are comparable to $a$, and their
first two normalized derivatives are bounded by $C_rh^r$.  The same is
true of $\tau_t^\mu$, $M_t^{-1}$, $\mathcal Q_t$, and $C_t$.
The sampling estimate \eqref{TR:eq:nodal-sampling} makes the range of
$C_t$ uniformly transverse to the area and translation gauges.  Thus,
on the area/translation complement $X$ at $t=0$, the pulled trace injections
\begin{equation*}
 B_t=\PiNG(T_{\gamma(t)}C_t)\circ\Phi_t,
 \qquad
 \widetilde B_t=\PiNG(\widetilde T_{\gamma(t)}C_t)\circ\Phi_t
\end{equation*}
are uniformly bounded below in the pulled $q_t$ norm.

Apply the normalized calculation
\eqref{TR:eq:normalized-t}--\eqref{TR:eq:E-normalized-jets} on this
quasiuniform material band.  Its first two material derivatives give
\begin{align}
 \left|\partial_t^r\{(t_i-s_i)\sin x\}\right|
 &\le C_rh^ra^3y(1-y),\label{TR:eq:material-E-value}\\
 \left|\partial_y\partial_t^r\{(t_i-s_i)\sin x\}\right|
 &\le C_rh^ra^3,\qquad 0\le r\le2.
 \label{TR:eq:material-E-derivative}
\end{align}
The derivative row uses the bound without an endpoint factor from
\eqref{TR:eq:weight-derivative}.  The coefficient-extraction and zero-wall
proofs then apply after each material differentiation.  Together with
the derivative bounds for $C_t$, this gives
\begin{align}
 \|\partial_t^rB_t\|+\|\partial_t^r\widetilde B_t\|
 &\le C_rh^r,
 \label{TR:eq:B-material}\\
 \|\partial_t^r(\widetilde B_t-B_t)\|
 &\le C_rh^ra^2,\qquad 0\le r\le2.
 \label{TR:eq:DeltaB-material}
\end{align}
All operator norms use the fixed energy norm on $X$ and the pulled
$q_t$ norm on the target.  Notice that the endpoint factors were proved
before differentiating, so no inverse minimum cell is present.

For completeness, the $q_t$-orthogonal projection onto
$\operatorname {ran}B_t$ is
\begin{equation}\label{TR:eq:Gram-projection}
 P_t=B_tG_t^{-1}B_t^{*t},\qquad
 G_t=B_t^{*t}B_t,
\end{equation}
where $*t$ is the adjoint in $q_t$.  The corresponding formula with
$\widetilde B_t$ defines $\widetilde P_t$.  Uniform trace coercivity
bounds both Gram inverses.  Differentiating
\eqref{TR:eq:Gram-projection} and
\begin{equation*}
 \partial_tG_t^{-1}=-G_t^{-1}(\partial_tG_t)G_t^{-1}
\end{equation*}
twice, using \eqref{TR:eq:form-material}, \eqref{TR:eq:B-material}, and
\eqref{TR:eq:DeltaB-material}, yields
\begin{align}
 \|\partial_t^rP_t\|+\|\partial_t^r\widetilde P_t\|
 &\le C_rh^r,\notag\\
 \|\partial_t^r(\widetilde P_t-P_t)\|
 &\le C_rh^ra^2,\qquad 0\le r\le2.
 \label{TR:eq:projection-material}
\end{align}
Lemma~\ref{TR:lem:projection-calculus}, with
$\varepsilon=a^2$ and $B=a^{3/2}$, now gives
\begin{equation}\label{TR:eq:R-second-h}
 |D^2R_{\rm tr}(\gamma)[\delta,\delta]|
 \le Ch^2a^5=Ca^3\|\delta\|_{\ell^\infty}^2.
\end{equation}
Finally
$\sum_i\gamma_i^2\delta_i^2\ge ca^2
\|\delta\|_{\ell^\infty}^2$, which proves
\eqref{TR:eq:transfer-hessian}.
\end{proof}

\begin{lemma}[Relative physical-trace Schur transfer]
\label{TR:lem:physical-transfer}
There is a modulus $\omega_{\rm tr}(S)\to0$ such that
\begin{equation}\label{TR:eq:physical-transfer}
 |\{\widetilde P_D(\alpha)-P_D(\alpha)\}
  -\{\widetilde P_D(\alpha_*)-P_D(\alpha_*)\}|
 \le\omega_{\rm tr}(S)\Jfour(\alpha).
\end{equation}
Moreover $\widetilde\G_\alpha$ and $\G_\alpha$ are uniformly equivalent,
with relative error $O(S^2)$.
\end{lemma}

\begin{proof}
The equivalence of the support metrics follows from
Lemma~\ref{TR:lem:trace-close}, the Bessel coercivity
$\qphys\asymp H^{1/2}$, and polarization; more precisely,
\begin{equation}\label{TR:eq:metric-equivalence}
 |\widetilde\G_\alpha(z)-\G_\alpha(z)|
 \le CS^2\G_\alpha(z).
\end{equation}
Lemma~\ref{TR:lem:trace-close} also gives a graph bijection
$J_\alpha:V_\alpha\to\widetilde V_\alpha$ satisfying
$\|J_\alpha-I\|\le CS^2$ in the $\qphys$ norm.  The elementary Hilbert
distance comparison therefore gives
\begin{equation}\label{TR:eq:transfer-absolute}
 |R_{\rm tr}(\alpha)|
 \le CS^2\qphys(b_\alpha)\le CS^5.
\end{equation}
The final source estimate follows by normalized same/adjacent/separated
Gagliardo summation for the continuous secant source:
$\qphys(b_\alpha)\le C\sum_i\alpha_i^4\le CS^3$.

Put $a=2\pi/N$, $d_i=\alpha_i-a$, and recall the exact identity
\begin{equation}\label{TR:eq:Jfour-identity}
 \Jfour(\alpha)=
 \sum_i d_i^2(\alpha_i^2+2a\alpha_i+3a^2).
\end{equation}
If $\Jfour(\alpha)\ge S^{9/2}$, apply
\eqref{TR:eq:transfer-absolute} at $\alpha$ and at
$\alpha_*=a\mathbf1$ to obtain
\begin{equation}\label{TR:eq:far-transfer}
 |R_{\rm tr}(\alpha)-R_{\rm tr}(\alpha_*)|
 \le CS^{1/2}\Jfour(\alpha).
\end{equation}

Suppose instead that $\Jfour(\alpha)<S^{9/2}$.  Then
$\|d\|_{\ell^\infty}^4\le C\Jfour(\alpha)$, so
$\|d\|_{\ell^\infty}/S\le CS^{1/8}$.  Since
$S=a+\max_i\alpha_i$, the whole segment
$\gamma(t)=a\mathbf1+td$ lies in a fixed quasiuniform band and
$a\asymp S$.  Cyclic invariance makes
$DR_{\rm tr}(a\mathbf1)[d]=0$, because the gradient at the regular fan
is constant and $\sum_i d_i=0$.  Taylor's formula and
Lemma~\ref{TR:lem:material-calculus} give
\begin{align}
 |R_{\rm tr}(\alpha)-R_{\rm tr}(\alpha_*)|
 &\le Ca\int_0^1(1-t)
       \sum_i\gamma_i(t)^2d_i^2\,dt\notag\\
 &\le CS\Jfour(\alpha),
 \label{TR:eq:near-transfer}
\end{align}
where the last inequality follows from
\eqref{TR:eq:Jfour-identity} and quasiuniformity.  Combining
\eqref{TR:eq:far-transfer} and \eqref{TR:eq:near-transfer} proves
\eqref{TR:eq:physical-transfer} with
\begin{equation}\label{TR:eq:transfer-modulus}
 \omega_{\rm tr}(S)=C(S^{1/2}+S).
\end{equation}
\end{proof}

\begin{theorem}[Physical disk endpoint D0]
\label{TR:thm:physical-D0}
For all sufficiently fine positive fans,
\begin{equation}\label{TR:eq:physical-D0}
 \qphys(b_\alpha+\PiNG\widetilde T_\alpha z)
 -\widetilde P_D(\alpha_*)
 \ge c_f\Jfour(\alpha)
   +c_s\widetilde\G_\alpha(z-\widetilde z_\alpha^\sharp),
\end{equation}
where $c_f,c_s>0$ are independent of all minimum cells.
\end{theorem}

\begin{proof}
Complete the physical support square exactly:
\begin{equation*}
 \qphys(b_\alpha+\PiNG\widetilde T_\alpha z)
 =\widetilde P_D(\alpha)
  +\widetilde\G_\alpha(z-\widetilde z_\alpha^\sharp).
\end{equation*}
The angular endpoint gap of stages 65 and 129, together with
Lemma~\ref{TR:lem:physical-transfer}, gives
$\widetilde P_D(\alpha)-\widetilde P_D(\alpha_*)\ge c_f\Jfour$ after
decreasing the fine-mesh threshold.  Uniform coercivity gives the support
constant $c_s$.
\end{proof}

\subsection{An exact fixed-disk scalar for the polygon}

Before the harmless area dilation, let $p_i(h)$ be the intersection of
the support lines with normals $\nu_i,\nu_{i+1}$.  It depends linearly on
$(h_i,h_{i+1})$.  On side $i$ define
\begin{equation}\label{TR:eq:material-boundary}
 X_{\alpha,z}(\theta)
 =(1-t_i(x))p_{i-1}(h_\alpha\mathbf1+z)
  +t_i(x)p_i(h_\alpha\mathbf1+z).
\end{equation}
As $\theta$ traverses $K_i$, this traverses the actual straight side;
hence its image is exactly $\partial P(\alpha,h_\alpha\mathbf1+z)$.
Scale normalization of the eigenvalue makes the omitted final area
dilation irrelevant.

Write
\begin{equation}\label{TR:eq:W}
 X_{\alpha,z}(\theta)=e_r(\theta)+W_{\alpha,z}(\theta).
\end{equation}
At $z=0$, \eqref{TR:eq:material-boundary} is the radial parametrization of the
tangential polygon.  Equations \eqref{TR:eq:endpoint-components}--
\eqref{TR:eq:Tphysical} therefore give the exact identity
\begin{equation}\label{TR:eq:normal-W}
 \PiNG(W_{\alpha,z}\cdot e_r)
 =b_\alpha+\PiNG\widetilde T_\alpha z
 =:\widetilde\tau(\alpha,z).
\end{equation}

Choose a fixed annular extension of $W$ which equals $W$ at $r=1$ and
vanishes for $r\le1/2$.  For $\|W\|_{W^{1,\infty}}$ small, the map
$F_W={\rm id}+EW$ is bi-Lipschitz.  Pulling stiffness and mass back by
$F_W$ defines the scale-normalized first eigenvalue
\begin{equation}\label{TR:eq:LambdaW}
 \mathcal L(W)=\frac{|F_W(\mathbb D)|}{\pi}
 \lambda_1(F_W(\mathbb D)).
\end{equation}

\begin{lemma}[Parametric fixed-chart analyticity]
\label{TR:lem:parametric-analytic}
$\mathcal L$ is real analytic on a fixed small ball of
$W^{1,\infty}(\T;\mathbb R^2)$.  It depends only on the boundary image,
and
\begin{equation}\label{TR:eq:parametric-Hessian}
 D\mathcal L(0)=0,
 \qquad
 \frac12D^2\mathcal L(0)[W,W]
 =\qphys\bigl(\PiNG(W\cdot e_r)\bigr).
\end{equation}
\end{lemma}

\begin{proof}
The pullback coefficients are rational analytic functions of $DF_W$ and
$\det DF_W$ on a small $L^\infty$ ball.  Uniform ellipticity, simplicity
of the first eigenvalue, and the Riesz projection give operator-norm
analyticity exactly as in the radial fixed chart.  Two extensions with the
same boundary image are related by a bi-Lipschitz change of variables and
therefore give the same eigenvalue and area.  The disk is critical for the
scale-normalized functional.  Reparametrization invariance removes the
tangential component from its second differential; the normal component
is the Fourier--Bessel disk Hessian.  This proves
\eqref{TR:eq:parametric-Hessian}.
\end{proof}

Define the nonlinear spectral remainder
\begin{equation}\label{TR:eq:N}
 \mathcal N(W)=\mathcal L(W)-\mathcal L(0)
 -\qphys\bigl(\PiNG(W\cdot e_r)\bigr).
\end{equation}
Since \eqref{TR:eq:material-boundary} parametrizes the actual polygon,
\begin{equation}\label{TR:eq:exact-value-split}
 \boxed{
 \lambda_1(P(\alpha,H_\alpha(z)))
 =\mathcal L(0)+\qphys(\widetilde\tau(\alpha,z))
  +\mathcal N(W_{\alpha,z}).}
\end{equation}
Here the left side is area normalized; equivalently one may use the
unnormalized support polygon because \eqref{TR:eq:LambdaW} already performs
the scale normalization.
\endgroup


\part{Complete cubic tensor and quadratic-bubble reduction}

\begingroup
\renewcommand{\Rone}{\mathcal R_1}
\section{Node CT: exact cubic tensor and retraction of CT-JM}
\label{module:CT}
\noindent\textit{Audited source: \nolinkurl{CT_complete_tensor.tex}.}
\subsection{Two functional spaces and the physical support quotient}
\label{CT:sec:functional-spaces}

The earlier working copies used the same symbol $\mathcal X$ for two
different objects.  We separate them here.  For an arbitrary boundary
vector field $W:\T\to\mathbb R^2$ let
\begin{equation}\label{CT:eq:Y-norm}
 \|W\|_{\Y}^{2}
 :=\sum_{n\in\mathbb Z}(1+|n|)|\widehat W(n)|^2.
\end{equation}
Thus $\Y=H^{1/2}(\T;\mathbb R^2)$ is the fixed space in which the disk
Taylor tensors are estimated.

The support space is a different, fan-dependent quotient.  Choose side
normals $\nu_i=e_r(\theta_i)$, put $\tau_i=e_\theta(\theta_i)$, and write
\begin{equation}\label{CT:eq:support-area-row}
 \mathsf Z_\alpha
 :=\left\{z\in\mathbb R^N:
 \sum_i\left(\tan\frac{\alpha_{i-1}}2+
              \tan\frac{\alpha_i}2\right)z_i=0\right\}
 \big/\{(q\cdot\nu_i)_i:q\in\mathbb R^2\}.
\end{equation}
For $x=\theta-\theta_i$ on
$K_i=[\theta_i-\alpha_{i-1}/2,\theta_i+\alpha_i/2]$, set
\begin{align}
 a_i^-&=\frac{z_i-z_{i-1}}{\sin\alpha_{i-1}}
          -z_i\tan\frac{\alpha_{i-1}}2,\notag\\
 a_i^+&=\frac{z_{i+1}-z_i}{\sin\alpha_i}
          +z_i\tan\frac{\alpha_i}2,\label{CT:eq:support-endpoint-jets}\\
 t_i(x)&=\frac{\tan x+\tan(\alpha_{i-1}/2)}
 {\tan(\alpha_{i-1}/2)+\tan(\alpha_i/2)}.\notag
\end{align}
The normal component of the genuine straight-side material velocity is
\begin{equation}\label{CT:eq:physical-normal-trace}
 v_\alpha(z)(\theta)
 =z_i\cos x+\bigl\{(1-t_i(x))a_i^-+t_i(x)a_i^+\bigr\}\sin x.
\end{equation}
Its tangential component, needed only in the curvature row, is
\begin{equation}\label{CT:eq:physical-tangential-trace}
 w_\alpha(z)(\theta)
 =-z_i\sin x+\bigl\{(1-t_i(x))a_i^-+t_i(x)a_i^+\bigr\}\cos x.
\end{equation}
Let $\PiNG$ delete the constant and first harmonics.  The physical Schur
form and the support Hilbert norm are, by definition,
\begin{equation}\label{CT:eq:Xalpha-definition}
 \Gtilde_\alpha(z)
 :=\qD(\PiNG v_\alpha(z)),
 \qquad
 \|z\|_{\X_\alpha}:=\Gtilde_\alpha(z)^{1/2}.
\end{equation}
The disk Bessel coercivity makes this a norm on the quotient
$\mathsf Z_\alpha$.  When a physical support slot is denoted by
$Z_\alpha(z)$, we use the unambiguous convention
\begin{equation}\label{CT:eq:support-X-convention}
 \|Z_\alpha(z)\|_{\X_\alpha}:=\|z\|_{\X_\alpha}.
\end{equation}
In particular, the comparison imported later by NR is not an additional
theorem: with the current notation it is the identity
\begin{equation}\label{CT:eq:exact-norm-interface}
 \boxed{\|Z_\alpha(z)\|_{\X_\alpha}^2
        =\Gtilde_\alpha(z).}
\end{equation}

For later bookkeeping put
$A_i(x)=(1-t_i(x))a_i^-+t_i(x)a_i^+$.  Direct rotation gives the Cartesian
identity
\begin{equation}\label{CT:eq:cartesian-support-trace}
 v_\alpha(z)e_r+w_\alpha(z)e_\theta
 =z_i\nu_i+A_i(x)\tau_i\qquad(x\in K_i).
\end{equation}
The endpoint formulas \eqref{CT:eq:support-endpoint-jets} make the two Cartesian
values at every common vertex identical.  They do \emph{not}, however, give
a ratio-free reconstruction of the full vector trace from its normal
component.

\begin{proposition}[Failure of full-vector reconstruction]
\label{CT:prop:no-vector-reconstruction}
There is no constant independent of $N$ for which
\[
 \|v_\alpha(z)e_r+w_\alpha(z)e_\theta\|_{H^{1/2}}
 \le C\|\PiNG v_\alpha(z)\|_{H^{1/2}}
\]
holds on all exact support classes.
\end{proposition}

\begin{proof}
Take an even regular fan, $\alpha_i=a=2\pi/N$, and the alternating support
mode $z_i=(-1)^i$, which satisfies the area row and has no translation
component.  The endpoint formulas give
\[
 a_i^-=z_i\cot(a/2),\qquad a_i^+=-z_i\cot(a/2).
\]
Thus the normal trace has a fixed normalized-cell profile of amplitude
$O(1)$, while the tangential coefficient $A_i$ runs between values of size
$\cot(a/2)\asymp a^{-1}$.  Scale invariance of the cell
$\dot H^{1/2}$ seminorm and summation over the $N$ cells give
\[
 \|\PiNG v_\alpha(z)\|_{H^{1/2}}^2\asymp N,
 \qquad
 \|v_\alpha(z)e_r+w_\alpha(z)e_\theta\|_{H^{1/2}}^2
 \asymp Na^{-2}.
\]
The squared ratio therefore diverges like $a^{-2}$.
\end{proof}

Consequently tangential material velocities may enter only through the
completed physical-coordinate chain rule in \eqref{CT:eq:physical-chain}; they
cannot be treated as an $H^{1/2}$ support slot.  This is exactly the
vertex-reparametrization issue isolated below.

\subsection{Fixed-domain coefficient ledger}

Let $E$ be a fixed annular extension from boundary vector fields to
$W^{1,\infty}(\mathbb D;\mathbb R^2)$, vanishing on $\{|x|\le1/2\}$.
For three boundary directions $V_i$ put
\[
 F_t(x)=x+\sum_{i=1}^3t_iEV_i(x),\qquad
 G_i=D(EV_i),\qquad G(t)=\sum_it_iG_i.
\]
Write
\begin{equation}\label{CT:eq:JC}
 J(t)=\det(I+G(t)),\qquad
 P(t)=(I+G(t))^{-1},\qquad
 C(t)=J(t)P(t)P(t)^T.
\end{equation}
If $I$ is an ordered list of distinct indices, a subscript $I$ denotes
$\partial_I|_{t=0}$, and $I_0\dot\cup\cdots\dot\cup I_r=I$ denotes an
ordered set partition, empty blocks being allowed.

All coefficient derivatives through cubic order are explicit.  Namely,
\begin{align}
 J_i&=\operatorname {tr}G_i,\notag\\
 J_{ij}&=(\operatorname {tr}G_i)(\operatorname {tr}G_j)
          -\operatorname {tr}(G_iG_j),\qquad J_{ijk}=0,
 \label{CT:eq:J-jets}\\
 P_i&=-G_i,\notag\\
 P_{ij}&=G_iG_j+G_jG_i,\notag\\
 P_{ijk}&=-\sum_{\pi\in S_3}
 G_{\pi(i)}G_{\pi(j)}G_{\pi(k)},
 \label{CT:eq:P-jets}\\
 C_I&=\sum_{I_0\dot\cup I_1\dot\cup I_2=I}
 J_{I_0}P_{I_1}P_{I_2}^{T}.
 \label{CT:eq:C-jets}
\end{align}
Here $J_\varnothing=1$ and $P_\varnothing=I$.  Formula
\eqref{CT:eq:C-jets}, rather than an ellipsis, is the metric/Jacobian ledger.

On $H_0^1(\mathbb D)$ define
\begin{align}
 A_t(u,v)&=\int_{\mathbb D}\nabla u\cdot C(t)\nabla v,
 &M_t(u,v)&=\int_{\mathbb D}J(t)uv,\label{CT:eq:forms}\\
 s(t)&=\pi^{-1}\int_{\mathbb D}J(t),
 &\widehat A_t&=s(t)A_t.\label{CT:eq:scale-form}
\end{align}
Thus the generalized eigenvalue of $(\widehat A_t,M_t)$ is the first
Dirichlet eigenvalue after normalization to area $\pi$.  Its scale rows are
\begin{equation}\label{CT:eq:scale-jets}
 \widehat A_I
 =\sum_{I_0\dot\cup I_1=I}s_{I_0}A_{I_1},
 \qquad
 A_I(u,v)=\int\nabla u\cdot C_I\nabla v,
 \qquad M_I(u,v)=\int J_Iuv.
\end{equation}
No area constraint has yet been imposed in \eqref{CT:eq:scale-jets}; in
particular, the scale terms have not been silently deleted.

To use a fixed inner product, let $S(t)$ be multiplication by
$J(t)^{-1/2}$ and set
\begin{equation}\label{CT:eq:H}
 H(t)=S(t)\widehat A_tS(t)
\end{equation}
in the form sense on $L^2(\mathbb D)$.  The mass/Jacobian jets are
\begin{align}
 S_i&=-\tfrac12J_i,\notag\\
 S_{ij}&=\tfrac34J_iJ_j-\tfrac12J_{ij},\notag\\
 S_{ijk}&=-\tfrac{15}{8}J_iJ_jJ_k
 +\tfrac34(J_{ij}J_k+J_{ik}J_j+J_{jk}J_i),
 \label{CT:eq:S-jets}
\end{align}
and the completed operator jets are
\begin{equation}\label{CT:eq:H-jets}
 H_I=\sum_{I_0\dot\cup I_1\dot\cup I_2=I}
 S_{I_0}\widehat A_{I_1}S_{I_2}.
\end{equation}
Equations \eqref{CT:eq:J-jets}--\eqref{CT:eq:H-jets} list every direct metric,
Jacobian, mass, and scale monomial through third order.

\subsection{The full symmetric state tensor}

Let $u_0$ be the positive normalized disk state, $H_0u_0=\lambda_0u_0$,
and put
\begin{equation}\label{CT:eq:R}
 Q=I-|u_0\rangle\langle u_0|,
 \qquad R=Q(H_0-\lambda_0)^{-1}Q,
 \qquad h_i=QH_iu_0.
\end{equation}
The inverse in \eqref{CT:eq:R} is positive because $\lambda_0$ is the first
simple eigenvalue.

\begin{proposition}[Complete third reduced-resolvent formula]
\label{CT:prop:full-tensor}
For arbitrary directions $V_1,V_2,V_3$, the full symmetric third
derivative is
\begin{align}
 \lambda_{123}={}&\langle u_0,H_{123}u_0\rangle\label{CT:eq:full-tensor}\\
 &-\sum_{\rm cyc(ij,k)}
 \bigl\{\langle QH_{ij}u_0,Rh_k\rangle
          +\langle h_k,RQH_{ij}u_0\rangle\bigr\}\notag\\
 &+\sum_{\pi\in S_3}
 \left\langle h_{\pi(1)},
 R(H_{\pi(2)}-\lambda_{\pi(2)}I)R
 h_{\pi(3)}\right\rangle.\notag
\end{align}
Here $\lambda_i=\langle u_0,H_iu_0\rangle$.  At the scale-normalized disk
$\lambda_i=0$ for every boundary direction.
\end{proposition}

\begin{proof}[Proof of Proposition~\ref{CT:prop:full-tensor}]
Choose the analytic normalized eigenvector $u(t)$ with
$\langle u(t),u(t)\rangle=1$.  The first differentiated state equation
gives
\begin{equation}\label{CT:eq:first-state}
 Qu_i=-Rh_i,
 \qquad \langle u_i,u_0\rangle=0.
\end{equation}
Differentiate once more, project by $Q$, and use the differentiated
normalization to eliminate the component along $u_0$.  Substitution in
$\lambda_i=\langle u,H_iu\rangle$ gives \eqref{CT:eq:full-tensor}.
For a repeated direction it reduces to
\[
 \lambda'''=\langle H'''\rangle
 -6\langle QH''u_0,Rh\rangle
 +6\langle h,R(H'-\lambda'I)Rh\rangle,
\]
which also checks all multiplicities.  Polarization proves the displayed
symmetric formula.  Criticality follows from dilation-normalized disk
stationarity.
\end{proof}

\begin{proposition}[Two energetic slots in the disk cubic tensor]
\label{CT:prop:cubic-tame}
The complete tensor of Proposition~\ref{CT:prop:full-tensor} extends to two
$\Y$ slots and one Lipschitz coefficient slot, and
\begin{equation}\label{CT:eq:cubic-tame}
 |D^3\mathcal L(0)[U,V,W]|
 \le C\sum_{\rm cyc}
 \|U\|_\Y\|V\|_\Y\|W\|_{W^{1,\infty}}.
\end{equation}
The statement concerns the complete vector tensor; no radial
reparametrization row has been separated from it.
\end{proposition}

\begin{proof}
For smooth traces use the exact physical chain rule
\eqref{CT:eq:physical-chain}.  The normal tensor is the null form
\eqref{CT:eq:null} plus the seven order-one rows
\eqref{CT:eq:seven-rows}.  Its multipliers have total order at most two and
obey the stated allocation by the dyadic argument of
Lemma~\ref{CT:lem:seven-row-audit}.

It remains to audit the curvature part.  The disk Hessian has symbol
$d^2(1+\mu_n)$ on $|n|\ge2$, hence order one with the standard discrete
differences.  The three terms in $\Gamma_{ij}$ have orders zero, one, and
one.  After symmetric polarization, their multipliers are finite
permutations of
\begin{equation*}
 (1+\mu_{n_1}),\qquad
 n_2(1+\mu_{n_1}),
 \qquad n_1+n_2+n_3=0,
\end{equation*}
up to the harmless shifts by one caused by $e_r,e_\theta$.  They have
total order at most two.  On each dyadic block put the two largest half
derivatives in $L^2$ and the remaining coefficient, together with its one
allowed derivative, in $L^\infty$.  The finite-rank area and translation
rows vanish by \eqref{CT:eq:geometric-kernel}.  Discrete Coifman--Meyer
summation proves \eqref{CT:eq:cubic-tame}.  Density defines the two energetic
slots and preserves extension independence.
\end{proof}

The uncancelled rows in \eqref{CT:eq:full-tensor} are now literally:
\begin{center}
\begin{longtable}{>{\raggedright\arraybackslash}p{.23\linewidth}
                  >{\raggedright\arraybackslash}p{.66\linewidth}}
\toprule
row & exact source\\
\midrule
\endfirsthead
\toprule
row & exact source\\
\midrule
\endhead
third direct row & $\langle H_{123}u_0,u_0\rangle$ with
\eqref{CT:eq:H-jets};\\
second-state rows & the six $H_{ij}$--$R$--$H_k$ pairings in the second
line of \eqref{CT:eq:full-tensor};\\
three first-state rows & the six $H_i$--$R$--$H_j$--$R$--$H_k$ pairings
in the last line;\\
state normalization & the three $-\lambda_iH_j$ insertions in that same
last line;\\
metric/Jacobian & the $C_I$ terms of \eqref{CT:eq:C-jets};\\
mass & the nonempty $S_I$ terms of \eqref{CT:eq:H-jets};\\
scale/area factor & the nonempty $s_I$ terms of
\eqref{CT:eq:scale-jets}.\\
\bottomrule
\end{longtable}
\end{center}
Thus the word ``state rows'' does not conceal a differentiated state, and
the geometric rows are finite set-partition sums.

\subsection{Physical coordinates, curvature, area, and translations}

Write a boundary vector direction as
\begin{equation}\label{CT:eq:components}
 V_i=f_i e_r+w_i e_\theta.
\end{equation}
The radial graph of the image of
$e_r+\sum t_iV_i$ has first jet $f_i$ and second jet
\begin{equation}\label{CT:eq:Gamma}
 \Gamma_{ij}=w_iw_j-w_if_j'-w_jf_i'.
\end{equation}
Indeed, take the modulus and argument of
$e^{i\theta}(1+\sum t_if_i+i\sum t_iw_i)$ and then express the modulus in
the new angular variable.  This gives \eqref{CT:eq:Gamma} before using any
invariance.

The exact area-and-centering chart adds a unique $\chi_{ij}$ in the span
of $1,\cos\theta,\sin\theta$ to the second radial jet.  Its coefficients
are
\begin{align}
 \chi_{ij}={}&c_{ij}^0+c_{ij}^1\cos\theta+c_{ij}^2\sin\theta,
 \label{CT:eq:chi}\\
 c_{ij}^0={}&-\frac1{2\pi}\int_\T
    (\Gamma_{ij}+f_if_j),\notag\\
 c_{ij}^1={}&-\frac1\pi\int_\T
    (\Gamma_{ij}+2f_if_j)\cos\theta,\notag\\
 c_{ij}^2={}&-\frac1\pi\int_\T
    (\Gamma_{ij}+2f_if_j)\sin\theta.\notag
\end{align}
These are obtained by twice differentiating
\[
 \tfrac12\int R^2=\pi,
 \qquad \tfrac13\int R^3(\cos\theta,\sin\theta)=0.
\]
They are the area row and the two translation-projection rows; there is no
unspecified finite-rank correction.

Let $\Lambda^{\rm n}(\rho)$ denote the scale-normalized scalar in the
normal radial chart.  Since $D\Lambda^{\rm n}(0)=0$, the exact cubic chain
rule is
\begin{multline}\label{CT:eq:physical-chain}
 D^3\mathcal L(0)[V_1,V_2,V_3]
 =D^3\Lambda^{\rm n}(0)[f_1,f_2,f_3]\\
 +\sum_{\rm cyc(ij,k)}D^2\Lambda^{\rm n}(0)
 [f_k,\Gamma_{ij}+\chi_{ij}].
\end{multline}
Scale and translation invariance imply
\begin{equation}\label{CT:eq:geometric-kernel}
 D^2\Lambda^{\rm n}(0)[g,\psi]=0,
 \qquad \psi\in\operatorname {span}\{1,\cos\theta,\sin\theta\}.
\end{equation}

The disk Hessian used throughout this node is the following literal
Fourier--Bessel form.  If
$d=\partial_ru_0(1)$ for the positive $L^2$-normalized disk state and
$\mathcal D_\circ$ is the multiplier in \eqref{CT:eq:DtN}, set
\begin{equation}\label{CT:eq:qD-definition}
 \qD(g,h):=d^2\int_\T
  (\PiNG g)(\mathcal D_\circ+1)(\PiNG h),
 \qquad \qD(g):=\qD(g,g).
\end{equation}
Thus
\begin{equation}\label{CT:eq:disk-Hessian}
 \boxed{\frac12D^2\Lambda^{\rm n}(0)[g,h]=\qD(g,h)},
 \qquad
 c\|\PiNG g\|_{H^{1/2}}^2
 \le \qD(g)\le
 C\|\PiNG g\|_{H^{1/2}}^2.
\end{equation}
Indeed, for a mean-zero diagonal direction the second boundary row is
$u_2|=d(f\mathcal D_\circ f+f^2/2)$.  Green's identity gives the raw
second Taylor coefficient
$d^2\int(f\mathcal D_\circ f+f^2/2)$.  Multiplication by the area factor
$1+\varepsilon^2(2\pi)^{-1}\int f^2$ adds
$j^2(2\pi)^{-1}\int f^2=(d^2/2)\int f^2$, because
$d^2=j^2/\pi$.  This proves the first identity by polarization.
Moreover, $\mu_{\pm1}=-1$ and
$\mu_n+1\asymp1+|n|$ for $|n|\ge2$, which proves the coercivity statement
and the geometric kernel \eqref{CT:eq:geometric-kernel}.

Consequently all $\chi_{ij}$ rows cancel exactly.  The surviving
$D^2\Lambda^{\rm n}[f_k,\Gamma_{ij}]$ rows are the curvature rows.  For
$V_1=V_2=fe_r$ and $V_3=we_\theta$, \eqref{CT:eq:physical-chain} becomes
\begin{equation}\label{CT:eq:tangent-special}
 D^3\mathcal L(0)[fe_r,fe_r,we_\theta]
 =-2D^2\Lambda^{\rm n}(0)[f,wf'],
\end{equation}
or, equivalently, $-4\qD(f,wf')$.  After multiplication by the Taylor
factor $1/2$ appropriate to the differentiated cubic term
$3\mathcal T_3$, the same row is $-2\qD(f,wf')$.  This distinction fixes
the multiplicity used downstream.

\begin{proposition}[Retraction of the curvature-as-remainder estimate]
\label{CT:prop:curvature-counterexample}
Let $f_\alpha$ be the normal fan field and let
$Z_\alpha(z)=v_\alpha(z)e_r+w_\alpha(z)e_\theta$ be a physical support
slot.  The formerly asserted estimate
\begin{equation}\label{CT:eq:curvature-absolute}
 \left|4\qD(f_\alpha,w_\alpha(z)f_\alpha')\right|
 \le CS^{5/2}\|Z_\alpha(z)\|_{\X_\alpha}.
\end{equation}
is false uniformly as $N\to\infty$.  Consequently the curvature row cannot
be placed by itself in the $o(1)$ mixed remainder.
\end{proposition}

\begin{proof}
Take $N=3M$, $a=2\pi/N$, and the three-periodic fan
$\alpha_{3j+r}=ap_r$ with $p=(1/6,8/3,1/6)$.  Choose a three-periodic
support vector satisfying the exact area row and the two translation rows;
after scaling it tends to $(17,-2,0)$.  In the material coordinate
$y=\theta/a$ one has
\[
 f_\alpha=a^2F_p+O(a^4),\qquad
 w_\alpha(z)f_\alpha'=G_{p,z}+O(a^2),\qquad
 \|Z_\alpha(z)\|_{\X_\alpha}^2\asymp a^{-1}.
\]
The high-frequency expansion
$\mathcal D_\circ+1=|D|+1+O(|D|^{-1})$ gives
\[
 \qD(f_\alpha,w_\alpha(z)f_\alpha')
 =2\pi d^2m_{p,z}a+O(a^2),
 \qquad m_{p,z}\ne0.
\]
Here $m_{p,z}$ is the mean Fourier pairing on the material period of length
three.  The factor $2\pi$ comes from its $2\pi/(3a)$ repetitions and the
change $d\theta=a\,dy$.  For a check independent of tail monotonicity, put
\[
 \widehat H(n)=\frac13\int_0^3H(y)e^{-2\pi iny/3}\,dy
\]
for every three-periodic material profile, and put
\[
 S_K=2\sum_{n=1}^K\frac{2\pi n}{3}
 \operatorname {Re}\{\overline{\widehat F_p(n)}
                         \widehat G_{p,z}(n)\}.
\]
On each of the three cells the two functions are polynomials of degree at
most two.  For $omega=2\pi n/3$ their integrals are evaluated by
\[
 I_0(l,u;\omega)=\frac{e^{-i\omega l}-e^{-i\omega u}}{i\omega},
 \qquad
 I_j(l,u;\omega)=
 \left[-\frac{x^je^{-i\omega x}}{i\omega}\right]_{l}^{u}
 +\frac{j}{i\omega}I_{j-1}(l,u;\omega).
\]
Substitution of the rational cell endpoints and coefficients, with
outward rational enclosures for $\pi$, sine, and cosine, gives at $K=4096$
\[
 2.23372<S_{4096}<2.23374.
\]
Moreover $F_p'$ has total variation $6$ and direct differentiation of the
three quadratic pieces, including their endpoint jumps, gives
$\operatorname {TV}(G_{p,z})<210$.  Twice integrating by parts in
$\widehat F_p$ and once in $\widehat G_{p,z}$ therefore yields
\[
 |m_{p,z}-S_K|
 \le\frac{\operatorname {TV}(F_p')\operatorname {TV}(G_{p,z})}
          {2\pi^2K},
 \qquad
 |m_{p,z}-S_{4096}|<0.01559.
\]
Thus $2.218<m_{p,z}<2.250$, in particular $m_{p,z}\ne0$.  The first ten
modes give the additional regression value $2.159716$, but are not used
for the rigorous tail separation.  The left side of
\eqref{CT:eq:curvature-absolute} is therefore of order $a$, whereas its right
side is $O(a^2)$.  This contradicts a uniform constant.
\end{proof}

The distributional curvature pairing remains part of the exact identity
\begin{multline}\label{CT:eq:curvature-as-difference}
 -4\qD(f_\alpha,w_\alpha(z)f_\alpha')\\
 =D^3\mathcal L(0)[B_\alpha,B_\alpha,Z_\alpha(z)]
  -D^3\Lambda^{\rm n}(0)
    [f_\alpha,f_\alpha,v_\alpha(z)].
\end{multline}
but Proposition~\ref{CT:prop:no-vector-reconstruction} and
Proposition~\ref{CT:prop:curvature-counterexample} show that its two terms may
not be estimated separately.  In fact even their sum is not a small
remainder, as proved below; it must be retained in a joint principal block.

\subsection{Extraction of the unique order-two symbol}

For real periodic functions define the same symmetric null form used by
the quadratic-bubble node,
\begin{equation}\label{CT:eq:null}
 \Cnull(f_1,f_2,f_3)
 =\frac13\sum_{\rm cyc}
 \int_\T f_1\{(\Dop f_2)(\Dop f_3)-f_2'f_3'\}.
\end{equation}
With $j=j_{0,1}$, the fixed nonzero disk solvability constant in our
integral convention is
\begin{equation}\label{CT:eq:kappa}
 \boxed{\kappa_{\mathbb D}=-\frac{6j^2}{\pi}}.
\end{equation}
Indeed, if $d=\partial_ru_0(1)$, then
$\partial_\lambda u_0(1)=d/(2j^2)$; the coefficient of the cubic boundary
solvability row is therefore $-2j^2$ times its angular mean.  The minus
sign comes from solving
\(
 (d/(2j^2))\,\delta\lambda+d\,R=0
\)
for the spectral increment.  Conversion from the Taylor coefficient to
$D^3$ supplies the factor $6$.  This sign is fixed by the outward radial
convention in \eqref{CT:eq:components}; it is not absorbed into the definition
of \(\Cnull\).

On the nongeometric subspace let
\begin{align}
 \mathcal D_\circ e^{in\theta}&=\mu_ne^{in\theta}\quad(n\ne0),
 &\mathcal D_\circ1&=0,\label{CT:eq:DtN}\\
 \mu_n&=j\frac{J_{|n|}'(j)}{J_{|n|}(j)}=|n|+r_n,
 &|\Delta^kr_n|&\le C_k(1+|n|)^{-1-k}.\notag
\end{align}
The Bessel equation also gives the exact second normal derivative
\begin{equation}\label{CT:eq:rr-exact}
 \partial_r^2\mathcal E_jg|_{r=1}
 =\Dop^2g+\mathcal B_1g,
 \qquad
 \mathcal B_1e^{in\theta}=(-\mu_n-j^2)e^{in\theta},
\end{equation}
where $\mathcal B_1$ has order one.

\begin{proposition}[Exact seven-row boundary formula]
\label{CT:prop:seven-row}
Let $f$ be real and $\PiZero f=f$, where $\PiZero$ deletes only the
constant mode.  Then
\begin{equation}\label{CT:eq:exact-boundary-cubic}
 D^3\Lambda^{\rm n}(0)[f,f,f]
 =\kappa_{\mathbb D}\,\mathfrak B(f),
\end{equation}
where
\begin{multline}\label{CT:eq:B-exact}
 \mathfrak B(f)
 =\int_\T f\,\mathcal D_\circ
       \left(f\mathcal D_\circ f+\frac12f^2\right)
 -\frac12\int_\T f^2(\Dop^2+\mathcal B_1)f\\
 +\frac{2-j^2}{6}\int_\T f^3.
\end{multline}
Consequently
\begin{equation}\label{CT:eq:seven-row}
 \mathfrak B(f)=\Cnull(f,f,f)+\sum_{a=1}^7\mathfrak r_a(f),
\end{equation}
with the seven literal lower rows
\begin{align}
 \mathfrak r_1(f)&=\int f\Dop(f\,r(D)f),
 &\mathfrak r_2(f)&=\int f\,r(D)(f\Dop f),\\
 \mathfrak r_3(f)&=\int f\,r(D)(f\,r(D)f),
 &\mathfrak r_4(f)&=\frac12\int f\Dop(f^2),\\
 \mathfrak r_5(f)&=\frac12\int f\,r(D)(f^2),
 &\mathfrak r_6(f)&=-\frac12\int f^2\mathcal B_1f,
 \label{CT:eq:seven-rows}\\
 \mathfrak r_7(f)&=\frac{2-j^2}{6}\int f^3.&&
\end{align}
Here $r(D)=\mathcal D_\circ-\Dop$ on nonzero modes and is zero on the
constant mode.  The polarized formulas are obtained by symmetric
polarization of \eqref{CT:eq:seven-rows}.
\end{proposition}

\begin{proof}
Write $u_0(r)=c_0J_0(jr)$, $d=\partial_ru_0(1)$, and expand without
factorials
\begin{equation*}
 u=u_0+\varepsilon u_1+\varepsilon^2u_2+
   \varepsilon^3u_3+O(\varepsilon^4),
 \qquad
 \lambda=j^2+\varepsilon^2\lambda_2+
   \varepsilon^3\lambda_3+O(\varepsilon^4).
\end{equation*}
The linear term vanishes because $f$ has zero mean.  The first two
Dirichlet rows are exactly
\begin{align}
 u_1|_{r=1}&=-df,
 &u_2|_{r=1}
 &=d\left(f\mathcal D_\circ f+\frac12f^2\right).
 \label{CT:eq:first-two-boundary-rows}
\end{align}
The nonzero angular modes of $\partial_ru_2$ are therefore
$d\mathcal D_\circ(f\mathcal D_\circ f+f^2/2)$.  Its constant mode contains
the second spectral correction and the state-normalization correction, but
it disappears after pairing with the mean-zero function $f$.  The Bessel
equation gives the exact identities
\begin{equation}\label{CT:eq:radial-identities}
 \partial_r^2u_1|_{r=1}
 =-d(\Dop^2+\mathcal B_1)f,
 \qquad
 \partial_r^3u_0(1)=(2-j^2)d.
\end{equation}

The third Dirichlet row is
\begin{equation*}
 u_3|_{r=1}
 =-f\partial_ru_2-\frac12f^2\partial_r^2u_1
   -\frac16f^3\partial_r^3u_0.
\end{equation*}
Choose the state normalization with $\langle u_1,u_0\rangle=0$.
Green's identity then gives
$\lambda_3=d\int_\T u_3|_{r=1}$.  Since the Rellich identity yields
$d^2=j^2/\pi$, multiplication by $3!=6$ gives
$-6d^2=\kappa_{\mathbb D}$.  Substitution of
\eqref{CT:eq:first-two-boundary-rows}--\eqref{CT:eq:radial-identities} proves
\eqref{CT:eq:exact-boundary-cubic}--\eqref{CT:eq:B-exact}.

Finally insert $\mathcal D_\circ=\Dop+r(D)$.  The two terms containing
only $\Dop$ are
\begin{equation*}
 \int f\Dop(f\Dop f)-\frac12\int f^2\Dop^2f
 =\Cnull(f,f,f).
\end{equation*}
The remaining terms are exactly \eqref{CT:eq:seven-rows}.
\end{proof}

\begin{lemma}[Seven-row multiplier audit]
\label{CT:lem:seven-row-audit}
For the symmetric polarization $\mathfrak r_a(f_1,f_2,f_3)$ of every row
in \eqref{CT:eq:seven-rows},
\begin{equation}\label{CT:eq:seven-row-estimate}
 |\mathfrak r_a(f_1,f_2,f_3)|
 \le C\sum_{p<q}\|f_p\|_{H^{1/2}}
                    \|f_q\|_{H^{1/2}}
                    \|f_\ell\|_{L^\infty}.
\end{equation}
The constant is uniform over $a=1,\ldots,7$.
\end{lemma}

\begin{proof}
The symbols of $r(D)$ and $\mathcal B_1$ satisfy
\begin{align*}
 |\Delta^kr_n|&\le C_k(1+|n|)^{-1-k},\\
 |\Delta^k(-\mu_n-j^2)|&\le C_k(1+|n|)^{1-k}.
\end{align*}
After polarization, the seven multipliers are finite permutations of
\begin{multline}\label{CT:eq:seven-symbols}
 |n_2|r_{n_3},\quad r_{n_2}|n_3|,\quad
 r_{n_2}r_{n_3},\quad |n_2|,\quad r_{n_2},\quad
 -\mu_{n_3}-j^2,\quad1,\\
 n_1+n_2+n_3=0.
\end{multline}
Each multiplier has total order at most one and satisfies
\begin{equation*}
 |\Delta^\gamma m(n_1,n_2,n_3)|
 \le C_\gamma(1+\max_i|n_i|)^{1-|\gamma|}.
\end{equation*}
On a dyadic block the two largest frequencies are comparable.  Put the
corresponding half derivatives in $L^2$ and the remaining low-frequency
factor in $L^\infty$.  Discrete Coifman--Meyer summation, followed by
Cauchy--Schwarz in the common high index, proves
\eqref{CT:eq:seven-row-estimate}.  The finitely many modes with
$|n|\le2$ are absorbed by the inhomogeneous norms.
\end{proof}

\begin{theorem}[Exact cubic split and literal remainder formula]
\label{CT:thm:complete-split}
For mean-zero normal directions, $\PiZero f_i=f_i$, one has the exact
identity
\begin{equation}\label{CT:eq:complete-split}
 \boxed{
 D^3\Lambda^{\rm n}(0)[f_1,f_2,f_3]
 =\kappa_{\mathbb D}\Cnull(f_1,f_2,f_3)
  +\Rone(f_1,f_2,f_3).}
\end{equation}
The remainder is not defined by an order symbol.  Its literal formula is
\begin{align}
 \Rone(f_1,f_2,f_3)={}&
 \langle u_0,H_{123}u_0\rangle
 -\sum_{\rm cyc(ij,k)}
 \bigl\{\langle QH_{ij}u_0,Rh_k\rangle
        +\langle h_k,RQH_{ij}u_0\rangle\bigr\}
 \notag\\
 &+\sum_{\pi\in S_3}
 \left\langle h_{\pi(1)},
 R(H_{\pi(2)}-\lambda_{\pi(2)}I)R
 h_{\pi(3)}\right\rangle
 -\kappa_{\mathbb D}\Cnull(f_1,f_2,f_3),
 \label{CT:eq:R-literal}
\end{align}
where every $H_I$ is the finite set-partition expression
\eqref{CT:eq:H-jets}, with \eqref{CT:eq:scale-jets}, \eqref{CT:eq:C-jets}, and
\eqref{CT:eq:J-jets} substituted.  Thus \eqref{CT:eq:R-literal}, unlike the old
tag names alone, specifies every scalar monomial and its multiplicity.
Equivalently, after the necessary direct/state cancellations have been
performed, the same remainder is the seven-row grouped expression
\begin{equation}\label{CT:eq:R-seven-polarized}
 \boxed{\Rone(f_1,f_2,f_3)
 =\kappa_{\mathbb D}\sum_{a=1}^7
   \mathfrak r_a(f_1,f_2,f_3),}
\end{equation}
where $\mathfrak r_a(f_1,f_2,f_3)$ denotes symmetric polarization of the
corresponding line of \eqref{CT:eq:seven-rows}.  Formula
\eqref{CT:eq:R-seven-polarized}, rather than the uncancelled raw monomials, is
the row decomposition used for analytic estimates.

For bookkeeping, expand \eqref{CT:eq:R-literal} fully and assign each scalar
monomial its first applicable tag in the ordered list
\begin{equation}\label{CT:eq:R-tag-priority}
 r(D),\quad \mathcal B_1,\quad C_I\ (I\ne\varnothing),\quad
 S_I\ (I\ne\varnothing),\quad s_I\ (I\ne\varnothing).
\end{equation}
This gives the following disjoint ledger:
\begin{equation}\label{CT:eq:R-ledger}
 \Rone=R_{\rm state}^{r}+R_{rr}+R_{\rm met}
       +R_{J/M}+R_{\rm scale}.
\end{equation}
Here:
\begin{enumerate}
\item $R_{\rm state}^{r}$ consists of the state monomials in
\eqref{CT:eq:full-tensor} containing at least one $r_n$ from
\eqref{CT:eq:DtN};
\item $R_{rr}$ consists of those containing $\mathcal B_1$ from
\eqref{CT:eq:rr-exact};
\item $R_{\rm met}$ consists of the remaining nonprincipal monomials
tagged by a nonempty $C_I$ in \eqref{CT:eq:C-jets};
\item $R_{J/M}$ consists of the remaining monomials whose first tag is a
nonempty $S_I$ in \eqref{CT:eq:H-jets};
\item $R_{\rm scale}$ consists of the remaining monomials whose first tag
is a nonempty $s_I$ in \eqref{CT:eq:scale-jets};
\end{enumerate}
The priority \eqref{CT:eq:R-tag-priority} makes the list disjoint, and the
set-partition formulae make it exhaustive.
\end{theorem}

\begin{proof}[Proof of Theorem~\ref{CT:thm:complete-split}]
Proposition~\ref{CT:prop:seven-row} proves the diagonal identity
\eqref{CT:eq:complete-split} together with
\eqref{CT:eq:R-seven-polarized}; symmetric polarization gives arbitrary
nongeometric slots.  Independently, Proposition~\ref{CT:prop:full-tensor}
gives the complete reduced-resolvent expression.  Subtracting the same
principal null form from that expression gives \eqref{CT:eq:R-literal}.
Hence the fixed-domain formula and the seven grouped boundary rows are two
exact representations of the same trilinear remainder.

Expanding \eqref{CT:eq:R-literal} and applying
\eqref{CT:eq:R-tag-priority} records the origin of every raw monomial in
\eqref{CT:eq:R-ledger}.  No estimate is taken before the raw metric and state
rows have been recombined into the seven rows
\eqref{CT:eq:R-seven-polarized}; this is the required order-two
direct/resolvent cancellation.
\end{proof}

\begin{corollary}[Order-one normal remainder]
\label{CT:cor:R-estimate}
For mean-zero normal slots,
\begin{equation}\label{CT:eq:R-estimate}
 |\Rone(f_1,f_2,f_3)|
 \le C\sum_{p<q}\|f_p\|_{H^{1/2}}
                    \|f_q\|_{H^{1/2}}
                    \|f_\ell\|_{L^\infty},
 \qquad \{\ell\}=\{1,2,3\}\setminus\{p,q\}.
\end{equation}
\end{corollary}

\begin{proof}
Use \eqref{CT:eq:R-seven-polarized} and sum
\eqref{CT:eq:seven-row-estimate} over the seven rows.
\end{proof}

For later fan inputs the constant scale mode must be retained explicitly.
For arbitrary normal slots write
\begin{equation*}
 g_i=\PiZero f_i,
 \qquad \psi_i=(I-\PiZero)f_i,
\end{equation*}
and, for $\varnothing\ne A\subset\{1,2,3\}$, let
$\xi_i^A=\psi_i$ if $i\in A$ and $\xi_i^A=g_i$ otherwise.  Define the
finite seven-term geometric remainder by
\begin{equation}\label{CT:eq:Rgeo}
 \Rgeo(f_1,f_2,f_3)
 :=\sum_{\varnothing\ne A\subset\{1,2,3\}}
 D^3\Lambda^{\rm n}(0)
 [\xi_1^A,\xi_2^A,\xi_3^A]
\end{equation}
and set
\begin{equation}\label{CT:eq:Rfull}
 \Rfull(f_1,f_2,f_3)
 :=\Rone(g_1,g_2,g_3)+\Rgeo(f_1,f_2,f_3).
\end{equation}
Multilinearity gives, with no omitted low mode,
\begin{equation}\label{CT:eq:normal-full-split}
 \boxed{
 D^3\Lambda^{\rm n}(0)[f_1,f_2,f_3]
 =\kappa_{\mathbb D}\Cnull(g_1,g_2,g_3)
  +\Rfull(f_1,f_2,f_3).}
\end{equation}
Every term of \eqref{CT:eq:Rgeo} contains a constant.
Apply Proposition~\ref{CT:prop:cubic-tame} with that smooth finite-mode slot
as the Lipschitz coefficient slot and the other two as energetic slots.
Together with Corollary~\ref{CT:cor:R-estimate}, this proves the same
order-one bound for $\Rfull$ as in \eqref{CT:eq:R-estimate}, with every $f_i$
replaced by its full inhomogeneous norm.  For the fan input,
\begin{equation}\label{CT:eq:fan-geometric-size}
 \|(I-\PiZero)f_\alpha\|_{W^{1,\infty}}\le CS^2,
\end{equation}
so the geometric pure and mixed rows have the lower majorants $CS^5$ and
$CS^{7/2}\|Z\|_\X$, respectively.

The physical curvature is a separate tensor, not a summand of the normal
remainder.  Define
\begin{equation}\label{CT:eq:K-physical}
 \mathcal K_{\rm curv}(V_1,V_2,V_3)
 :=\sum_{\rm cyc(ij,k)}D^2\Lambda^{\rm n}(0)
     [f_k,\Gamma_{ij}].
\end{equation}
Set
\begin{equation}\label{CT:eq:Rphys}
 \Rphys(V_1,V_2,V_3)
 :=\Rfull(f_1,f_2,f_3)
   +\mathcal K_{\rm curv}(V_1,V_2,V_3).
\end{equation}
Then \eqref{CT:eq:physical-chain} and the geometric kernel give the typed
identity
\begin{equation}\label{CT:eq:physical-complete-split}
 \boxed{
 D^3\mathcal L(0)[V_1,V_2,V_3]
 =\kappa_{\mathbb D}\Cnull(f_1,f_2,f_3)
  +\Rphys(V_1,V_2,V_3).}
\end{equation}
The $\chi_{ij}$ rows are absent from this formula because they vanish
exactly, not because they have been assigned to $\Rone$.

For clarity, the cancellations used above are recorded separately:
\begin{center}
\begin{tabular}{>{\raggedright\arraybackslash}p{.27\linewidth}
                >{\raggedright\arraybackslash}p{.61\linewidth}}
\toprule
family & cancellation or surviving order\\
\midrule
two principal state derivatives & self-adjointness and one integration by
parts combine them into \eqref{CT:eq:null};\\
mixed-sign order-two interactions & cancel inside
$(\Dop f_2)(\Dop f_3)-f_2'f_3'$;\\
direct metric versus second state & their order-two pieces reconstruct the
boundary acceleration in Proposition~\ref{CT:prop:seven-row}; after the
combination only the seven order-one rows \eqref{CT:eq:seven-rows} remain;\\
state normalization & zero because $\lambda_i=0$ at the normalized disk;\\
mass and scale & no tangential multiplier, hence order zero;\\
curvature & total order two in general; its ambient two-energetic-slot bound
does not imply a radial support bound.  Both the separate remainder estimate
and the CT--JM small-sum estimate are false; the row is retained inside the
honest JP--JC joint scalar audit;\\
area and translations & exact Hessian kernel
\eqref{CT:eq:geometric-kernel}.\\
\bottomrule
\end{tabular}
\end{center}

We next define the scalar and support functionals used in the relative
argument.  Put $\alpha_*=a\mathbf1$,
$\gamma(t)=\alpha_*+t(\alpha-\alpha_*)$, and write
$f_t=f_{\gamma(t)}$.  Fix at $t=0$ a complement to the area and two
translation rows.  Let
\begin{equation}\label{CT:eq:constraint-transport}
 C_t:\mathsf Z_{\alpha_*}\longrightarrow\mathsf Z_{\gamma(t)}
\end{equation}
be the analytic constraint transport obtained by adding the unique linear
combination of the scale and two translation vectors that satisfies the
three rows at $\gamma(t)$.  Its $3\times3$ Gram matrix is a uniform
perturbation of the regular-fan matrix on the material polydisc, so $C_t$
and $C_t^{-1}$ are uniformly bounded.  For a class $z$ in the fixed
complement define
\begin{align}
 F_1(t)&:=\Rfull(f_t,f_t,f_t),\label{CT:eq:F1-definition}\\
 \Psi_{1,n}(t)[z]
 &:=\Rfull(f_t,f_t,v_{\gamma(t)}(C_tz)),
 \label{CT:eq:Psi1n-definition}\\
\Psi_{1,t}(t)[z]
 &:=-4\qD
 (f_t,w_{\gamma(t)}(C_tz)f_t').
 \label{CT:eq:Psi1t-definition}
\end{align}
The last line is exactly the two curvature placements in
\eqref{CT:eq:physical-chain}; it is not included a second time in
$\Psi_{1,n}$.

\begin{lemma}[Material Cauchy bounds for the normal order-one rows]
\label{CT:lem:R-Cauchy}
In the uniformly comparable material regime
\eqref{CT:eq:uniformly-comparable}, $F_1$ and $\Psi_{1,n}$ extend
holomorphically to the line polydisc of radius $c/\delta(t)$, where
\begin{equation*}
 \delta(t)=\max_i\frac{|\alpha_i-a|}{\gamma_i(t)}.
\end{equation*}
Their scalar and dual majorants are
\begin{equation}\label{CT:eq:R-complex-majorants}
 |F_1|\le CS^5,\qquad
 \|\Psi_{1,n}\|_{\X_{\gamma(t)}^*}\le CS^{7/2}.
\end{equation}
Consequently, for
$E(t)=\sum_i\gamma_i(t)^2(\alpha_i-a)^2$,
\begin{align}
 |F_1''(t)|&\le CSE(t),\label{CT:eq:F1-Cauchy}\\
 \|\Psi_{1,n}'(t)\|_{\X_{\gamma(t)}^*}
 &\le CS^{3/2}E(t)^{1/2}.\label{CT:eq:Psi1n-Cauchy}
\end{align}
\end{lemma}

\begin{proof}
Split every polar cell at its normal and pull both halves to $[0,1]$.
The secant profile, the traces
\eqref{CT:eq:physical-normal-trace}--\eqref{CT:eq:physical-tangential-trace},
and the constraint matrix in \eqref{CT:eq:constraint-transport} are rational
analytic functions of the incident widths.  Their denominators stay in a
fixed sector on the stated polydisc.  The normal estimate
\eqref{CT:eq:R-estimate} together with the finite geometric expansion
\eqref{CT:eq:Rgeo} gives both bounds in
\eqref{CT:eq:R-complex-majorants}.  Same and adjacent cells use
the normalized jets.  On separated cells, subtracting both cell means
makes the first-separating dyadic annuli summable.  No bound contains the
number of cells.

One-variable Cauchy on
$\gamma(t)+\zeta(\alpha-\alpha_*)$ gives respectively
$CS^5\delta(t)^2$ and $CS^{7/2}\delta(t)$.  If $k$ realizes $\delta(t)$, comparability gives
\begin{equation*}
 E(t)\ge\gamma_k(t)^2(\alpha_k-a)^2
       \asymp S^4\delta(t)^2.
\end{equation*}
This proves \eqref{CT:eq:F1-Cauchy}--\eqref{CT:eq:Psi1n-Cauchy}.
\end{proof}

\begin{lemma}[Relative closure of the pure cubic row]
\label{CT:lem:R-relative}
Use the material notation of Section~\ref{CT:sec:material}.  Let $f_\alpha$
be the exact normal fan input and let
$Z_\alpha(z)$ be one physical support slot.  On the exact scale/translation
quotient there is a modulus $\varepsilon_{1}\to0$ such that
\begin{equation}
 |\Rfull(f_\alpha,f_\alpha,f_\alpha)
  -\Rfull(f_{a\mathbf1},f_{a\mathbf1},f_{a\mathbf1})|
 \le\varepsilon_1(S)\Jfour(\alpha).\label{CT:eq:R-relative-pure}
\end{equation}
Equivalently, in the typed physical notation \eqref{CT:eq:Rphys}, the left
side is
\begin{equation}
 \left|\Rphys(B_\alpha,B_\alpha,B_\alpha)
       -\Rphys(B_{a\mathbf1},B_{a\mathbf1},B_{a\mathbf1})\right|.
\end{equation}
\end{lemma}

\begin{proof}
The full remainder estimate following \eqref{CT:eq:Rfull},
\eqref{CT:eq:fan-geometric-size}, and the fan sizes give absolute majorants
$CS^5$ for the pure row.

Fix $0<\beta<1$.  If $\Jfour\ge S^{5-\beta}$, these bounds divided by the
right side of \eqref{CT:eq:R-relative-pure} gives $CS^\beta$.
In the complementary regime the exact quartic Bregman identity
\eqref{CT:eq:J4-bregman-comparability} and the quantitative estimate following
it give
$\|\alpha-a\mathbf1\|_\infty/S\le2S^{(1-\beta)/2}$, so for small $S$ the
entire segment is uniformly comparable and lies in the material polydisc
\eqref{CT:eq:material-polydisc}.
Lemma~\ref{CT:lem:R-Cauchy} gives the required material derivative bound.
Cyclic symmetry kills the first pure derivative at the regular fan.
Integrate with \eqref{CT:eq:action}.  This proves
\eqref{CT:eq:R-relative-pure} with no cell-count factor.
\end{proof}

\begin{proposition}[Retraction of CT--JM]
\label{CT:prop:joint-mixed-counterexample}
There is no modulus $\varepsilon_{\rm JM}(S)\to0$ for which
\begin{equation}\label{CT:eq:false-R-relative-mixed}
 \left|\Rphys(B_\alpha,B_\alpha,Z_\alpha(z))\right|
 \le\varepsilon_{\rm JM}(S)\Jfour(\alpha)^{1/2}
       \|Z_\alpha(z)\|_{\X_\alpha}
\end{equation}
holds uniformly on the exact scale/translation quotient.
\end{proposition}

\begin{proof}
Use the three-periodic family from
Proposition~\ref{CT:prop:curvature-counterexample}.  Its seven normal rows
satisfy, by the displayed order-one multiplier estimate,
\[
 \Rfull(f_\alpha,f_\alpha,v_\alpha(z))=O(a^3),
\]
whereas
\[
 -4\qD(f_\alpha,w_\alpha(z)f_\alpha')
 =-8\pi d^2m_{p,z}a+O(a^2),\qquad m_{p,z}\ne0.
\]
For this family $S\asymp a$,
$\Jfour(\alpha)\asymp a^3$, and
$\|Z_\alpha(z)\|_{\X_\alpha}\asymp a^{-1/2}$; hence the product on the
right of \eqref{CT:eq:false-R-relative-mixed} is comparable to $a$.  Dividing
by that product gives a positive limiting lower bound, contradicting
$\varepsilon_{\rm JM}(S)\to0$.
\end{proof}

The surviving leading support forcing in the differentiated cubic Taylor
row is therefore
\begin{equation}\label{CT:eq:effective-leading-forcing}
 2\qD(f_\alpha,v_\alpha(z))
 -2\qD(f_\alpha,w_\alpha(z)f_\alpha')
 =2\qD\bigl(f_\alpha,
     v_\alpha(z)-w_\alpha(z)f_\alpha'\bigr).
\end{equation}
For a polar graph $R_\alpha$, the corresponding genuine polygon-normal
trace is
\begin{equation}\label{CT:eq:effective-normal-trace}
 v_{\alpha,{\rm eff}}(z)
 =\frac{v_\alpha(z)-(R_\alpha'/R_\alpha)w_\alpha(z)}
        {\sqrt{1+(R_\alpha'/R_\alpha)^2}}.
\end{equation}
Equations \eqref{CT:eq:effective-leading-forcing}--
\eqref{CT:eq:effective-normal-trace}, together with all remaining cubic rows,
are the exact input exported to the JP/JC joint scalar audit.  No estimate
in this node replaces that joint principal analysis.

\begin{remark}[Sign and geometric tests]
The Fourier symbol in \eqref{CT:eq:null} vanishes when the two differentiated
frequencies have opposite signs and equals twice their product when they
have the same sign.  A constant in any slot gives zero by Parseval.  The
full tensor, including \eqref{CT:eq:physical-chain}, vanishes on dilation and
translation paths; the principal and lower rows need not vanish there
separately.
\end{remark}

\subsection{Uniform analytic tail on material cells}
\label{CT:sec:material}

Let
\begin{equation}\label{CT:eq:fan}
 \alpha_i>0,\qquad \sum_i\alpha_i=2\pi,\qquad
 a=2\pi/N,\qquad S=a+\max_i\alpha_i,
\end{equation}
and let $B_\alpha$ be the physical fan boundary field.

\begin{theorem}[Unconditional fixed-chart and natural-order facts]
\label{CT:thm:fixed-chart-unconditional}
On a fixed sufficiently small real or complex $W^{1,\infty}$ ball, the
pulled-back stiffness and mass forms, their compact solution operator, and
the scale-normalized simple first-eigenvalue branch $\mathcal L(W)$ are
analytic in operator norm.  On that complexified chart its bilinear
Hessian, and hence at every real near-disk active polygon the genuine
polygon-preserving Hessian on continuous physical-side-affine fields,
satisfies
\begin{equation}\label{CT:eq:polygon-natural-Hessian}
 |D^2\mathcal L(W)[U,V]|
 \le C\|U\|_{H^{1/2}(\partial P)}
       \|V\|_{H^{1/2}(\partial P)}.
\end{equation}
Here the trace norms on the complexified chart are transported to the
underlying real boundary; these norms are uniformly equivalent throughout
the fixed chart.
The constant is independent of the number, the minimum size, and the
ratios of the sides.
\end{theorem}

\begin{proof}
For the first assertion use the annular extension in the fixed-domain
ledger.  The coefficients $J$, $P$, $C$, $S$, and $H$ in
\eqref{CT:eq:JC}--\eqref{CT:eq:H} are convergent rational power series in
$D(EW)$ on a fixed $L^\infty$ ball.  Uniform ellipticity, the disk spectral
gap, the resolvent identity, and the Riesz projection therefore give the
operator-norm analytic simple branch.  This applies directly to polygonal
boundary traces and does not round their corners.

For the second assertion, extend a physical affine-side trace by a uniform
right inverse from $H^{1/2}(\partial P;\mathbb R^2)$ to
$H^1(P;\mathbb R^2)$.  The radial bi-Lipschitz maps between the disk and
the polygons give a uniform norm for this inverse.  Convexity and the
common inner and outer radii give uniform bounds for the eigenvalue, its
spectral gap, the ground state, and its gradient.  At a supporting line
the barrier $s(D-s)$ gives
$u(x)\le C\operatorname {dist}(x,\partial P)$; a scaled interior estimate
then gives $\|\nabla u\|_\infty\le C$.

In the exact second-variation formula, the first coefficient jet is linear
and the second is quadratic in the two extension gradients.  The direct
rows are bounded by the product of their $H^1$ norms; the two state
forcings have the same $H^{-1}$ bounds, and the reduced resolvent is
uniformly coercive.  On the smaller complex ball the real part of the
pulled-back form remains uniformly coercive and the same spectral contour
separates the analytic simple branch; the identical energy estimates
therefore hold for the complex bilinear Hessian.  Polarization proves
\eqref{CT:eq:polygon-natural-Hessian}.  Approximation by Lipschitz extensions
and invariance under changes with zero boundary trace justify the use of
the energy extensions.
\end{proof}

Put
\begin{equation}\label{CT:eq:R4-definition}
 \mathcal R_4(W)=\mathcal L(W)-\mathcal L(0)
 -\frac12D^2\mathcal L(0)[W,W]
 -\frac16D^3\mathcal L(0)[W,W,W].
\end{equation}

\begin{corollary}[Unconditional tame scalar fourth remainder]
\label{CT:cor:tame-scalar-fourth}
On a fixed smaller real or complex chart ball,
\begin{equation}\label{CT:eq:tame-scalar-fourth}
 \boxed{
 |\mathcal R_4(W)|
 \le C\|W\|_{H^{1/2}}^2\|W\|_{W^{1,\infty}}^2.}
\end{equation}
The constant is independent of the number, minimum size, and ratios of the
physical sides.  This is a scalar statement and makes no assertion about
the retired support-covector row of CT--M4.
\end{corollary}

\begin{proof}
The scale-normalized disk branch has $D\mathcal L(0)=0$.  For fixed $W$ put
\[
 g(\zeta)=D^2\mathcal L(\zeta W)[W,W].
\]
Theorem~\ref{CT:thm:fixed-chart-unconditional} makes $g$ analytic for
$|\zeta|<r_0/\|W\|_{W^{1,\infty}}$ and bounds it there by
$C\|W\|_{H^{1/2}}^2$.  Cauchy's coefficient estimate, after decreasing the
chart so that this radius is greater than $2$, gives for $0\le t\le1$
\[
 |g(t)-g(0)-tg'(0)|
 \le C\|W\|_{H^{1/2}}^2\|W\|_{W^{1,\infty}}^2t^2.
\]
Twice integrating the scalar Taylor formula yields
\[
 \mathcal R_4(W)
 =\int_0^1(1-t)\{g(t)-g(0)-tg'(0)\}\,dt,
\]
and proves \eqref{CT:eq:tame-scalar-fourth}.
\end{proof}

Call a material line $h_i(t)=a+td_i$ uniformly comparable when
\begin{equation}\label{CT:eq:uniformly-comparable}
 \frac a2\le h_i(t)\le\frac{3a}{2}
 \qquad(1\le i\le N,\ 0\le t\le1).
\end{equation}
For such a line put
\[
 E(t)=\sum_i h_i(t)^2d_i^2.
\]

\begin{remark}[Retired CT--M4 proposal: not an active obligation]
\label{CT:obl:material-fourth}
Historically, this proposal assumed that CT--JP supplied a no-ratio Hilbert
form $\Gjp_\alpha$ on the exact support quotient and wrote
$\|z\|_{\Xjp_\alpha}=\Gjp_\alpha(z)^{1/2}$.  After splitting every physical
cell at its normal and pulling the two halves to fixed unit intervals,
prove the following four estimates with a constant independent of the
number, minimum size, and ratios of the cells and moduli
$\omega_4,\omega_{4,s}\downarrow0$ at the origin:
\begin{align}
 |\mathcal R_4(B_\alpha)|&\le CS^5,\notag\\
 \|D\mathcal R_4(B_\alpha)\|_{(\Xjp_\alpha)^*}&\le CS^{7/2},
 \label{CT:eq:M4-size}\\
 \left|\frac{d^2}{dt^2}\mathcal R_4(B_{h(t)})\right|
 &\le \omega_4(S) E(t),\notag\\
 \left\|\frac d{dt}\left\{D\mathcal R_4(B_{h(t)})
     \circ\mathcal T_t\right\}\right\|_{(\Xjp_\alpha)^*}
 &\le \omega_{4,s}(S)E(t)^{1/2},
 \label{CT:eq:M4-path}
\end{align}
Here the varying support fibers are identified by the uniformly bounded
constraint-preserving material transport
\[
 \mathcal T_t:\Xjp_\alpha\longrightarrow\Xjp_{h(t)},
 \qquad
 \|\mathcal T_t\|+\|\mathcal T_t^{-1}\|\le C,
 \qquad \mathcal T_1=I.
\]
It is obtained by the fixed half-cell pullback followed by the exact
area/translation projection.  In the second estimate of
\eqref{CT:eq:M4-path}, the derivative is the total derivative of the displayed
fixed-fiber covector and therefore includes the variation of the physical
support field as the fan changes.  The estimates concern the completed direct-plus-state
functional; its order-two rows may not be estimated separately.
\end{remark}

\begin{remark}[Why the old CT--M4 proposal was insufficient]
Theorem~\ref{CT:thm:fixed-chart-unconditional} proves scalar analyticity and
an absolute $H^{1/2}$ Hessian bound, but neither gives the vanishing moduli
in \eqref{CT:eq:M4-path}.  Those moduli express the cancellation between the
direct metric row and the reduced-resolvent row at the natural boundary
order.  Historical rounded Calder\'on arguments estimated those rows
separately and therefore do not prove this obligation.

Corollary~\ref{CT:cor:tame-scalar-fourth} now closes the first, purely scalar
bound in \eqref{CT:eq:M4-size}.  It does not close the support-covector bound
there, either relative path estimate in \eqref{CT:eq:M4-path}, or the moving
support-fiber comparison.  Thus it is sufficient for JP--BU but does not
resurrect CT--M4.

The former fifth ``support Hessian'' row is not part of CT--M4.  It used the
wrong radial Schur metric and cannot control the physical support problem.
Its correct replacement is the exact support Hessian in JP--EH, integrated
without a post-cubic support remainder.

The former CT--AH hypothesis---bounded holomorphy of
$W\mapsto D^2\mathcal L(W)$ in
$\operatorname {Bil}(H^{1/2}\times H^{1/2})$ on a uniform complex
$W^{1,\infty}$ ball---does not by itself imply CT--M4, because the ambient
full-vector norm need not be uniformly comparable to the joint support
metric.  CT--AH is no longer a vertex of the active DAG.
\end{remark}

The no-ratio physical sizes are
\begin{equation}\label{CT:eq:fan-sizes}
 \|B_\alpha\|_\Y^2\le CS^3,\qquad
 \|B_\alpha\|_{W^{1,\infty}}\le CS.
\end{equation}
Split each material cell at its normal and pull both halves to $[0,1]$.
The secant, support, Jacobian, and finite constraint coefficients are
rational analytic functions of the incident widths.  Hence, on a fan
satisfying \eqref{CT:eq:uniformly-comparable}, they extend to
\begin{equation}\label{CT:eq:material-polydisc}
 |\zeta_i-\alpha_i|<c\alpha_i,\qquad
 \sum_i\zeta_i=2\pi,
\end{equation}
with the same bounds as \eqref{CT:eq:fan-sizes}.  The denominator sector and
the chart radius are independent of $N$.

\begin{lemma}[Archived consequences of the retired CT--M4 proposal]
\label{CT:lem:Cauchy}
Assume the retired CT--M4 proposal \ref{CT:obl:material-fourth}.  Let
\begin{align*}
 F_4(\alpha)&=\mathcal R_4(B_\alpha),\\
 \Psi_4(\alpha)[Z]&=D\mathcal R_4(B_\alpha)[Z].
\end{align*}
At the real fan $\alpha$,
\begin{equation}\label{CT:eq:tail-majorants}
 |F_4|\le CS^5,\qquad
 \|\Psi_4\|_{(\Xjp_\alpha)^*}\le CS^{7/2}.
\end{equation}
If $h_i(t)=a+t d_i$ is uniformly comparable in the sense of
\eqref{CT:eq:uniformly-comparable} and
$E(t)=\sum_i h_i(t)^2d_i^2$, then
\begin{equation}\label{CT:eq:Cauchy-bounds}
 |\partial_t^2F_4(h(t))|\le \omega_4(S)E(t),
 \qquad
 \left\|\partial_t\{\Psi_4(h(t))\circ\mathcal T_t\}
     \right\|_{(\Xjp_\alpha)^*}
 \le \omega_{4,s}(S)E(t)^{1/2}.
\end{equation}
\end{lemma}

\begin{proof}
The two bounds in \eqref{CT:eq:tail-majorants} are
\eqref{CT:eq:M4-size}; the two bounds in \eqref{CT:eq:Cauchy-bounds} are
\eqref{CT:eq:M4-path}.  The identification of the varying dual spaces is the
fixed material pullback required in CT--M4, so no derivative of a moving
corner or hidden sum over material coordinates is introduced.
\end{proof}

The quartic action identities
\begin{equation}\label{CT:eq:action}
 \Jfour(\alpha)=12\int_0^1(1-t)E(t)\,dt,
 \qquad
 \int_0^1E(t)\,dt\le C\Jfour(\alpha)
\end{equation}
give the relative estimates as follows.  At the regular fan, cyclic
symmetry makes $DF_4(a\mathbf1)$ a constant material covector, so it
vanishes on $\sum_i d_i=0$.  The same symmetry makes
$\Psi_4(a\mathbf1)$ a constant support covector, which vanishes on the
area/translation quotient.  Taylor's formula and \eqref{CT:eq:Cauchy-bounds}
therefore give
\[
 |F_4(\alpha)-F_4(a\mathbf1)|
 \le \omega_4(S)\int_0^1(1-t)E(t)\,dt
 \le C\omega_4(S)\Jfour(\alpha).
\]
Apply the fundamental theorem of calculus to the fixed-fiber covector
$\Psi_4(h(t))\circ\mathcal T_t$.  The regular-fan value vanishes on the
transported quotient, and Cauchy--Schwarz gives
\[
 |\Psi_4(\alpha)[Z]|
 \le \omega_{4,s}(S)\left(\int_0^1E(t)\,dt\right)^{1/2}
       \|Z\|_{\Xjp_\alpha}
 \le C\omega_{4,s}(S)\Jfour(\alpha)^{1/2}\|Z\|_{\Xjp_\alpha}.
\]
These are the pure and mixed tail bounds in the uniformly comparable
regime.  The following elementary dichotomy verifies that this regime and
the absolute regime exhaust all positive fans.  The exact Bregman identity
is
\begin{equation}\label{CT:eq:J4-bregman-comparability}
 \Jfour(\alpha)=\sum_i d_i^2
  (\alpha_i^2+2a\alpha_i+3a^2).
\end{equation}
Let $m=\max_i\alpha_i$, so $S=a+m$, and put
$\rho=\Jfour(\alpha)^{1/2}/S^2$.  Since $m\ge a$ and the summand at an
index where $\alpha_i=m$ is at least $(m-a)^2m^2$,
\[
 \frac{m-a}{S}\le2\rho.
\]
If $\rho\le1/8$, then $a=(S-(m-a))/2\ge3S/8$; applying
\eqref{CT:eq:J4-bregman-comparability} at every index gives
\[
 \max_i\frac{|d_i|}{S}
 \le\frac{8}{3\sqrt3}\rho<2\rho,
 \qquad
 \max_i\frac{|d_i|}{a}\le\frac{16}{3}\rho.
\]
Fix $0<\beta<1$.  If $\Jfour<S^{5-\beta}$, then
$\rho<S^{(1-\beta)/2}$; for all sufficiently small $S$ the last bound is
at most $1/2$, and the whole segment is uniformly comparable by convexity.
In the complementary case $\Jfour\ge S^{5-\beta}$,
\eqref{CT:eq:tail-majorants} gives respectively the relative factors
$CS^\beta$ and $CS^{1+\beta/2}$.  Thus the two cases are exhaustive and
both are $o(1)$ multiples of the required defects.
Thus, under CT--M4, the quartic-and-higher Taylor tail has the uniform
material bounds required by CT.  This deduction uses no arbitrary-order
tensor estimate.

\subsection{Literal fan-to-bubble dictionary}

Choose normal angles $\theta_i$ with
$\theta_{i+1}-\theta_i=\alpha_i$ and set
\begin{equation}\label{CT:eq:cells}
 K_i=[\theta_i-\alpha_{i-1}/2,\theta_i+\alpha_i/2].
\end{equation}
On $K_i$, $x=\theta-\theta_i$, the exact normal fan input is
\begin{equation}\label{CT:eq:secant}
 f_\alpha(\theta)=h_\alpha\sec x-1,
 \qquad
 h_\alpha=\left(\frac{\pi}{\sum_j\tan(\alpha_j/2)}\right)^{1/2}.
\end{equation}
Define, on the same cell and with no change of nodes,
\begin{align}
 c_\alpha&=h_\alpha-1,\notag\\
 q_\alpha(\theta)&=\tfrac12x^2,\label{CT:eq:bubble}\\
 r_\alpha(\theta)&=(h_\alpha-1)\tfrac12x^2
 +h_\alpha(\sec x-1-\tfrac12x^2).\notag
\end{align}
Then the coordinate dictionary is the literal equality
\begin{equation}\label{CT:eq:dictionary}
 \boxed{f_\alpha=c_\alpha+q_\alpha+r_\alpha.}
\end{equation}
The function $q_\alpha$ is continuous at the midpoint vertices because
both adjacent values are $\alpha_i^2/8$, and
\begin{equation}\label{CT:eq:q-distribution}
 q_\alpha''=1-\sum_i\alpha_i
 \delta_{\theta_i+\alpha_i/2}.
\end{equation}
Taylor's formula on each normalized half-cell gives
\begin{equation}\label{CT:eq:r-bounds}
 \|r_\alpha\|_\infty\le CS^4,\qquad
 \|r_\alpha\|_{W^{1,\infty}}\le CS^3,\qquad
 \|r_\alpha\|_{H^{1/2}}^2\le CS^7.
\end{equation}
Constants vanish in \eqref{CT:eq:null}, including after polarization, because
$\int(\Dop g)(\Dop h)=\int g'h'$.  The tame estimate for the null form and
\eqref{CT:eq:r-bounds} place every term containing $r_\alpha$ in the
order-one/tail remainder already controlled above.  The normal component
of the physical support slot is exactly the trace $v_\alpha(z)$ used in
PT.  Therefore the two surviving principal rows of the third derivative
are precisely
\begin{equation}\label{CT:eq:terminal-rows}
 \kappa_{\mathbb D}\Cnull(q_\alpha,q_\alpha,q_\alpha),
 \qquad
 \kappa_{\mathbb D}\Cnull(q_\alpha,q_\alpha,v_\alpha(z)),
\end{equation}
with the same nodes, normalization, and projection as the quadratic-bubble
interface; the common fixed coefficient $\kappa_{\mathbb D}$ remains
outside both terminal estimates.

For downstream Taylor assembly one must use
$\mathcal T_3=(3!)^{-1}D^3\mathcal L(0)$.  Thus the corresponding pure
Taylor row and differentiated mixed Taylor row are
\begin{equation}\label{CT:eq:terminal-taylor-rows}
 \boxed{
 \frac{\kappa_{\mathbb D}}6
   \Cnull(q_\alpha,q_\alpha,q_\alpha),
 \qquad
 \frac{\kappa_{\mathbb D}}2
   \Cnull(q_\alpha,q_\alpha,v_\alpha(z)).}
\end{equation}
The second coefficient is $3(\kappa_{\mathbb D}/6)$, because
differentiating $\mathcal T_3(W,W,W)$ in one support slot produces three
equal placements.

\subsection{CT reduction theorem and audit}

\begin{theorem}[Exact CT core]
\label{CT:thm:CT}
Unconditionally, the fixed-disk third derivative is the sum of
$\kappa_{\mathbb D}$ times the null form \eqref{CT:eq:null} and the explicit
normal order-one ledger \eqref{CT:eq:R-ledger}.  The full
physical-coordinate tensor is \eqref{CT:eq:physical-complete-split}; its area
and translation rows vanish exactly.  The radial null-form outputs are
\eqref{CT:eq:terminal-rows}, with Taylor factors
\eqref{CT:eq:terminal-taylor-rows}; the nonzero curvature support symbol
\eqref{CT:eq:effective-leading-forcing} is retained separately and exported
to the JP--JC audit.  All unconditional stated constants are independent
of the number and minimum size of active cells.  The old material-tail
proposal is not used by the active DAG.
\end{theorem}

\begin{proof}
Combine Proposition~\ref{CT:prop:full-tensor}, Theorem~\ref{CT:thm:complete-split},
Lemma~\ref{CT:lem:R-relative}, Propositions
\ref{CT:prop:curvature-counterexample} and
\ref{CT:prop:joint-mixed-counterexample}, and \eqref{CT:eq:dictionary}.  The
None of these statements imports an estimate for either terminal null-form
row, and the active proof imports no premise from the archived calculation.
\end{proof}

\begin{center}
\begin{longtable}{>{\raggedright\arraybackslash}p{.30\linewidth}
                  >{\raggedright\arraybackslash}p{.58\linewidth}}
\toprule
audit item & disposition\\
\midrule
\endfirsthead
\toprule
audit item & disposition\\
\midrule
\endhead
state rows & all three families displayed in \eqref{CT:eq:full-tensor};\\
metric, mass, and scale & generated exhaustively by
\eqref{CT:eq:J-jets}--\eqref{CT:eq:H-jets};\\
area and translation projection & explicit coefficients
\eqref{CT:eq:chi}, then exact kernel cancellation;\\
curvature/tangential support & both the separate bound and CT--JM are
retracted; \eqref{CT:eq:effective-leading-forcing} is retained for JP--JC;\\
principal symbol & same-sign null form \eqref{CT:eq:null};\\
all lower rows & disjoint ledger \eqref{CT:eq:R-ledger} and
order-one estimate \eqref{CT:eq:R-estimate};\\
quartic pure/linear tail & old CT--M4 proposal and its consequences are
retained only as an archive; the active exact-support route uses neither;\\
bubble output & literal equality \eqref{CT:eq:dictionary}, with identical
cells and midpoint atoms.\\
\bottomrule
\end{longtable}
\end{center}
\endgroup

\begingroup
\section{Node QB: quadratic-bubble reduction}
\label{module:QB}
\noindent\textit{Audited source: \nolinkurl{QB_bubble_reduction.tex}.}
\subsection{Quadratic polar spline}

Let
\begin{equation}\label{QB:eq:fan}
 \alpha_i>0,\qquad \sum_i\alpha_i=2\pi,
 \qquad a=\frac{2\pi}{N},\qquad S=\max_i\alpha_i+a.
\end{equation}
Choose normal angles $\theta_i$ with
$\theta_{i+1}-\theta_i=\alpha_i$ and put
\begin{equation}\label{QB:eq:vertex}
 \varphi_i=\theta_i+\frac{\alpha_i}{2}.
\end{equation}
The polar cell of side $i$ is
\begin{equation}\label{QB:eq:cell}
 K_i=[\varphi_{i-1},\varphi_i]
 =[\theta_i-\alpha_{i-1}/2,\theta_i+\alpha_i/2].
\end{equation}
Define the continuous periodic quadratic spline
\begin{equation}\label{QB:eq:q}
 q_\alpha(\theta)=\frac12(\theta-\theta_i)^2,
 \qquad \theta\in K_i.
\end{equation}
Continuity follows because the two values at $\varphi_i$ are both
$\alpha_i^2/8$.

The exact radial fan graph is
\begin{equation}\label{QB:eq:f}
 f_\alpha(\theta)=h_\alpha\sec(\theta-\theta_i)-1,
 \qquad \theta\in K_i,
\end{equation}
with
$h_\alpha=(\pi/\sum_j\tan(\alpha_j/2))^{1/2}$.

\begin{lemma}[Secant/quadratic profile split]
\label{QB:lem:split}
There is a constant $c_\alpha$ and a continuous remainder $r_\alpha$ such
that
\begin{equation}\label{QB:eq:split}
 f_\alpha=c_\alpha+q_\alpha+r_\alpha,
\end{equation}
and, without a minimum-angle assumption,
\begin{equation}\label{QB:eq:r-size}
 \|r_\alpha\|_\infty\le CS^4,
 \qquad \|r_\alpha\|_{W^{1,\infty}}\le CS^3,
 \qquad \|r_\alpha\|_{H^{1/2}}^2\le CS^7.
\end{equation}
Also
\begin{equation}\label{QB:eq:q-size}
 \|q_\alpha\|_\infty\le CS^2,
 \qquad \|q_\alpha\|_{W^{1,\infty}}\le CS,
 \qquad \|q_\alpha\|_{H^{1/2}}^2\le CS^3.
\end{equation}
\end{lemma}

\begin{proof}
Take $c_\alpha=h_\alpha-1$ and write, on a cell,
\begin{equation*}
 r_\alpha(x)
 =(h_\alpha-1)\frac{x^2}{2}
 +h_\alpha\left(\sec x-1-\frac{x^2}{2}\right).
\end{equation*}
Since $h_\alpha-1=O(\sum_i\alpha_i^3)=O(S^2)$ and
$\sec x-1-x^2/2=O(x^4)$ through one normalized derivative, the first two
estimates in \eqref{QB:eq:r-size} follow.  Same and adjacent half-cells in the
Slobodeckij seminorm cost $C\alpha_i^8$.  On separated half-cells subtract
both means and use the first-separating dyadic annulus.  Thus
\begin{equation*}
 \|r_\alpha\|_{H^{1/2}}^2
 \le C\sum_i\alpha_i^8
 \le CS^4\sum_i\alpha_i^4\le CS^7.
\end{equation*}
The proof of \eqref{QB:eq:q-size} is identical with local amplitude
$\alpha_i^2$ and derivative amplitude $\alpha_i$.
\end{proof}

\subsection{The higher secant profile is harmless}

Recall the symmetric null form of stage 138,
\begin{multline}\label{QB:eq:C}
 \Cnull(g_1,g_2,g_3)\\
 =\frac13\sum_{\rm cyc}\int_\T g_1
 \{(\Dop g_2)(\Dop g_3)-g_2'g_3'\}\,d\theta.
\end{multline}
It has total order two and the tame estimate
\begin{equation}\label{QB:eq:C-tame}
 |\Cnull(g_1,g_2,g_3)|
 \le C\sum_{p<q}\|g_p\|_{H^{1/2}}\|g_q\|_{H^{1/2}}
 \|g_\ell\|_{W^{1,\infty}}.
\end{equation}

\begin{lemma}[Constants disappear]
\label{QB:lem:constant}
For every $g$ and every constant $c$,
\begin{equation}\label{QB:eq:constant}
 \Cnull(g+c,g+c,g+c)=\Cnull(g,g,g).
\end{equation}
The same is true after polarization whenever a constant occupies any one
slot.
\end{lemma}

\begin{proof}
Terms in which a differentiated slot is constant vanish.  If the
undifferentiated slot is constant, Parseval gives
$\|\Dop g\|_2^2=\|g'\|_2^2$.  Polarization proves the last assertion.
\end{proof}

\begin{theorem}[Secant-to-quadratic transfer]
\label{QB:thm:transfer}
There are moduli $\varepsilon_q(S),\varepsilon_{q,s}(S)\to0$ such that
\begin{align}
 &|\{\Cnull(f_\alpha^3)-\Cnull(q_\alpha^3)\}
 -\{\Cnull(f_{a\mathbf1}^3)-\Cnull(q_{a\mathbf1}^3)\}|
 \le\varepsilon_q(S)\Jfour(\alpha),
 \label{QB:eq:transfer-fan}\\
 &|\Cnull(f_\alpha,f_\alpha,v_\alpha(z))
 -\Cnull(q_\alpha,q_\alpha,v_\alpha(z))|
 \le\varepsilon_{q,s}(S)\Jfour(\alpha)^{1/2}
 \widetilde{\mathfrak G}_\alpha(z)^{1/2}.
 \label{QB:eq:transfer-mixed}
\end{align}
Here $v_\alpha(z)$ is the normal component of the physical support trace.
\end{theorem}

\begin{proof}
By Lemma~\ref{QB:lem:constant}, only $r_\alpha$ occurs in the difference.
Equations \eqref{QB:eq:C-tame}--\eqref{QB:eq:q-size} give the absolute estimates
\begin{equation}\label{QB:eq:transfer-absolute}
 |\Cnull(f_\alpha^3)-\Cnull(q_\alpha^3)|\le CS^6,
\end{equation}
and
\begin{equation}\label{QB:eq:transfer-mixed-absolute}
 |\Cnull(f_\alpha,f_\alpha,v)-\Cnull(q_\alpha,q_\alpha,v)|
 \le CS^{9/2}\|v\|_{H^{1/2}}.
\end{equation}
For example, an energetic $(q,r)$ pair and a coefficient $q$ cost
$S^{3/2}S^{7/2}S=S^6$; the other placements are no larger.

Use the far threshold $\Jfour\ge S^{6-\beta}$, $0<\beta<1$.
Equations \eqref{QB:eq:transfer-absolute}--
\eqref{QB:eq:transfer-mixed-absolute} give an $o(1)$ multiple of the required
right sides.  In the complementary regime,
\begin{equation*}
 \|\alpha-a\mathbf1\|_\infty/S
 \le CS^{(2-\beta)/4},
\end{equation*}
so the whole segment is in a fixed complex material polydisc.  Pull every
half-cell to $[0,1]$ before differentiating.  Cauchy's estimate applied to
the two absolute majorants gives
\begin{equation*}
 |F_q''(t)|\le CS^2E(t),
 \qquad \|\Psi_q'(t)\|\le CS^{5/2}E(t)^{1/2}.
\end{equation*}
Cyclic symmetry kills the regular first fan derivative and the regular
support functional on the exact quotient.  Integrate with the weighted and
unweighted quartic actions.  Physical trace equivalence replaces
$\|v_\alpha(z)\|_{H^{1/2}}$ by
$C\widetilde{\mathfrak G}_\alpha(z)^{1/2}$.
\end{proof}

\subsection{Exact midpoint-atom formula}

Let
\begin{equation}\label{QB:eq:mu}
 \mu_\alpha=\sum_i\alpha_i\delta_{\varphi_i}.
\end{equation}

\begin{lemma}[Distributional spline identity]
\label{QB:lem:distribution}
One has
\begin{equation}\label{QB:eq:distribution}
 q_\alpha''=1-\mu_\alpha
\end{equation}
on $\T$.  With Fourier convention
$\widehat g(n)=(2\pi)^{-1}\int_\T g(\theta)e^{-in\theta}\,d\theta$, put
\begin{equation}\label{QB:eq:M}
 M_\alpha(n)=\sum_i\alpha_i e^{-in\varphi_i}.
\end{equation}
Then, for $n\ne0$,
\begin{equation}\label{QB:eq:qhat}
 \widehat q_\alpha(n)=\frac{M_\alpha(n)}{2\pi n^2}.
\end{equation}
Moreover the midpoint cancellation is the exact identity
\begin{equation}\label{QB:eq:midpoint-cancellation}
 M_\alpha(n)=\sum_i\alpha_i e^{-in\varphi_i}
 \left\{1-\operatorname{sinc}\left(\frac{n\alpha_i}{2}\right)\right\},
 \qquad n\ne0.
\end{equation}
\end{lemma}

\begin{proof}
Inside every open cell, $q_\alpha''=1$.  At the vertex $\varphi_i$ the
one-sided derivatives are $\alpha_i/2$ and $-\alpha_i/2$, so the derivative
jump is $-\alpha_i$.  This proves \eqref{QB:eq:distribution}; Fourier
transformation gives \eqref{QB:eq:qhat}.

The normal intervals
$I_i=[\theta_i,\theta_{i+1}]$ partition the circle and have midpoint
$\varphi_i$.  Hence, for $n\ne0$,
\begin{equation*}
 0=\int_\T e^{-in\theta}\,d\theta
 =\sum_i\alpha_i e^{-in\varphi_i}
 \operatorname{sinc}\left(\frac{n\alpha_i}{2}\right).
\end{equation*}
Subtract this equality from \eqref{QB:eq:M}.
\end{proof}

\begin{theorem}[Exact discrete cubic scalar]
\label{QB:thm:discrete}
The pure quadratic-bubble null form is
\begin{equation}\label{QB:eq:discrete}
 \boxed{
 \Cnull(q_\alpha,q_\alpha,q_\alpha)
 =\frac1{\pi^2}\operatorname{Re}
 \sum_{p,q\ge1}
 \frac{M_\alpha(p)M_\alpha(q)
 \overline{M_\alpha(p+q)}}{pq(p+q)^2}.}
\end{equation}
The series is interpreted by symmetric Abel limits and is absolutely
convergent after the midpoint cancellation \eqref{QB:eq:midpoint-cancellation}
is inserted sourcewise.
\end{theorem}

\begin{proof}
For a real $g$, Fourier expansion of \eqref{QB:eq:C} shows that only triples
with two frequencies of the same sign survive.  If $p,q>0$ and the third
frequency is $-(p+q)$, the unsymmetrized coefficient is $2pq$.
The positive and negative triples are conjugate, so
\begin{equation*}
 \Cnull(g,g,g)=8\pi\operatorname{Re}
 \sum_{p,q\ge1}pq\,\widehat g(p)\widehat g(q)
 \overline{\widehat g(p+q)}.
\end{equation*}
Insert \eqref{QB:eq:qhat}; since
$8\pi/(2\pi)^3=1/\pi^2$, this gives \eqref{QB:eq:discrete}.
Abel regularization $r^{p+q}$ justifies rearrangement.  Formula
\eqref{QB:eq:midpoint-cancellation} gives two powers of $n\alpha_i$ at low
frequency and the trivial bound at high frequency; dyadic source
allocation then makes the Abel limit independent of the order of
summation.
\end{proof}

\subsection{Final discrete obligations}

\begin{obligation}[DF: midpoint cubic Bregman estimate]
\label{QB:obl:DF}
For the explicit scalar on the right of \eqref{QB:eq:discrete}, prove
\begin{equation}\label{QB:eq:DF}
 |\Cnull(q_\alpha^3)-\Cnull(q_{a\mathbf1}^3)|
 \le CS\Jfour(\alpha).
\end{equation}
Every cumulative phase produced by differentiating $\varphi_i$ must first
be summed cyclically; the midpoint factor in
\eqref{QB:eq:midpoint-cancellation} is not to be differentiated row by row.
\end{obligation}

\begin{obligation}[DM: polarized midpoint/support estimate]
\label{QB:obl:DM}
For the polarization of \eqref{QB:eq:discrete} with one slot replaced by the
physical normal trace $v_\alpha(z)$, prove
\begin{equation}\label{QB:eq:DM}
 |\Cnull(q_\alpha,q_\alpha,v_\alpha(z))|
 \le\varepsilon_{\rm DM}(S)\Jfour(\alpha)^{1/2}
 \widetilde{\mathfrak G}_\alpha(z)^{1/2},
 \qquad \varepsilon_{\rm DM}(S)\to0.
\end{equation}
\end{obligation}

\begin{proposition}[Equivalence]
\label{QB:prop:equivalence}
DF and the secant-to-quadratic transfer imply the pure principal row NF.
DF, DM, and the same transfer imply the mixed principal row NM.  These are
exactly the two terminal principal inputs exported by QB; the conversion
to CF, CM, NR, and J0 is performed only in the current
\texttt{NR\_assembly.tex} node.
\end{proposition}

\begin{proof}
Apply Theorem~\ref{QB:thm:transfer} to the pure and mixed null forms and absorb
its moduli into the DF and DM moduli.  No downstream tensor or value
statement is imported here.
\end{proof}
\endgroup


\part{The pure midpoint row}

\begingroup
\section{Foundation F144: Hilbert stress and the exact midpoint identity}
\label{module:F144}
\noindent\textit{Audited source: \nolinkurl{convex_largeN_polya_szego_stage144_NR_Hilbert_midpoint_quadrature_reduction.tex}.}
\subsection{A second exact form of the cubic null structure}

Use the periodic Hilbert transform convention
\begin{equation}\label{F144:eq:H}
 \widehat{\Hil f}(n)=-i\operatorname {sgn}(n)\widehat f(n),
 \qquad |D|=\Hil\partial_\theta.
\end{equation}
The symmetric null form is the one of stage 138; on the diagonal,
\begin{equation}\label{F144:eq:C}
 \Cnull(f,f,f)=\int_\T f\{(|D|f)^2-(f')^2\}.
\end{equation}

\begin{lemma}[Hilbert stress identity]
\label{F144:lem:stress}
For every smooth real periodic $f$,
\begin{equation}\label{F144:eq:stress}
 \boxed{\Cnull(f,f,f)=\int_\T f''(\Hil f)^2.}
\end{equation}
The identity extends by Abel regularization whenever the right side is a
well-defined distributional pairing.
\end{lemma}

\begin{proof}
For a real function $g$ the Hilbert product identity is
\begin{equation}\label{F144:eq:product}
 (\Hil g)^2-g^2=\Hil(2g\Hil g).
\end{equation}
Indeed it follows first on trigonometric polynomials from
$\Hil^2=-I$ and then by density.  Put $g=f'$.  By skew-adjointness of
$\Hil$ and commutation with $\partial_\theta$,
\begin{align*}
 \Cnull(f,f,f)
 &=\int f\Hil(2f'\Hil f')
 =-2\int (\Hil f)f'(\Hil f)'\\
 &=-\int f'\{(\Hil f)^2\}'
 =\int f''(\Hil f)^2.
\end{align*}
Abel regularization proves the distributional extension.
\end{proof}

Polarization of \eqref{F144:eq:stress} also gives, for real $f,g,v$,
\begin{equation}\label{F144:eq:polarized-stress}
 \Cnull(f,g,v)=\frac13\int_\T
 \{f''\Hil g\Hil v+g''\Hil f\Hil v+v''\Hil f\Hil g\}.
\end{equation}
In particular,
\begin{equation}\label{F144:eq:q-q-v}
 3\Cnull(q,q,v)
 =2\int q''\Hil q\Hil v+\int v''(\Hil q)^2.
\end{equation}
This will be the smaller starting point for the mixed row DM.

\subsection{The cubic scalar is a midpoint quadrature error}

Let $I_i=[\theta_i,\theta_{i+1}]$ have length $h_i>0$ and midpoint
$c_i$.  On $I_i$ put
\begin{equation}\label{F144:eq:q}
 q_\alpha(\theta_i+x)=\frac12\min\{x^2,(h_i-x)^2\}.
\end{equation}
Then, as in stage 139,
\begin{equation}\label{F144:eq:qdd}
 q_\alpha''=1-\mu_\alpha,
 \qquad \mu_\alpha=\sum_i h_i\delta_{c_i}.
\end{equation}

\begin{theorem}[Exact Hilbert midpoint formula]
\label{F144:thm:midpoint}
Put $w_\alpha=\Hil q_\alpha$.  Then
\begin{equation}\label{F144:eq:midpoint}
 \boxed{
 \Cnull(q_\alpha,q_\alpha,q_\alpha)
 =\int_\T(q_\alpha-\overline q_\alpha)^2
 -\sum_i h_i|w_\alpha(c_i)|^2.}
\end{equation}
For the regular fan $a\mathbf1$, every sampled value is zero and
\begin{equation}\label{F144:eq:regular}
 \Cnull(q_{a\mathbf1}^3)=\frac{Na^5}{720}
 =\frac{\pi a^4}{360}.
\end{equation}
\end{theorem}

\begin{proof}
Insert \eqref{F144:eq:qdd} in \eqref{F144:eq:stress}.  The Hilbert transform is an
$L^2$ isometry on mean-zero functions, so
\begin{equation*}
 \int w_\alpha^2=\int(q_\alpha-\overline q_\alpha)^2.
\end{equation*}
This proves \eqref{F144:eq:midpoint}.  At the regular fan, reflection symmetry
about each midpoint makes $q$ even and $\Hil q$ odd, hence
$w(c_i)=0$.  Direct integration on one cell gives
\begin{align}
 \int_{I_i}q_\alpha&=\frac{h_i^3}{24},&
 \int_{I_i}q_\alpha^2&=\frac{h_i^5}{320}.
 \label{F144:eq:cell-integrals}
\end{align}
Using $Na=2\pi$ in the resulting variance gives \eqref{F144:eq:regular}.
\end{proof}

Thus the apparently cubic two-frequency obstruction is the difference of
one explicit local scalar and one nonnegative one-frequency sampling row.
\endgroup

\begingroup
\section{Foundation F148: the exact cumulative-free material matrix}
\label{module:F148}
\noindent\textit{Audited source: \nolinkurl{convex_largeN_polya_szego_stage148_NR_quadrature_primitive_material_matrix.tex}.}
\subsection{Clausen material derivative as a length gradient}

Let $I_j=[\theta_j,\theta_{j+1}]$ be a cyclic interval partition of
$\T$, with length $h_j$ and midpoint $c_j$.  Put
\begin{equation}\label{F148:eq:clausen}
 \Phi(x)=\sum_{n\ge1}\frac{\sin(nx)}{n^2},\qquad
 K(x)=\Phi'(x)=-\log\left|2\sin\frac x2\right|.
\end{equation}
The sampled conjugate of the quadratic fan spline is
\begin{equation}\label{F148:eq:w}
 w_i=\frac1\pi\sum_jh_j\Phi(c_i-c_j).
\end{equation}
Along a positive affine path, let $h_j'=d_j$, $\sum_jd_j=0$, and let
$s_j=c_j'$.  Stage 144 gives
\begin{equation}\label{F148:eq:material-old}
 W_i:=\dot w_i=\frac1\pi\sum_{j\ne i}
 \{d_j\Phi(c_i-c_j)+h_j(s_i-s_j)K(c_i-c_j)\}.
\end{equation}

Fix the target $i$.  Relabel cyclically and choose a partition endpoint
$\theta_0$ as a cut; write $\theta_0+2\pi$ for its second copy.  The cut
may be chosen at circular distance at least $\pi/3$ from $c_i$.  This
target-dependent relabelling changes neither \eqref{F148:eq:w} nor
\eqref{F148:eq:material-old}.

For $k<i$ define
\begin{align}
 \mathscr E_{ik}^-&=
 \sum_{j<k}h_jK(c_i-c_j)+\frac{h_k}{2}K(c_i-c_k)
 -\int_{\theta_0}^{c_k}K(c_i-y)\,dy. \label{F148:eq:Eminus}
\end{align}
For $k>i$ define
\begin{align}
 \mathscr E_{ik}^+&=
 \int_{c_k}^{\theta_0+2\pi}K(c_i-y)\,dy
 -\frac{h_k}{2}K(c_i-c_k)-\sum_{j>k}h_jK(c_i-c_j).
 \label{F148:eq:Eplus}
\end{align}
The two arcs in these definitions do not contain $c_i$ in their
interiors.  Finally put
\begin{align}
 E_i^-&=\sum_{j<i}h_jK(c_i-c_j)
       -\int_{\theta_0}^{\theta_i}K(c_i-y)\,dy,\notag\\
 E_i^+&=\sum_{j>i}h_jK(c_i-c_j)
       -\int_{\theta_{i+1}}^{\theta_0+2\pi}K(c_i-y)\,dy,
 \label{F148:eq:E-self}
\end{align}
and define the matrix
\begin{equation}\label{F148:eq:matrix}
 \mathscr E_{ik}=\begin{cases}
 \mathscr E_{ik}^-,&k<i,\\
 (E_i^--E_i^+)/2,&k=i,\\
 \mathscr E_{ik}^+,&k>i.
 \end{cases}
\end{equation}

\begin{theorem}[Exact cumulative-free material matrix]
\label{F148:thm:matrix}
For every zero-sum length velocity $d$,
\begin{equation}\label{F148:eq:W-matrix}
 \boxed{W_i=\frac1\pi\sum_k d_k\mathscr E_{ik}.}
\end{equation}
Thus no node velocity $s_j$, root velocity, or cumulative sum of the
$d_j$ occurs in the final coefficient matrix.
\end{theorem}

\begin{proof}
Temporarily keep the cut fixed.  Its length coordinates give
\begin{equation*}
 \frac{\partial c_j}{\partial h_k}=
 \begin{cases}0,&j<k,\\1/2,&j=k,\\1,&j>k.
 \end{cases}
\end{equation*}
If $M_{ik}=\pi\,\partial w_i/\partial h_k$, differentiation of
\eqref{F148:eq:w} gives
\begin{equation}\label{F148:eq:M}
 M_{ik}=\Phi(c_i-c_k)+\sum_jh_j
 \left(\frac{\partial c_i}{\partial h_k}
       -\frac{\partial c_j}{\partial h_k}\right)K(c_i-c_j),
\end{equation}
where the self coefficient of $K(0)$ is zero before the Abel limit.

For $k<i$, equation \eqref{F148:eq:M} becomes
\begin{equation*}
 M_{ik}=\Phi(c_i-c_k)+\sum_{j<k}h_jK(c_i-c_j)
                    +\frac{h_k}{2}K(c_i-c_k).
\end{equation*}
The fundamental theorem of calculus gives exactly
\begin{equation*}
 \int_{\theta_0}^{c_k}K(c_i-y)\,dy
 =\Phi(c_i-\theta_0)-\Phi(c_i-c_k).
\end{equation*}
After rotating the origin to the cut, the first term on the right is
$\Phi(c_i)$.  Hence
$M_{ik}=\Phi(c_i)+\mathscr E_{ik}^-$.  The same calculation on the
right arc gives
$M_{ik}=\Phi(c_i)+\mathscr E_{ik}^+$ for $k>i$.

For $k=i$, \eqref{F148:eq:M} reads
\begin{equation*}
 M_{ii}=\frac12\sum_{j<i}h_jK(c_i-c_j)
       -\frac12\sum_{j>i}h_jK(c_i-c_j).
\end{equation*}
The exact integrals over the two sides of $I_i$ give
$M_{ii}=\Phi(c_i)+(E_i^--E_i^+)/2$.  Therefore
\begin{equation*}
 M_{ik}=\Phi(c_i)+\mathscr E_{ik}\qquad\hbox{for every }k.
\end{equation*}
Since $\sum_kd_k=0$, the common Clausen term disappears and
\eqref{F148:eq:W-matrix} follows.  This also proves that the result is
independent of the chosen cut.
\end{proof}

\subsection{The matrix is a cotangent integral of a cellwise sawtooth}

On each cell define
\begin{equation}\label{F148:eq:eta}
 \eta(y)=\begin{cases}
 -(y-\theta_j),&\theta_j<y<c_j,\\
 \theta_{j+1}-y,&c_j<y<\theta_{j+1}.
 \end{cases}
\end{equation}
Then, distributionally on $I_j$,
\begin{equation}\label{F148:eq:eta-properties}
 d\eta=h_j\delta_{c_j}-\mathbf1_{I_j}\,dy,\qquad
 |\eta|\le h_j/2,\qquad
 \int_{I_j}\eta=0.
\end{equation}
Moreover
\begin{equation}\label{F148:eq:Kprime}
 K'(x)=-\frac12\cot\frac x2=-\pi\kappa(x),
 \qquad \kappa(x)=\frac1{2\pi}\cot\frac x2.
\end{equation}

\begin{proposition}[Exact sawtooth representation]
\label{F148:prop:sawtooth}
For $k<i$ and $k>i$, respectively,
\begin{align}
 \mathscr E_{ik}^-&=
 \int_{\theta_0}^{c_k}\eta(y)K'(c_i-y)\,dy,
 \label{F148:eq:saw-left}\\
 \mathscr E_{ik}^+&=-
 \int_{c_k}^{\theta_0+2\pi}\eta(y)K'(c_i-y)\,dy.
 \label{F148:eq:saw-right}
\end{align}
The terminal half-atom at $c_k$ closes the cumulative mass to zero.
Likewise $E_i^-$ and $E_i^+$ are the corresponding full-cell integrals
on the two sides of $I_i$.  Hence every matrix entry is assembled from
the same cellwise mean-zero sawtooth and one cotangent frequency.
\end{proposition}

\begin{proof}
On the left arc, the signed midpoint-quadrature measure consists of all
full atoms before $I_k$, half of the atom at $c_k$, and minus Lebesgue
measure on the arc.  Its cumulative distribution is exactly $\eta$ and
vanishes at both endpoints.  Stieltjes integration by parts therefore
gives
\begin{equation*}
 \int K(c_i-y)\,d\eta(y)
 =-\int\eta(y)\frac d{dy}K(c_i-y)\,dy
 =\int\eta(y)K'(c_i-y)\,dy,
\end{equation*}
which is \eqref{F148:eq:saw-left}.  The right arc has the opposite orientation,
giving \eqref{F148:eq:saw-right}.  The full-side statements are identical.
\end{proof}

\begin{remark}[Both former cancellations are now identities]
The endpoint vanishing in stage 147 is encoded by the fact that $\eta$
is zero at every cell endpoint.  The cumulative-velocity cancellation is
encoded by $\int_{I_j}\eta=0$ and the disappearance of the common
$\Phi(c_i)$ in Theorem~\ref{F148:thm:matrix}.  Thus neither cancellation needs
to be recovered after a triangle inequality.
\end{remark}
\endgroup

\begingroup
\section{Node DF0A: intrinsic variance and the far stratum}
\label{module:DF0A}
\noindent\textit{Audited source: \nolinkurl{DF0_variance.tex}.}
\subsection{The intrinsic equality-stratum variance}

Let
\begin{equation}\label{DF0A:eq:data}
 h_i>0,\qquad \sum_i h_i=2\pi=:L,
 \qquad a=\frac{2\pi}{N},\qquad
 S=\max_i h_i+a.
\end{equation}
The quadratic cell bubble and its sampled Hilbert conjugate are
\begin{align}
 q_\alpha(\theta_i+x)
 &=\frac12\min\{x^2,(h_i-x)^2\},\qquad 0\le x\le h_i,
 \label{DF0A:eq:q}\\
 T(\alpha)&=\sum_i h_i|\Hil q_\alpha(c_i)|^2,
 \qquad c_i=\theta_i+h_i/2.\label{DF0A:eq:T}
\end{align}
Put $A_k=\sum_i h_i^k$ and define
\begin{equation}\label{DF0A:eq:D}
 \ell^2:=\frac{A_3}{L},
 \qquad
 \boxed{\mathfrak D(\alpha)
 :=\sum_i h_i(h_i^2-\ell^2)^2
 =A_5-\frac{A_3^2}{L}.}
\end{equation}

\begin{lemma}[Pair-variance formula and equality strata]
\label{DF0A:lem:pair}
One has
\begin{equation}\label{DF0A:eq:pair}
 \mathfrak D(\alpha)
 =\frac1{2L}\sum_{i,j}h_ih_j(h_i^2-h_j^2)^2.
\end{equation}
Consequently $\mathfrak D\ge0$, and it vanishes exactly when all
positive cell lengths are equal.  After zero-cell deletion, these are
precisely the collapsed-uniform strata.
\end{lemma}

\begin{proof}
Expand the double sum and use $\sum_i h_i=L$ and
$\sum_i h_i^3=A_3$.  Equality in a weighted variance holds precisely
when $h_i^2$ is constant on its positive support.
\end{proof}

Recall the quartic Bregman charge
\begin{equation}\label{DF0A:eq:J}
 \Jfour(\alpha)=\sum_i(h_i-a)^2(h_i^2+2ah_i+3a^2).
\end{equation}

\begin{proposition}[The intrinsic variance is paid by the DAG charge]
\label{DF0A:prop:D-J}
For every positive fan,
\begin{equation}\label{DF0A:eq:D-J}
 \boxed{\mathfrak D(\alpha)\le S\Jfour(\alpha).}
\end{equation}
The estimate is independent of the minimum cell and all adjacent ratios.
\end{proposition}

\begin{proof}
The number $\ell^2=A_3/L$ is the weighted least-squares minimizer of
$m\mapsto\sum_i h_i(h_i^2-m)^2$.  Hence, using the competitor $m=a^2$,
\begin{align*}
 \mathfrak D(\alpha)
 &\le\sum_i h_i(h_i^2-a^2)^2\\
 &=\sum_i h_i(h_i-a)^2(h_i+a)^2\\
 &\le S\sum_i(h_i-a)^2(h_i^2+2ah_i+3a^2)
 =S\Jfour(\alpha).
\end{align*}
\end{proof}

\subsection{Exact relation with the cubic null form}

Stage 144 proves
\begin{align}
 \Cnull(q_\alpha^3)&=V(\alpha)-T(\alpha),\label{DF0A:eq:C-VT}\\
 V(\alpha)&=\frac{A_5}{320}-\frac{A_3^2}{1152\pi}.
 \label{DF0A:eq:V}
\end{align}

\begin{proposition}[Intrinsic defect identity]
\label{DF0A:prop:defect}
For every positive fan,
\begin{equation}\label{DF0A:eq:defect}
 \boxed{
 \Cnull(q_\alpha^3)-\frac1{720}\sum_i h_i^5
 =\frac1{576}\mathfrak D(\alpha)-T(\alpha).}
\end{equation}
In particular the sharp inequality
\begin{equation}\label{DF0A:eq:sharp-candidate}
 \Cnull(q_\alpha^3)\ge\frac1{720}\sum_i h_i^5
\end{equation}
is equivalent to $T\le\mathfrak D/576$.
\end{proposition}

\begin{proof}
Since $L=2\pi$,
\[
 \frac1{320}-\frac1{720}=\frac1{576},
 \qquad
 \frac{A_3^2}{1152\pi}=\frac1{576}\frac{A_3^2}{L}.
\]
Insert these two equalities in \eqref{DF0A:eq:C-VT}--\eqref{DF0A:eq:V}.
\end{proof}

\begin{remark}[Logical status of the sharp constant]
Equation \eqref{DF0A:eq:defect} is proved; inequality
\eqref{DF0A:eq:sharp-candidate} is not used as an established fact.  The
unconditional proof only needs an unspecified universal constant in the
following smaller endpoint leaf.
\end{remark}

\begin{obligation}[AVS: adaptive equality-stratum variance sampling]
\label{DF0A:obl:AVS}
Prove
\begin{equation}\label{DF0A:eq:AVS}
 \boxed{T(\alpha)\le C\mathfrak D(\alpha)}
\end{equation}
for every positive fine fan.  By Proposition~\ref{DF0A:prop:D-J}, AVS implies
HQ and hence DF.  Unlike the deleted HMS/QPS route, AVS vanishes on every
collapsed-uniform endpoint stratum and contains no path velocity.
\end{obligation}

\subsection{A global absolute sampling bound}

We use the following mesh-free sampling theorem in the precise special
case needed here.

\begin{lemma}[Adaptive truncated-Hilbert sampling]
\label{DF0A:lem:sampling}
Let $\{I_i\}$ be any interval partition of the circle, with lengths
$h_i$ and midpoints $c_i$.  If $f\in L^2(\T)$ satisfies
\begin{equation}\label{DF0A:eq:self-zero}
 \operatorname {pv}\int_{I_i}\frac1{2\pi}
 \cot\frac{c_i-y}{2}\,f(y)\,dy=0
 \qquad\text{for every }i,
\end{equation}
then
\begin{equation}\label{DF0A:eq:sampling}
 \sum_i h_i|\Hil f(c_i)|^2\le C\|f\|_2^2.
\end{equation}
The constant is independent of all interval ratios.
\end{lemma}

\begin{proof}
The left side of \eqref{DF0A:eq:self-zero} is the part removed by symmetric
truncation at radius $h_i/2$.  For every $x\in I_i$, the
scale-restricted Cotlar inequality gives
\[
 |\Hil f(c_i)|\le C\{M(\Hil f)(x)+Mf(x)\}.
\]
Indeed, split at radius $h_i$; moving the centre from $c_i$ to $x$
changes the far part by a Hardy--Littlewood average, and the intervening
annulus is another such average.  Square, integrate on $I_i$, sum the
disjoint cells, and use the $L^2$ boundedness of $M$ and $\Hil$.
\end{proof}

\begin{theorem}[Unconditional absolute endpoint bound]
\label{DF0A:thm:absolute}
For every positive fan,
\begin{equation}\label{DF0A:eq:absolute}
 \boxed{T(\alpha)\le C A_5=C\sum_i h_i^5.}
\end{equation}
\end{theorem}

\begin{proof}
On its own cell, $q_\alpha$ is even about $c_i$, while the Hilbert kernel
is odd.  Thus \eqref{DF0A:eq:self-zero} holds with $f=q_\alpha$.  Moreover
direct integration of the two half-bubbles gives
\[
 \|q_\alpha\|_2^2=\sum_i\frac{h_i^5}{320}.
\]
Apply Lemma~\ref{DF0A:lem:sampling}.
\end{proof}

\subsection{All fans away from the equality strata are closed}

\begin{theorem}[Far-stratum AVS and HQ]
\label{DF0A:thm:far}
Fix $0<\delta<1$.  If
\begin{equation}\label{DF0A:eq:far}
 \mathfrak D(\alpha)\ge\delta A_5,
\end{equation}
then
\begin{equation}\label{DF0A:eq:far-conclusion}
 T(\alpha)\le C_\delta\mathfrak D(\alpha)
 \le C_\delta S\Jfour(\alpha).
\end{equation}
Thus AVS, HQ, and the pure cubic row DF are unconditional on the complete
far-stratum region, including arbitrary minimum cells and ratios.
\end{theorem}

\begin{proof}
Theorem~\ref{DF0A:thm:absolute} and \eqref{DF0A:eq:far} give the first inequality;
Proposition~\ref{DF0A:prop:D-J} gives the second.  Stage 144 then gives
HQ $\Rightarrow$ DF.
\end{proof}

\subsection{Quantitative selection of the nearest collapsed-uniform stratum}

It remains to treat $\mathfrak D<\delta A_5$.  This condition has a
strong geometric consequence which we now make exact.

\begin{lemma}[Concentration of physical angular mass]
\label{DF0A:lem:concentration}
Assume $0<\delta<1$ and
$\mathfrak D<\delta A_5$.  Then
\begin{equation}\label{DF0A:eq:D-small}
 \mathfrak D\le\frac{\delta}{1-\delta}L\ell^4.
\end{equation}
For $0<\rho<1$ let
\begin{equation}\label{DF0A:eq:bad}
 \mathcal B_\rho=\{i:|h_i/\ell-1|\ge\rho\},
 \qquad B_\rho=\sum_{i\in\mathcal B_\rho}h_i.
\end{equation}
Then
\begin{equation}\label{DF0A:eq:bad-mass}
 \boxed{
 \frac{B_\rho}{L}
 \le\frac{\delta}{(1-\delta)\rho^2}.}
\end{equation}
\end{lemma}

\begin{proof}
Equation \eqref{DF0A:eq:D} gives $A_5=\mathfrak D+L\ell^4$;
this proves \eqref{DF0A:eq:D-small}.  On a bad cell,
$|h_i^2-\ell^2|\ge\rho\ell^2$.  Hence
$\rho^2\ell^4B_\rho\le\mathfrak D$, and
\eqref{DF0A:eq:bad-mass} follows.
\end{proof}

Let $\mathcal G_\rho$ be the complementary good set and
$M=|\mathcal G_\rho|$.  If
$\delta<(1-\delta)\rho^2$, Lemma~\ref{DF0A:lem:concentration} ensures
$M\ge1$.  Define the collapsed-uniform reference fan $\beta$ by
\begin{equation}\label{DF0A:eq:beta}
 \beta_i=\begin{cases}b:=L/M,&i\in\mathcal G_\rho,\\0,&i\in\mathcal B_\rho.
 \end{cases}
\end{equation}
After deleting its zero cells, $\beta$ is the regular $M$-fan; therefore
all its Hilbert midpoint samples vanish.

\begin{lemma}[Quantitative distance to the selected stratum]
\label{DF0A:lem:reference}
Put $\varepsilon=B_\rho/L$.  Then
\begin{equation}\label{DF0A:eq:b-ell}
 \frac{1-\rho}{1-\varepsilon}
 \le\frac b\ell\le
 \frac{1+\rho}{1-\varepsilon}.
\end{equation}
If $\rho+\varepsilon\le1/4$, then
\begin{equation}\label{DF0A:eq:reference-charge}
 \boxed{
 b^3\sum_{i\in\mathcal G_\rho}(h_i-b)^2
 +\sum_{i\in\mathcal B_\rho}h_i(h_i^2-\ell^2)^2
 \le C_\rho\mathfrak D(\alpha).}
\end{equation}
\end{lemma}

\begin{proof}
Since every good length lies between $(1-\rho)\ell$ and
$(1+\rho)\ell$,
\[
 M(1-\rho)\ell\le L-B_\rho
 \le M(1+\rho)\ell.
\]
Using $b=L/M$ gives \eqref{DF0A:eq:b-ell}.

Write $e_i=h_i-\ell$ on the good set.  There
\[
 h_i(h_i^2-\ell^2)^2\ge c_\rho\ell^3e_i^2.
\]
Also
\[
 M(b-\ell)=\sum_{i\in\mathcal G_\rho}e_i+B_\rho.
\]
Therefore Cauchy--Schwarz, $M\asymp L/\ell$, and
$\rho^2\ell^4B_\rho\le\mathfrak D$ give
\begin{align*}
 b^3\sum_{\mathcal G_\rho}(h_i-b)^2
 &\le C_\rho\ell^3\sum_{\mathcal G_\rho}e_i^2
   +C_\rho\frac{\ell^3B_\rho^2}{M}\\
 &\le C_\rho\mathfrak D
   +C_\rho\ell^4\frac{B_\rho^2}{L}
 \le C_\rho\mathfrak D.
\end{align*}
The bad term in \eqref{DF0A:eq:reference-charge} is a sub-sum of
$\mathfrak D$.
\end{proof}

The preceding lemmas reduce the near-stratum region to the following
strictly localized endpoint estimate.

\begin{obligation}[ELCA: equality-stratum first-separation square]
\label{DF0A:obl:ELCA}
For the good/bad decomposition and collapsed-uniform reference
\eqref{DF0A:eq:beta}, prove directly from the endpoint kernel-difference
identity of stage 152 that
\begin{multline}\label{DF0A:eq:ELCA}
 T(\alpha)\\
 \le C_\rho\left\{
 b^3\sum_{i\in\mathcal G_\rho}(h_i-b)^2
 +\sum_{i\in\mathcal B_\rho}h_i(h_i^2-\ell^2)^2
 \right\}.
\end{multline}
The regular $M$-fan is used only after zero-cell deletion.  A maximal
consecutive bad block is kept intact until its first separating scale;
the two endpoint kernels may not be estimated separately inside that
block.  Good--good rectangles use the signed variable-grid discrete
Calder\'on--Zygmund square, while every rectangle meeting a bad block is
paid by the second term in \eqref{DF0A:eq:ELCA}.  Descendants are squared
before ancestors are summed.
\end{obligation}

\begin{proposition}[ELCA is the only remaining pure endpoint row]
\label{DF0A:prop:ELCA}
ELCA implies AVS in the near-stratum region.  Together with
Theorem~\ref{DF0A:thm:far}, it implies AVS, HQ, and DF for every positive fan.
\end{proposition}

\begin{proof}
Lemma~\ref{DF0A:lem:reference} bounds the right side of \eqref{DF0A:eq:ELCA} by
$C_\rho\mathfrak D$.  This is AVS.  The far and near regions exhaust all
fans.  Proposition~\ref{DF0A:prop:D-J} and stage 144 then give HQ and DF.
\end{proof}
\endgroup

\begingroup
\section{Node DF0B: the comparable-grid discrete square}
\label{module:DF0B}
\noindent\textit{Audited source: \nolinkurl{DF0_good_grid.tex}.}
\subsection{Quasiuniform data and the exact material matrix}

Let $M\ge2$, $b=2\pi/M$, and let
\begin{equation}\label{DF0B:eq:quasi}
 h_i>0,\qquad \sum_i h_i=2\pi,
 \qquad (1-\rho)b\le h_i\le(1+\rho)b,
 \qquad 0<\rho<\frac14.
\end{equation}
Put $d_i=h_i-b$ and $h_i(t)=b+td_i$.  Then
\begin{equation}\label{DF0B:eq:path-quasi}
 (1-\rho)b\le h_i(t)\le(1+\rho)b,
 \qquad |d_i|\le\rho b,\qquad \sum_i d_i=0.
\end{equation}

For a positive current partition $g=(g_i)$ with endpoints $\theta_i$ and
midpoints $c_i$, the assembled foundation F148 gives the exact material
derivative of
\begin{equation}\label{DF0B:eq:w}
 w_i=\Hil q_g(c_i)=\frac1\pi\sum_jg_j\Phi(c_i-c_j),
 \qquad \Phi'(x)=-\log\left|2\sin\frac x2\right|,
\end{equation}
in the form
\begin{equation}\label{DF0B:eq:matrix}
 \dot w_i=\frac1\pi\sum_kd_k\mathcal E_{ik}(g).
\end{equation}
For completeness, choose for each target $i$ a cut $\theta_0$ at circular
distance at least $\pi/3$ from $c_i$.  With
$K=\Phi'$ and cyclic indices relabelled at that cut, the entries are
\begin{align}
 \mathcal E_{ik}^-&=
 \sum_{j<k}g_jK(c_i-c_j)+\frac{g_k}{2}K(c_i-c_k)
 -\int_{\theta_0}^{c_k}K(c_i-y)\,dy, &&k<i,\notag\\
 \mathcal E_{ik}^+&=
 \int_{c_k}^{\theta_0+2\pi}K(c_i-y)\,dy
 -\frac{g_k}{2}K(c_i-c_k)-\sum_{j>k}g_jK(c_i-c_j), &&k>i,
 \label{DF0B:eq:E-explicit}\\
 \mathcal E_{ii}&=\frac12\left\{
 \sum_{j<i}g_jK(c_i-c_j)-\int_{\theta_0}^{\theta_i}K(c_i-y)\,dy
 \right.\notag\\
 &\hspace{34mm}\left.
 -\sum_{j>i}g_jK(c_i-c_j)
 +\int_{\theta_{i+1}}^{\theta_0+2\pi}K(c_i-y)\,dy
 \right\}.\notag
\end{align}
Set $\mathcal E_{ik}=\mathcal E_{ik}^-$ for $k<i$ and
$\mathcal E_{ik}=\mathcal E_{ik}^+$ for $k>i$.  F148 proves
\eqref{DF0B:eq:matrix} by differentiating the Clausen formula: every raw
coefficient equals a common $\Phi(c_i)$ plus the displayed quadrature
primitive, and the common term vanishes because $\sum_kd_k=0$.

Every entry also has the following signed sawtooth representation.  With
\begin{equation}\label{DF0B:eq:eta}
 \eta(y)=\begin{cases}
 -(y-\theta_j),&\theta_j<y<c_j,\\
 \theta_{j+1}-y,&c_j<y<\theta_{j+1},
 \end{cases}
\end{equation}
one has on every complete cell
\begin{equation}\label{DF0B:eq:eta-properties}
 |\eta|\le g_j/2,\qquad \int_{I_j}\eta=0,
\end{equation}
and, on the left or right arc selected by an antipodal cut,
\begin{equation}\label{DF0B:eq:saw}
 \mathcal E_{ik}=\pm\int_{\mathrm{arc}(i,k)}
 \eta(y)K'(c_i-y)\,dy,
 \qquad K'(x)=-\frac12\cot\frac x2.
\end{equation}
The terminal half-cell is part of the integral; it is not separated from
the complete-cell primitive.

\subsection{The variable-grid discrete CZ lemma}

Write $|i-k|_M$ for cyclic index distance.  Since adding a constant to
every entry of one row does not change \eqref{DF0B:eq:matrix} on
$\sum d_k=0$, choose the representative
$\widetilde{\mathcal E}$ whose one-sided primitives are centred at an
antipodal cut.  Thus the two terminal values at the cut are opposite (or
zero when there is a single antipodal index).  We do not subtract a row
average, which could destroy the far-index size.

\begin{lemma}[Signed primitive estimates on a comparable mesh]
\label{DF0B:lem:kernel}
Assume
\begin{equation}\label{DF0B:eq:current-comparable}
 \kappa b\le g_j\le\kappa^{-1}b
\end{equation}
with fixed $0<\kappa<1$.  Then, after choosing the three usual shifted
antipodal cuts, the normalized matrix
$A=b^{-1}\widetilde{\mathcal E}$ satisfies
\begin{align}
 |A_{ik}|&\le\frac{C_\kappa}{1+|i-k|_M},\label{DF0B:eq:size}\\
 |A_{i,k+1}-A_{ik}|+|A_{i+1,k}-A_{ik}|
 &\le\frac{C_\kappa}{(1+|i-k|_M)^2}
 \label{DF0B:eq:smooth}
\end{align}
away from the two harmless cut indices.  At a cut the two one-sided
values obey \eqref{DF0B:eq:size}.  Moreover, for every cyclic index interval
$J$,
\begin{multline}\label{DF0B:eq:testing}
 \sum_{i\in J}\left|\sum_{k\in J}A_{ik}\right|^2
 +\sum_{k\in J}\left|\sum_{i\in J}A_{ik}\right|^2
 \le C_\kappa|J|.
\end{multline}
All constants are independent of $M$ and of cumulative displacement of
the variable grid from the regular grid.
\end{lemma}

\begin{proof}
Quasiuniformity converts physical distance to index distance:
\[
 \operatorname {dist}(c_i,c_k)\asymp_\kappa
 b(1+|i-k|_M)
\]
on the selected arc.  In \eqref{DF0B:eq:saw}, a complete cell has both the
endpoint cancellation and the zero mean in
\eqref{DF0B:eq:eta-properties}.  Subtracting $K'(c_i-c_j)$ on that cell and
using $|(K')'(x)|\le C/\operatorname {dist}(x,2\pi\mathbb Z)^2$
gives
\[
 \left|\int_{I_j}\eta(y)K'(c_i-y)\,dy\right|
 \le \frac{C_\kappa b}{(1+|i-j|_M)^2}.
\]
The terminal half-cell is estimated together with its endpoint atom and
costs $C_\kappa b/(1+|i-k|_M)$.  Summing from the antipode proves
\eqref{DF0B:eq:size}.  Advancing the terminal index by one cell leaves exactly
one complete zero-mean cell, proving the $k$-difference in
\eqref{DF0B:eq:smooth}.  Moving the target by one comparable cell and
subtracting the old kernel value gives the $i$-difference.  The cut
values have physical distance comparable with one and satisfy the same
size estimate directly.

For testing, retain the left/right signs from the antipodally centred
primitive.  Summation by parts with \eqref{DF0B:eq:smooth} shows that a target
at index distance $r$ from the nearer boundary of $J$ contributes at
most
\[
 C_\kappa\left(1+\log\frac{|J|+1}{r+1}\right).
\]
The square of this profile is summable:
\[
 \sum_{r=0}^{|J|-1}
 \left(1+\log\frac{|J|+1}{r+1}\right)^2\le C|J|.
\]
This proves the first test.  For the adjoint test, reverse the source and
target roles in \eqref{DF0B:eq:saw}; the complete-cell zero mean and the
terminal endpoint term are unchanged, so the same argument proves the
second test.  No absolute harmonic row sum is used.
\end{proof}

\begin{lemma}[Finite cyclic discrete $T1$ square]
\label{DF0B:lem:T1}
Any cyclic matrix satisfying \eqref{DF0B:eq:size}--\eqref{DF0B:eq:testing} obeys
\begin{equation}\label{DF0B:eq:T1}
 \|Ax\|_{\ell^2}\le C_\kappa\|x\|_{\ell^2}.
\end{equation}
\end{lemma}

\begin{proof}
Use the three shifted dyadic grids on the cyclic index set and expand
$x$ and a test vector $y$ in normalized Haar functions.  If their
support intervals are disjoint, subtract the kernel at the two interval
centres.  Equation \eqref{DF0B:eq:smooth} gives
\[
 |\langle Ah_I,h_J\rangle|
 \le C_\kappa
 \frac{|I|^{1/2}|J|^{1/2}\min\{|I|,|J|\}}
 {(\operatorname {dist}(I,J)+\max\{|I|,|J|\})^2}.
\]
The two-sided annular sum is Schur summable.  For nested supports, split
the larger Haar function into its constant child value plus a mean-zero
remainder.  The remainder has the same separated estimate; the constant
parts form the two paraproducts.  Their coefficients are Carleson because
\eqref{DF0B:eq:testing} is exactly the square testing condition for $A1_J$
and $A^*1_J$.  The dyadic Carleson embedding theorem bounds both
paraproducts.  The three grids place every pair of cyclic intervals in
one of these separated or nested configurations with comparable parent.
Summing the Haar coefficients proves \eqref{DF0B:eq:T1}.
\end{proof}

\begin{theorem}[Quasiuniform quadrature-primitive square]
\label{DF0B:thm:QPS}
Under \eqref{DF0B:eq:current-comparable}, for every zero-sum direction $d$,
\begin{equation}\label{DF0B:eq:QPS}
 \boxed{
 \sum_i\left|\sum_kd_k\mathcal E_{ik}(g)\right|^2
 \le C_\kappa b^2\sum_kd_k^2.}
\end{equation}
\end{theorem}

\begin{proof}
The chosen row gauge does not change the action on $d$.  Apply
Lemmas~\ref{DF0B:lem:kernel}--\ref{DF0B:lem:T1} to
$A=b^{-1}\widetilde{\mathcal E}$ and rescale.
\end{proof}

\subsection{Endpoint integration inside one equality component}

Along \eqref{DF0B:eq:path-quasi}, Theorem~\ref{DF0B:thm:QPS} and
\eqref{DF0B:eq:matrix} give
\begin{equation}\label{DF0B:eq:material-L2}
 \sum_i|\dot w_i(t)|^2\le C_\rho b^2\sum_kd_k^2.
\end{equation}
At $t=0$ the regular $M$-fan has $w_i(0)=0$.

\begin{theorem}[Complete quasiuniform AVS]
\label{DF0B:thm:AVS}
Every fan satisfying \eqref{DF0B:eq:quasi} obeys
\begin{equation}\label{DF0B:eq:AVS}
 \boxed{
 T(\alpha)=\sum_i h_i|\Hil q_\alpha(c_i)|^2
 \le C_\rho b^3\sum_i(h_i-b)^2
 \le C_\rho\Dvar(\alpha).}
\end{equation}
Hence HQ and DF are closed on the entire quasiuniform component, with no
adjacent-ratio derivative or cumulative-position hypothesis.
\end{theorem}

\begin{proof}
Cauchy--Schwarz in $t$, \eqref{DF0B:eq:quasi}, and
\eqref{DF0B:eq:material-L2} give
\begin{align*}
 T(\alpha)
 &\le(1+\rho)b\int_0^1\sum_i|\dot w_i(t)|^2\,dt\\
 &\le C_\rho b^3\sum_i d_i^2.
\end{align*}
For the last inequality in \eqref{DF0B:eq:AVS}, use the pair formula of stage
153.  On \eqref{DF0B:eq:quasi},
\[
 h_ih_j(h_i^2-h_j^2)^2\asymp_\rho
 b^4(h_i-h_j)^2.
\]
Since $2\pi=Mb$ and $\sum_i d_i=0$,
\[
 \Dvar(\alpha)
 \asymp_\rho\frac{b^4}{Mb}
 \sum_{i,j}(d_i-d_j)^2
 =2b^3\sum_i d_i^2
\]
up to constants depending only on $\rho$.  Stage 153 gives
AVS $\Rightarrow$ HQ $\Rightarrow$ DF.
\end{proof}

The same square is needed at the collapsed locations at which a bad
block will later be opened.  Write $z_i=\theta_i$ for the cell endpoints.

\begin{corollary}[Quasiuniform endpoint samples]
\label{DF0B:cor:endpoints}
Under \eqref{DF0B:eq:quasi},
\begin{equation}\label{DF0B:eq:endpoint-samples}
 \boxed{
 b\sum_i|\Hil q_\alpha(z_i)|^2
 \le C_\rho b^3\sum_i(h_i-b)^2.}
\end{equation}
\end{corollary}

\begin{proof}
Repeat the material calculation with the target carried by the endpoint
$z_i(t)$ instead of the midpoint $c_i(t)$.  The exact length gradient is
again a signed quadrature primitive.  At the target endpoint the two
incident terminal half-cells must be kept together; their quadratic
profiles and first derivatives vanish at their common endpoint.  Thus
there is no boundary logarithm.  Advancing either source or target by one
cell leaves one complete mean-zero sawtooth, so
\eqref{DF0B:eq:size}--\eqref{DF0B:eq:testing} hold for the endpoint matrix with the
same constants.  Lemma~\ref{DF0B:lem:T1} therefore gives the analogue of
\eqref{DF0B:eq:material-L2}.  At the regular fan the spline is even about every
endpoint, so all endpoint Hilbert samples vanish.  Integration in $t$
proves \eqref{DF0B:eq:endpoint-samples}.
\end{proof}

\begin{remark}[Why stage 152 does not contradict this theorem]
The false HMS/QPS family has half of its current cells of length
$\varepsilon a$ while their fixed straight-path directions have size
$a$.  It leaves every fixed quasiuniform component.  In
Theorem~\ref{DF0B:thm:AVS}, both the current lengths and the entire connecting
segment remain comparable with $b$, and $|d_i|\le\rho b$.
\end{remark}

\subsection{Effect on the ELCA ledger}

Use the good/bad selection of stage 153.  On a maximal collection of
good cells whose zero-cell deletion contains no intervening bad block,
the current effective mesh is comparable with the selected stratum
length $b$.  The preceding theorem closes every good--good
first-separation rectangle.

\begin{corollary}[Only bad-block rectangles remain]
\label{DF0B:cor:remaining}
In ELCA of stage 153, all rows whose source and target shadows avoid the
bad set satisfy
\begin{equation}\label{DF0B:eq:good-good}
 T_{\mathrm{good\text{-}good}}
 \le C_\rho b^3\sum_{i\in\mathcal G_\rho}(h_i-b)^2.
\end{equation}
Thus the remaining pure endpoint ledger consists only of rectangles for
which the source shadow or target shadow meets a maximal bad block.  Such
a block must be retained as a whole until its first separating scale and
paid by
\begin{equation}\label{DF0B:eq:bad-charge}
 \sum_{i\in\mathcal B_\rho}h_i(h_i^2-\ell^2)^2.
\end{equation}
\end{corollary}

\begin{proof}
Apply Theorem~\ref{DF0B:thm:AVS} on each good component after the common
collapsed-uniform reference has been fixed.  At its two boundary cuts,
keep the terminal half-cell with the adjacent bad block; those boundary
rectangles are not included in \eqref{DF0B:eq:good-good}.  Disjoint good
components have disjoint physical mass, so their squares sum.  The
reference-charge estimate of stage 153 then gives the displayed right
side.  Every omitted rectangle meets a bad block by construction.
\end{proof}
\endgroup

\begingroup
\section{Node ZG: strong ZG-T1 retracted; structured ZG-COM closed}
\label{module:ZG}
\noindent\textit{Audited source: \nolinkurl{ZG_zero_gap.tex}.}
\subsection{Intrinsic cut and source decomposition}

Let $I_i=[\theta_i,\theta_{i+1}]$ have length $h_i>0$ and midpoint
$c_i$.  Put
\begin{equation}\label{ZG:eq:q}
 q_h=\sum_iq_i,
 \qquad
 q_i(\theta_i+x)=\frac12\min\{x^2,(h_i-x)^2\}
 \quad(0\le x\le h_i),
\end{equation}
where every $q_i$ is extended by zero outside $I_i$.  Define
\begin{equation}\label{ZG:eq:T}
 T(h)=\sum_i h_i|\Hil q_h(c_i)|^2.
\end{equation}

As in stage 153, let
\begin{equation}\label{ZG:eq:D}
 \ell^2=\frac1{2\pi}\sum_i h_i^3,
 \qquad
 \Dvar(h)=\sum_i h_i(h_i^2-\ell^2)^2.
\end{equation}
The far region $\Dvar\ge\delta\sum_i h_i^5$ is already closed by the
absolute adaptive sampling estimate.  In the complementary region choose
$0<\rho<1/8$ and set
\begin{equation}\label{ZG:eq:GB}
 \mathcal G=\{i:|h_i/\ell-1|<\rho\},
 \qquad \mathcal B=\{1,\ldots,N\}\setminus\mathcal G.
\end{equation}
Put
\begin{equation}\label{ZG:eq:bE}
 M=|\mathcal G|,\qquad b=\frac{2\pi}{M},\qquad
 E=\sum_{i\in\mathcal B}h_i.
\end{equation}
The near-stratum selection gives $b\asymp_\rho\ell$ and $E$ smaller than
a fixed multiple of $2\pi$.  We fix the stratum threshold so that
$E\le\pi/4$.  If $E$ is larger, the absolute sampling bound
$T\le C\sum_i h_i^5$, together with \eqref{ZG:eq:budget} below and
$b^4\le Cb^4E$, already belongs to the closed far row.  The following
static budget is the only
place where the definition of the bad set is used:
\begin{equation}\label{ZG:eq:budget}
 b^3\sum_{i\in\mathcal G}(h_i-b)^2
 +b^4E+\sum_{i\in\mathcal B}h_i^5
 \le C_\rho\Dvar(h).
\end{equation}
Indeed, the good term is the reference-charge estimate of stage 153.
On the bad set, with $r=h_i/\ell$,
\[
 |r-1|\ge\rho,
 \qquad r^5+r\le C_\rho r(r^2-1)^2.
\]
Multiplication by $\ell^5$, summation, and $b\asymp_\rho\ell$ give
\[
 b^4E+\sum_{i\in\mathcal B}h_i^5
 \le C_\rho\sum_{i\in\mathcal B}
 h_i(h_i^2-\ell^2)^2,
\]
which proves \eqref{ZG:eq:budget} without invoking the retracted block
lemma of stage 155.

Split
\begin{equation}\label{ZG:eq:split}
 q_{\rm G}=\sum_{i\in\mathcal G}q_i,
 \qquad q_{\rm B}=\sum_{i\in\mathcal B}q_i.
\end{equation}
Then
\begin{equation}\label{ZG:eq:Tsplit}
 T(h)\le2T_{\rm G}+2T_{\rm B},
 \qquad
 T_X:=\sum_i h_i|\Hil q_X(c_i)|^2.
\end{equation}

\subsection{The bad source is an absolute square}

We recall the special adaptive sampling fact proved in stage 153.  If an
interval partition has midpoints $c_i$ and $f\in L^2(\T)$ satisfies
\begin{equation}\label{ZG:eq:selfzero}
 \operatorname {pv}\int_{I_i}\frac1{2\pi}
 \cot\frac{c_i-y}{2}f(y)\,dy=0
 \quad\text{for every }i,
\end{equation}
then
\begin{equation}\label{ZG:eq:adaptive}
 \sum_i h_i|\Hil f(c_i)|^2\le C\|f\|_2^2.
\end{equation}
It follows by comparing the truncation at radius $h_i/2$ with the
maximal Hilbert transform and integrating the Cotlar inequality over the
disjoint cells.

\begin{lemma}[Bad-source square]\label{ZG:lem:bad}
For the decomposition \eqref{ZG:eq:split},
\begin{equation}\label{ZG:eq:bad}
 \boxed{
 T_{\rm B}\le C\sum_{j\in\mathcal B}h_j^5.}
\end{equation}
\end{lemma}

\begin{proof}
If $i\in\mathcal G$, the function $q_{\rm B}$ vanishes on $I_i$.  If
$i\in\mathcal B$, its restriction to $I_i$ is the even bubble $q_i$,
while the cotangent kernel centred at $c_i$ is odd.  Thus
\eqref{ZG:eq:selfzero} holds for $f=q_{\rm B}$ on every target cell.
Equation \eqref{ZG:eq:adaptive} and the exact cell integral
\[
 \|q_{\rm B}\|_2^2
 =\sum_{j\in\mathcal B}\int_{I_j}q_j^2
 =\frac1{320}\sum_{j\in\mathcal B}h_j^5
\]
prove \eqref{ZG:eq:bad}.
\end{proof}

\begin{remark}
This single application of \eqref{ZG:eq:adaptive} is the decisive cleaning
step.  All bad midpoints are now inside the source $q_{\rm B}$ and are
paid by their individual $h_j^5$ energies.  A maximal bad block is never
asked to pack its internal midpoint set at every dyadic scale.
\end{remark}

\subsection{The good skeleton with zero-source gaps}

Delete the bad bubbles but retain their intervals as gaps.  Thus
$q_{\rm G}$ consists of comparable good bubbles, while a maximal
consecutive bad block $B_\nu$ is a zero-source gap of length
\begin{equation}\label{ZG:eq:gaps}
 e_\nu=\sum_{i\in B_\nu}h_i,
 \qquad \sum_\nu e_\nu=E.
\end{equation}
Only the two endpoints $\partial B_\nu$ are singular locations for the
good source.

For a target interval $Q$ of physical length $L_Q$ (either a member of
one of the shifted grids or a node of the balanced target tree below),
define the endpoint weight
\begin{equation}\label{ZG:eq:Omega-end}
 \Omega_{\rm end}(B,Q)=
 \begin{cases}
 \displaystyle \frac e{L_Q}
 \left(1+\frac{\operatorname {dist}(B,Q)}{L_Q}\right)^{-3},
 &e\le L_Q,\\[3mm]
 \displaystyle \frac {L_Q}e
 \left(1+\frac{\operatorname {dist}(\partial B,Q)}{L_Q}\right)^{-3},
 &e>L_Q.
 \end{cases}
\end{equation}

\begin{lemma}[Genuine endpoint packing]\label{ZG:lem:end-packing}
For disjoint gaps, on each shifted dyadic grid and on every balanced
target interval tree used below,
\begin{equation}\label{ZG:eq:end-packing}
 \sup_Q\sum_\nu\Omega_{\rm end}(B_\nu,Q)
 +\sup_\nu\sum_Q\Omega_{\rm end}(B_\nu,Q)\le C.
\end{equation}
\end{lemma}

\begin{proof}
For $e\le L_Q$, disjointness bounds the total gap length in the $n$th
distance annulus by $C2^nL_Q$; the factor
$L_Q^{-1}2^{-3n}$ is summable.  For $e>L_Q$, sort gaps by
$e\asymp2^mL_Q$.  The endpoints themselves need not be
$2^mL_Q$-separated: two large gaps can be separated by one good cell.
What is separated is the interior mass of the disjoint gaps.  Hence an
annulus of length $C2^nL_Q$ meets at most
$C(1+2^{n-m})$ such gaps.  Multiplication by
$2^{-m}2^{-3n}$ is summable in $m,n$.

With a gap fixed, the coarse scale sum is
$\sum_{L_Q\ge e}e/L_Q<\infty$.  At fine scales there are only two
singular points, so the annular sum is bounded at each scale and
$\sum_{L_Q<e}L_Q/e<\infty$.  This proves \eqref{ZG:eq:end-packing}.  The
argument would fail with all internal bad midpoints in place of the two
endpoints; Lemma~\ref{ZG:lem:bad} is what removed them.
\end{proof}

We next isolate the precise zero-gap extension of the variable-grid
square.  Only the good midpoint atoms occur in the opening operator:
\begin{equation}\label{ZG:eq:tau}
 d\tau_{\rm G}=\sum_{i\in\mathcal G}h_i\delta_{c_i},
 \qquad d\mu=dx.
\end{equation}
Thus the source and target measures are different.  In particular the
one-weight matrix theorem of stage 154 cannot be quoted here.  We give
the continuous-to-atomic operator and its square estimate explicitly.

The exact identity behind the next estimate is worth recording.  With
$\Phi(x)=\sum_{n\ge1}\sin(nx)/n^2$, twice integrating
$q_j''=\mathbf1_{I_j}-h_j\delta_{c_j}$ gives
\begin{equation}\label{ZG:eq:Peano-H}
 \Hil q_{\rm G}(x)=\frac1\pi\sum_{j\in\mathcal G}
 \left\{h_j\Phi(x-c_j)-\int_{I_j}\Phi(x-y)\,dy\right\}.
\end{equation}
Each bracket has zero mass and zero first moment.  Formula
\eqref{ZG:eq:Peano-H}, rather than separate endpoint kernels, is used until
the first source--target separating scale.

Put $B=\bigcup_\nu B_\nu$, $G^*=\bigcup_{j\in\mathcal G}I_j$, and
\begin{equation}\label{ZG:eq:sigma}
 \gamma=\frac EM,
 \qquad
 \sigma=\mathbf1_B-
 \sum_{j\in\mathcal G}\frac{\gamma}{h_j}\mathbf1_{I_j}.
\end{equation}
Then $\int_\T\sigma=0$.  Let $u$ be the periodic primitive of $\sigma$
with zero mean and set
\begin{equation}\label{ZG:eq:collapse-map}
 R_t(x)=x-tu(x),\qquad 0\le t\le1.
\end{equation}
On a good cell $R_t'=1+t\gamma/h_j$; on a gap $R_t'=1-t$.
Consequently $R_1$ collapses every gap and maps $I_j$ affinely onto an
interval of length $h_j+\gamma$.  In particular
$R_1(c_j)=\widehat c_j$.

For $f\in L^2(\mu)$ put
\[
 f_0=f-\frac1{2\pi}\int_\T f.
\]
For an oriented cyclic arc from $x$ to $z$, let
$\epsilon_{x,z}$ be its signed indicator.  If $V_f$ is the periodic
zero-mean primitive of $f_0$, then
\begin{equation}\label{ZG:eq:oriented-primitive}
 V_f(z)-V_f(x)=\int_\T\epsilon_{x,z}(s)f_0(s)\,ds.
\end{equation}
For $y\in I_j$, $j\in\mathcal G$, write $m(y)=c_j$ and
$a_t(y)=R_t'(y)$.  Define
\begin{align}
 P_t(x,y)&=\Phi(R_t x-R_t m(y))-\Phi(R_t x-R_t y),
 \label{ZG:eq:Pt}\\
 (\mathcal T_tf)(x)
 &=\frac1\pi\int_{G^*}
 \Big[-f_0(y)P_t(x,y)\notag\\
 &\hspace{13mm}+a_t(y)\{\Phi'(R_tx-R_tm(y))
       [V_f(m(y))-V_f(x)]\notag\\
 &\hspace{35mm}-\Phi'(R_tx-R_ty)
       [V_f(y)-V_f(x)]\}\Big],dy .
 \label{ZG:eq:Tt}
\end{align}
This is a linear map from $L^2(\mu)$ to functions on the good target
points and annihilates constants.  On the zero-mean subspace,
\eqref{ZG:eq:oriented-primitive} gives the raw kernel
\begin{align}
 K_t(x,s)=\frac1\pi\Big[&-\mathbf1_{G^*}(s)P_t(x,s)
 \notag\\
 &+\int_{G^*}a_t(y)
 \{\Phi'(R_tx-R_tm(y))\epsilon_{x,m(y)}(s)\notag\\
 &\hspace{35mm}-\Phi'(R_tx-R_ty)\epsilon_{x,y}(s)\}\,dy\Big].
 \label{ZG:eq:Kt}
\end{align}
For arbitrary $f$ the actual kernel is the source-centred version
\begin{equation}\label{ZG:eq:Kt-centred}
 K_t^0(x,s)=K_t(x,s)-\frac1{2\pi}\int_\T K_t(x,r)\,dr,
 \qquad
 \mathcal T_tf(x)=\int_\T K_t^0(x,s)f(s)\,ds.
\end{equation}
Below we suppress the superscript $0$.  The subtracted row constant has
zero source difference.  Integrating \eqref{ZG:eq:Kt} in $s$, using
$\int_\T\Phi=0$, and pairing the two terminal half-cells gives
\begin{equation}\label{ZG:eq:row-gauge}
 \left|\frac1{2\pi}\int_\T K_t(x,s)\,ds\right|
 +\left|\partial_x\frac1{2\pi}\int_\T K_t(x,s)\,ds\right|
 \le C_\rho b^2.
\end{equation}
Thus the centred kernel obeys the same size, difference, and testing
bounds (with a changed constant).

\begin{proposition}[The strong ZG--$T1$ square is false]
\label{ZG:prop:Peano-counterexample}
There is no constant independent of $M$ for which
\begin{equation}\label{ZG:eq:false-Peano-square}
 \sum_{i\in\mathcal G}h_i|\mathcal T_0f(c_i)|^2
 \le C b^4\|f\|_{L^2(\T)}^2
\end{equation}
holds for every $f$.  In fact, on the uniform gap-free fan the norm of
$\mathcal T_0:L^2(dx)\to L^2(\tau_{\rm G})$ is at least a constant
multiple of $b$.
\end{proposition}

\begin{proof}
Take $h_j=b=2\pi/M$, $c_j=(j+\tfrac12)b$, and
\[
 f_M(\theta)=(2\pi)^{-1/2}e^{i(M-1)\theta}.
\]
Put $n=M-1$ and $q=e^{-ib}$.  Cellwise integration by parts in
\eqref{ZG:eq:Tt}, followed by cyclic summation, gives for the unnormalised
mode $e^{in\theta}$
\begin{equation}\label{ZG:eq:counterexample-exact}
 \mathcal T_0(e^{in\cdot})(c_0)
 =\frac{B_M+C_M}{\pi i(M-1)},
\end{equation}
where
\begin{align}
 B_M&=-(1-q)\sum_{j=0}^{M-1}q^j\Phi(jb),
 \label{ZG:eq:counterexample-B}\\
 C_M&=b e^{-ib/2}\sum_{j=1}^{M-1}
       \Phi'(jb)(1-q^j).
 \label{ZG:eq:counterexample-C}
\end{align}
Here is the calculation before summing.  Write $x=c_0$, $m=c_j$,
$d=x-m=-jb$, $F(y)=e^{iny}$, and
$P(r)=\Phi(d)-\Phi(d-r)$, and denote by
$[\mathcal T_0F(x)]_{I_j}$ the $I_j$-integral in \eqref{ZG:eq:Tt}.
Integration by parts of $-FP$, using
$F=(F/in)'$, cancels its interior $F\Phi'$ term against the last
primitive term in \eqref{ZG:eq:Tt} and gives
\begin{multline}\label{ZG:eq:counterexample-cell}
 \pi[\mathcal T_0F(x)]_{I_j}
 =\frac1{in}\Big\{
 -[F(m+r)P(r)]_{-b/2}^{b/2}\\
 +F(x)[\Phi(d+b/2)-\Phi(d-b/2)]
 +b\Phi'(d)[F(m)-F(x)]\Big\}.
\end{multline}
For $j=0$ the last product is understood by continuity and is zero.
Now
\[
 F(c_j)=-e^{-ib/2}q^j,
 \qquad \Phi(-z)=-\Phi(z),
 \qquad \Phi'(-z)=\Phi'(z).
\]
Substitution in \eqref{ZG:eq:counterexample-cell}, followed by one cyclic
index shift in the endpoint terms, gives
\eqref{ZG:eq:counterexample-B}--\eqref{ZG:eq:counterexample-C} and hence
\eqref{ZG:eq:counterexample-exact}.

The Riemann sum in \eqref{ZG:eq:counterexample-B} satisfies
\[
 b\sum_{j=0}^{M-1}e^{-ijb}\Phi(jb)
 \longrightarrow\int_0^{2\pi}e^{-ix}\Phi(x)\,dx=-i\pi.
\]
Since $1-e^{-ib}=ib+O(b^2)$, it follows that $B_M\to-\pi$.
Although $\Phi'$ has logarithmic endpoint singularities, the function
\[
 x\longmapsto\Phi'(x)(1-e^{-ix})
\]
extends continuously by zero at both endpoints.  Hence
\begin{align*}
 C_M&\longrightarrow
 \int_0^{2\pi}\Phi'(x)(1-e^{-ix})\,dx\\
 &= -\int_0^{2\pi}\Phi'(x)e^{-ix}\,dx=-\pi;
\end{align*}
the last equality is the first cosine coefficient of
$\Phi'(x)=\sum_{k\ge1}\cos(kx)/k$.  Thus
\begin{equation}\label{ZG:eq:counterexample-asymptotic}
 |\mathcal T_0(e^{in\cdot})(c_0)|\sim\frac2M.
\end{equation}
Translation covariance gives the same modulus at all $M$ target atoms.
The normalisation of $f_M$ cancels the total target mass $2\pi$, and so
\[
 \|\mathcal T_0f_M\|_{L^2(\tau_{\rm G})}
 =\frac{|B_M+C_M|}{\pi(M-1)}\sim\frac2M.
\]
Since $b^2=4\pi^2/M^2$, the ratio of the two sides at the norm level is
asymptotic to $M/(2\pi^2)$.  This contradicts
\eqref{ZG:eq:false-Peano-square}.  If a nonempty deleted block is required,
insert one gap of length $o(M^{-3})$ and distribute that length over the
good cells.  For each fixed $M$ the finite cell integrals above depend
continuously on those endpoints, so the same contradiction persists.
\end{proof}

\begin{remark}[Consequence for the former proof route]
The former ZG--$T1$ obligation asserted \eqref{ZG:eq:false-Peano-square} for
all $f$ and all $t$.  Proposition~\ref{ZG:prop:Peano-counterexample} shows
that no completion of its proposed Haar paragraph can prove that claim.
In particular, the previously stated kernel/testing package cannot have
all the asserted uniform consequences.  None of those claims is used
below.  The opening argument only needs the special density $\sigma$ and
the time integral appearing in the exact deformation identity.
\end{remark}

\begin{lemma}[Endpoint Calder\'on estimate for a small change of variable]
\label{ZG:lem:small-change}
Let $S:\T\to\T$ be an orientation-preserving, degree-one, piecewise
affine map, and choose its lift in the form $S(x)=x+v(x)$.  Suppose
\begin{equation}\label{ZG:eq:small-change-hypothesis}
 \|v'\|_\infty\le\eta<1.
\end{equation}
For $f\in W^{1,\infty}(\T)$ put
\[
 \mathfrak C_Sf=(\Hil f)\circ S-\Hil(f\circ S).
\]
Then
\begin{align}
 \|\mathfrak C_Sf\|_2
 &\le C_\eta\|S'-1\|_2\|f\|_\infty,
 \label{ZG:eq:small-change-H0}\\
 \|(\mathfrak C_Sf)'\|_2
 &\le C_\eta\|S'-1\|_2\|f'\|_\infty.
 \label{ZG:eq:small-change-H1}
\end{align}
The same constants hold for limits of such maps with nonnegative
derivative, provided each approximating map satisfies
\eqref{ZG:eq:small-change-hypothesis}.
\end{lemma}

\begin{proof}
We include the endpoint argument because replacing the right side of
\eqref{ZG:eq:small-change-H0} by an arbitrary-input $L^2$ norm is precisely
the false step retracted above.  First smooth $v$ and $f$; the estimates
obtained below are uniform under this smoothing.

On a lifted interval, changing variables in the first Hilbert transform
and writing
\[
 \Delta_v(x,y)=\frac{v(x)-v(y)}{x-y}
\]
gives, for the singular part of the cotangent kernel,
\begin{equation}\label{ZG:eq:calderon-kernel}
 \frac{1+v'(y)}{Sx-Sy}-\frac1{x-y}
 =\frac{v'(y)-\Delta_v(x,y)}
 {(x-y)(1+\Delta_v(x,y))}.
\end{equation}
For the moment suppose $\|v'\|_\infty\le\eta_0$, where the universal
$\eta_0>0$ will be chosen below.  Expand the last denominator
geometrically.  It is therefore enough to estimate
\begin{equation}\label{ZG:eq:calderon-Am}
 A_m(v,g)(x)=\operatorname {pv}\!\int
 g(y)\frac{(v'(y)-\Delta_v(x,y))
                 \Delta_v(x,y)^m}{x-y}\,dy,
 \qquad m\ge0,
\end{equation}
where in the application $g=f\circ S$.  The following two endpoint
multiplier bounds are the elementary Calder\'on estimate:
\begin{align}
 \|A_m(v,g)\|_2
 &\le C^{m+1}\|v'\|_2\|v'\|_\infty^m\|g\|_\infty,
 \label{ZG:eq:Am-H0}\\
 \|\partial_xA_m(v,g)\|_2
 &\le C^{m+1}\|v'\|_2\|v'\|_\infty^m\|g'\|_\infty.
 \label{ZG:eq:Am-H1}
\end{align}
Here is a direct frequency verification.  Use
\[
 \Delta_v(x,y)=\int_0^1v'(y+s(x-y))\,ds
\]
in \eqref{ZG:eq:calderon-Am}, expand into Fourier series, and divide the
frequencies into the regions in which one of the $m+1$ copies of $v'$
has largest modulus.  The difference
$v'(y)-\Delta_v(x,y)$ makes the multiplier vanish unless such a
$v'$-frequency dominates the frequency of $g$.  On each region the
multiplier and all its discrete differences are bounded by $C^{m+1}$.
Smooth dyadic resolution consequently writes that piece as
\[
 \sum_r P_rv'\,
 T_r(P_{\le r+3}v',\ldots,P_{\le r+3}v',P_{\le r+3}g),
\]
up to a uniformly summable family of translates; the kernels of $T_r$
have uniformly bounded $L^1$ norm.  The square-function inequality for
the sole high-frequency factor and the uniform $L^\infty$ bounds for
smooth low-frequency projections prove \eqref{ZG:eq:Am-H0}.  After applying
$\partial_x$, summation by parts in the Fourier frequency of $g$ moves
the output frequency onto $g'$; on the nonzero multiplier support its
coefficient is bounded by the same largest $v'$-frequency.  The identical
decomposition proves \eqref{ZG:eq:Am-H1}.  This also shows that the constants
are independent of the number of affine pieces.

Choose $\eta_0$ so that $C\eta_0<1/2$.  Summing
\eqref{ZG:eq:Am-H0}--\eqref{ZG:eq:Am-H1} with the geometric coefficients gives
the assertions for the flat singular kernel whenever
$\|v'\|_\infty\le\eta_0$.  The
difference $\frac12\cot(t/2)-1/t$ is smooth.  Subtracting $g(x)$ (constants
are annihilated) and integrating once by parts reduces its two estimates
to convolution with an $L^1$ kernel, and gives the same right sides.
Thus \eqref{ZG:eq:small-change-H0}--\eqref{ZG:eq:small-change-H1} hold in the
small-distortion case.  Notice that
\[
 \|g\|_\infty=\|f\|_\infty,
 \qquad \|g'\|_\infty\le(1+\eta_0)\|f'\|_\infty.
\]

It remains to pass from $\eta_0$ to the stated fixed $\eta<1$ without
using the divergent series above.  For $0\le s\le1$ define the degree-one
map $S_s$ by
\begin{equation}\label{ZG:eq:logarithmic-interpolation}
 S_s'(x)=\frac{S'(x)^s}{(2\pi)^{-1}\int_\T S'(y)^s\,dy},
 \qquad S_0=\operatorname {id}.
\end{equation}
Up to a harmless rotation, $S_1=S$.  Partition $[0,1]$ into a number
$L=L(\eta,\eta_0)$ of equal pieces and set
$T_j=S_{(j+1)/L}\circ S_{j/L}^{-1}$.  Since
$1-\eta\le S'\le1+\eta$, differentiating
\eqref{ZG:eq:logarithmic-interpolation} in $s$ gives
\begin{equation}\label{ZG:eq:interpolation-cost}
 \|T_j'-1\|_\infty\le\eta_0,
 \qquad
 \sum_{j=0}^{L-1}\|T_j'-1\|_2
 \le C_\eta\|S'-1\|_2.
\end{equation}
Indeed, on this compact interval the functions $\log t$ and $t-1$ are
uniformly comparable, and change of variables by every $S_s$ has
Jacobian bounded above and below by constants depending only on $\eta$.

For degree-one maps the commutator has the exact cocycle
\begin{equation}\label{ZG:eq:change-cocycle}
 \mathfrak C_{S_2\circ S_1}f
 =\mathfrak C_{S_1}(f\circ S_2)
   +(\mathfrak C_{S_2}f)\circ S_1.
\end{equation}
Iterate \eqref{ZG:eq:change-cocycle} over the $T_j$.  All intermediate
Jacobians and inverse Jacobians are bounded by $C_\eta$; composition
preserves the $L^\infty$ norm and multiplies the derivative norm by at
most $C_\eta$.  Apply the already proved small-distortion estimates to
each summand and use \eqref{ZG:eq:interpolation-cost}.  This proves
\eqref{ZG:eq:small-change-H0}--\eqref{ZG:eq:small-change-H1} for every fixed
$\eta<1$.

Approximation and weak lower semicontinuity in $H^1$ remove the smoothing
and prove the lemma.
\end{proof}

Put $R=R_1$ and $\widehat h_j=h_j+\gamma$ on the good indices.  After
deleting the gaps, the $\widehat h_j$ form a positive partition of the
circle and remain comparable with $b$.  Let $\widehat q$ be its complete
quadratic spline.  Set $p=\widehat q\circ R$ and define the monotone-pullback
commutator
\begin{equation}\label{ZG:eq:collapse-commutator}
 \mathfrak C_R\widehat q(x)
 :=\Hil\widehat q(Rx)-\Hil p(x).
\end{equation}
Changing variables $z=R(y)$ in the first Hilbert transform gives the
entirely explicit kernel formula
\begin{multline}\label{ZG:eq:collapse-commutator-kernel}
 \mathfrak C_R\widehat q(x)=\frac1{2\pi}\int_\T
 \widehat q(Ry)\Big\{R'(y)
 \cot\frac{Rx-Ry}{2}\\
 -\cot\frac{x-y}{2}\Big\}\,dy ,
\end{multline}
with principal values taken symmetrically.  Here $R'=0$ on every gap and
$R'=1+\gamma/h_j$ on $I_j$.

\begin{theorem}[Collapsed-bubble commutator square]
\label{ZG:thm:collapse-commutator}
One has
\begin{equation}\label{ZG:eq:collapse-commutator-square}
 \boxed{
 \sum_{i\in\mathcal G}h_i
 |\mathfrak C_R\widehat q(c_i)|^2
 \le C_\rho b^4E.}
\end{equation}
The constant must be independent of the number, length, and subdivision
of the gaps.
\end{theorem}

\begin{proof}
We resolve the collapse into geometrically summable, nondegenerate
changes of variable.  If $B_\nu$ is a gap of length $e_\nu$, define the
level-$k$ geometry by
\begin{equation}\label{ZG:eq:dyadic-geometry}
 e_\nu^{(k)}=2^{-k}e_\nu,
 \qquad h_j^{(k)}=h_j+(1-2^{-k})\gamma,
 \qquad k\ge0.
\end{equation}
These lengths sum to $2\pi$.  Let $S_k$ be the affine map from the
level-$k$ circle to the level-$(k+1)$ circle which preserves the cyclic
order of all cells.  Thus
\begin{equation}\label{ZG:eq:Sk-derivative}
 S_k'=\frac12\quad\hbox{on the gaps},
 \qquad
 S_k'=1+\frac{2^{-k-1}\gamma}{h_j^{(k)}}
       \quad\hbox{on the $j$th good cell}.
\end{equation}
Put $R_0=\operatorname {id}$ and $R_{k+1}=S_k\circ R_k$.  Then
$R_k$ converges uniformly to the collapse $R$.

The good lengths are mutually comparable and their average is
$g=(2\pi-E)/M$.  Since $E\le\pi/4$ and $\rho<1/8$, it follows directly
from \eqref{ZG:eq:Sk-derivative} that
\begin{equation}\label{ZG:eq:Sk-uniform}
 \|S_k'-1\|_\infty\le\frac12,
 \qquad
 1\le R_k'|_{I_j}\le C_\rho .
\end{equation}
Moreover, using $\sum_j(h_j^{(k)})^{-1}\le C_\rho M/b$ and
$\gamma=E/M$, we have the exact layer computation
\begin{align}
 \|S_k'-1\|_2^2
 &=\frac14,2^{-k}E
   +2^{-2k-2}\gamma^2
       \sum_{j\in\mathcal G}\frac1{h_j^{(k)}}
 \le C_\rho2^{-k}E.                         \label{ZG:eq:Sk-L2}
\end{align}

Let $R_{k,\infty}$ be the monotone tail collapse from the level-$k$
geometry to the final complete fan, and put
\begin{equation}\label{ZG:eq:pk}
 p_k=\widehat q\circ R_{k,\infty}.
\end{equation}
Then $p_0=p$, $p_k=p_{k+1}\circ S_k$, and $p_k$ vanishes on the
level-$k$ gaps.  On a good cell $R_{k,\infty}$ is affine with derivative
$\widehat h_j/h_j^{(k)}$.  Since all these lengths are comparable with
$b$, the explicit quadratic bubble gives, uniformly in $k$,
\begin{equation}\label{ZG:eq:pk-bounds}
 \|p_k\|_\infty\le C_\rho b^2,
 \qquad \|p_k'\|_\infty\le C_\rho b.
\end{equation}

On the original good set define
\[
 G_k=(\Hil p_k)\circ R_k.
\]
The identities $R_{k+1}=S_k\circ R_k$ and
$p_k=p_{k+1}\circ S_k$ give the exact commutator cocycle
\begin{equation}\label{ZG:eq:commutator-cocycle}
 G_{k+1}-G_k=(\mathfrak C_{S_k}p_{k+1})\circ R_k.
\end{equation}
Lemma~\ref{ZG:lem:small-change}, \eqref{ZG:eq:Sk-L2},
\eqref{ZG:eq:pk-bounds}, and the change of variables on the good cells in
\eqref{ZG:eq:Sk-uniform} therefore imply
\begin{align}
 \|G_{k+1}-G_k\|_{L^2(G)}
 &\le C_\rho b^2 2^{-k/2}E^{1/2},
 \label{ZG:eq:Gk-H0}\\
 \|(G_{k+1}-G_k)'\|_{L^2(G)}
 &\le C_\rho b 2^{-k/2}E^{1/2}.
 \label{ZG:eq:Gk-H1}
\end{align}
Both series are summable.  Since $G_0=\Hil p$ and
$G_k\to(\Hil\widehat q)\circ R$ in $H^1$ on the union of the original
good cells, their sum yields
\begin{equation}\label{ZG:eq:commutator-continuous-H1}
 \|\mathfrak C_R\widehat q\|_{L^2(G)}^2
 +b^2\|(\mathfrak C_R\widehat q)'\|_{L^2(G)}^2
 \le C_\rho b^4E.
\end{equation}

For $F\in H^1(I_i)$,
\[
 h_i|F(c_i)|^2
 \le2\int_{I_i}|F|^2+h_i^2\int_{I_i}|F'|^2.
\]
Apply this to $F=\mathfrak C_R\widehat q$, sum over the disjoint good
cells, and use $h_i\asymp_\rho b$ together with
\eqref{ZG:eq:commutator-continuous-H1}.  This proves
\eqref{ZG:eq:collapse-commutator-square}.
\end{proof}

\begin{lemma}[Static zero-gap opening square]\label{ZG:lem:opening}
Let
$\widehat c_j$ be the midpoints of the preceding collapsed partition and put
\begin{equation}\label{ZG:eq:reference-vector}
 U_i=\Hil\widehat q(\widehat c_i),\qquad i\in\mathcal G.
\end{equation}
Then
\begin{equation}\label{ZG:eq:opening-square}
 \boxed{
 \sum_{i\in\mathcal G}h_i
 \left|\Hil q_{\rm G}(c_i)-U_i\right|^2
 \le C_\rho b^4E.}
\end{equation}
The constant is independent of the number and subdivision of the gaps.
\end{lemma}

\begin{proof}
For $x=c_i$ define
\begin{equation}\label{ZG:eq:Wt}
 W_i(t)=\frac1\pi\sum_{j\in\mathcal G}\int_{I_j}a_t(y)
 \{\Phi(R_tc_i-R_tc_j)-\Phi(R_tc_i-R_ty)\}\,dy.
\end{equation}
Changing variables $z=R_t(y)$ shows that $W_i(t)$ is the Hilbert
sample of the quadratic bubbles on the images $R_t(I_j)$, evaluated at
their transported midpoint.  Hence
\[
 W_i(0)=\Hil q_{\rm G}(c_i),
 \qquad W_i(1)=\Hil\widehat q(\widehat c_i)=U_i.
\]
Differentiating \eqref{ZG:eq:Wt} under the integral and using $u'=\sigma$
gives the exact identity
\begin{equation}\label{ZG:eq:opening-identity}
 \dot W_i(t)=\mathcal T_t\sigma(c_i).
\end{equation}
There is no omitted endpoint term: the gap motion is contained in the
primitive differences in \eqref{ZG:eq:Tt}.

It remains to separate the part of the opening error already controlled
by elementary Sobolev sampling.  On a good cell put
$a_j=1+\gamma/h_j=\widehat h_j/h_j$.  The quadratic bubble scales
exactly, so
\begin{equation}\label{ZG:eq:pullback-bubble}
 p|_{I_j}=a_j^2q_j,
 \qquad p=0\quad\hbox{on every gap}.
\end{equation}
Consequently $r:=q_{\rm G}-p$ belongs to $H^1(\T)$ and, since
$h_j\asymp_\rho b$ and the near stratum has bounded small $E$,
\begin{align}
 \|r\|_2^2
 &\le C_\rho E^2\sum_{j\in\mathcal G}h_j^5
 \le C_\rho b^4E,\label{ZG:eq:amplitude-H0}\\
 \|r'\|_2^2
 &\le C_\rho E^2\sum_{j\in\mathcal G}h_j^3
 \le C_\rho b^2E.\label{ZG:eq:amplitude-H1}
\end{align}
Indeed $|a_j^2-1|\le C_\rho E$, while
$Mb=2\pi$.  For every $F\in H^1(I_j)$ the fundamental theorem of
calculus and Cauchy--Schwarz give
\[
 h_j|F(c_j)|^2
 \le2\int_{I_j}|F|^2+h_j^2\int_{I_j}|F'|^2.
\]
Apply this with $F=\Hil r$, sum the good cells, and use the $L^2$
isometry of $\Hil$ and its commutation with differentiation.  Equations
\eqref{ZG:eq:amplitude-H0}--\eqref{ZG:eq:amplitude-H1} yield
\begin{equation}\label{ZG:eq:amplitude-trace}
 \sum_{i\in\mathcal G}h_i|\Hil r(c_i)|^2
 \le C_\rho b^4E.
\end{equation}
Finally, \eqref{ZG:eq:collapse-commutator} gives the exact decomposition
\[
 \Hil q_{\rm G}(c_i)-\Hil\widehat q(Rc_i)
 =\Hil r(c_i)-\mathfrak C_R\widehat q(c_i).
\]
Combine \eqref{ZG:eq:amplitude-trace} with
\eqref{ZG:eq:collapse-commutator-square} to prove
\eqref{ZG:eq:opening-square}.  Integrating
\eqref{ZG:eq:opening-identity} also shows that this static difference is
$-\int_0^1\mathcal T_t\sigma(c_i)\,dt$; hence ZG--COM closes exactly the
structured integrated opening-flow square, without a pointwise-in-$t$
operator estimate.

\end{proof}

The restriction $i\in\mathcal G$ in Lemma~\ref{ZG:lem:opening} is
essential.  On a bad midpoint the collapsed reference value is an
endpoint trace, not zero.  It must be retained in the endpoint layer,
rather than incorrectly charged to the opening density alone.

\begin{lemma}[Normalized-bubble bad-target trace]\label{ZG:lem:gap-target}
With the same notation,
\begin{equation}\label{ZG:eq:gap-target}
 \boxed{
 \sum_{i\in\mathcal B}h_i|\Hil q_{\rm G}(c_i)|^2
 \le C_\rho\left\{
 b^3\sum_{j\in\mathcal G}(\widehat h_j-b)^2+b^4E\right\}.}
\end{equation}
The constant is independent of the number and subdivision of the gaps.
\end{lemma}

\begin{proof}
Put
\[
 g:=\frac{2\pi-E}{M}=b-\gamma.
\]
Thus $\widehat h_j-b=h_j-g$, and $g\asymp_\rho b$.  On $0\le s\le1$
write
\[
 \psi(s)=\frac12\min\{s^2,(1-s)^2\},
 \qquad \overline\psi=\int_0^1\psi(s)\,ds=\frac1{24}.
\]
For a good cell $I_j=[\theta_j,\theta_{j+1}]$ define the normalized
bubble
\[
 p_j(\theta_j+h_js)=g^2\psi(s),
 \qquad p_{\rm G}=\sum_{j\in\mathcal G}p_j.
\]
The amplitude error
\[
 f:=q_{\rm G}-p_{\rm G}
 =\sum_{j\in\mathcal G}(h_j^2-g^2)
   \psi\!\left(\frac{\,\cdot-\theta_j\,}{h_j}\right)\mathbf1_{I_j}
\]
is even on every good target cell and is zero on every bad target cell.
It therefore satisfies the self-zero hypothesis \eqref{ZG:eq:selfzero}.
Since $\int_0^1\psi^2=1/320$ and $h_j,g\asymp_\rho b$,
\begin{align}
 \sum_i h_i|\Hil f(c_i)|^2
 &\le C\|f\|_2^2\notag\\
 &=\frac1{320}\sum_{j\in\mathcal G}
 h_j(h_j^2-g^2)^2
 \le C_\rho b^3\sum_{j\in\mathcal G}(h_j-g)^2.
 \label{ZG:eq:amplitude-error}
\end{align}

It remains to estimate $p_{\rm G}$.  Decompose it as
\begin{equation}\label{ZG:eq:normalized-decomposition}
 p_{\rm G}=g^2\overline\psi\,\mathbf1_{G^*}+r_{\rm G},
 \qquad
 r_{\rm G}|_{I_j}
 =g^2\left\{\psi\!\left(\frac{\,\cdot-\theta_j\,}{h_j}\right)
             -\overline\psi\right\}.
\end{equation}
Because $\Hil\mathbf1=0$,
$\Hil\mathbf1_{G^*}=-\Hil\mathbf1_B$.  The function $\mathbf1_B$ is
constant on every target cell, hence also satisfies
\eqref{ZG:eq:selfzero}.  The adaptive sampling estimate gives
\begin{equation}\label{ZG:eq:mean-gap}
 \sum_i h_i|\Hil(g^2\overline\psi\,\mathbf1_{G^*})(c_i)|^2
 \le Cg^4\|\mathbf1_B\|_2^2
 =Cg^4E.
\end{equation}

We now prove the same bound for $r_{\rm G}$ on the bad targets.  Fix a
maximal gap $B_\nu=(a,a+e)$ and $x\in B_\nu$.  Put
\[
 r_-=x-a,\qquad r_+=a+e-x.
\]
The two good cells adjacent to $a$ and $a+e$ have lengths comparable
with $g$.  On either adjacent cell the profile
$\psi-\overline\psi$ has the common endpoint value
$-\overline\psi$; its difference from that endpoint is $\psi$ and
vanishes quadratically.  Since the periodic Hilbert kernel
$K(t)=(2\pi)^{-1}\cot(t/2)$ is odd, pairing the two adjacent cells gives
\begin{equation}\label{ZG:eq:paired-log}
 \left|\int_{I_- \cup I_+}K(x-y)r_{\rm G}(y)\,dy\right|
 \le C_\rho g^2
 \left(1+\left|\log\frac{r_-}{r_+}\right|\right).
\end{equation}
Indeed, the two constant endpoint pieces are
\[
 -g^2\overline\psi
 \left\{\int_0^{h_-}K(r_-+u)\,du
       -\int_0^{h_+}K(r_++u)\,du\right\}.
\]
Lift the gap and its two adjacent cells to the real arc; this is
unambiguous because $E\le\pi/4$ and the fan is sufficiently fine.  Then
\[
 \int_0^hK(t+u)\,du
 =\frac1\pi\log\frac{\sin((t+h)/2)}{\sin(t/2)}
\]
exactly.  On this arc $\sin(v/2)\asymp v$, while
$h_\pm\asymp_\rho g$.  Moreover
$u\mapsto\log(1+e^{-u})$ is $1$-Lipschitz.  Applying this with
$u=\log(t/h)$ shows that the difference of the two logarithms is bounded
by $C_\rho+|\log(r_-/r_+)|$, which proves the constant part of
\eqref{ZG:eq:paired-log}.  For the two remaining pieces use
$\psi(u/h)\le C\min\{(u/h)^2,((h-u)/h)^2\}$; their integrals against
$|K(t+u)|$ are uniformly bounded.

Every other good cell $J$ is at distance at least $c_\rho g$ from $x$.
On $J$ the function $r_{\rm G}$ has zero mass and zero first moment
about the midpoint.  Two Taylor subtractions and
$|K''(t)|\le C\operatorname {dist}(t,2\pi\mathbb Z)^{-3}$ give
\[
 \left|\int_JK(x-y)r_{\rm G}(y)\,dy\right|
 \le \frac{C_\rho g^5}{\operatorname {dist}(x,J)^3}.
\]
The $k$th good cell in either cyclic direction is at distance at least
$c_\rho kg$; summing $k^{-3}$ and combining with
\eqref{ZG:eq:paired-log} yields
\begin{equation}\label{ZG:eq:gap-pointwise}
 |\Hil r_{\rm G}(x)|
 \le C_\rho g^2
 \left(1+\left|\log\frac{r_-}{r_+}\right|\right).
\end{equation}

Finally, for an arbitrary subdivision of $(a,a+e)$ into target cells
with midpoints $c_i$,
\begin{equation}\label{ZG:eq:midpoint-log-square}
 \sum_{I_i\subset B_\nu}h_i
 \left(1+\left|\log
 \frac{c_i-a}{a+e-c_i}\right|\right)^2
 \le Ce.
\end{equation}
To see this, scale to $(0,1)$.  On either half interval the displayed
profile is monotone, so the value at a cell midpoint is bounded by twice
its integral over the half of that cell nearer the corresponding
endpoint.  A cell crossing $1/2$ contributes $O(h_i)$.  The remaining
claim is the elementary finite integral
\[
 \int_0^1\left(1+\left|\log\frac{s}{1-s}\right|\right)^2ds<\infty.
\]
Equations \eqref{ZG:eq:gap-pointwise}--\eqref{ZG:eq:midpoint-log-square},
summed over the disjoint gaps, prove
\begin{equation}\label{ZG:eq:oscillation-gap}
 \sum_{i\in\mathcal B}h_i|\Hil r_{\rm G}(c_i)|^2
 \le C_\rho g^4E.
\end{equation}
Combining \eqref{ZG:eq:amplitude-error}, \eqref{ZG:eq:mean-gap}, and
\eqref{ZG:eq:oscillation-gap}, and using $g\asymp_\rho b$ and
$h_j-g=\widehat h_j-b$, proves \eqref{ZG:eq:gap-target}.
\end{proof}

We can now state the corrected replacement for the stage 155 opening
square.

\begin{lemma}[Zero-gap skeleton square]\label{ZG:lem:skeleton}
Let good source
cells have lengths comparable with $b$, let disjoint
zero-source gaps have total length $E$, and let the target measure put
the actual cell mass at every good or bad midpoint.  If the configuration
is obtained from the collapsed regular $M$-fan by varying the good
lengths to $h_i$ and opening the gaps, then
\begin{equation}\label{ZG:eq:skeleton}
 \boxed{
 \sum_i h_i|\Hil q_{\rm G}(c_i)|^2
 \le C_\rho\left\{
 b^3\sum_{j\in\mathcal G}(h_j-b)^2+b^4E\right\}.}
\end{equation}
The constant is independent of the number of gaps, the number of
target cells inside a gap, and every gap ratio.
\end{lemma}

\begin{proof}
Let $\gamma=E/M$ and $\widehat h_j=h_j+\gamma$.  Then
$\sum_{j\in\mathcal G}\widehat h_j=2\pi$ and
\[
 \sum_{j\in\mathcal G}(\widehat h_j-b)^2
 =\sum_{j\in\mathcal G}(h_j-b)^2-\frac{E^2}{M}
 \le\sum_{j\in\mathcal G}(h_j-b)^2.
\]
The complete fan $\widehat h$ is comparable with $b$.  The variable-grid
$T1$ theorem of stage 154, applied on that complete fan, gives
\[
 \sum_{i\in\mathcal G}h_i
 |\Hil\widehat q(\widehat c_i)|^2
 \le C_\rho b^3
 \sum_{j\in\mathcal G}(h_j-b)^2,
\]
because $h_i\asymp_\rho b$.  This is precisely the square of the
reference vector on the good target atoms.  Lemma~\ref{ZG:lem:opening}
bounds the good-target opening error by $C_\rho b^4E$, and
Lemma~\ref{ZG:lem:gap-target} bounds all bad target atoms.  Since
\[
 \sum_{j\in\mathcal G}(\widehat h_j-b)^2
 \le\sum_{j\in\mathcal G}(h_j-b)^2,
\]
the inequality $|x+y|^2\le2|x|^2+2|y|^2$ proves
\eqref{ZG:eq:skeleton}.
\end{proof}

\subsection{AVS, HQ, and DF}

\begin{theorem}[Adaptive variance sampling]
\label{ZG:thm:AVS}
For every
positive sufficiently fine fan,
\begin{equation}\label{ZG:eq:AVS}
 \boxed{
 T(h)\le C\Dvar(h)\le CS\Jfour(h).}
\end{equation}
\end{theorem}

\begin{proof}
The far-stratum region is Theorem 4.1 of stage 153.  In the near region,
combine \eqref{ZG:eq:Tsplit}, Lemma~\ref{ZG:lem:bad}, and
Lemma~\ref{ZG:lem:skeleton}, and then apply the static budget
\eqref{ZG:eq:budget}.  This gives $T\le C_\rho\Dvar$.  Stage 153 proves
$\Dvar\le S\Jfour$.
\end{proof}

\begin{corollary}[Pure midpoint row DF]\label{ZG:cor:DF}
The pure cubic Bregman row DF holds unconditionally.
\end{corollary}

\begin{proof}
Stage 144 gives the exact identity
\[
 \Cnull(q_h^3)-\Cnull(q_{a\mathbf1}^3)
 =\frac1{720}\left(\sum_i h_i^5-Na^5\right)
  +\frac1{576}\Dvar(h)-T(h),
 \qquad a=\frac{2\pi}{N}.
\]
The linear part of the fifth-power Bregman difference cancels because
$\sum_i(h_i-a)=0$.  Taylor's formula and $h_i,a\le S$ give
\[
 0\le\sum_i h_i^5-Na^5
 \le CS\sum_i(h_i-a)^2(h_i^2+2ah_i+3a^2)
 =CS\Jfour(h).
\]
Also $\Dvar\le S\Jfour$ by stage 153, and
$0\le T\le C\Dvar$ by Theorem~\ref{ZG:thm:AVS}.  Taking absolute values in
the exact identity proves DF with no omitted sign case.
\end{proof}
\endgroup


\part{The mixed support row and exact support-Hessian assembly}

\begingroup
\section{Node PT: physical material transport and the mixed row DM}
\label{module:PT}
\noindent\textit{Audited source: \nolinkurl{PT_material_transport.tex}.}
\subsection{The row and the two regimes}

Let
\begin{equation}\label{PT:eq:fan}
 h_i>0,\qquad \sum_i h_i=2\pi,\qquad
 a=\frac{2\pi}{N},\qquad S=a+\max_i h_i,
\end{equation}
and put
\begin{equation}\label{PT:eq:J}
 \Jfour(h)=\sum_i h_i^4-Na^4.
\end{equation}
On $I_i=[\theta_i,\theta_{i+1}]$ define
\begin{equation}\label{PT:eq:q}
 q_h(\theta_i+x)=\frac12\min\{x^2,(h_i-x)^2\},
 \qquad0\le x\le h_i.
\end{equation}

Let $Z_h$ be the exact area tangent space of support increments in a fixed
translation gauge, and let $v_h(z)$ be the physical-affine normal trace.
On
\[
 K_i=[\theta_i-h_{i-1}/2,\theta_i+h_i/2],\qquad x=\theta-\theta_i,
\]
it is
\begin{align}
 v_h(z)(\theta)
 &=\PiNG\left[z_i\cos x+
 \{(1-t_i(x))A_i^-+t_i(x)A_i^+\}\sin x\right],
 \label{PT:eq:trace}\\
 t_i(x)&=\frac{\tan x+\tan(h_{i-1}/2)}
 {\tan(h_{i-1}/2)+\tan(h_i/2)},\notag\\
 A_i^-&=\frac{z_i-z_{i-1}}{\sin h_{i-1}}
       -z_i\tan\frac{h_{i-1}}2,\qquad
 A_i^+=\frac{z_{i+1}-z_i}{\sin h_i}
       +z_i\tan\frac{h_i}2.\notag
\end{align}
The exact support metric satisfies
\begin{equation}\label{PT:eq:metric}
 c\|v_h(z)\|_{H^{1/2}}^2
 \le\Gtilde_h(z)\le C\|v_h(z)\|_{H^{1/2}}^2.
\end{equation}

The symmetric null form is
\begin{equation}\label{PT:eq:null}
 \Cnull(f_1,f_2,f_3)
 =\frac13\sum_{\rm cyc}\int_\T f_1
 \{(|D|f_2)(|D|f_3)-f_2'f_3'\}\,d\theta.
\end{equation}
The row DM is
\begin{equation}\label{PT:eq:DM-target}
 |\Cnull(q_h,q_h,v_h(z))|
 \le\varepsilon_{\rm DM}(S)\Jfour(h)^{1/2}
       \Gtilde_h(z)^{1/2},
 \qquad\varepsilon_{\rm DM}(S)\longrightarrow0.
\end{equation}

Fix $\beta=1/2$.  As in stage 161, split into
\begin{align}
 \Jfour(h)&\ge S^{5-\beta}&&\text{(far)},\label{PT:eq:far-regime}\\
 \Jfour(h)&<S^{5-\beta}&&\text{(near)}.\label{PT:eq:near-regime}
\end{align}
The fully polarized one-Lipschitz estimate
\begin{equation}\label{PT:eq:one-slot}
 |\Cnull(f,g,w)|
 \le C\|f'\|_\infty
       \|g\|_{\dot H^{1/2}}\|w\|_{\dot H^{1/2}}
\end{equation}
follows from the same-sign Fourier formula by putting the low derivative
block in $L^\infty$; when the Lipschitz slot is high, the remaining
separation factor is $2^{-|j-k|/2}$.  Since
\begin{equation}\label{PT:eq:q-size}
 \|q_h'\|_\infty\le S,
 \qquad \|q_h\|_{\dot H^{1/2}}\le CS^{3/2},
\end{equation}
the far regime gives
\begin{equation}\label{PT:eq:far-bound}
 |\Cnull(q_h,q_h,v_h(z))|
 \le CS^{1/4}\Jfour(h)^{1/2}\Gtilde_h(z)^{1/2}.
\end{equation}

In the near regime set
\begin{equation}\label{PT:eq:path}
 d_i=h_i-a,\qquad h_i(t)=a+td_i.
\end{equation}
The exact identity
\begin{equation}\label{PT:eq:Jidentity}
 \Jfour(h)=\sum_i d_i^2(h_i^2+2ah_i+3a^2)
\end{equation}
gives $|d_i|^4\le C\Jfour(h)$.  In the near regime,
$\max_i|d_i|\le CS^{(5-\beta)/4}=CS^{9/8}=o(S)$ as $S\to0$.
Since $S=a+\max_i h_i$ and $|h_i-a|\le\max_j|d_j|$, both $a$ and every
$h_i(t)$ therefore lie between fixed positive multiples of $S$, for all
sufficiently small $S$.  Hence
\begin{equation}\label{PT:eq:comparable}
 cS\le a\le h_i(t)\le CS,
 \qquad
 \delta(t):=\max_i\frac{|d_i|}{h_i(t)}
 \le C\frac{\max_i|d_i|}{S}.
\end{equation}

\subsection{Lipschitz covariance of the same-sign null form}

For a periodic Lipschitz function $u$, put
$F_\varepsilon(x)=x+\varepsilon u(x)$ and
\begin{equation}\label{PT:eq:transport-derivative}
 \mathfrak D_u\Cnull(f,g,w)
 =\frac d{d\varepsilon}\bigg|_{\varepsilon=0}
 \Cnull(f\circ F_\varepsilon^{-1},
        g\circ F_\varepsilon^{-1},
        w\circ F_\varepsilon^{-1}).
\end{equation}

\begin{lemma}[Critical Lipschitz covariance]\label{PT:lem:covariance}
For real $f\in W^{1,\infty}(\T)$ and real
$g,w\in H^{1/2}(\T)$,
\begin{equation}\label{PT:eq:covariance}
 \boxed{
 |\mathfrak D_u\Cnull(f,g,w)|
 \le C\|u'\|_\infty\|f'\|_\infty
       \|g\|_{H^{1/2}}\|w\|_{H^{1/2}}.}
\end{equation}
The polarized estimate is uniform for piecewise-affine $u$.
\end{lemma}

\begin{proof}
First take trigonometric polynomials.  The same-sign representation of
\eqref{PT:eq:null} is a sum, over $p,q>0$, of the three placements of
\begin{equation}\label{PT:eq:same-sign}
 pq\,\widehat f_p\widehat g_q\overline{\widehat w_{p+q}}
\end{equation}
and their conjugates.  Since
\[
 \frac d{d\varepsilon}\bigg|_0
 (f\circ F_\varepsilon^{-1})=-uf',
\]
differentiating the three slots in \eqref{PT:eq:same-sign} produces a
quadrilinear multiplier with the deformation frequency $\ell$.  At
$\ell=0$ the three contributions add to zero: this is exactly invariance
under a common translation.  Therefore the summed multiplier is divisible
by $\ell$, and $\ell\widehat u_\ell$ is a Fourier coefficient of $u'$.

For completeness, this divisibility has a pointwise symbol bound.  Put
\[
 \tau(r,s,t)=\sum_{\rm cyc}(|s||t|+st),
 \qquad r+s+t=0.
\]
Thus $\tau$ is twice the product of the two same-sign frequencies.  If
$\ell+k+m+n=0$, direct substitution in the six possible sign orders gives
\begin{multline}\label{PT:eq:symbol-bound}
 \big|k\tau(k+\ell,m,n)+m\tau(k,m+\ell,n)\\
       +n\tau(k,m,n+\ell)\big|
 \le3|\ell|\,|k|\,|m|^{1/2}|n|^{1/2}.
\end{multline}
For example, on a region with $m,n>0$ the only changes occur when
$k+\ell$, $m+\ell$, or $n+\ell$ crosses zero; on each resulting interval
the left side is affine in $\ell$ after its zero value at $\ell=0$ is
removed, and its two endpoint slopes are bounded by the right side.
The other sign orders are permutations.  This proves
\eqref{PT:eq:symbol-bound} without a smoothness assumption on $u$.

We spell out the endpoint summation, since the pointwise bound alone is
not an $L^\infty\times L^\infty\times L^2\times L^2$ theorem.  Let
$P_j$ be cyclic Littlewood--Paley projections and write
\[
 U_j=P_j(u'),\quad F_j=P_j(f'),\quad
 G_j=|D|^{1/2}P_jg,\quad W_j=|D|^{1/2}P_jw.
\]
Split each dyadic block further at the six sign boundaries
\[
 k=0,\quad m=0,\quad n=0,\quad
 k+\ell=0,\quad m+\ell=0,\quad n+\ell=0.
\]
On every resulting interval the divided symbol
\begin{equation}\label{PT:eq:divided-symbol}
 \frac{k\tau(k+\ell,m,n)+m\tau(k,m+\ell,n)
       +n\tau(k,m,n+\ell)}
      {\ell k|m|^{1/2}|n|^{1/2}}
\end{equation}
and each discrete difference, multiplied by the corresponding dyadic
scale, are bounded.  At a boundary the one-sided limits agree with the
formula in the adjacent sign chamber; if one of $k,m,n$ is zero, the
undivided symbol is zero.  Thus no boundary atom or zero-frequency row is
lost.

Let $J$ be the largest dyadic index and recall that the frequency sum
forces the two largest indices to differ by at most $3$.  Up to exchanging
$g,w$ and $u',f'$, all flag rows have the following bounds:
on the chambers where one of $G,W$ is $r$ generations below a high
$U$ or $F$ block, direct substitution in the same-sign product gives an
extra $2^{-r/2}$; when both $G,W$ are lower by $r,s$ and $U,F$ are high,
the divided symbol is bounded by $C2^{-(r+s)/2}$.  For example, if
$k>0$, $m,n>0$, and $\ell<0$ are in the latter configuration, then
$k+\ell=-(m+n)$ and the numerator in
\eqref{PT:eq:divided-symbol} is $6kmn$, so division gives
$C|mn|^{1/2}/|\ell|$.  The other five sign orders are identical after a
permutation.  Thus all flag rows have the following bounds:
\begin{center}
\begin{tabular}{c|c}
two high blocks & bound after the order-zero kernel is summed\\
\hline
$G_J,W_J$ &
$C\|U_{\le J}\|_\infty\|F_{\le J}\|_\infty
  \|G_J\|_2\|W_{J+O(1)}\|_2$\\
$U_J,G_J$ &
$C\sum_{r\ge0}2^{-r/2}\|U_J\|_\infty
  \|F_{\le J}\|_\infty\|G_{J+O(1)}\|_2\|W_{J-r}\|_2$\\
$F_J,G_J$ &
$C\sum_{r\ge0}2^{-r/2}\|U_{\le J}\|_\infty
  \|F_J\|_\infty\|G_{J+O(1)}\|_2\|W_{J-r}\|_2$\\
$U_J,F_J$ &
$C\sum_{r,s\ge0}2^{-(r+s)/2}\|U_J\|_\infty
  \|F_{J+O(1)}\|_\infty\|G_{J-r}\|_2\|W_{J-s}\|_2$.
\end{tabular}
\end{center}
The factors $2^{-r/2}$ are the unused half derivative in
\eqref{PT:eq:symbol-bound}.  The same table holds when a shifted frequency
touches a sign boundary, using the one-sided divided symbol above.

Each Littlewood--Paley block is bounded in $L^\infty$ by the full
$L^\infty$ norm.  Sum first in $r,s$ and then in $J$.  Young's inequality
for the geometric sequences and Cauchy--Schwarz for the $G,W$ square
functions give
\[
 C\|u'\|_\infty\|f'\|_\infty
 \left(\sum_j\|G_j\|_2^2\right)^{1/2}
 \left(\sum_j\|W_j\|_2^2\right)^{1/2}.
\]
The finitely many zero and unit frequencies are bounded by the
inhomogeneous $H^{1/2}$ norms.  This is \eqref{PT:eq:covariance}.

Approximate a Lipschitz $u$ by smooth functions with uniformly bounded
derivative.  The estimate just proved is stable under this approximation,
so it applies in particular to piecewise-affine velocities.
\end{proof}

\begin{remark}
Lemma~\ref{PT:lem:covariance} measures $u'$ and not $u$.  Consequently a
large cumulative displacement of grid nodes and a common rotation do not
enter the estimate.  This is the precise replacement for the invalid
bare-kernel order used in stage 161.
\end{remark}

\subsection{Exact material fan identity and support transport}

Refine the fan mesh by the nodes $\theta_i(t)$ and the midpoints
$c_i(t)=(\theta_i(t)+\theta_{i+1}(t))/2$.  Let
$F_t:\T_0\to\T_t$ be the unique orientation-preserving map which is
affine on every reference half-cell and sends
\(\theta_i(0),c_i(0)\) to \(\theta_i(t),c_i(t)\).  Set
\begin{equation}\label{PT:eq:finite-material-map}
 G_{t,\varepsilon}=F_{t+\varepsilon}\circ F_t^{-1},
 \qquad
 u_t=\partial_\varepsilon G_{t,\varepsilon}\big|_{\varepsilon=0}.
\end{equation}
On the two halves of $I_i(t)$ one has
\begin{equation}\label{PT:eq:uprime}
 u_t'=\frac{d_i}{h_i(t)},
 \qquad \|u_t'\|_\infty=\delta(t).
\end{equation}

\begin{lemma}[Exact fan transport]\label{PT:lem:q-transport}
On both halves of $I_i(t)$,
\begin{equation}\label{PT:eq:q-transport}
 q_{h(t+\varepsilon)}\circ G_{t,\varepsilon}
 =\left(1+\varepsilon\frac{d_i}{h_i(t)}\right)^2q_{h(t)}.
\end{equation}
Consequently, if a dot denotes the derivative of the pulled profile,
\begin{equation}\label{PT:eq:qdot}
 \|\dot q_t\|_{H^{1/2}}
 \le C\delta(t)S^{3/2}.
\end{equation}
\end{lemma}

\begin{proof}
On either half-cell, $q_h$ is one half of the squared distance from its
nearest endpoint.  The affine map multiplies that distance by
$1+\varepsilon d_i/h_i(t)$, proving \eqref{PT:eq:q-transport}.

Thus $\dot q_t=2(d_i/h_i(t))q_{h(t)}$ on $I_i(t)$.  Each cell piece
vanishes quadratically at its outer endpoints.  Its zero extension has
$H^{1/2}$ norm squared at most $Ch_i(t)^4|d_i/h_i(t)|^2$ by scaling the
fixed unit-cell profile.  To sum the pieces, use comparability from
\eqref{PT:eq:comparable}.  The $H^{1/2}$ pairing of two nonadjacent cell
profiles at cyclic index distance $r$ is bounded by
$CS^4|a_i a_j|/(1+r)^2$, where $a_i=d_i/h_i(t)$; adjacent pairs are
bounded by the two-cell unit-profile norm.  Schur summation of
$(1+r)^{-2}$ gives
\[
 \left\|\sum_i2a_iq_i\right\|_{H^{1/2}}^2
 \le C\sum_i S^4|a_i|^2.
\]
Hence
\[
 \|\dot q_t\|_{H^{1/2}}^2
 \le C\delta(t)^2\sum_i h_i(t)^4
 \le C\delta(t)^2S^3,
\]
because $\sum_i h_i(t)=2\pi$ and $h_i(t)\le CS$.
\end{proof}

We next record the only finite-element statement about the physical
support trace.  The estimate is on its natural support space, not on an
arbitrary ambient target.

\begin{lemma}[Stable physical material transport]\label{PT:lem:v-transport}
Along every comparable path \eqref{PT:eq:comparable} there is a linear
transport $z_1\mapsto z_t\in Z_{h(t)}$, preserving the exact area and
translation quotient, such that, with
\begin{equation}\label{PT:eq:vtilde}
 \widetilde v_t=v_{h(t)}(z_t)\circ F_t,
\end{equation}
interpreted infinitesimally on the common refined grid,
\begin{align}
 \sup_{0\le t\le1}\|v_{h(t)}(z_t)\|_{H^{1/2}}
 &\le C\|v_h(z_1)\|_{H^{1/2}},\label{PT:eq:v-uniform}\\
 \|\dot{\widetilde v}_t\|_{H^{1/2}}
 &\le C\delta(t)\|v_{h(t)}(z_t)\|_{H^{1/2}}.
 \label{PT:eq:vdot}
\end{align}
\end{lemma}

\begin{proof}
We include the discrete endpoint estimate which prevents a hidden
mesh-frequency loss.  On a comparable cyclic mesh define
\begin{equation}\label{PT:eq:discrete-H}
 \|z\|_{\mathfrak h^{1/2}_S}^2
 =S\sum_i|z_i|^2
 +\sum_{i\ne j}\frac{|z_i-z_j|^2}{(1+|i-j|_{\rm cyc})^2}.
\end{equation}

First let $I_hz$ be the affine interpolant with nodal values $z_i$.
For non-neighbouring cells $I_i,I_j$ at cyclic distance $r$,
comparability gives
\[
 |x-y|\asymp S(1+r),\qquad |I_i|\asymp|I_j|\asymp S.
\]
The corresponding Gagliardo rectangle is therefore comparable with
$|\bar z_i-\bar z_j|^2/(1+r)^2$, where $\bar z_i$ is the cell average.
The two-cell rectangles give $|z_{i+1}-z_i|^2$.  Since
\[
 |\bar z_i-z_i|\le\tfrac12|z_{i+1}-z_i|,
\]
the triangle inequality converts nodal differences to average
differences and conversely, with an error paid by these neighbouring
rectangles.  The summability of $(1+r)^{-2}$ pays each such error a
bounded number of times.
Together with
$\|I_hz\|_2^2\asymp S\sum_i|z_i|^2$, summing the rectangles in both
directions proves
\begin{equation}\label{PT:eq:affine-discrete}
 \|I_hz\|_{H^{1/2}}^2\asymp
 \|z\|_{\mathfrak h^{1/2}_S}^2.
\end{equation}

Now pull every patch $K_i$ in \eqref{PT:eq:trace} to $[0,1]$.  With
$w_i=(h_{i-1}+h_i)/2$, $p_i=h_{i-1}/(h_{i-1}+h_i)$, and
\[
 Z=z_i,\qquad A=w_iA_i^-,\qquad B=w_iA_i^+,
\]
the normalized trace is
\[
 Z\cos\xi+\{(1-y)A+yB\}\frac{\sin\xi}{w_i},
 \qquad \xi=w_i(y-p_i),
\]
where replacing $t_i$ by $y$ changes the coefficient functions by
$O(S^2)$ in $C^1$.  Comparability places $p_i$ in a fixed compact
subinterval of $(0,1)$.  As $w_i\to0$, the local three-node map tends to
\begin{equation}\label{PT:eq:limiting-local-trace}
 (z_{i-1},Z,z_{i+1})\longmapsto
 Z+(y-p_i)\left\{
 (1-y)\frac{Z-z_{i-1}}{2p_i}
 +y\frac{z_{i+1}-Z}{2(1-p_i)}\right\}.
\end{equation}
This limiting map is injective: evaluation at $y=p_i$ first gives
$Z=0$, and division by $y-p_i$ then gives an affine function whose two
endpoint coefficients force $z_{i-1}=z_{i+1}=0$.  Finite-dimensional
norm equivalence, compactness in $p_i$, and an $O(S^2)$ perturbation
therefore give uniform upper and lower bounds for the local map and its
inverse.  The physical trace and the affine interpolant are consequently
related by uniformly invertible three-point finite-element maps.
Applying the preceding rectangle argument to these local maps gives
\begin{equation}\label{PT:eq:discrete-equivalence}
 \|v_h(z)\|_{H^{1/2}}^2
 \asymp\|z\|_{\mathfrak h^{1/2}_S}^2
\end{equation}
on the exact quotient.  Indeed, the area row controls the constant mode,
while support translations are exactly the first two harmonics and are
removed by $\PiNG$.  Hence no low-mode kernel remains in either direction.

We next differentiate the actual coefficients.  Put
\[
 b_i^- =\frac{h_{i-1}}2,\quad b_i^+=\frac{h_i}2,\quad
 D_i=\tan b_i^-+\tan b_i^+,\quad
 N_i(x)=\tan x+\tan b_i^-.
\]
For the pulled local coordinate $x=x(t)$, direct differentiation gives
\begin{align}
 \dot t_i(x)
 &=\frac{\{\sec^2x\,\dot x+\tfrac12\sec^2b_i^-d_{i-1}\}D_i
       -N_i(x)\{\tfrac12\sec^2b_i^-d_{i-1}
                  +\tfrac12\sec^2b_i^+d_i\}}
       {D_i^2},\label{PT:eq:tdot}\\
 \dot A_i^-
 &=\frac{\dot z_i-\dot z_{i-1}}{\sin h_{i-1}}
   -\frac{(z_i-z_{i-1})\cos h_{i-1}}{\sin^2h_{i-1}}d_{i-1}\\
 &\hspace{12mm}-\dot z_i\tan\frac{h_{i-1}}2
   -\frac12z_i\sec^2\frac{h_{i-1}}2d_{i-1},
 \label{PT:eq:Adot-minus}\\
 \dot A_i^+
 &=\frac{\dot z_{i+1}-\dot z_i}{\sin h_i}
   -\frac{(z_{i+1}-z_i)\cos h_i}{\sin^2h_i}d_i\\
 &\hspace{12mm}+\dot z_i\tan\frac{h_i}2
   +\frac12z_i\sec^2\frac{h_i}2d_i.
 \label{PT:eq:Adot-plus}
\end{align}
Moreover
\begin{align}
 \partial_t\cos x&=-\dot x\sin x,\qquad
 \partial_t\sin x=\dot x\cos x,\label{PT:eq:frame-dots}\\
 |\dot x|&\le C\delta(t)S,\qquad
 |\dot t_i|\le C\delta(t).
\end{align}
Substitution of \eqref{PT:eq:tdot}--\eqref{PT:eq:frame-dots} in
\eqref{PT:eq:trace} shows that every pulled trace derivative is a bounded
three-point finite-element combination of $\dot z$ and
\begin{equation}\label{PT:eq:transport-ledger}
 S a_i z_i,\qquad
 a_i(z_{i+1}-z_i),\qquad
 a_{i-1}(z_i-z_{i-1}),
 \qquad |a_i|\le\delta(t),
\end{equation}
and their one-cell shifts.  There is no unweighted term $a_i z_i$: the
trigonometric frame contributes $d_i=S a_i$, while every derivative of a
normalized slope is paired with a nodal difference.  The two one-sided
wall terms cancel because the physical affine vector field is continuous.
Differentiating $\PiNG$ adds only the modes $1,\cos\theta,\sin\theta$ and
their pulled derivatives $u_t\sin\theta,u_t\cos\theta$.  After subtracting
the irrelevant common rotation, periodic Poincar\'e gives
$\|u_t\|_\infty\le C\|u_t'\|_\infty$; these fixed-dimensional rows obey
the same bound as \eqref{PT:eq:transport-ledger} and create no factor of $N$.

The multiplier bound can be proved without a Fourier endpoint theorem.
Write $Dz_i=z_{i+1}-z_i$ and let $K(i,j)=(1+|i-j|_{\rm cyc})^{-2}$.
For $w_i=Sa_iz_i$,
\begin{align*}
 S\sum_i|w_i|^2&\le C\|a\|_\infty^2S\sum_i|z_i|^2,\\
 \sum_{i\ne j}K(i,j)|w_i-w_j|^2
 &\le2S^2\|a\|_\infty^2
      \sum_{i\ne j}K(i,j)|z_i-z_j|^2\\
 &\quad+8S^2\|a\|_\infty^2
      \sum_j|z_j|^2\sum_iK(i,j).
\end{align*}
Since $\sup_j\sum_iK(i,j)<\infty$ and $S$ is bounded, this is controlled
by \eqref{PT:eq:discrete-H}.  For $w_i=a_iDz_i$ use
\[
 S\sum_i|w_i|^2
 \le S\|a\|_\infty^2\sum_i|Dz_i|^2
 \le C\|a\|_\infty^2\|z\|_{\mathfrak h^{1/2}_S}^2
\]
and
\[
 |w_i-w_j|^2
 \le2\|a\|_\infty^2|Dz_i-Dz_j|^2
     +8\|a\|_\infty^2|Dz_j|^2
\]
and
\[
 |Dz_i-Dz_j|^2
 \le2|z_{i+1}-z_{j+1}|^2+2|z_i-z_j|^2.
\]
The last row is controlled by the $|i-j|_{\rm cyc}=1$ part of
\eqref{PT:eq:discrete-H}.  These estimates prove
\begin{equation}\label{PT:eq:ledger-bound}
 \|Sa z+a\Delta z\|_{\mathfrak h^{1/2}_S}
 \le C\|a\|_{\ell^\infty}
       \|z\|_{\mathfrak h^{1/2}_S}.
\end{equation}

It remains to define, rather than merely invoke, the constraint
transport.  Let
\begin{equation}\label{PT:eq:area-row}
 L_tz=\sum_i\ell_i(t)z_i,\qquad
 \ell_i(t)=\tan\frac{h_{i-1}(t)}2+\tan\frac{h_i(t)}2,
\end{equation}
and let
\[
 \tau_t^1=(\cos\theta_i(t))_i,\qquad
 \tau_t^2=(\sin\theta_i(t))_i.
\]
These are the exact two support-translation vectors and satisfy
$L_t\tau_t^\mu=0$.  The matrix
\begin{equation}\label{PT:eq:translation-Gram}
 M_t^{\mu\nu}=\sum_i\ell_i(t)\tau_{t,i}^\mu\tau_{t,i}^\nu
\end{equation}
is uniformly comparable with the identity.  Define
\begin{align}
 \mathcal Q_tz
 &=z-\sum_{\mu,\nu=1}^2(M_t^{-1})_{\mu\nu}
 \left(\sum_i\ell_i(t)z_i\tau_{t,i}^\nu\right)\tau_t^\mu,
 \label{PT:eq:translation-projection}\\
 C_tz
 &=\mathcal Q_t\left(z-\frac{L_tz}{L_t\mathbf1}\mathbf1\right).
 \label{PT:eq:constraint-transport}
\end{align}
If $z_1$ is in the chosen exact gauge, set
\begin{equation}\label{PT:eq:z-transport}
 z_t=C_tz_1.
\end{equation}
Then $L_tz_t=0$, $z_t$ is orthogonal to both translation rows in the
$\ell(t)$-weighted Gram product, and $C_1z_1=z_1$.

The explicit rows give
\begin{align}
 \ell_i&\asymp S,\qquad
 |\dot\ell_i|\le C\delta(t)S,\label{PT:eq:constraint-row-bounds}\\
 \|\dot\tau_t^\mu\|_{\mathfrak h^{1/2}_S}
 &\le C\delta(t),\qquad
 |\dot M_t|+|\partial_tM_t^{-1}|\le C\delta(t).
\end{align}
Indeed, $|\dot\theta_i|=|\sum_{k<i}d_k|\le C\delta(t)$ and adjacent
differences of $\dot\theta_i$ equal $d_i$.  Differentiating
\eqref{PT:eq:translation-projection}--\eqref{PT:eq:constraint-transport}, and
using the bounded area and translation functionals furnished by
\eqref{PT:eq:discrete-equivalence}, first yields
\begin{equation}\label{PT:eq:Ct-raw-bounds}
 \|C_t\|\le C,
 \qquad
 \|\dot C_tz_1\|_{\mathfrak h^{1/2}_S}
 \le C\delta(t)\|z_1\|_{\mathfrak h^{1/2}_S}.
\end{equation}
On $Z_{h(1)}$ the terminal map $C_1$ is the identity.  Hence
\[
 \|(C_t-C_1)|_{Z_{h(1)}}\|
 \le C\int_t^1\delta(s)\,ds
 \le CS^{1/8}.
\]
For sufficiently small $S$ this norm is at most $1/2$.  The Neumann
series therefore proves that $C_t|_{Z_{h(1)}}$ is injective onto its
image and that its inverse is uniformly bounded.  Combining this inverse
bound with \eqref{PT:eq:Ct-raw-bounds} gives
\begin{equation}\label{PT:eq:Ct-bounds}
 \|C_t\|+\|(C_t|_{Z_{h(1)}})^{-1}\|\le C,\qquad
 \|\dot C_tz_1\|_{\mathfrak h^{1/2}_S}
 \le C\delta(t)\|C_tz_1\|_{\mathfrak h^{1/2}_S}.
\end{equation}

Equations \eqref{PT:eq:discrete-equivalence},
\eqref{PT:eq:ledger-bound}, and \eqref{PT:eq:Ct-bounds} now give
\[
 \|\dot{\widetilde v}_t\|_{H^{1/2}}
 \le C\delta(t)\|v_{h(t)}(z_t)\|_{H^{1/2}}.
\]
The finite maps $F_t$ and $F_t^{-1}$ have uniformly bounded Lipschitz
constants by \eqref{PT:eq:comparable}, so their pullbacks are uniformly
bounded on $H^{1/2}$.  Finally,
\begin{equation}\label{PT:eq:delta-integral}
 \int_0^1\delta(t)\,dt
 \le C\frac{\max_i|d_i|}{S}\le CS^{1/8}
\end{equation}
in the near regime.  Gronwall's inequality proves
\eqref{PT:eq:v-uniform}, and the displayed derivative estimate is
\eqref{PT:eq:vdot}.
\end{proof}

\subsection{The near material estimate}

\begin{proposition}[Comparable Lipschitz material bound]
\label{PT:prop:material}
Along the near path \eqref{PT:eq:path} and the transport of
Lemma~\ref{PT:lem:v-transport},
\begin{equation}\label{PT:eq:material}
 \boxed{
 \left|\frac d{dt}
 \Cnull(q_{h(t)},q_{h(t)},v_{h(t)}(z_t))\right|
 \le CS^{5/2}\delta(t)\Gtilde_h(z_1)^{1/2}.}
\end{equation}
\end{proposition}

\begin{proof}
Pull all three functions by the common refined-grid map.  The derivative
splits into the common-coordinate row and three profile rows.
Lemma~\ref{PT:lem:covariance}, \eqref{PT:eq:q-size},
\eqref{PT:eq:uprime}, and \eqref{PT:eq:v-uniform} bound the coordinate row by
the right side of \eqref{PT:eq:material}.

For either fan-profile row use the polarized estimate
\eqref{PT:eq:one-slot}, placing the Lipschitz derivative on the unchanged
fan slot, and then use \eqref{PT:eq:qdot}:
\[
 |\Cnull(q_t,\dot q_t,v_t)|
 \le C S\,(\delta(t)S^{3/2})\|v_t\|_{H^{1/2}}.
\]
For the physical-profile row, use \eqref{PT:eq:one-slot},
\eqref{PT:eq:q-size}, and \eqref{PT:eq:vdot}:
\[
 |\Cnull(q_t,q_t,\dot{\widetilde v}_t)|
 \le C S\,S^{3/2}\delta(t)\|v_t\|_{H^{1/2}}.
\]
Metric equivalence \eqref{PT:eq:metric} completes the proof.
\end{proof}

\begin{lemma}[Regular quotient cancellation]\label{PT:lem:regular}
For the uniform fan and every exact support class,
\begin{equation}\label{PT:eq:regular}
 \Cnull(q_{a\mathbf1},q_{a\mathbf1},v_{a\mathbf1}(z))=0.
\end{equation}
\end{lemma}

\begin{proof}
The fan spline and the physical support construction commute with the
cyclic shift.  Hence the left side is a cyclic-invariant linear functional
of the support coefficients, and therefore a multiple of $\sum_i z_i$.
At the regular tangential polygon the exact area tangent condition is
$\sum_i z_i=0$; translations have already been removed.
\end{proof}

\begin{theorem}[Unconditional mixed row DM]\label{PT:thm:DM}
For every sufficiently fine positive fan and every exact support class,
\begin{equation}\label{PT:eq:DM}
 \boxed{
 |\Cnull(q_h,q_h,v_h(z))|
 \le C(S^{1/4}+S^{1/2})\Jfour(h)^{1/2}
       \Gtilde_h(z)^{1/2}.}
\end{equation}
In particular \eqref{PT:eq:DM-target} holds independently of every minimum
cell and adjacent ratio.
\end{theorem}

\begin{proof}
Use \eqref{PT:eq:far-bound} in the far regime.  In the near regime transport
the final support class to the regular quotient and integrate
\eqref{PT:eq:material}.  Equations \eqref{PT:eq:comparable} and
\eqref{PT:eq:Jidentity} give
\[
 \Jfour(h)^{1/2}\ge cS\left(\sum_i d_i^2\right)^{1/2}
 \ge cS\max_i|d_i|.
\]
Therefore
\begin{align*}
 |\Cnull(q_h,q_h,v_h(z))|
 &\le CS^{5/2}\frac{\max_i|d_i|}{S}
       \Gtilde_h(z)^{1/2}\\
 &\le CS^{1/2}\Jfour(h)^{1/2}\Gtilde_h(z)^{1/2},
\end{align*}
where Lemma~\ref{PT:lem:regular} supplies the zero initial value.  Combining
the regimes proves \eqref{PT:eq:DM}.
\end{proof}

\begin{corollary}[Handoff to the current NR assembly]\label{PT:cor:J0}
Theorem~\ref{PT:thm:DM} supplies the mixed terminal input DM.  The exact flux
audit shows, however, that DM does not imply the former JE--FM clause,
which remains unavailable; its former period-three material falsification
is retracted after the polygonal corner finite-part audit.  NR and active-chart J0 are
assembled conditionally in \texttt{JC\_joint\_comparison.tex} and
\texttt{NR\_assembly.tex}.  At the current audit DF, the CT core, and
JP--DE are closed, while JC--PL is the unique open premise and JP--JC is
its dependent exact output.  No historical CT
proposition is imported in that assembly, and none of these conclusions
uses a minimum fan, polar, or conformal cell.
\end{corollary}
\endgroup

\begingroup
\section{Node JP: exact support Hessian and the closed disk endpoint}
\label{module:JP}
\noindent\textit{Audited source: \nolinkurl{JP_joint_principal.tex}.}
\subsection{Fixed-normal support chart}

Fix cyclic unit normals $\nu_i$ with exterior angles
$\alpha_i\in(0,\pi)$ and $\sum_i\alpha_i=2\pi$.  For a support vector
$h=(h_i)$ put
\begin{equation}\label{JP:eq:P}
 P_\alpha(h)=\{x\in\mathbb R^2:x\cdot\nu_i<h_i\ \hbox{for all }i\},
 \qquad
 \Phi_\alpha(h)=\frac{|P_\alpha(h)|}{\pi}\lambda_1(P_\alpha(h)).
\end{equation}
We work in an active star on which all sides persist.  Translation vectors
$(b\cdot\nu_i)_i$ are quotiented exactly.  Although $\Phi_\alpha$ is also
dilation invariant, at a nonstationary point its Hessian does not descend
to an abstract dilation quotient.  We therefore fix an affine scale section
below.

Let $p_i(h)$ be the intersection of the two support lines with normals
$\nu_i,\nu_{i+1}$.  Solving the two linear equations shows that
\begin{equation}\label{JP:eq:vertex-linear}
 p_i(h+tz)=p_i(h)+t p_i(z).
\end{equation}
Consequently the physical-affine material parametrization of every side is
affine in $t$.  Its normal velocity on side $S_i(t)$ is exactly
\begin{equation}\label{JP:eq:constant-normal}
 V(t)\cdot\nu_i=z_i,\qquad \partial_\tau(V\cdot\nu_i)=0.
\end{equation}
This is the regular finite-energy replacement for the discontinuous
piecewise-constant function obtained by writing $z_i$ in the radial angular
coordinate.  That raw angular function is not itself an $H^{1/2}$ trace.

The area $A_\alpha(h)=|P_\alpha(h)|$ is a quadratic polynomial in $h$.
Indeed, by \eqref{JP:eq:vertex-linear},
\begin{equation}\label{JP:eq:area-jets}
\begin{aligned}
 A_\alpha(h)&=\frac12\sum_i\det(p_i(h),p_{i+1}(h)),\\
 A_\alpha'[h;z]
 &=\frac12\sum_i\{\det(p_i(z),p_{i+1}(h))
                 +\det(p_i(h),p_{i+1}(z))\},\\
 A_\alpha''[z,z]&=\sum_i\det(p_i(z),p_{i+1}(z)).
\end{aligned}
\end{equation}
Equivalently, direct differentiation of the side lengths gives the
ratio-sensitive but exact cyclic formula
\begin{equation}\label{JP:eq:area-second-explicit}
\begin{aligned}
 A_\alpha''[z,z]
 =\sum_i z_i\{&z_{i-1}\csc\alpha_{i-1}
               +z_{i+1}\csc\alpha_i\\
              &-z_i(\cot\alpha_{i-1}+\cot\alpha_i)\}.
\end{aligned}
\end{equation}
Equivalently,
\begin{equation}\label{JP:eq:area-jump}
 A_\alpha''[z,z]=\frac12\sum_i\left\{
 \tan\frac{\alpha_i}{2}(z_i+z_{i+1})^2
 -\cot\frac{\alpha_i}{2}(z_{i+1}-z_i)^2\right\}.
\end{equation}
Thus the area row contains a negative jump channel; it may not be bounded
from below independently of the endpoint part of $\Dred$.

Let $h_0=h_\alpha\mathbf1$ be the canonical tangential support vector,
$A_0=A_\alpha(h_0)$, and define the linear scale functional
\begin{equation}\label{JP:eq:scale-section}
 s_\alpha(h)=\frac{A_\alpha'[h_0;h]}{2A_0}.
\end{equation}
Euler's identity for the quadratic area gives $s_\alpha(h_0)=1$, while
translation invariance of area makes $s_\alpha$ vanish on translation
vectors.  Fix once and for all a translation complement
$\mathsf Y_\alpha$ which contains $h_0$, and put
\begin{equation}\label{JP:eq:Xalpha}
 \mathsf X_\alpha=\{z\in\mathsf Y_\alpha:s_\alpha(z)=0\}.
\end{equation}
Then $h_0+tz$ stays in the affine section $s_\alpha=1$ for every
$z\in\mathsf X_\alpha$.  All support variables, Hessians, dual norms, and
Riesz maps used in JP--SC/JP--JE/JP--FS refer to this fixed section.

To match the global wrapper, let $\mathcal U_{S_0,\rho_0}$ denote the
normalized radial endpoint chart: $S=a+\max_i\alpha_i\le S_0$, all sides
are active, and, after the fixed scale/translation normalization, the
radial graph of $P_\alpha(h)$ is within $\rho_0$ of the unit circle in
$W^{1,\infty}$.  The radii $S_0,\rho_0$ are universal and may be
decreased.  No smallness of the affine-side parametrization relative to
$P_\alpha(h_0)$ is imposed: such a condition would incorrectly exclude a
positive side whose length is much smaller than its normal gap.  The
ratio-free compatibility between the radial endpoint chart, the whole
fixed-support segment, and the support metric is proved below as JP--NS.

\begin{proposition}[The radial chart does not imply a labelled affine-side chart]
\label{JP:prop:radial-not-affine}
There are active area-$\pi$ convex $N$-gons $Q_N$ with positive normal
fans such that
\[
 \max_i\alpha_i+\frac{2\pi}{N}\longrightarrow0,
 \qquad \|r_{Q_N}-1\|_{W^{1,\infty}(\mathbb S^1)}\longrightarrow0,
\]
whereas the derivative of the labelled physical-affine displacement from
the tangential polygon with the same fan tends to infinity.  Thus no
ratio-free radial-to-affine-side chart implication is available.
\end{proposition}

\begin{proof}
Put $a=2\pi/N$, $\varepsilon=a^4$,
$b=(2\pi-2\varepsilon)/(N-2)$, and $t=a^2$.  Choose consecutive normals
$-\varepsilon,0,\varepsilon$ and make all remaining gaps equal to $b$.
For
\[
 c=\left\{\frac{\pi}
 {2\tan(\varepsilon/2)+(N-2)\tan(b/2)}\right\}^{1/2}
 =1+O(a^2),
\]
the constant support vector $c\mathbf1$ gives the area-$\pi$ tangential
polygon.  Keep the central support equal to $c$, change the two adjacent
supports to $c\cos\varepsilon+t\sin\varepsilon$, and leave all other
supports equal to $c$.  The central side then has endpoints $(c,-t)$ and
$(c,t)$ and length $2t$, whereas its tangential reference length is
$2c\tan(\varepsilon/2)$.

Writing
\[
 \delta=\sin\varepsilon
 \left(t-c\tan\frac\varepsilon2\right)
 =t\varepsilon(1+o(1)),
\]
the two incident side lengths change by
$-\delta(\cot\varepsilon+\cot b)=O(a^2)$ from positive lengths
asymptotic to $a$, the next two change by $\delta\csc b$, and all others
are unchanged.  Hence all sides remain active.  The affine scale-section
factor is $1+O(bt\varepsilon)$, so scale and area normalization do not
alter the following asymptotics.  On the central labelled side the
tangential derivative of the displacement is at least
\[
 \frac{t}{(1+o(1))c\tan(\varepsilon/2)}-1
 \sim\frac{2t}{\varepsilon}=\frac2{a^2}\longrightarrow\infty.
\]

Every changed vertex is nevertheless displaced by $O(t)$; the two
central displacements are exactly
$t-c\tan(\varepsilon/2)$, and the outer displacements are
$O(\delta/\sin b)$.  On a side with normal $\theta_i$ the radial graph is
$r(\theta)=h_i/\cos(\theta-\theta_i)$ and
$|r'|=r|\tan(\theta-\theta_i)|$.  Thus, after the harmless area
normalization,
$\|r_{Q_N}-1\|_{W^{1,\infty}}\le Ca\to0$.  This proves the claim.
\end{proof}

\subsection{Exact eigenvalue and scale-normalized Hessians}

Let $u_t>0$ be the $L^2$-normalized first eigenfunction of
$P_\alpha(h+tz)$, let $\lambda_t$ be its eigenvalue, and put
$p_t=-\partial_{\nu_t}u_t$ on the open sides.  If $g$ is boundary data,
write $\Dred_{P_t}(g)$ for the reduced Helmholtz energy: solve the
Eulerian state equation with trace $-g$, impose $L^2$ orthogonality to
$u_t$, and evaluate
\(
 \int_{P_t}|\nabla\psi|^2-\lambda_t\int_{P_t}\psi^2.
\)
This is the same reduced form as in the polygon-preserving volume Hessian.

\begin{theorem}[JP--EH: exact fixed-normal support Hessian]
\label{JP:thm:EH}
For every active support line $h(t)=h+tz$,
\begin{align}
 \lambda_t'&=-\sum_i z_i\int_{S_i(t)}p_t^2\,ds,
 \label{JP:eq:lambda-first}\\
 \lambda_t''&=2\Dred_{P_t}(z_i\partial_{\nu_t}u_t).
 \label{JP:eq:lambda-second}
\end{align}
Hence, with $A_t=A_\alpha(h(t))$ and
\(
 F_t(z)=\sum_i z_i\int_{S_i(t)}p_t^2\,ds,
\)
one has the exact identity
\begin{equation}\label{JP:eq:Phi-second}
 \boxed{
 \frac{d^2}{dt^2}\Phi_\alpha(h(t))
 =\frac1\pi\left\{
 2A_t\Dred_{P_t}(z_i\partial_{\nu_t}u_t)
 -2A_t'[z]F_t(z)+\lambda_t A_\alpha''[z,z]
 \right\}.}
\end{equation}
No stationarity, minimum side, adjacent-ratio, moving radial knot, or
Taylor truncation is used in this formula.
\end{theorem}

\begin{proof}
Hadamard's formula and \eqref{JP:eq:constant-normal} give
\eqref{JP:eq:lambda-first}.  Differentiate it on the moving polygon.  If
$\psi_t$ is the Eulerian eigenfunction derivative, the material derivative
of the boundary flux is
\[
 \dot p_t=-\partial_{\nu_t}\psi_t+V_\tau\partial_\tau p_t.
\]
The differentiated arclength, normal, and transport rows combine, after
one sidewise integration by parts, into
\[
 \lambda_t''=2\Dred_{P_t}((V\cdot\nu_t)\partial_{\nu_t}u_t)
 +2\sum_i\int_{S_i(t)}p_t^2V_\tau
          \partial_\tau(V\cdot\nu_i)\,ds.
\]
The endpoint terms vanish by the convex-corner expansion of the ground
state.  The second row is zero by \eqref{JP:eq:constant-normal}, proving
\eqref{JP:eq:lambda-second}.  Finally differentiate
$\Phi_\alpha=A_t\lambda_t/\pi$ twice, use
$\lambda_t'=-F_t(z)$ and \eqref{JP:eq:area-jets}, and obtain
\eqref{JP:eq:Phi-second}.
\end{proof}

\subsection{Why the disk cubic truncation cannot be the global metric}

For reference, let
\begin{equation}\label{JP:eq:disk-cubic-candidate}
 H^{(3)}_\alpha(z,z)
 =D^2\mathcal L(0)[Z_\alpha,Z_\alpha]
  +D^3\mathcal L(0)[B_\alpha,Z_\alpha,Z_\alpha]
\end{equation}
be the formerly proposed disk-based support Hessian.  The complete physical
chain rule is understood in \eqref{JP:eq:disk-cubic-candidate}; in particular
the two curvature terms are not omitted.

\begin{proposition}[Alternating-mesh obstruction to the cubic metric]
\label{JP:prop:cubic-counterexample}
There are even cyclic fans with fixed adjacent ratio and admissible
translation/area-normalized support directions for which
$H^{(3)}_\alpha(z,z)<0$.  Consequently no positive no-ratio metric can be
obtained by proving coercivity of \eqref{JP:eq:disk-cubic-candidate}.
\end{proposition}

\begin{proof}
Let $N=2M$, $a=2\pi/N$, take alternating exterior angles
$a/10,19a/10$, and set $z_{2j}=1$, $z_{2j+1}=-1$.  The alternating
symmetry gives the two translation constraints and the exact area row.
After the material rescaling $y=\theta/a$, one period has length $2$.
Put $p=(1/10,19/10)$, $c_0=0$, $c_1=1/10$, and on the cell
$x=y-c_r\in[-p_{r-1}/2,p_r/2]$ define
\[
 F_r=\frac{x^2}{2}-\frac{p_0^3+p_1^3}{48},\qquad
 W_0=-\frac{360}{19}-\frac{400}{19}x,\qquad
 W_1=-\frac{360}{19}+\frac{400}{19}x,
\]
\[
 V_r=z_r+xW_r,\qquad U_r=xW_r,\qquad
 K_r=W_r^2-2W_rV_r'.
\]
For period-two Fourier coefficients
$\widehat h(n)=\frac12\int_0^2h(y)e^{-\pi iny}\,dy$, write
\(
 \langle g,h\rangle_2
 =\sum_{n\ne0}\pi|n|\overline{\widehat g(n)}\widehat h(n).
\)
The small-cell expansions, in the physical angular coordinate, are
\[
 f=a^2F+O(a^4),\quad w=a^{-1}W+O(a),\quad
 v=V+O(a^2),\quad f'=aF'+O(a^3),
\]
\[
 w^2-2wv'=a^{-2}K+O(1),\qquad wf'=U+O(a^2).
\]
The exact physical-chain identity is
\[
 D^3\mathcal L(0)[B,Z,Z]
 =D^3\Lambda^n[f,v,v]
  +D^2\Lambda^n[f,w^2-2wv']
  +2D^2\Lambda^n[v,-wf'],
\]
and $D^2\Lambda^n=2q_{\mathbb D}$.  The normal cubic is $O(1)$;
rescaling the other three pairings therefore gives
\[
 H^{(3)}_\alpha(z,z)=\frac{4\pi d^2}{a}E+O(1),\qquad
 E=\langle V,V\rangle_2+\langle F,K\rangle_2
       -2\langle V,U\rangle_2,
\]
where $d>0$ is the disk normalization constant.  Exact rational-polynomial
Fourier integration through mode $4096$, with outward intervals of radius
$10^{-12}$ including the trigonometric Taylor and rounding remainders,
gives
\begin{align*}
 A_{4096}&=46.5557803349116641,\\
 B_{4096}&=22.9387363887007744,\\
 C_{4096}&=35.1293696664596509,
\end{align*}
and hence $E_{4096}=A_{4096}+B_{4096}-2C_{4096}
=-0.7642226093068633$.  Integration by parts in each Fourier coefficient
and Cauchy--Schwarz give the explicit tail bound
\[
 |E-E_K|\le
 \frac{\operatorname{TV}(V')^2}{4\pi^3K^2}
 +\frac{\operatorname{TV}(F')\operatorname{TV}(K)}{2\pi^2K}
 +\frac{\operatorname{TV}(V')\operatorname{TV}(U)}{\pi^2K}.
\]
Here
\begin{align*}
 \operatorname{TV}(F')&=4,&
 \operatorname{TV}(V')&=1600/19,\\
 \operatorname{TV}(U)&=476/19,&
 \operatorname{TV}(K)&=523200/361.
\end{align*}
The reproducible certificate is
\nolinkurl{tests/check_jp_cubic_counterexample.py}.  At $K=4096$ the
right-hand side is less than $0.123892271$, so
$E<-0.640330338$.  Therefore \eqref{JP:eq:disk-cubic-candidate} is negative
for all sufficiently small $a$.  Reflection through a common cell vertex
makes $B$ even and $Z$ odd, so the old linear forcing in this direction and
$D^3\mathcal L(0)[Z,Z,Z]$ vanish.  The negative sign is therefore intrinsic
to the candidate Hessian.  This is an obstruction to the truncation, not
to the exact Hessian \eqref{JP:eq:Phi-second}.
\end{proof}

\begin{corollary}[Stationary specialization]
\label{JP:cor:stationary}
If $h$ is stationary for $\Phi_\alpha$ under the fixed-normal support
variations, then at $t=0$
\begin{equation}\label{JP:eq:stationary-specialization}
 D^2\Phi_\alpha(h)[z,z]
 =\frac{2A}{\pi}\Dred_P(z_i\partial_\nu u)
 +\frac\lambda\pi\left\{A_\alpha''[z,z]
       -\frac{2A_\alpha'[z]^2}{A}\right\}.
\end{equation}
On the area tangent row $A_\alpha'[z]=0$, the last square vanishes.
\end{corollary}

\begin{proof}
Stationarity and \eqref{JP:eq:lambda-first} give
$F_0(z)=\lambda A_\alpha'[z]/A$.  Substitute this in
\eqref{JP:eq:Phi-second}.
\end{proof}

\subsection{Exact value integration and the reduced frontier}

Let $h_\alpha\mathbf1$ denote the canonical tangential support vector and
let $\alpha_*=a\mathbf1$, $a=2\pi/N$.  Put
\begin{equation}\label{JP:eq:forcing}
 \ell_\alpha=D_h\Phi_\alpha(h_\alpha\mathbf1)
\end{equation}
restricted to $\mathsf X_\alpha$.  Ordinary one-variable Taylor
integration, using the genuine Hessian rather than its disk cubic
truncation, gives the identity
\begin{equation}\label{JP:eq:exact-integration}
 \boxed{
 \Phi_\alpha(h_\alpha\mathbf1+z)-\Phi_\alpha(h_\alpha\mathbf1)
 =\ell_\alpha[z]
 +\int_0^1(1-t)
 D_h^2\Phi_\alpha(h_\alpha\mathbf1+tz)[z,z]\,dt.}
\end{equation}

The exact Hadamard variable is the constant side-normal speed $z_i$, not
the radial component of the full material vertex field.  Let
$\Pi_{\rm ng}$ delete the constant and first Fourier harmonics.  We import
from the CT core the literal disk Bessel form
\[
 q_D(f)=d^2\int_{\mathbb T}(\Pi_{\rm ng}f)
       (\mathcal D_\circ+1)(\Pi_{\rm ng}f),
 \qquad
 q_D(f)\asymp\|\Pi_{\rm ng}f\|_{H^{1/2}}^2,
\]
where $d=\partial_ru_0(1)$ and
$\mathcal D_\circ e^{in\theta}
=jJ_{|n|}'(j)J_{|n|}(j)^{-1}e^{in\theta}$ for $n\ne0$.
Let $\theta_{i+1}-\theta_i=\alpha_i$ and define the periodic fan-normal sine
interpolant by
\begin{equation}\label{JP:eq:normal-speed-interpolant}
 (\mathcal I_\alpha z)(\theta)
 =\frac{\sin(\theta_{i+1}-\theta)z_i
         +\sin(\theta-\theta_i)z_{i+1}}{\sin\alpha_i},\qquad
 \theta_i\le\theta\le\theta_{i+1}.
\end{equation}
Put
\begin{equation}\label{JP:eq:normal-speed-metric}
 \mathfrak G_\alpha^{\rm ns}(z)
 :=q_D(\Pi_{\rm ng}\mathcal I_\alpha z),
 \qquad z\in\mathsf X_\alpha.
\end{equation}
This is the metric interface \(\mathrm{JP\text{-}NS}\).

\begin{proposition}[Closed chart entry for the normal-speed metric]
\label{JP:prop:normal-speed-chart-entry}
The form $\mathfrak G_\alpha^{\rm ns}$ is a Hilbert form on
$\mathsf X_\alpha$.  There is a universal $C$ such that
\begin{equation}\label{JP:eq:normal-speed-chart-entry}
 P_\alpha(h_0+z)\in\mathcal U_{S_0,\rho_0}
 \quad\Longrightarrow\quad
 h_0+[0,1]z\ \hbox{is active},\qquad
 \mathfrak G_\alpha^{\rm ns}(z)^{1/2}\le C(S_0+\rho_0),
\end{equation}
and every polygon on the segment has a radial graph $1+\rho_t$ satisfying
\begin{equation}\label{JP:eq:radial-path-entry}
 \sup_{0\le t\le1}\|\rho_t\|_{W^{1,\infty}(\mathbb S^1)}
 \le C(S_0+\rho_0).
\end{equation}
No minimum angle or adjacent-ratio bound is used.
\end{proposition}

\begin{proof}
If $\mathfrak G_\alpha^{\rm ns}(z)=0$, then
$\mathcal I_\alpha z$ contains only the constant and first Fourier modes.
On every open normal-angle interval it satisfies
$(\mathcal I_\alpha z)''+\mathcal I_\alpha z=0$.  Applying this equation
to a function of the form $c+b\cdot e_r(\theta)$ forces $c=0$.  Hence
$z_i=b\cdot\nu_i$ is an exact support translation, which is zero in the
chosen complement $\mathsf X_\alpha$.  This proves positive definiteness.
It also shows that \eqref{JP:eq:normal-speed-metric} descends exactly through
the translation quotient, rather than merely after choosing the complement.

Let $P_i(z)$ be the displacement of the vertex shared by sides $i$ and
$i+1$ from its position in $P_\alpha(h_0)$.  If a star-shaped boundary is
$r(\varphi)e_r(\varphi)$, its one-sided outer normal is parallel to
$r e_r-r'e_\varphi$.  Hence $\|r-1\|_{W^{1,\infty}}\le\rho_0<1/2$
puts the polar direction of each endpoint vertex within $C\rho_0$ of
its adjacent normal cone.  The tangential reference has
$\|r_0-1\|_{W^{1,\infty}}\le CS_0$, and the cone width is at most $S_0$.
The two vertex radii are bounded above and below universally.  Therefore
\begin{equation}\label{JP:eq:vertex-displacement-radial}
 \max_i|P_i(z)|\le C(S_0+\rho_0),
\end{equation}
with no inverse angle or side length.  At that vertex the
constant-normal-speed identity gives
\[
 z_i=P_i(z)\cdot\nu_i,\qquad
 z_{i+1}=P_i(z)\cdot\nu_{i+1}.
\]
Solving these two projection equations gives the exact identity
\[
 (\mathcal I_\alpha z)(\theta)=P_i(z)\cdot e_r(\theta),
 \qquad \theta_i\le\theta\le\theta_{i+1}.
\]
Therefore
\[
 \int_{\theta_i}^{\theta_{i+1}}
 \left\{|\mathcal I_\alpha z|^2
       +|\partial_\theta\mathcal I_\alpha z|^2\right\}\,d\theta
 =\alpha_i|P_i(z)|^2.
\]
Summation gives
$\|\mathcal I_\alpha z\|_{H^1}\le C(S_0+\rho_0)$, and the
Bessel equivalence $q_D\asymp H^{1/2}$ proves the metric bound.  For fixed
normals every side length is a linear function of the support numbers.
It is positive at both $h_0$ and $h_0+z$; therefore it remains positive on
their whole affine segment.  This proves activity of every labelled side.

Every intermediate vertex is the affine combination of its endpoint
vertices.  Its radial and tangential coordinates relative to either
incident normal therefore obey the same $C(S_0+\rho_0)$ bounds.  On each
side the exact formulas
\[
 r_t(\theta)=\frac{h_i(t)}{\cos(\theta-\theta_i)},
 \qquad
 r_t'(\theta)=r_t(\theta)\tan(\theta-\theta_i)
\]
then give~\eqref{JP:eq:radial-path-entry}, uniformly in $t$.  In view of
Proposition~\ref{JP:prop:radial-not-affine}, this radial/metric path statement,
not a labelled affine-side derivative bound, is the ratio-free interface
used downstream.
\end{proof}

\begin{proposition}[Exact area identity in the normal-speed chart]
\label{JP:prop:normal-speed-area}
For $f=\mathcal I_\alpha z$ one has
\begin{equation}\label{JP:eq:normal-speed-area}
 \boxed{A_\alpha''[z,z]
 =\int_{\mathbb T}\bigl(f^2-|f'|^2\bigr)\,d\theta.}
\end{equation}
Moreover, on each interval $I_i=[\theta_i,\theta_{i+1}]$,
\begin{equation}\label{JP:eq:normal-speed-L2}
\begin{aligned}
 2\int_{I_i}f^2\,d\theta
 ={}&\frac{\alpha_i-\sin\alpha_i}{4\sin^2(\alpha_i/2)}
       (z_{i+1}-z_i)^2\\
 &+\frac{\alpha_i+\sin\alpha_i}{4\cos^2(\alpha_i/2)}
       (z_{i+1}+z_i)^2.
\end{aligned}
\end{equation}
Thus the two displayed coefficients are exactly $2\int_{I_i}f^2$.  After
the Hessian factor $\lambda_\alpha/\pi$, this is the polygonal local
$L^2$ row converging to the $L^2$ part of $2q_D$.  The low-mode projection
and the nonlocal order-one row of $q_D$ remain global.
\end{proposition}

\begin{proof}
Write $s=\theta-(\theta_i+\theta_{i+1})/2$ and
\[
 X_i=\frac{z_i+z_{i+1}}{2\cos(\alpha_i/2)},
 \qquad
 Y_i=\frac{z_{i+1}-z_i}{2\sin(\alpha_i/2)}.
\]
The sine interpolation formula is exactly
$f=X_i\cos s+Y_i\sin s$.  Orthogonality of the even and odd summands on
$[-\alpha_i/2,\alpha_i/2]$ gives \eqref{JP:eq:normal-speed-L2}.  The same
calculation gives
\[
 \int_{I_i}(f^2-|f'|^2)\,d\theta
 =2z_iz_{i+1}\csc\alpha_i
  -(z_i^2+z_{i+1}^2)\cot\alpha_i.
\]
Summing in $i$ is precisely \eqref{JP:eq:area-second-explicit}, which proves
\eqref{JP:eq:normal-speed-area}.
\end{proof}

Henceforth the support Hilbert form used by JP--SC, JP--FS, and NR is fixed
to be
\begin{equation}\label{JP:eq:JP-metric-fixed}
 \Gjp_\alpha:=\mathfrak G_\alpha^{\rm ns}.
\end{equation}

\begin{remark}[Archived stronger JP--SC route]
\label{JP:rem:SC}
For every fixed $\delta\in(0,1)$ there are
$S_0,\rho_0,\eta_0>0$, independent of the number, minimum size, and ratios
of the sides, such that the following holds for every fan with
$S=a+\max_i\alpha_i\le S_0$.  For all
$y,\xi\in\mathsf X_\alpha$ with the
entire segment $h_0+[0,1]y$ in the physical active support chart and
$\mathfrak G_\alpha^{\rm ns}(y)^{1/2}\le\eta_0$,
\begin{equation}\label{JP:eq:SC}
 D_h^2\Phi_\alpha(h_\alpha\mathbf1+y)[\xi,\xi]
 \ge (2-\delta)\mathfrak G_\alpha^{\rm ns}(\xi).
\end{equation}
The exported Hilbert constant is $c=2-\delta$; it may not be silently
decreased before JP--FS.
The proof must use the complete expression \eqref{JP:eq:Phi-second}.  A lower
bound for its reduced-energy diagonal alone is insufficient unless the
area and mixed-flux rows have first been Schur-completed.
\end{remark}

\begin{remark}[Why stationary O--H does not discharge JP--SC]
At the base point $h_\alpha\mathbf1$, the definition of the fixed scale
section gives $A_\alpha'[\xi]=0$, so the exact base Hessian is
\begin{equation}\label{JP:eq:base-support-hessian}
 D_h^2\Phi_\alpha(h_\alpha\mathbf1)[\xi,\xi]
 =\frac1\pi\left\{
 2A_\alpha\Dred_{P_\alpha^0}(\xi_i\partial_\nu u_\alpha)
 +\lambda_\alpha A_\alpha''[\xi,\xi]\right\}.
\end{equation}
The stationary O--H comparison is neither stated at this generally
nonstationary tangential polygon nor relative to the scalar support metric.
Its available error has the form
$\omega(S)\|V\|_{H^{1/2}(\mathbb T;\mathbb R^2)}^2$.  This full-vector
norm is not uniformly controlled by \eqref{JP:eq:normal-speed-metric}: the
vertex solution contains the tangential coefficient
\[
 \frac{\xi_{i+1}-\xi_i\cos\alpha_i}{\sin\alpha_i},
\]
which has an inverse-angle amplification that the scale-critical
fan-normal interpolation does not.  A rate-free $\omega(S)\to0$ therefore
cannot be absorbed into $\mathfrak G_\alpha^{\rm ns}$.  A proof of JP--SC
needs a support-compatible,
endpoint-resummed relative lower bound for the complete form
\eqref{JP:eq:base-support-hessian}, followed by control of the mixed-flux row
away from the base.  Invoking stationary O--H or its full-vector norm is
not such a bound.
\end{remark}

\begin{remark}[Why a pointwise corner-amplitude estimate is insufficient]
The terminal jump--cotangent cancellation in the stationary O--H ledger
uses sidewise stationary flux normalization.  At the generally
nonstationary tangential polygon one has only the global Rellich average,
not
\(
 \int_{S_i}p^2=(\lambda/\pi)|S_i|
\)
for each side.  If $\beta_i$ denotes the normalized leading flux amplitude
at a corner of exterior angle $\alpha_i$, separating that corner before
its descendants leaves a schematic residual
\[
 (\beta_i^2-1)\alpha_i^{-1}(z_{i+1}-z_i)^2.
\]
The corresponding sharp transition in the JP--NS form costs only
$O((z_{i+1}-z_i)^2(1+|\log\alpha_i|))$.  Consequently even a uniform
$\beta_i=1+o(1)$ is not an absorbable pointwise error.  The amplitude loss
must be retained with all descendant annuli and LCA pairs until their
positive energy is recovered.  This is why \eqref{JP:eq:graded-DtN-blocker}
is a completed operator inequality rather than a cornerwise estimate.
\end{remark}

The fan-normal candidate reduces JP--SC to two precise estimates.  In the
pure two-ray wedge model, combining the endpoint reduced energy with the
area row gives exactly the local integral identity
\eqref{JP:eq:normal-speed-L2}.
Both coefficients are nonnegative, and the common coefficient is positive.
What is not supplied by this local calculation is the ratio-free global
normalization of all terminal rays, annuli, and LCA pairs.

\begin{remark}[Archived JP--S0 sufficient estimate]
\label{JP:rem:S0}
Prove a modulus $\omega_0(S)\to0$, independent of all minimum angles and
adjacent ratios, such that
\begin{equation}\label{JP:eq:S0}
 \boxed{
 D_h^2\Phi_\alpha(h_0)[z,z]
 \ge(2-\omega_0(S))\mathfrak G_\alpha^{\rm ns}(z),
 \qquad z\in\mathsf X_\alpha.}
\end{equation}
The complete two-ray, annular, and LCA ledger must be retained, and its
amplified jump must be combined with $\lambda_0A_\alpha''[z,z]$ before
either term is estimated.
\end{remark}

At the base scale section, \eqref{JP:eq:normal-speed-area} rewrites this
obligation without any hidden geometric row as
\begin{equation}\label{JP:eq:graded-DtN-blocker}
 \begin{aligned}
 \frac1\pi\left\{
 2A_\alpha\Dred_{P_\alpha^0}(z_i\partial_\nu u_\alpha)
 +\lambda_\alpha\int_{\mathbb T}(f^2-|f'|^2)\right\}
 &\ge(2-\omega_0(S))q_D(\Pi_{\rm ng}f),\\
 f&=\mathcal I_\alpha z.
 \end{aligned}
\end{equation}
The same/adjacent two-ray part is \eqref{JP:eq:normal-speed-L2}; the genuinely
missing part is the no-ratio nonlocal terminal-ray, annular, and LCA
normalization in \eqref{JP:eq:graded-DtN-blocker}.  None of H0, QI, DF, or DM
states this polygonal reduced-DtN operator inequality.

\begin{remark}[Archived JP--ST sufficient estimate]
\label{JP:rem:ST}
Along every active fixed-scale support segment in the physical chart,
prove a modulus $\omega_1(\rho,S)\to0$ such that
\begin{equation}\label{JP:eq:ST}
 D_h^2\Phi_\alpha(h_0+y)[\xi,\xi]
 \ge\min\left\{D_h^2\Phi_\alpha(h_0)[\xi,\xi],
                 2\mathfrak G_\alpha^{\rm ns}(\xi)\right\}
       -\omega_1(\rho,S)\mathfrak G_\alpha^{\rm ns}(\xi).
\end{equation}
All reduced-energy, mixed-flux, and area rows in
\eqref{JP:eq:Phi-second} are included.  No two-sided absolute Hessian
transport is required: a large positive endpoint coefficient need only
remain above the disk threshold, not be approximated to additive
$\mathfrak G_\alpha^{\rm ns}$ accuracy.
\end{remark}

\begin{proposition}[Typed sufficient cut for JP--SC]
\label{JP:prop:SC-cut}
JP--S0 and JP--ST, together with
Proposition~\ref{JP:prop:normal-speed-chart-entry}, imply JP--SC with
$\Gjp_\alpha=\mathfrak G_\alpha^{\rm ns}$ and
$c=2-\delta$ after decreasing $S_0$ and $\rho_0$ so that
$\omega_0(S)+\omega_1(\rho,S)\le\delta$.
\end{proposition}

\begin{proof}
Choose the thresholds so that
$\omega_0(S)+\omega_1(\rho,S)\le\delta$.  Then
\eqref{JP:eq:S0} makes the minimum in \eqref{JP:eq:ST} at least
$(2-\omega_0(S))\mathfrak G_\alpha^{\rm ns}(\xi)$.  Hence
\eqref{JP:eq:S0}--\eqref{JP:eq:ST} give
$D_h^2\Phi_\alpha(h_0+y)[\xi,\xi]
\ge(2-\delta)\mathfrak G_\alpha^{\rm ns}(\xi)$.  The remaining chart and activity
requirements are exactly \eqref{JP:eq:normal-speed-chart-entry}.
\end{proof}

\begin{remark}[Numerical warning is not a counterexample]
Finite-element prechecks on alternating fans have not found a negative
exact support Hessian.  The quotient relative to the older physical-radial
trace deteriorates as one angle ratio decreases at fixed modest $N$; that
calculation may also underresolve the nearly flat corner exponent and is
not a certified asymptotic.  On the same samples, the quotient relative to
$\mathfrak G_\alpha^{\rm ns}$ does not deteriorate.  These are design
diagnostics only: \eqref{JP:eq:S0} still requires a ratio-free proof for the
complete endpoint ledger, and no node status follows from the precheck.
\end{remark}

\begin{remark}[Status of the candidate cut]
JP--S0 and JP--ST are typed sufficient subcontracts; they are not
additional active DAG vertices.  A replacement proof using another
no-ratio support metric would have to replace JP--NS and the JP--SC
interface together.  Conversely, a rigorous degeneration of JP--S0 would
reject only this particular completion route, not the target theorem.
\end{remark}

Let $\|\ell_\alpha\|_{(\Gjp_\alpha)^*}$ be the dual norm.  We first remove
the part of the former JP--BR contract which is already a consequence of
the closed pure-fan nodes.

\begin{theorem}[JP--BU: exact tangential-fan baseline upper bound]
\label{JP:thm:BU}
There are universal $C,S_0>0$ such that, for every positive fan with
$S=a+\max_i\alpha_i\le S_0$,
\begin{equation}\label{JP:eq:fan-upper}
 \boxed{
 \Phi_\alpha(h_\alpha\mathbf1)
       -\Phi_{\alpha_*}(h_{\alpha_*}\mathbf1)
 \le C\Jfour(\alpha).}
\end{equation}
The constant is independent of the smallest angle and all adjacent ratios.
\end{theorem}

\begin{proof}
The exact disk-normalized expansion furnished by the closed TR, CT core,
QB, and DF nodes splits the pure tangential-fan value into its quadratic
disk-source row, its complete cubic row, and a fourth-and-higher scalar
tail $F_4$; explicitly, if $\Delta_2$ and $\Delta_3$ denote the complete
degree-two and degree-three scalar rows, then
\begin{align*}
 \Phi_\alpha(h_\alpha\mathbf1)
 -\Phi_{\alpha_*}(h_{\alpha_*}\mathbf1)
 &=[\Delta_2(\alpha)-\Delta_2(\alpha_*)]
   +[\Delta_3(\alpha)-\Delta_3(\alpha_*)]\\
 &\quad+[F_4(\alpha)-F_4(\alpha_*)].
\end{align*}
The relative source transfer gives the required one-sided quadratic
Bregman estimate, while the closed pure cubic row is absolute:
\begin{align*}
 \Delta_2(\alpha)-\Delta_2(\alpha_*)&\le C\Jfour(\alpha),\\
 |\Delta_3(\alpha)-\Delta_3(\alpha_*)|&\le C\Jfour(\alpha).
\end{align*}
It remains only to record why the scalar tail needs no support-direction
CT--M4 estimate.

Put $d_i=\alpha_i-a$ and
$\gamma_i(t)=a+td_i$.  The unconditional tame scalar fourth-remainder
corollary of the closed CT core (obtained by applying Cauchy's coefficient
estimate to the complex natural-order Hessian) gives, on the real chart
and on its material complex polydisc,
\begin{equation}\label{JP:eq:F4-absolute}
 |F_4(\zeta)|
 \le C\|B_\zeta\|_{H^{1/2}}^2
        \|B_\zeta\|_{W^{1,\infty}}^2
 \le CS^5\le CS^4.
\end{equation}
Choose a fixed sufficiently small $\varepsilon>0$.  If
$\Jfour(\alpha)\ge\varepsilon S^4$, applying
\eqref{JP:eq:F4-absolute} at the two endpoints gives the desired bound with
constant $2C/\varepsilon$.  If
$\Jfour(\alpha)<\varepsilon S^4$, the exact identity
\begin{equation}\label{JP:eq:J4-action}
 \Jfour(\alpha)
 =12\int_0^1(1-t)E(t)\,dt,
 \qquad E(t)=\sum_i\gamma_i(t)^2d_i^2,
\end{equation}
and
\(
 \Jfour=\sum_i d_i^2(\alpha_i^2+2a\alpha_i+3a^2)
\)
show, after decreasing $\varepsilon$, that
$a\asymp S$ and the whole segment is uniformly comparable.  The material
complex radius in the direction $d$ is at least
$ca/\|d\|_{\ell^\infty}$.  Cauchy's estimate and
\eqref{JP:eq:F4-absolute} therefore give
\[
 |F_4''(t)|
 \le CS^4\frac{\|d\|_{\ell^\infty}^2}{a^2}
 \le C E(t).
\]
Cyclic symmetry makes $F_4'(0)$ constant in the fan coordinates, hence
$F_4'(0)[d]=0$ because $\sum_i d_i=0$.  Taylor integration and
\eqref{JP:eq:J4-action} yield
\[
 |F_4(\alpha)-F_4(\alpha_*)|
 \le C\int_0^1(1-t)E(t)\,dt
 \le C\Jfour(\alpha).
\]
Adding the three rows proves \eqref{JP:eq:fan-upper}.  Only the pure scalar
tail was used here; none of the retired support-linear CT--M4 estimates is
being reintroduced.
\end{proof}

The remaining forcing admits a coordinate-free finite-dimensional
description.  Let
\begin{equation}\label{JP:eq:side-masses}
 P_\alpha^0=P_\alpha(h_\alpha\mathbf1),\qquad
 L_i=|S_i|,\qquad
 m_i=\int_{S_i}(-\partial_\nu u_\alpha)^2\,ds,
\end{equation}
where $|P_\alpha^0|=\pi$ and $u_\alpha$ is $L^2$-normalized.

\begin{proposition}[Exact flux representation of the support forcing]
\label{JP:prop:flux-forcing}
On the fixed affine scale section, equivalently
\(
 \sum_iL_i z_i=0
\)
at $h_\alpha\mathbf1$, one has
\begin{equation}\label{JP:eq:exact-flux-forcing}
 \boxed{
 \ell_\alpha[z]
 =-\sum_i m_i z_i
 =-\sum_i\left(m_i-\frac{\lambda_\alpha}{\pi}L_i\right)z_i.}
\end{equation}
Moreover, with
\begin{equation}\label{JP:eq:centered-flux}
 r_{\alpha,i}=m_i-\frac{\lambda_\alpha}{\pi}L_i,
\end{equation}
one has
\begin{equation}\label{JP:eq:flux-moments}
 \sum_i r_{\alpha,i}=0,
 \qquad
 \sum_i r_{\alpha,i}\nu_i=0.
\end{equation}
Thus $r_\alpha$ is an exact covector on the scale/translation quotient.
\end{proposition}

\begin{proof}
Hadamard's formula and the first area variation give
\[
 D_h\Phi_\alpha(h_\alpha\mathbf1)[z]
 =\frac1\pi\left\{
 \lambda_\alpha\sum_iL_i z_i-\pi\sum_i m_i z_i\right\}.
\]
The first term vanishes on the fixed scale section.  Rellich's identity
on the tangential polygon and tangentiality of its boundary give
\[
 h_\alpha\sum_i m_i=2\lambda_\alpha,
 \qquad
 \sum_iL_i=\frac{2\pi}{h_\alpha}.
\]
These identities prove the first equation in \eqref{JP:eq:flux-moments} and
allow the centered form in \eqref{JP:eq:exact-flux-forcing}.  Translation
invariance in Hadamard's formula gives
$\sum_i m_i\nu_i=0$, while the boundary normal-balance identity is
$\sum_iL_i\nu_i=0$.  Their difference proves the vector equation in
\eqref{JP:eq:flux-moments}.
\end{proof}

We now isolate the analytic input which is actually missing from the old
TR/DM route.  On the side-normal cell
$K_i=[\theta_i-\alpha_{i-1}/2,\theta_i+\alpha_i/2]$, put
\begin{equation}\label{JP:eq:disk-source-B}
 B_\alpha=\Pi_{\rm ng}\{h_\alpha\sec(\theta-\theta_i)-1\},
 \qquad K_\alpha z=\Pi_{\rm ng}\mathcal I_\alpha z.
\end{equation}
Let $\langle\cdot,\cdot\rangle_D$ be the bilinear disk Bessel form and
define
\begin{align}
 p_\alpha[z]&=2\langle B_\alpha,K_\alpha z\rangle_D,
 &\Gjp_\alpha(z)&=q_D(K_\alpha z),\label{JP:eq:JE-forms}\\
 P_D^{\rm ns}(\alpha)&=
 \inf_{z\in\mathsf X_\alpha}q_D(B_\alpha+K_\alpha z),
 &\Delta_\alpha&=
 \Phi_\alpha(h_\alpha\mathbf1)
 -\Phi_{\alpha_*}(h_{\alpha_*}\mathbf1).
 \label{JP:eq:D0ns-definition}
\end{align}
Cyclic symmetry gives $p_{\alpha_*}=0$ on the regular scale/translation
quotient.

We first close the disk part of the endpoint.  This is not a transfer from
the old physical-radial space: the normal-cell partition admits its own
exact half-strip square.

For $I_i=[\theta_i,\theta_{i+1}]$, let
\begin{equation}\label{JP:eq:normal-linear-spline}
 (\mathcal I_\alpha^0z)(\theta_i+y)
 =z_i+\frac{y}{\alpha_i}(z_{i+1}-z_i),
 \qquad 0\le y\le\alpha_i,
\end{equation}
and let $G_\alpha^{\rm e}$ be the edge-quadratic source
\begin{equation}\label{JP:eq:edge-source-on-normal-cell}
 G_\alpha^{\rm e}(\theta_i+y)
 =\frac12\min\{y,\alpha_i-y\}^2.
\end{equation}
This is precisely $x^2/2$ on the two side cells incident to the polygonal
vertex direction $\theta_i+\alpha_i/2$.  Put
\begin{equation}\label{JP:eq:q-plus-normal}
 q_+(f)=\sum_{n\ne0}(|n|+1)|\widehat f_n|^2,
 \qquad
 L_\alpha^0z=\frac12\sum_i(\alpha_{i-1}+\alpha_i)z_i,
\end{equation}
and
\begin{equation}\label{JP:eq:P-plus-normal-linear}
 P_{+,0}^{\rm ns}(\alpha)
 =\inf_{L_\alpha^0z=0}q_+(G_\alpha^{\rm e}+\mathcal I_\alpha^0z).
\end{equation}
The missing constant in \eqref{JP:eq:q-plus-normal} is understood as the
global mean minimization.

\begin{lemma}[Exact normal-cell half-strip and variance squares]
\label{JP:lem:normal-cell-square}
Let $\mathcal N_i$ be the least Dirichlet energy in the Neumann half-strip
$I_i\times(0,\infty)$ with prescribed bottom trace.  Then
\begin{equation}\label{JP:eq:normal-cell-strip-square}
 \boxed{
 \mathcal N_i(G_\alpha^{\rm e}+\mathcal I_\alpha^0z)
 =\frac{\zeta(3)}{\pi^3}
 \left\{\frac{\alpha_i^4}{4}
       +7(z_{i+1}-z_i)^2\right\}.}
\end{equation}
Moreover, after allowing an independent local mean $C_i$,
\begin{equation}\label{JP:eq:normal-cell-L2-square}
 \boxed{
 \inf_{C_i\in\mathbb R}\int_{I_i}
 |G_\alpha^{\rm e}+\mathcal I_\alpha^0z-C_i|^2
 =\frac{\alpha_i^5}{720}
  +\frac{\alpha_i}{12}(z_{i+1}-z_i)^2.}
\end{equation}
Consequently
\begin{equation}\label{JP:eq:normal-homogeneous-gap}
 P_{+,0}^{\rm ns}(\alpha)-P_{+,0}^{\rm ns}(\alpha_*)
 \ge \frac{\zeta(3)}{8\pi^4}\Jfour(\alpha).
\end{equation}
The estimate is pointwise in the coefficient vector before imposing the
area and translation rows.
\end{lemma}

\begin{proof}
Write $w=\alpha_i$, $t=y/w$, and
\[
 g(t)=\frac12\min\{t,1-t\}^2,
 \qquad \rho=\frac{z_{i+1}-z_i}{w^2}.
\]
Modulo a constant, the normalized bottom trace is
$w^2h(t)$ with $h(t)=g(t)+\rho t$.  Its distributional second derivative
is $h''=1-\delta_{1/2}$.  Twice integrating the cosine coefficient gives,
for $k\ge1$,
\[
 2\int_0^1h(t)\cos(k\pi t)\,dt
 =\frac{2\{\rho((-1)^k-1)+\cos(k\pi/2)\}}{(k\pi)^2}.
\]
Thus the even modes contain only the source and the odd modes only the
slope.  Since
\[
 \sum_{k\ {\rm even}}k^{-3}=\frac18\zeta(3),
 \qquad
 \sum_{k\ {\rm odd}}k^{-3}=\frac78\zeta(3),
\]
conformal rescaling of the half-strip gives
\eqref{JP:eq:normal-cell-strip-square}.

The function $g$ is symmetric about $1/2$, so it is orthogonal in variance
to $t-1/2$.  Direct integration gives
\[
 \int_0^1g=\frac1{24},\qquad
 \int_0^1g^2=\frac1{320},\qquad
 \operatorname {Var}(g)=\frac1{720},\qquad
 \operatorname {Var}(t)=\frac1{12}.
\]
This proves \eqref{JP:eq:normal-cell-L2-square}.

Restriction of the periodic harmonic extension yields
\[
 2\pi\sum_{n\ne0}|n||\widehat f_n|^2
 \ge\sum_i\mathcal N_i(f).
\]
For the $L^2$ row, independent cell means only decrease the infimum relative
to one global mean.  At the regular fan with $z=0$, reflection at every
normal line makes both restrictions equalities and all local means agree.
Subtracting the regular value, using \eqref{JP:eq:normal-cell-strip-square}--
\eqref{JP:eq:normal-cell-L2-square}, and applying convexity of $x^5$ gives
\[
 P_{+,0}^{\rm ns}(\alpha)-P_{+,0}^{\rm ns}(\alpha_*)
 \ge\frac{\zeta(3)}{8\pi^4}
       \left(\sum_i\alpha_i^4-Na^4\right)
 +\frac1{1440\pi}
       \left(\sum_i\alpha_i^5-Na^5\right),
\]
which proves \eqref{JP:eq:normal-homogeneous-gap}.
\end{proof}

Let $Z_\alpha=\ker L_\alpha$, where
\begin{equation}\label{JP:eq:normal-exact-area-row}
 L_\alpha z=\sum_i\left\{
 \tan\frac{\alpha_{i-1}}2+\tan\frac{\alpha_i}2\right\}z_i,
\end{equation}
and define the exact-sine, edge-source value
\begin{equation}\label{JP:eq:P-plus-normal-sine}
 P_{+,\sin}^{\rm e}(\alpha)
 =\inf_{z\in Z_\alpha}
 q_+(G_\alpha^{\rm e}+\mathcal I_\alpha z).
\end{equation}
Translations may be left in $Z_\alpha$: the sine interpolant sends them
exactly to the first harmonics, so the minimizer cancels those modes.  This
is equivalent to applying $\Pi_{\rm ng}$ and then taking the fixed
translation complement.

\begin{lemma}[Relative linear-to-sine area graph]
\label{JP:lem:normal-sine-graph}
There is a modulus $\omega_{\sin}(S)\to0$, independent of all minimum
angles and adjacent ratios, such that
\begin{equation}\label{JP:eq:normal-sine-relative}
 \left|\{P_{+,\sin}^{\rm e}(\alpha)-P_{+,0}^{\rm ns}(\alpha)\}
 -\{P_{+,\sin}^{\rm e}(\alpha_*)-P_{+,0}^{\rm ns}(\alpha_*)\}\right|
 \le\omega_{\sin}(S)\Jfour(\alpha).
\end{equation}
One may take $\omega_{\sin}(S)=C(S^{1/4}+S^{1/2})$.
\end{lemma}

\begin{proof}
On $I_i$, Taylor's formula for the two endpoint basis functions gives
\begin{align*}
 \|\mathcal I_\alpha z-\mathcal I_\alpha^0z\|_{L^\infty(I_i)}
 &\le C\alpha_i^2(|z_i|+|z_{i+1}|),\\
 \|\partial_\theta(\mathcal I_\alpha z-\mathcal I_\alpha^0z)
        \|_{L^\infty(I_i)}
 &\le C\alpha_i(|z_i|+|z_{i+1}|).
\end{align*}
For every constant $c$, direct integration of a linear function gives the
ratio-free nodal sampling estimate
\[
 \sum_i(\alpha_{i-1}+\alpha_i)|z_i-c|^2
 \le C\|\mathcal I_\alpha^0z-c\|_{L^2}^2.
\]
If $L_\alpha^0z=0$, then
$\int_{\mathbb T}\mathcal I_\alpha^0z=L_\alpha^0z=0$, so take $c=0$.
Summing the displayed jets, applying the sampling estimate, and using
$\|e\|_{H^{1/2}}^2\le C\|e\|_{L^2}\|e\|_{H^1}$ gives
\begin{equation}\label{JP:eq:normal-linear-sine-close}
 \|\mathcal I_\alpha z-\mathcal I_\alpha^0z\|_{H^{1/2}}
 \le CS^{3/2}\|\mathcal I_\alpha^0z\|_{H^{1/2}}.
\end{equation}
Indeed, the squared $L^2$ and $H^1$ errors are bounded respectively by
$CS^4\|\mathcal I_\alpha^0z\|_{L^2}^2$ and
$CS^2\|\mathcal I_\alpha^0z\|_{L^2}^2$.

For $z\in\ker L_\alpha^0$, put
\[
 C_\alpha z=z-\frac{L_\alpha z}{L_\alpha\mathbf1}\mathbf1.
\]
Then $C_\alpha z\in Z_\alpha$, and the analogous formula using
$L_\alpha^0$ is its inverse.  Since
\[
 L_\alpha z-L_\alpha^0z
 =\sum_i\left(\tan\frac{\alpha_i}{2}-\frac{\alpha_i}{2}\right)
       (z_i+z_{i+1}),
\]
the sampling estimate and
$|\tan(u/2)-u/2|\le Cu^3$ show that the scalar correction in $C_\alpha$
has size at most $CS^2\|\mathcal I_\alpha^0z\|_{H^{1/2}}$.
The sine trace of the constant nodal vector has uniformly bounded
$H^{1/2}$ norm.  Combining this fact with
\eqref{JP:eq:normal-linear-sine-close} yields
\begin{equation}\label{JP:eq:normal-sine-graph-close}
 \|\mathcal I_\alpha C_\alpha z-\mathcal I_\alpha^0z\|_{H^{1/2}}
 \le CS^{3/2}\|\mathcal I_\alpha^0z\|_{H^{1/2}},
\end{equation}
with constants independent of the smallest interval and all adjacent
ratios.

The source bound
$q_+(G_\alpha^{\rm e})\le C\sum_i\alpha_i^4\le CS^3$
follows by splitting the Slobodeckij integral into same-cell,
adjacent-cell, and dyadically separated-cell pieces.  The elementary
Hilbert-space fact
\[
 \left|\operatorname {dist}(b,V_1)^2
       -\operatorname {dist}(b,V_0)^2\right|
 \le C\eta\|b\|^2
\]
whenever $V_1$ is the graph of a bijection $J:V_0\to V_1$ with
$\|J-I\|\le\eta<1/2$, applied to
\eqref{JP:eq:normal-sine-graph-close}, gives
\begin{equation}\label{JP:eq:normal-sine-absolute}
 |R_{\sin}(\alpha)|\le CS^{9/2},
 \qquad R_{\sin}=P_{+,\sin}^{\rm e}-P_{+,0}^{\rm ns}.
\end{equation}

It remains to make the comparison relative to $\Jfour$.  In a
quasiuniform material band
$(1-\eta_0)a\le\gamma_i\le(1+\eta_0)a$, pull every interval affinely to
the real fan and complexify
$\gamma(t)=\gamma+t\delta$ for
$|t|<ca/\|\delta\|_{\ell^\infty}$.  The pulled-back Gagliardo kernel, the
$L^2$ row, both interpolation graphs, the area row, and the quadratic
source are holomorphic.  The sampling coercivity above keeps the complex
Gram inverse uniformly bounded after the numerical constant $c$ is
decreased.  Thus the proof of \eqref{JP:eq:normal-sine-absolute} applies to
the complex bilinear Schur formula and gives
$|R_{\sin}(\gamma(t))|\le Ca^{9/2}$.  Cauchy's estimate and
$\sum_i\gamma_i^2\delta_i^2\ge ca^2
\|\delta\|_{\ell^\infty}^2$ give
\[
 |D^2R_{\sin}(\gamma)[\delta,\delta]|
 \le Ca^{1/2}\sum_i\gamma_i^2\delta_i^2,
\]

Put $d_i=\alpha_i-a$.  The exact identities
\[
 \Jfour=\sum_i d_i^2(\alpha_i^2+2a\alpha_i+3a^2)
 =12\int_0^1(1-t)\sum_i(a+td_i)^2d_i^2\,dt
\]
imply $\max_i|d_i|^4\le C\Jfour$.  If
$\Jfour\ge S^{17/4}$, subtract the two bounds
\eqref{JP:eq:normal-sine-absolute}.  Otherwise, $a\asymp S$ and
$\max_i|d_i|/a\le CS^{1/16}$, so the whole radial segment is
quasiuniform.  Cyclic symmetry kills
$DR_{\sin}(a\mathbf1)[d]$ because $\sum_i d_i=0$.  Taylor integration and
the Hessian bound give the $CS^{1/2}\Jfour$ estimate in the near regime;
the far regime gives $CS^{1/4}\Jfour$.  This proves
\eqref{JP:eq:normal-sine-relative}.
\end{proof}

\begin{lemma}[Normal-sine Bessel and secant transfer]
\label{JP:lem:normal-bessel-secant}
With $\widehat P_D^{\rm ns}=(2\pi d^2)^{-1}P_D^{\rm ns}$, there is a
modulus $\omega_D(S)\to0$ such that
\begin{equation}\label{JP:eq:normal-bessel-secant-relative}
 \left|\{\widehat P_D^{\rm ns}(\alpha)-P_{+,\sin}^{\rm e}(\alpha)\}
 -\{\widehat P_D^{\rm ns}(\alpha_*)-P_{+,\sin}^{\rm e}(\alpha_*)\}\right|
 \le\omega_D(S)\Jfour(\alpha).
\end{equation}
\end{lemma}

\begin{proof}
Put
\[
 W_\alpha=\Pi_{\rm ng}\mathcal I_\alpha Z_\alpha,
 \qquad
 \widehat q_D=(2\pi d^2)^{-1}q_D.
\]
Because exact support translations belong to $Z_\alpha$ and are carried by
$\mathcal I_\alpha$ to the first harmonics, minimizing $q_+$ over the
unprojected sine space is the same as minimizing it over $W_\alpha$ after
applying $\Pi_{\rm ng}$.

We first construct the required no-ratio quasi-interpolant.  For
$\psi\in H_{\rm ng}^{3/2}:=\Pi_{\rm ng}H^{3/2}(\mathbb T)$, let
$\phi=P_{\le M}\psi$, $M=\lfloor S^{-1}\rfloor$, sample $phi$ at the
normal angles, and use the linear interpolant
\eqref{JP:eq:normal-linear-spline}.  Fourier truncation gives
\begin{equation}\label{JP:eq:normal-QI-tail}
 \|\psi-\phi\|_{H^{1/2}}\le CS\|\psi\|_{H^{3/2}},
 \qquad
 \|\phi''\|_{L^2}\le CS^{-1/2}\|\psi\|_{H^{3/2}}.
\end{equation}
Taylor's integral formula on each interval, followed by summation, gives
\[
 \|\phi-\mathcal I_\alpha^0\phi|_{\{\theta_i\}}\|_{L^2}
 \le CS^2\|\phi''\|_{L^2},\qquad
 \|\phi-\mathcal I_\alpha^0\phi|_{\{\theta_i\}}\|_{H^1}
 \le CS\|\phi''\|_{L^2}.
\]
The constants are independent of cell ratios because each point belongs to
one normal interval and the estimate is charged to that same interval.
Interpolation between $L^2$ and $H^1$, the tail
\eqref{JP:eq:normal-QI-tail}, the exact mean correction in the homogeneous
area row, and the graph \eqref{JP:eq:normal-sine-graph-close} therefore
produce
$Q_\alpha\psi\in\Pi_{\rm ng}\mathcal I_\alpha Z_\alpha$ with
\begin{equation}\label{JP:eq:normal-QI}
 \|\psi-Q_\alpha\psi\|_{H^{1/2}}
 \le CS\|\psi\|_{H^{3/2}}.
\end{equation}
No overlap or ratio constant occurs.

Write the normalized disk multiplier as
\[
 \widehat q_D(f)
 =q_+(f)-k(f),\qquad
 k(f)=\sum_{|n|\ge2}k_{|n|}|\widehat f_n|^2,
 \qquad 0<k_n<\frac4n.
\]
Let $f_\alpha$ be the $q_+$-minimizing residual of
$\Pi_{\rm ng}G_\alpha^{\rm e}+W_\alpha$.  If
$g\in H_{\rm ng}^{1/2}$ and $(|D|+1)\psi=g$, Galerkin orthogonality and
\eqref{JP:eq:normal-QI} give
\[
 |\langle f_\alpha,g\rangle|
 =|q_+(f_\alpha,\psi-Q_\alpha\psi)|
 \le CS\|f_\alpha\|_{H^{1/2}}\|g\|_{H^{1/2}}.
\]
Taking the dual supremum proves the negative-norm gain
\[
 \|f_\alpha\|_{H^{-1/2}}
 \le CS\|f_\alpha\|_{H^{1/2}}.
\]
Since $q_+(f_\alpha)\le q_+(G_\alpha^{\rm e})\le CS^3$, the exact
Hilbert Schur identity
\begin{align*}
 R_B(\alpha)&:=P_{+,\sin}^{\rm e}(\alpha)-P_D^{\rm e}(\alpha)\\
 &=k(f_\alpha,f_\alpha)
  +\langle g_\alpha,
     (\widehat q_D|_{W_\alpha})^{-1}g_\alpha\rangle,
 &g_\alpha(v)&=k(f_\alpha,v),
\end{align*}
where
$P_D^{\rm e}=\inf_{v\in W_\alpha}\widehat q_D
(\Pi_{\rm ng}G_\alpha^{\rm e}+v)$, gives
\begin{equation}\label{JP:eq:normal-bessel-absolute}
 0\le R_B(\alpha)\le CS^5.
\end{equation}

For the source replacement, the exact decomposition on every side-normal
cell is
\[
 h_\alpha\sec x-1=(h_\alpha-1)+G_\alpha^{\rm e}+E_\alpha,
 \qquad E_\alpha=h_\alpha(\sec x-1)-\frac{x^2}{2}.
\]
Here $x$ is measured from the side normal, so $G_\alpha^{\rm e}=x^2/2$
on both incident half-cells.  Taylor's formula through one normalized
derivative gives
\[
 |E_\alpha|+w|E_\alpha'|\le CS^2w^2
\]
on a side cell of total angular width $w$.  Splitting the Gagliardo
integral into same, adjacent, and dyadically separated cells yields the
ratio-free estimate
\begin{equation}\label{JP:eq:normal-secant-remainder}
 \|E_\alpha\|_{H^{1/2}}
 \le CS^2\|G_\alpha^{\rm e}\|_{H^{1/2}}.
\end{equation}
The constant term $h_\alpha-1$ disappears under $\Pi_{\rm ng}$.
Coercivity of $\widehat q_D$ and the Hilbert distance inequality therefore
give
\begin{equation}\label{JP:eq:normal-secant-absolute}
 |R_{\rm sec}(\alpha)|\le CS^5,
 \qquad R_{\rm sec}=\widehat P_D^{\rm ns}-P_D^{\rm e}.
\end{equation}

We finally turn both absolute corrections into relative ones.  In a
quasiuniform material band, pull all intervals to a fixed real fan and put
$h=\|\delta\|_{\ell^\infty}/a$.  The first two material derivatives of
$q_+$ are bounded by $Ch^r$ in $H^{1/2}$, while those of $k$ are bounded by
$Ch^r$ in $H^{-1/2}$; the latter follows from
$|\Delta^rk_n|\le C_r(1+n)^{-1-r}$ and discrete summation by parts.
Differentiating the Galerkin equation and the negative-norm estimate gives,
for $r=0,1,2$,
\[
 \|\partial_t^rf_t\|_{H^{1/2}}\le Ch^ra^{3/2},
 \qquad
 \|\partial_t^rf_t\|_{H^{-1/2}}\le Ch^ra^{5/2}.
\]
Twice differentiating the displayed Schur identity gives
\[
 |D^2R_B(\gamma)[\delta,\delta]|
 \le Ca\sum_i\gamma_i^2\delta_i^2.
\]
For $R_{\rm sec}$, affine normalization of the two source profiles gives
\[
 \|\partial_t^rG_{\gamma(t)}^{\rm e}\|_{H^{1/2}}
 \le Ch^ra^{3/2},\qquad
 \|\partial_t^rE_{\gamma(t)}\|_{H^{1/2}}
 \le Ch^ra^{7/2},\qquad r=0,1,2.
\]
Differentiating the two coercive Schur equations in
\eqref{JP:eq:normal-secant-absolute} consequently gives the same Hessian
bound for $R_{\rm sec}$.  Hence their sum
$R_D=\widehat P_D^{\rm ns}-P_{+,\sin}^{\rm e}$ satisfies
\begin{equation}\label{JP:eq:normal-RD-hessian}
 |D^2R_D(\gamma)[\delta,\delta]|
 \le Ca\sum_i\gamma_i^2\delta_i^2.
\end{equation}

If $\Jfour\ge S^{9/2}$, subtract the endpoint bounds
\eqref{JP:eq:normal-bessel-absolute} and
\eqref{JP:eq:normal-secant-absolute}; this gives
$CS^{1/2}\Jfour$.  Otherwise the exact fourth-defect identity used in the
preceding lemma makes the radial segment quasiuniform.  Cyclic symmetry
removes the first derivative of $R_D$ at the regular fan, and Taylor's
formula with \eqref{JP:eq:normal-RD-hessian} gives $CS\Jfour$.  This proves
\eqref{JP:eq:normal-bessel-secant-relative} with
$\omega_D(S)=C(S^{1/2}+S)$.
\end{proof}

\begin{theorem}[Closed fan-normal disk endpoint and principal forcing]
\label{JP:thm:normal-disk-endpoint}
There are universal $S_D,\kappa,C_p>0$ such that every positive fan with
$S\le S_D$ satisfies
\begin{align}
 P_D^{\rm ns}(\alpha)-P_D^{\rm ns}(\alpha_*)
 &\ge2\kappa\Jfour(\alpha),
 \label{JP:eq:D0ns}\\
 \|p_\alpha\|_{(\Gjp_\alpha)^*}
 &\le C_p\Jfour(\alpha)^{1/2}.
 \label{JP:eq:JE-P}
\end{align}
Both constants are independent of all minimum angles and adjacent ratios.
\end{theorem}

\begin{proof}
Combine \eqref{JP:eq:normal-homogeneous-gap},
\eqref{JP:eq:normal-sine-relative}, and
\eqref{JP:eq:normal-bessel-secant-relative}, then decrease $S_D$ so that the
two moduli consume at most half of the homogeneous margin.  Multiplication
by $2\pi d^2$ gives \eqref{JP:eq:D0ns} for a universal $\kappa>0$.

The closed exact source envelope from F129 gives
\[
 q_D(B_\alpha)-q_D(B_{\alpha_*})\le C\Jfour(\alpha).
\]
Hilbert completion and $p_{\alpha_*}=0$ give
\[
 \frac14\|p_\alpha\|_{(\Gjp_\alpha)^*}^2
 =\{q_D(B_\alpha)-q_D(B_{\alpha_*})\}
  -\{P_D^{\rm ns}(\alpha)-P_D^{\rm ns}(\alpha_*)\}.
\]
Drop the favorable second brace and take square roots to obtain
\eqref{JP:eq:JE-P}.
\end{proof}

\begin{theorem}[Closed exact fan-normal value matching]
\label{JP:thm:normal-value-matching}
There is a modulus $\varepsilon_V(S)\to0$, independent of all minimum
angles and adjacent ratios, such that every sufficiently fine positive fan
satisfies
\begin{equation}\label{JP:eq:JE-V}
 \boxed{
 \left|\Delta_\alpha-
 \{q_D(B_\alpha)-q_D(B_{\alpha_*})\}\right|
 \le\varepsilon_V(S)\Jfour(\alpha).}
\end{equation}
\end{theorem}

\begin{proof}
Let $W_\alpha=B_\alpha e_r$ be the radial physical boundary field of the
area-normalized tangential polygon; its normal scalar trace at the disk is
the source $B_\alpha$ in \eqref{JP:eq:disk-source-B}.  The exact
disk-normalized scalar expansion used in JP--BU may be retained
at its natural orders rather than weakened to a one-sided bound.  It is
\begin{equation}\label{JP:eq:value-matching-split}
 \Delta_\alpha
 =\{q_D(B_\alpha)-q_D(B_{\alpha_*})\}
  +\{\Delta_3(\alpha)-\Delta_3(\alpha_*)\}
  +\{F_4(\alpha)-F_4(\alpha_*)\},
\end{equation}
where $\Delta_3$ is the complete scalar cubic row, including its factor
$1/6$, and $F_4$ is the exact fourth-and-higher scalar remainder.

The closed CT cubic identity splits $\Delta_3$ into its principal null form
and the complete normal/geometric order-one remainder.  CT's relative
closure theorem controls the latter by $o(\Jfour)$.  The QB
secant-to-quadratic transfer and the closed DF theorem control the principal
null form by the same scale.  Consequently there is a modulus
$\varepsilon_3(S)\to0$ with
\begin{equation}\label{JP:eq:value-cubic-relative}
 |\Delta_3(\alpha)-\Delta_3(\alpha_*)|
 \le\varepsilon_3(S)\Jfour(\alpha).
\end{equation}
This is the pure scalar row, so no support-metric comparison is involved.

It remains to sharpen the tail estimate already used in JP--BU.  On the
real chart and its material complex disk, the tame scalar remainder gives
\begin{equation}\label{JP:eq:value-tail-absolute}
 |F_4(\zeta)|
 \le C\|W_\zeta\|_{H^{1/2}}^2
       \|W_\zeta\|_{W^{1,\infty}}^2
 \le CS^5.
\end{equation}
Put $d_i=\alpha_i-a$ and $\gamma_i(t)=a+td_i$.  If
$\Jfour\ge S^{9/2}$, applying \eqref{JP:eq:value-tail-absolute} at the two
endpoints gives
\[
 |F_4(\alpha)-F_4(\alpha_*)|
 \le CS^{1/2}\Jfour(\alpha).
\]
If $\Jfour<S^{9/2}$, the exact fourth-defect identity implies
$a\asymp S$ and keeps the whole radial segment quasiuniform.  The material
complex radius in direction $d$ is at least
$ca/\|d\|_{\ell^\infty}$.  Cauchy's estimate and
\eqref{JP:eq:value-tail-absolute} give, with
$E(t)=\sum_i\gamma_i(t)^2d_i^2$,
\[
 |F_4''(t)|
 \le CS^5\frac{\|d\|_{\ell^\infty}^2}{a^2}
 \le CS E(t).
\]
Cyclic symmetry makes $F_4'(0)$ constant in the fan coordinates and hence
$F_4'(0)[d]=0$.  Taylor integration and
$\Jfour=12\int_0^1(1-t)E(t)\,dt$ give the near estimate
$CS\Jfour$.  Together with \eqref{JP:eq:value-cubic-relative} and
\eqref{JP:eq:value-matching-split}, this proves \eqref{JP:eq:JE-V} with
$\varepsilon_V(S)=\varepsilon_3(S)+C(S^{1/2}+S)$.
\end{proof}

\begin{proposition}[Exact forcing-residual covector]
\label{JP:prop:forcing-residual-split}
On the polar side cell $K_i$, put
\[
 f_\alpha(\theta)=h_\alpha\sec(\theta-\theta_i)-1
\]
and let
$Z_\alpha(z)=v_\alpha(z)e_r+w_\alpha(z)e_\theta$ be the continuous
physical-affine support displacement.  In the notation of the complete CT
tensor, let $\kappa_{\mathbb D}\mathcal C+\mathcal R_{\rm full}$ be the
normal third-derivative split and set
\[
 \Psi_{4,\alpha}[z]
 :=D\mathcal R_4(f_\alpha e_r)[Z_\alpha(z)].
\]
Then the exact residual $e_\alpha=\ell_\alpha-p_\alpha$ is
\begin{equation}\label{JP:eq:exact-forcing-residual}
\boxed{
\begin{aligned}
 e_\alpha[z]={}&
 2\langle f_\alpha,
       v_\alpha(z)-\mathcal I_\alpha z\rangle_D\\
 &+\frac{\kappa_{\mathbb D}}2
   \mathcal C(f_\alpha,f_\alpha,v_\alpha(z))
  +\frac12\mathcal R_{\rm full}
   (f_\alpha,f_\alpha,v_\alpha(z))\\
 &-2\langle f_\alpha,
       w_\alpha(z)f_\alpha'\rangle_D
  +\Psi_{4,\alpha}[z].
\end{aligned}}
\end{equation}
Every term is evaluated after the exact area/translation projection.  The
formula is an identity of covectors on $\mathsf X_\alpha$; in particular,
it contains no norm comparison and no asymptotic remainder notation.
\end{proposition}

\begin{proof}
The boundary of the tangential polygon is the disk boundary displaced by
$W_\alpha=f_\alpha e_r$.  Because vertices and side parametrizations are
affine in the support numbers, differentiation in the support direction is
the single physical slot $Z_\alpha(z)$.  The exact scalar Taylor identity
therefore gives
\[
 \ell_\alpha[z]
 =D^2\mathcal L(0)[W_\alpha,Z_\alpha(z)]
 +\frac12D^3\mathcal L(0)
       [W_\alpha,W_\alpha,Z_\alpha(z)]
 +\Psi_{4,\alpha}[z].
\]
The disk Hessian is
$D^2\mathcal L(0)[W_\alpha,Z_\alpha(z)]
=2\langle f_\alpha,v_\alpha(z)\rangle_D$.
The complete physical-coordinate CT identity gives
\[
\begin{aligned}
 \frac12D^3\mathcal L(0)
 [W_\alpha,W_\alpha,Z_\alpha(z)]
 ={}&\frac{\kappa_{\mathbb D}}2
   \mathcal C(f_\alpha,f_\alpha,v_\alpha(z))\\
 &+\frac12\mathcal R_{\rm full}
   (f_\alpha,f_\alpha,v_\alpha(z))
 -2\langle f_\alpha,w_\alpha(z)f_\alpha'\rangle_D.
\end{aligned}
\]
The last coefficient is the curvature multiplicity after differentiating
the cubic Taylor coefficient: the raw third derivative contributes
$-4\langle f_\alpha,w_\alpha(z)f_\alpha'\rangle_D$ and is multiplied by
$1/2$.  Finally
$p_\alpha[z]=2\langle f_\alpha,\mathcal I_\alpha z\rangle_D$ because the
disk form annihilates the constant and first harmonics.  Subtraction proves
\eqref{JP:eq:exact-forcing-residual}.
\end{proof}

\begin{remark}[Retraction of the period-three material forcing transfer]
\label{JP:rem:period-three-material-forcing}
The former proposition which identified the three-cell material constant
$R_3$ with the genuine polygonal covector
$e_\alpha=\ell_\alpha-p_\alpha$ is withdrawn.  The revision-22
Schwarz--Christoffel audit finds a same-order corner finite part already in
the support diagonal, so the material-to-polygon transfer used below does
not justify the mixed forcing coefficient either.  The displayed material
series remains a useful diagnostic, but it proves neither truth nor
falsity of the former JE--FM estimate
\begin{equation}\label{JP:eq:JE-FM}
 \|e_\alpha\|_{(\mathfrak G_\alpha^{\rm ns})^*}
 \le\varepsilon(S)\Jfour(\alpha)^{1/2},
 \qquad \varepsilon(S)\longrightarrow0.
\end{equation}
JE--FM is absent from the active
DAG, and no active implication uses this archived calculation.
\end{remark}

\iffalse
For a three-periodic material function define
\[
 \widehat h(n)=\frac13\int_0^3h(y)e^{-2\pi iny/3}\,dy,
 \qquad
 \langle f,g\rangle_3
 =\sum_{n\ne0}\left|\frac{2\pi n}{3}\right|
       \widehat f(-n)\widehat g(n).
\]
On a normal cell of length $l_i$, let $G_l$ be the edge quadratic
source, let $S_Z$ equal $Z_i$ and $Z_{i+1}$ on the two half-cells, and let
$I_l^0Z$ be the linear interpolant.  If $\varphi_i$ is the cell midpoint,
direct integration gives, for $\omega=2\pi n/3\ne0$,
\begin{align}
 \widehat G_l(\omega)
 &=\frac1{3\omega^2}\sum_i l_i e^{-i\omega\varphi_i},\label{JP:eq:G-Fourier}\\
 \widehat{(S_Z-I_l^0Z)}(\omega)
 &=\frac1{3i\omega}\sum_i(Z_{i+1}-Z_i)e^{-i\omega\varphi_i}
 \left(1-\operatorname{sinc}\frac{\omega l_i}{2}\right).
 \label{JP:eq:step-sine-Fourier}
\end{align}
We now give the alias calculation, including the realification.  Put
\[
 \kappa=\frac{2\pi}{3},\qquad
 D_i^+=e^{i\kappa i},\qquad Z_i^-=e^{-i\kappa i},
\]
\[
 \Delta=e^{-i\kappa}-1,\qquad
 A=\frac1{e^{i\kappa}-1}+\frac12=-\frac{i}{2\sqrt3}.
\]
For the complexified material family,
\[
 l_i(t)=1+tD_i^+,\qquad
 \varphi_i(t)=i+\frac12+tAD_i^+.
\]
Let $H_t=S_{Z^-}-I_{l(t)}^0Z^-$ and
$\omega_m=2\pi(m-\frac13)$.  Direct substitution in
\eqref{JP:eq:G-Fourier}--\eqref{JP:eq:step-sine-Fourier} gives
\begin{align}
 \widehat H_0(\omega_m)
 &=\frac{\Delta e^{-i\omega_m/2}}{i\omega_m}
 \left(1-\operatorname{sinc}\frac{\omega_m}{2}\right),
 \label{JP:eq:R3-H-first}\\
 \left.\partial_t\right|_0\widehat G_t(-\omega_m)
 &=\frac{e^{i\omega_m/2}}{\omega_m^2}
 \left(1+\frac{\omega_m}{2\sqrt3}\right).
 \label{JP:eq:R3-G-first}
\end{align}
The second alias family comes from the regular source and the material
derivative of $H_t$.  For $m\ne0$,
\begin{align}
 \widehat G_0(2\pi m)&=\frac{(-1)^m}{(2\pi m)^2},
 \label{JP:eq:R3-G-second}\\
 \left.\partial_t\right|_0\widehat H_t(2\pi m)
 &=\frac{\Delta(-1)^m}{i}
 \left(-\frac1{2\sqrt3}-\frac{(-1)^m}{2\pi m}\right).
 \label{JP:eq:R3-H-second}
\end{align}
Consequently the complex bilinear derivative is
\begin{equation}\label{JP:eq:R3-alias-complex}
 \left.\partial_t\right|_0\langle G_t,H_t\rangle_3
 =\frac{\Delta}{i}R_3,
\end{equation}
where the two sums below are taken with the same symmetric cutoff:
\begin{align}
 R_3={}&
 \sum_{m\in\mathbb Z}
 \frac{\operatorname {sgn}\omega_m}{\omega_m^2}
 \left(1+\frac{\omega_m}{2\sqrt3}\right)
 \left(1-\operatorname{sinc}\frac{\omega_m}{2}\right)
 \notag\\
 &+\sum_{m\ne0}\frac1{2\pi|m|}
 \left(-\frac1{2\sqrt3}-\frac{(-1)^m}{2\pi m}\right).
 \label{JP:eq:R3-alias-sum}
\end{align}
This joint cutoff is essential: the logarithmic terms cancel between the
two alias families.

For $x_m=m-\frac13$,
\[
 1-\operatorname{sinc}\frac{\omega_m}{2}
 =1+\frac{(-1)^m\sqrt3}{2\pi x_m}.
\]
Expanding \eqref{JP:eq:R3-alias-sum} gives
\begin{equation}\label{JP:eq:R3-four-sums}
 R_3=\frac{S_1}{4\pi^2}
 +\frac{S_0}{4\pi\sqrt3}
 +\frac{\sqrt3S_3}{8\pi^3}
 +\frac{S_2}{8\pi^2},
\end{equation}
where, with the same symmetric convention,
\[
\begin{aligned}
 S_1&=\sum_{m\in\mathbb Z}
 \frac{\operatorname {sgn}x_m}{x_m^2},&
 S_0&=\sum_{m\in\mathbb Z}\frac1{|x_m|}
      -\sum_{m\ne0}\frac1{|m|},\\
 S_3&=\sum_{m\in\mathbb Z}
 \frac{(-1)^m\operatorname {sgn}x_m}{x_m^3},&
 S_2&=\sum_{m\in\mathbb Z}
 \frac{(-1)^m\operatorname {sgn}x_m}{x_m^2}.
\end{aligned}
\]
The remaining odd summand satisfies
$\sum_{m\ne0}(-1)^m/(|m|m)=0$.  Pairing the positive and negative
indices and then grouping by residue classes modulo $3$ and $6$ yields
\begin{equation}\label{JP:eq:R3-residue-sums}
 S_1=-9L(2,\chi_3),\qquad S_0=3\log3,\qquad
 S_3=\frac{39}{2}\zeta(3),\qquad
 S_2=-\frac{27}{2}L(2,\chi_3).
\end{equation}
Equations \eqref{JP:eq:R3-four-sums}--\eqref{JP:eq:R3-residue-sums} give
\eqref{JP:eq:R3-closed}.

The real family in the proposition is obtained from
\[
 D=\operatorname {Re}D^+,\qquad
 Z=\operatorname {Re}(e^{-i\kappa}Z^-).
\]
Real bilinearity and \eqref{JP:eq:R3-alias-complex} therefore give the
separate realification identity
\begin{equation}\label{JP:eq:R3-alias-real}
 \left.\partial_t\right|_0
 \langle G_{\alpha(t)},S_Z-I_{\alpha(t)}^0Z\rangle_3
 =\frac12\operatorname {Re}
 \left(e^{-i\kappa}\frac{\Delta}{i}\right)R_3
 =\frac{\sqrt3}{2}R_3.
\end{equation}

The sign in \eqref{JP:eq:R3-closed} is rigorous: write
\[
 L(2,\chi_3)=\sum_{k\ge0}
 \frac{6k+3}{(3k+1)^2(3k+2)^2}.
\]
After the first $1000$ terms its tail lies between zero and
$1/(18\cdot999^2)$.  Integral tails after $1000$ terms for $\zeta(3)$,
the positive series
\[
\log3=2\sum_{j\ge0}\frac{2^{-(2j+1)}}{2j+1},
\]
the rational enclosures
\[
 \frac{333}{106}<\pi<\frac{355}{113},\qquad
 \frac{1732050807}{10^9}<\sqrt3<\frac{1732050808}{10^9}
\]
give the displayed rational interval for $R_3$.

On each side cell, in the material coordinate
$x\in[-1/2,1/2]$, one has, also after one $t$ derivative,
\begin{equation}\label{JP:eq:R3-cell-expansions}
\begin{gathered}
 f_\alpha=a^2\left(\frac{x^2}{2}-\frac1{24}\right)+O(a^4),\\
 v_\alpha(Z)-w_\alpha(Z)f_\alpha'=S_Z+O(a^2),
 \qquad \mathcal I_\alpha Z=I^0Z+O(a^2).
\end{gathered}
\end{equation}
The corresponding material norms satisfy
\begin{equation}\label{JP:eq:R3-material-scales}
 \|f_\alpha\|_\infty=O(a^2),\quad
 \|f_\alpha'\|_\infty=O(a),\quad
 \|f_\alpha\|_{H^{1/2}}=O(a^{3/2}),\quad
 \|v_\alpha(Z)\|_{H^{1/2}}=O(a^{-1/2}).
\end{equation}
Thus the unique $at$ term in
\eqref{JP:eq:exact-forcing-residual} is
$2q_D(f_\alpha,S_Z-I^0Z)$.  Indeed, the high-frequency symbol
\[
 \mathcal D_\circ+1=|D|+1+O(|D|^{-1})
\]
and the $2\pi/(3a)$ material periods show that the $|D|$ row is
$4\pi d_\circ^2at$ times \eqref{JP:eq:R3-alias-real}; the $+1$ row is
$O(a^2|t|)$, while the order-minus-one row and the errors in
\eqref{JP:eq:R3-cell-expansions} are $O(a^3|t|)$.

For completeness, the one-Lipschitz null-form estimate gives
\[
 |\mathcal C(f,f,v)|
 \le C\|f'\|_\infty
       \|f\|_{H^{1/2}}\|v\|_{H^{1/2}}=O(a^2).
\]
At $t=0$ the period-one source is Fourier-orthogonal to the period-three
mode, so along the family
\begin{equation}\label{JP:eq:R3-lower-orders}
 \mathcal C(f_t,f_t,v_t)=O(a^2|t|),\qquad
 \mathcal R_{\rm full}(f_t,f_t,v_t)=O(a^3|t|).
\end{equation}
We record the special-family fourth-tail estimate rather than invoking
the retracted general support-covector row CT--M4.  Write the analytic
disk expansion as
\[
 \Psi_{4,\alpha(t)}[Z]
 =\sum_{k\ge4}\frac1{(k-1)!}
 D^k\mathcal L(0)
 [f_{\alpha(t)}^{\,k-1},Z_{\alpha(t)}(Z)].
\]
After rescaling each of the three material cells to unit length,
$f=a^2F_{a,t}$ and $f'=a\,\partial_xF_{a,t}$, with the profiles
holomorphic and uniformly bounded in $t$ on a fixed complex
neighbourhood.  On the
physical circle the relevant full-vector bounds are
\begin{equation}\label{JP:eq:R3-full-vector-scales}
\begin{gathered}
 \|f_{\alpha(t)}e_r\|_{H^{1/2}}\le Ca^{3/2},\qquad
 \|f_{\alpha(t)}e_r\|_{W^{1,\infty}}\le Ca,\\
 \|Z_{\alpha(t)}(Z)\|_{H^{1/2}}\le Ca^{-3/2}.
\end{gathered}
\end{equation}
To see the last scale, the normal component has period $3a$ and size one,
while the
tangential endpoint coefficients are $O(a^{-1})$; the piecewise-affine
full vector therefore has $H^{1/2}$ size $O(a^{-3/2})$.  We do not place
this support slot in $W^{1,\infty}$, where its derivative is
$O(a^{-2})$.

The fixed-chart analytic Hessian estimate is bilinear in two
$H^{1/2}$ slots and analytic in the remaining
$W^{1,\infty}$ slots.  Applying its Cauchy bound with one copy of $f$
and the support direction in the two Hilbert slots gives, uniformly for
$|t|\le1/2$,
\begin{equation}\label{JP:eq:R3-special-tail-bound}
 \frac1{(k-1)!}
 \left|D^k\mathcal L(0)
 [f_{\alpha(t)}^{\,k-1},Z_{\alpha(t)}(Z)]\right|
 \le C_0C_1^k
 \|f e_r\|_{H^{1/2}}\|Z_\alpha(Z)\|_{H^{1/2}}
 \|f e_r\|_{W^{1,\infty}}^{k-2}
 \le C_0C_1^k a^{k-2}.
\end{equation}
For $C_1a<1/2$, summing \eqref{JP:eq:R3-special-tail-bound} gives
$|\Psi_4(t)|\le Ca^2$.

The derivative in $t$ is taken only after passing to the fixed
three-cell material representation.  There the side-to-material maps,
Jacobians, and tensor coefficients are rational holomorphic functions of
$t$ on a disk $|t|<r$ with $r>0$ independent of $a$.  Thus the
\emph{scalar} function $t\mapsto\Psi_4(t)$ is holomorphic on that disk,
and the preceding $Ca^2$ bound holds uniformly there.  One-variable
Cauchy gives $|\Psi_4'(t)|\le C_ra^2$ for $|t|<r/2$.  This argument does
not differentiate a moving-knot physical field in $H^{1/2}$.
The value at $t=0$ vanishes by cyclic Fourier orthogonality, hence
\[
 \Psi_{4,\alpha(t)}[Z]=O(a^2|t|).
\]
The nonlinear material geometry contributes $O(at^2)$.  Combining these
estimates with \eqref{JP:eq:R3-alias-real} first proves the asserted
expansion for the unprojected $Z$.

It remains to check the exact quotient projection.  From
\eqref{JP:eq:scale-section} and the tangential side lengths,
\begin{align}
 s_{\alpha(t)}(Z)
 &=\frac{h_{\alpha(t)}}{2\pi}
 \sum_i\tan\frac{\alpha_i(t)}2\,(Z_i+Z_{i+1})\notag\\
 &=-\frac{h_{\alpha(t)}}4t+O(a^2|t|+t^2).
 \label{JP:eq:R3-scale-projection}
\end{align}
Thus, for a translation vector $\tau_{\alpha(t)}$,
\begin{equation}\label{JP:eq:R3-projected-Z}
 Z_{\alpha(t)}
 =Z+\frac t4\mathbf1+\tau_{\alpha(t)}
   +O(a^2|t|+t^2)\mathbf1.
\end{equation}
Both covectors annihilate translations, and scale invariance gives
$\ell_\alpha[\mathbf1]=0$; in contrast, $p_\alpha[\mathbf1]$ need not
vanish.  On the present quasiuniform family the material-cell expansions
\[
 B_{\alpha(t)}=a^2F_{a,t}+O(a^4),\qquad
 K_{\alpha(t)}\mathbf1=a^2C_{a,t}+O(a^4)
\]
have uniformly bounded three-periodic profiles.  Since $q_D$ has
material order one,
\begin{equation}\label{JP:eq:R3-scale-p}
 |p_{\alpha(t)}[\mathbf1]|
 =2|\langle B_{\alpha(t)},K_{\alpha(t)}\mathbf1\rangle_D|
 \le Ca^3.
\end{equation}
Equations \eqref{JP:eq:R3-projected-Z}--\eqref{JP:eq:R3-scale-p} change
$e_\alpha[Z]$ by only $O(a^3|t|+a^3t^2)$, which is contained in
$O(a^2|t|+at^2)$.  This proves \eqref{JP:eq:period-three-e} for the exact
projected direction.

Finally
\begin{align*}
 \Jfour(\alpha(t))&=6\pi a^3t^2+O(a^3|t|^3),\\
 \mathfrak G_{\alpha(t)}^{\rm ns}(Z_{\alpha(t)})
 &=\frac{2\pi d_\circ^2}{a}Q_3+O(1),
 \qquad Q_3=\langle I^0Z,I^0Z\rangle_3>0.
\end{align*}
Hence the quotient of the left side of \eqref{JP:eq:period-three-e} by
$\Jfour^{1/2}(\mathfrak G^{\rm ns})^{1/2}$ tends in absolute value to
$|d_\circ|R_3/\sqrt{Q_3}>0$.  Taking $t=a$ keeps all angles positive and
gives $\Jfour\asymp a^5<S^{9/2}$, so the counterexample lies in the near
branch and has adjacent ratios tending to one.
\end{proof}
\fi

The closed disk Schur, forcing, and value rows are retained as JP--DE.
The unproved separated flux interface is absent from the active DAG; this
does not remove the fan-normal metric.  The replacement is the fixed-loss
joint scalar comparison in
\texttt{JC\_joint\_comparison.tex}.

\begin{remark}[Retraction of the period-three joint-loss transfer]
\label{JP:rem:period-three-material-joint-loss}
The former period-three joint-loss proposition is also withdrawn.  Its
coefficient $E_3$ is the material cell value, while the true polygonal
support finite part is $P(1/3)/2=9\log3/(4\pi)$.  Moreover the material
constant $R_3$ has not been transferred through the mixed polygonal corner
layer.  Consequently the limit formerly labelled
\textup{``joint-period3-loss''} is not a theorem about the genuine
polygonal residual.  The active JC--PL formulation already uses a fixed
absorbable loss and does not depend on this archived diagnostic.
\end{remark}

\iffalse
The exact regular-fan leading terms furnished by the displayed material
profiles are
\begin{align*}
 D_h^2\Phi_{\alpha_*}[Z,Z]
 &=\frac{4\pi d_\circ^2}{a}E_3+O(1),\\
 \mathfrak G_{\alpha_*}^{\rm ns}(Z)
 &=\frac{2\pi d_\circ^2}{a}Q_3+O(1).
\end{align*}
Here is an explicit certificate for \eqref{JP:eq:E3Q3-bounds}.  Put
\[
 \omega_n=\frac{2\pi n}{3},\qquad
 J_k(\omega)=\int_{-1/2}^{1/2}x^ke^{-i\omega x}\,dx.
\]
Then
\begin{align}
 J_0(\omega)&=\frac{2\sin(\omega/2)}{\omega},\label{JP:eq:J0-recurrence}\\
 J_k(\omega)&=
 -\frac{2^{-k}e^{-i\omega/2}-(-2)^{-k}e^{i\omega/2}}{i\omega}
 +\frac{k}{i\omega}J_{k-1}(\omega).\label{JP:eq:Jk-recurrence}
\end{align}
For a cell polynomial $P_r(x)=\sum_{j=0}^2p_{rj}x^j$,
\begin{equation}\label{JP:eq:poly-Fourier}
 \widehat P(n)=\frac13\sum_{r=0}^2e^{-i\omega_nr}
 \sum_{j=0}^2p_{rj}J_j(\omega_n).
\end{equation}
The three cell polynomials are
\[
\begin{array}{c|ccc}
r&0&1&2\\ \hline
V_r&-\frac12-\frac34x+\frac32x^2&
 -\frac12+\frac34x+\frac32x^2&1-3x^2\\
U_r&-\frac34x+\frac32x^2&
 \frac34x+\frac32x^2&-3x^2\\
K_r&-\frac9{16}+\frac92x-\frac{27}{4}x^2&
 -\frac9{16}-\frac92x-\frac{27}{4}x^2&-27x^2.
\end{array}
\]
For the linear interpolant, put
\[
 L_0(\omega)=\frac{1-e^{-i\omega}}{i\omega},\qquad
 L_1(\omega)=-\frac{e^{-i\omega}}{i\omega}
              +\frac{L_0(\omega)}{i\omega};
\]
then
\begin{equation}\label{JP:eq:I-Fourier}
 \widehat I(n)=\frac13\sum_{r=0}^2e^{-i\omega_nr}
 \{Z_rL_0(\omega_n)+(Z_{r+1}-Z_r)L_1(\omega_n)\}.
\end{equation}
Define
\[
\begin{aligned}
 A_3&=\sum_{n\ne0}|\omega_n||\widehat V(n)|^2,&
 B_3&=\sum_{n\ne0}|\omega_n|
       \overline{\widehat F(n)}\widehat K(n),\\
 C_3&=\sum_{n\ne0}|\omega_n|
       \overline{\widehat V(n)}\widehat U(n),&
 Q_3&=\sum_{n\ne0}|\omega_n||\widehat I(n)|^2.
\end{aligned}
\]
Thus $E_3=A_3+B_3-2C_3$.  Symmetric summation over
$0<|n|\le4096$, using
\eqref{JP:eq:J0-recurrence}--\eqref{JP:eq:I-Fourier}, gives the outward
intervals
\[
\begin{aligned}
0.76393070&<A_K<0.76393075,&
-0.26168518&<B_K<-0.26168512,\\
-0.13865959&<C_K<-0.13865952,&
0.56698459&<Q_K<0.56698465.
\end{aligned}
\]
Only $0,\pm1/2,\pm\sqrt3/2,\pm1$ occur in the trigonometric factors, so
these are finite rational interval calculations using rational outward
bounds for $\pi$ and $\sqrt3$.  Finally
\[
\operatorname {TV}(F')=6,\quad \operatorname {TV}(V')=15,\quad
\operatorname {TV}(I')=6,\quad
\operatorname {TV}(K)\le\frac{135}{4},\quad
\operatorname {TV}(U)\le\frac{15}{2}.
\]
Twice integrating the continuous profiles and once integrating the BV
profiles gives
\[
\begin{aligned}
0\le A_3-A_K&<1.63\cdot10^{-7},&
0\le Q_3-Q_K&<2.60\cdot10^{-8},\\
|B_3-B_K|&<0.002505,&
|C_3-C_K|&<0.001392.
\end{aligned}
\]
These intervals imply
$0.56698456<Q_3<0.56698468$ and
$0.2072<E_3-Q_3<0.2179$, hence
\eqref{JP:eq:E3Q3-bounds}.

We finish by checking the projection, active interval, and every Taylor
remainder on the stated path $t=a$.  Equations
\eqref{JP:eq:R3-projected-Z}--\eqref{JP:eq:R3-scale-p}, the analytic
three-cell material pullback, and the two leading terms at the start of
the proof give, uniformly for $|s|\le c_0a^2$,
\begin{align}
 \mathscr R_{\alpha(a)}(sZ_{\alpha(a)})
 ={}&
 \{2\pi\sqrt3\,d_\circ^2R_3a^2+O(a^3)\}s\notag\\
 &+\left\{\frac{2\pi d_\circ^2}{a}(E_3-Q_3)+O(1)\right\}s^2
 +O(a^{-2}|s|^3).
 \label{JP:eq:joint-period3-Taylor}
\end{align}
Here the $O(1)s^2$ term includes the material change
$H_{\alpha(a)}-H_{\alpha_*}$ and the scale projection; the latter uses
$K_\alpha\mathbf1=O(a^2)$ in material amplitude.  The last term follows
by contracting the third support derivative on the fixed three material
cells: its circular multipliers have total order at most two, hence its
norm on three-periodic support data is $O(a^{-2})$.

Side lengths are affine in the support numbers.  Their base size is
comparable to $a$, while a support increment $sZ_{\alpha(a)}$ changes
them by $O(|s|/a)$.  Choosing $c_0$ sufficiently small therefore keeps
the entire interval $|s|\le c_0a^2$ active.  The minimizer of the
quadratic part of \eqref{JP:eq:joint-period3-Taylor} is
\begin{equation}\label{JP:eq:joint-period3-minimizer}
 s_*=-\frac{\sqrt3R_3}{2(E_3-Q_3)}a^3+O(a^4),
\end{equation}
so $|s_*|\ll a^2$.  At this point the linear-coefficient error,
the $O(1)s^2$ term, the projection error, and the support cubic remainder
are respectively $O(a^6)$, $O(a^6)$, $O(a^6)$, and $O(a^7)$.
Since
\[
 \Jfour(\alpha(a))=6\pi a^5+O(a^6),
\]
substitution of \eqref{JP:eq:joint-period3-minimizer} into
\eqref{JP:eq:joint-period3-Taylor} gives the upper bound in
\eqref{JP:eq:joint-period3-loss}.  For the matching lower bound, throughout
$|s|\le c_0a^2$ the cubic term divided by the positive leading quadratic
term is $O(|s|/a)\le O(c_0a)$.  Hence outside
$|s|\le Ma^3$, with $M$ fixed sufficiently large, the positive quadratic
term dominates both the linear term and every remainder.  On the
remaining interval put $s=a^3\sigma$ and divide
\eqref{JP:eq:joint-period3-Taylor} by $a^5$.  The result converges uniformly
on bounded $\sigma$ to
\[
 2\pi\sqrt3\,d_\circ^2R_3\sigma
 +2\pi d_\circ^2(E_3-Q_3)\sigma^2.
\]
Its minimum, divided by
$\Jfour(\alpha(a))/a^5\to6\pi$, is exactly the right side of
\eqref{JP:eq:joint-period3-loss}.  This proves the matching lower bound and
the asserted limit of the local infimum.
\end{proof}
\fi

\begin{remark}[Why the closed TR/DM metric cannot be transferred]
TR completes its square in
$\widetilde G_\alpha(z)=q_D(\Pi_{\rm ng}\widetilde T_\alpha z)$, not in
$\Gjp_\alpha$.  These metrics have no no-ratio comparison in the needed
direction.  Indeed, take adjacent angles $\varepsilon\ll a$ and support
data $z_0=0$, $z_1=z_2=1$, return to zero across ordinary cells, and remove
the scale and translation modes.  The sine trace has only one sharp
transition, so
\[
 \Gjp_\alpha(z)\le C\{1+\log(a/\varepsilon)\}.
\]
On the adjacent physical cell, however,
 $a_1^-=(\sin\varepsilon)^{-1}-\tan(\varepsilon/2)\asymp\varepsilon^{-1}$.
On a subinterval of length comparable to $a$ this gives
$|\widetilde T_\alpha z|\ge ca/\varepsilon$.  Removing three low modes
does not remove this localized mass, and hence
\[
 \widetilde G_\alpha(z)\ge c\frac{a^3}{\varepsilon^2},
 \qquad
 \frac{\widetilde G_\alpha(z)}{\Gjp_\alpha(z)}\longrightarrow\infty.
\]
Thus the closed DM estimate in the old physical-radial metric cannot imply
the fan-normal dual estimate \eqref{JP:eq:JE-FM}.  This is a strict interface
mismatch, not a missing citation.
\end{remark}

\begin{remark}[Diagnostic check of the new Schur sign]
The reproducible script
\texttt{tests/check\_jp\_joint\_endpoint.py} samples the exact
finite-dimensional Schur complement in \eqref{JP:eq:D0ns-definition} on the
area/translation quotient.  For regular-nearby, alternating-tiny,
one-tiny-angle, and fixed-seed random fans with
$N=8,12,16,24$, the sampled ratio
\[
 \frac{P_D^{\rm ns}(\alpha)-P_D^{\rm ns}(\alpha_*)}
      {\Jfour(\alpha)}
\]
lies between $0.0189$ and $0.0333$; the sampled metric remains positive on
the quotient.  This independently checks the sign and normalization used
in Theorem~\ref{JP:thm:normal-disk-endpoint}; the script itself is not a proof
and its finite resolution does not test the mixed polygonal corner finite
part isolated in the JC node.
\end{remark}

\begin{proposition}[Archived conditional JP--FS algebra]
\label{JP:prop:FS}
Assume JP--SC with $c=2-\delta$, $0<\delta\le1/2$, and JP--JE.  After
decreasing the common fan and chart thresholds,
\begin{equation}\label{JP:eq:FS}
 \boxed{
 \Delta_\alpha
 -\frac1{2c}\|\ell_\alpha\|_{(\Gjp_\alpha)^*}^2
 \ge\kappa\Jfour(\alpha).}
\end{equation}
Equivalently, with $\mathcal Q_\alpha=c\Gjp_\alpha$, the dual term is
\(
 \frac12\sup_{z\ne0}\frac{(\sum_i r_{\alpha,i}z_i)^2}
 {\mathcal Q_\alpha(z)}
\).
\end{proposition}

\begin{proof}
Write $\|\cdot\|_* =\|\cdot\|_{(\Gjp_\alpha)^*}$ and
$e_\alpha=\ell_\alpha-p_\alpha$, and enlarge the two moduli to the common
one
\[
 \varepsilon(S):=\max\{\varepsilon_V(S),
                         \varepsilon_{\rm JE}(S)\}\longrightarrow0.
\]
Hilbert completion and
$p_{\alpha_*}=0$ give
\[
 P_D^{\rm ns}(\alpha)-P_D^{\rm ns}(\alpha_*)
 =q_D(B_\alpha)-q_D(B_{\alpha_*})
  -\frac14\|p_\alpha\|_*^2.
\]
For $c=2-\delta$ and $\delta\le1/2$,
\begin{align*}
 \frac1{2c}\|\ell_\alpha\|_*^2-\frac14\|p_\alpha\|_*^2
 &\le\left\{
 \frac{\delta C_p^2}{6}
 +\frac{2C_p\varepsilon(S)
        +\varepsilon(S)^2}{3}\right\}\Jfour(\alpha).
\end{align*}
Indeed,
$1/(2c)-1/4\le\delta/6$, $1/(2c)\le1/3$, and then use
\eqref{JP:eq:JE-P}--\eqref{JP:eq:JE-FM} and Cauchy--Schwarz.  Combining this
with \eqref{JP:eq:D0ns} and \eqref{JP:eq:JE-V} yields
\begin{align*}
 \Delta_\alpha-\frac1{2c}\|\ell_\alpha\|_*^2
 \ge\Bigl[2\kappa-\varepsilon(S)
 -\frac{\delta C_p^2}{6}
 -\frac{2C_p\varepsilon(S)
        +\varepsilon(S)^2}{3}\Bigr]\Jfour(\alpha).
\end{align*}
Choose $\delta$ and then $S_0,\rho_0$ so that the total subtracted bracket
is at most $\kappa$.  This proves \eqref{JP:eq:FS}.  Finally
\eqref{JP:eq:exact-flux-forcing} gives the displayed supremum identity.
\end{proof}

\begin{proposition}[Archived exact reduction of JP--BR]
\label{JP:prop:BR-reduction}
JP--BU and JP--FS give the former JP--BR interface, namely
\eqref{JP:eq:fan-upper} and
\begin{equation}\label{JP:eq:BR}
 \Phi_\alpha(h_\alpha\mathbf1)
 -\Phi_{\alpha_*}(h_{\alpha_*}\mathbf1)
 -\frac1{2c}\|\ell_\alpha\|_{(\Gjp_\alpha)^*}^2
 \ge c_0\Jfour(\alpha).
\end{equation}
\end{proposition}

\begin{proof}
Equation \eqref{JP:eq:exact-flux-forcing} and
$\mathcal Q_\alpha=c\Gjp_\alpha$ give the exact identity
\[
 \frac1{2c}\|\ell_\alpha\|_{(\Gjp_\alpha)^*}^2
 =\frac12\sup_{z\ne0}
 \frac{(\sum_i r_{\alpha,i}z_i)^2}{\mathcal Q_\alpha(z)}.
\]
Now apply JP--BU and JP--FS.
\end{proof}

\begin{remark}[Why the closed DM row did not save JE--FM]
In disk coordinates the first nonzero part of the exact covector is
\[
 2q_D\bigl(f_\alpha,
       v_\alpha-w_\alpha f_\alpha'\bigr).
\]
DM controls the null-form contribution in the $v_\alpha$ slot, but not the
curvature placement $-2q_D(f_\alpha,w_\alpha f_\alpha')$.  The certified
three-periodic CT family makes this curvature placement comparable to
$a$, while
$\Jfour(\alpha)^{1/2}\Gjp_\alpha(z)^{1/2}$ is also comparable to $a$.
Thus the archived material model does not behave as an $o(1)$ perturbation
of the disk Schur margin.  The old stationary O--H determinant cannot be
inserted at the nonstationary tangential polygon either.  After the
corner-transfer retraction in
Remark~\ref{JP:rem:period-three-material-forcing}, this is only a diagnostic,
not a proof that the genuine separate covector estimate is false.  The
fixed-loss scalar JP--JC, rather than an unproved citation to TR or DM, is
the active replacement.
\end{remark}

\begin{theorem}[Archived separated completion]
\label{JP:thm:completion}
Assume JP--SC and JP--JE.  Let $\mathcal R_\alpha$ be the Riesz map of
$\Gjp_\alpha$ and set
\begin{equation}\label{JP:eq:corrector}
 z_\alpha^\sharp=-c^{-1}\mathcal R_\alpha^{-1}\ell_\alpha.
\end{equation}
Then, for every normalized
$P_\alpha(h_0+z)\in\mathcal U_{S_0,\rho_0}$,
\begin{align}
 \Gjp_\alpha(z_\alpha^\sharp)&\le C\Jfour(\alpha),
 \label{JP:eq:corrector-bound}\\
 \Phi_\alpha(h_\alpha\mathbf1+z)
 -\Phi_{\alpha_*}(h_{\alpha_*}\mathbf1)
 &\ge c_0\Jfour(\alpha)
 +\frac c2\Gjp_\alpha(z-z_\alpha^\sharp).
 \label{JP:eq:joint-gap}
\end{align}
Thus JP--SC plus JP--JE implies the active-chart coercive node directly;
there is no support-direction CT--M4 or post-cubic remainder.
\end{theorem}

\begin{proof}
Proposition~\ref{JP:prop:FS} supplies JP--FS from the two hypotheses; together
with the closed JP--BU theorem, Proposition~\ref{JP:prop:BR-reduction} supplies
the baseline Schur margin.
By JP--SC and \eqref{JP:eq:exact-integration},
\[
 \Phi_\alpha(h_\alpha\mathbf1+z)-\Phi_\alpha(h_\alpha\mathbf1)
 \ge \ell_\alpha[z]+\frac c2\Gjp_\alpha(z).
\]
Hilbert completion gives
\[
 \ell_\alpha[z]+\frac c2\Gjp_\alpha(z)
 =\frac c2\Gjp_\alpha(z-z_\alpha^\sharp)
 -\frac1{2c}\|\ell_\alpha\|_{(\Gjp_\alpha)^*}^2.
\]
Proposition~\ref{JP:prop:BR-reduction} supplies \eqref{JP:eq:BR}; add it to
obtain \eqref{JP:eq:joint-gap}.  Equations
\eqref{JP:eq:fan-upper} and \eqref{JP:eq:BR} imply
$\|\ell_\alpha\|_{(\Gjp_\alpha)^*}^2\le C\Jfour(\alpha)$; use
\eqref{JP:eq:corrector} to obtain \eqref{JP:eq:corrector-bound}.
\end{proof}

\subsection{Audit disposition}

\begin{center}
\begin{longtable}{>{\raggedright\arraybackslash}p{.34\linewidth}
                  >{\raggedright\arraybackslash}p{.54\linewidth}}
\toprule
item & disposition\\
\midrule
fixed-normal vertex/support path & exact affine identity
\eqref{JP:eq:vertex-linear};\\
genuine side-normal velocity & constant $z_i$, equation
\eqref{JP:eq:constant-normal};\\
JP--EH & closed by Theorem~\ref{JP:thm:EH};\\
JP--NS & closed by Proposition~\ref{JP:prop:normal-speed-chart-entry}; the
downstream metric is fixed by \eqref{JP:eq:JP-metric-fixed};\\
raw $q_D$ norm of the piecewise-constant true normal & rejected as
ill-typed because vertex jumps have logarithmically divergent
$H^{1/2}$ energy;\\
disk cubic candidate & rigorously noncoercive on the alternating family of
Proposition~\ref{JP:prop:cubic-counterexample};\\
JP--SC & archived stronger route; estimates \eqref{JP:eq:S0} and
\eqref{JP:eq:ST} are not active premises;\\
JP--BU & closed by Theorem~\ref{JP:thm:BU};\\
exact support forcing & centered side-flux covector
\eqref{JP:eq:exact-flux-forcing}--\eqref{JP:eq:flux-moments};\\
$D0_{\rm ns}$ and JE--P & closed by
Theorem~\ref{JP:thm:normal-disk-endpoint};\\
JE--V & closed by Theorem~\ref{JP:thm:normal-value-matching};\\
JE--FM & unavailable; its former period-three material falsification is
retracted by Remark~\ref{JP:rem:period-three-material-forcing}, and the row is
absent from the active DAG;\\
vanishing joint comparison & former material counterexample retracted by
Remark~\ref{JP:rem:period-three-material-joint-loss}; the active route uses a
fixed absorbable loss in JP--JC and does not assume this statement;\\
JP--FS and JP--BR & archived conditional algebra; unnecessary on the
JP--JC route;\\
support post-cubic and support-linear fourth tails & removed from the DAG
by the exact value identity \eqref{JP:eq:exact-integration}.\\
\bottomrule
\end{longtable}
\end{center}
\endgroup

\begingroup
\section{Node JC: honest joint endpoint comparison}
\label{module:JC}
\noindent\textit{Audited source: \nolinkurl{JC_joint_comparison.tex}.}
\subsection{Two honest scalars on one fixed fan}

Fix a positive normal fan $\alpha$, let $h_\alpha\mathbf1$ be its
area-$\pi$ tangential support vector, and work on the fixed affine
scale/translation section $\mathsf X_\alpha$ of the JP node.  Put
\begin{equation}\label{JC:eq:polygon-scalar}
 \Phi_\alpha(h)=\frac{|P_\alpha(h)|}{\pi}
                    \lambda_1(P_\alpha(h)).
\end{equation}
The closed JP--NS and disk-endpoint rows define
\begin{equation}\label{JC:eq:disk-data}
 B_\alpha=\Pi_{\rm ng}
 \{h_\alpha\sec(\theta-\theta_i)-1\},\qquad
 K_\alpha z=\Pi_{\rm ng}\mathcal I_\alpha z,
\end{equation}
and
\begin{equation}\label{JC:eq:metric-forcing}
 \Gjp_\alpha(z)=q_D(K_\alpha z),\qquad
 p_\alpha[z]=2\langle B_\alpha,K_\alpha z\rangle_D.
\end{equation}
Let $\mathcal R_\alpha$ be the Riesz map of $\Gjp_\alpha$ and set
\begin{equation}\label{JC:eq:disk-corrector}
 z_\alpha^D=-\frac12\mathcal R_\alpha^{-1}p_\alpha.
\end{equation}

\begin{proposition}[Exact disk completion]\label{JC:prop:disk-completion}
On $\mathsf X_\alpha$ one has
\begin{align}
 q_D(B_\alpha+K_\alpha z)
 &=P_D^{\rm ns}(\alpha)
   +\Gjp_\alpha(z-z_\alpha^D),\label{JC:eq:disk-completion}\\
 P_D^{\rm ns}(\alpha)
 &=q_D(B_\alpha)-\frac14
       \|p_\alpha\|_{(\Gjp_\alpha)^*}^2,\label{JC:eq:disk-schur}\\
 \Gjp_\alpha(z_\alpha^D)
 &=\frac14\|p_\alpha\|_{(\Gjp_\alpha)^*}^2.
 \label{JC:eq:corrector-size}
\end{align}
At the regular fan $\alpha_*$, cyclic symmetry gives
$p_{\alpha_*}=0$ and $z_{\alpha_*}^D=0$.
\end{proposition}

\begin{proof}
Expand the fixed quadratic form:
\[
 q_D(B_\alpha+K_\alpha z)
 =q_D(B_\alpha)+p_\alpha[z]+\Gjp_\alpha(z).
\]
The definition \eqref{JC:eq:disk-corrector} is exactly
$2\mathcal R_\alpha z_\alpha^D=-p_\alpha$.  Completing the Hilbert square
proves all three identities.  The regular assertion follows because the
regular source is cyclic while every vector in the quotient is orthogonal
to the cyclic mode.
\end{proof}

Define the fixed-fan joint residual
\begin{equation}\label{JC:eq:joint-residual}
\boxed{
 \mathscr R_\alpha(z):=
 \Phi_\alpha(h_\alpha\mathbf1+z)
 -\Phi_\alpha(h_\alpha\mathbf1)
 -\{q_D(B_\alpha+K_\alpha z)-q_D(B_\alpha)\}.}
\end{equation}
Both terms in \eqref{JC:eq:joint-residual} are genuine scalar functions of
$z$.  The disk term is the fixed Fourier--Bessel quadratic form; it is not
the eigenvalue of a moving disk domain.

\begin{proposition}[Exact integrated form]\label{JC:prop:integrated-form}
Let $\ell_\alpha=D_h\Phi_\alpha(h_\alpha\mathbf1)$ and
$e_\alpha=\ell_\alpha-p_\alpha$.  Whenever the fixed-normal support
segment is active,
\begin{equation}\label{JC:eq:integrated-form}
 \mathscr R_\alpha(z)=e_\alpha[z]
 +\int_0^1(1-t)
 \{D_h^2\Phi_\alpha(h_\alpha\mathbf1+tz)[z,z]
       -2\Gjp_\alpha(z)\}\,dt.
\end{equation}
\end{proposition}

\begin{proof}
Apply the fundamental theorem of calculus twice to the polygonal scalar.
For the disk scalar use
$q_D(B+Kz)-q_D(B)=p[z]+\Gjp(z)$.  Subtraction gives
\eqref{JC:eq:integrated-form} exactly.
\end{proof}

\subsection{Exact Schur reduction and support-time cancellation}

Let $g_\alpha(\cdot,\cdot)$ be the bilinear form associated with
$\Gjp_\alpha$, put $d_\alpha=z_\alpha^D$, and abbreviate
\begin{equation}\label{JC:eq:Hessian-family}
 H_{\alpha,y}=D_h^2\Phi_\alpha(h_\alpha\mathbf1+y).
\end{equation}
Equation~\eqref{JC:eq:disk-corrector} is equivalently
\begin{equation}\label{JC:eq:p-d-relation}
 p_\alpha[\xi]=-2g_\alpha(d_\alpha,\xi).
\end{equation}
Define the completed covector
\begin{equation}\label{JC:eq:completed-covector}
 b_\alpha[\xi]
 :=\ell_\alpha[\xi]+g_\alpha(d_\alpha,\xi)
 =\ell_\alpha[\xi]-\frac12p_\alpha[\xi]
\end{equation}
and the pathwise Schur action
\begin{equation}\label{JC:eq:pathwise-action}
 \mathscr A_\alpha(z):=
 b_\alpha[z]
 +\int_0^1(1-t)
 \{H_{\alpha,tz}[z,z]-g_\alpha(z,z)\}\,dt
 +\frac12g_\alpha(d_\alpha,d_\alpha)
 +\frac\kappa2\Jfour(\alpha).
\end{equation}

\begin{proposition}[JC--ALG: exact completed-action equivalence]
\label{JC:prop:JC-ALG}
For every active fixed-fan segment,
\begin{equation}\label{JC:eq:JC-equivalence}
 \boxed{
 \mathscr A_\alpha(z)
 =\mathscr R_\alpha(z)
  +\frac12\Gjp_\alpha(z-z_\alpha^D)
  +\frac\kappa2\Jfour(\alpha).}
\end{equation}
Consequently the former JP--JC inequality is equivalent, without a loss
of constants, to $\mathscr A_\alpha(z)\ge0$.
\end{proposition}

\begin{proof}
Insert the integrated identity~\eqref{JC:eq:integrated-form} and expand
\[
 \frac12g_\alpha(z-d_\alpha,z-d_\alpha)
 =\frac12g_\alpha(z,z)-g_\alpha(z,d_\alpha)
   +\frac12g_\alpha(d_\alpha,d_\alpha).
\]
By~\eqref{JC:eq:p-d-relation},
$-g_\alpha(z,d_\alpha)=p_\alpha[z]/2$.  Hence the linear row is
$(\ell_\alpha-p_\alpha)[z]+p_\alpha[z]/2=b_\alpha[z]$.
Moreover
\[
 \int_0^1(1-t)(-2g_\alpha(z,z))\,dt
 +\frac12g_\alpha(z,z)=-\frac12g_\alpha(z,z),
\]
which is exactly the integral of $-g_\alpha(z,z)$ with weight $1-t$.
The constant rows agree with~\eqref{JC:eq:pathwise-action}, proving
\eqref{JC:eq:JC-equivalence}.
\end{proof}

The same identity has the following equivalent form, which displays the
positive disk corrector square before any estimate:
\begin{equation}\label{JC:eq:joint-Peano-block}
 \boxed{
 \mathscr A_\alpha(z)
 =\frac12\Gjp_\alpha(z-d_\alpha)
  +e_\alpha[z]
  +\int_0^1(1-t)
    \{H_{\alpha,tz}[z,z]-2\Gjp_\alpha(z)\}\,dt
  +\frac\kappa2\Jfour(\alpha).}
\end{equation}
Indeed, $b_\alpha=e_\alpha-g_\alpha(d_\alpha,\cdot)$ and
$\int_0^1(1-t)g_\alpha(z,z)\,dt=\Gjp_\alpha(z)/2$.
Thus the period-three component of $e_\alpha$ is not to be bounded as an
isolated covector: it must be paid jointly by the jump surplus in
$H_{\alpha,tz}-2\Gjp_\alpha$ and the fan row.

\begin{lemma}[JC--TL: support-time logarithm cancellation]
\label{JC:lem:support-time-log}
Let $\beta>0$, $0<r\le1$, and
\begin{equation}\label{JC:eq:I-beta}
 I_\beta(r):=\int_0^1(1-t)
 \left\{1-\bigl((1-t)+tr\bigr)^{2\beta}\right\}\,dt.
\end{equation}
Then
\begin{align}
 0\le I_\beta(r)
 &\le\frac\beta3\log\frac1r,\label{JC:eq:log-bound}\\
 I_\beta(r)
 &\le\frac{\beta}{2(1+\beta)}\le\frac\beta2.\label{JC:eq:absolute-bound}
\end{align}
In particular,
\begin{equation}\label{JC:eq:min-bound}
 \boxed{I_\beta(r)\le\frac\beta3
 \min\left\{\log\frac1r,\frac32\right\}.}
\end{equation}
All constants are independent of $r$.
\end{lemma}

\begin{proof}
For $0<x\le1$, $1-x^{2\beta}\le2\beta(-\log x)$.
Weighted arithmetic--geometric mean gives
$(1-t)+tr\ge r^t$ and hence
\[
 -\log((1-t)+tr)\le t\log(1/r).
\]
Since $\int_0^1t(1-t)\,dt=1/6$, these two inequalities prove
\eqref{JC:eq:log-bound}.  Also $(1-t)+tr\ge1-t$, so monotonicity of the
power gives
\[
 I_\beta(r)\le I_\beta(0)
 =\frac12-\frac1{2+2\beta}
 =\frac\beta{2(1+\beta)}.
\]
This proves~\eqref{JC:eq:absolute-bound}; taking the smaller bound gives
\eqref{JC:eq:min-bound}.
\end{proof}

\begin{lemma}[JC--TL$^{\min}$: switches of the shortest incident side]
\label{JC:lem:support-time-min}
Let $L_-$ and $L_+$ be positive affine functions on $[0,1]$, put
$m(t):=\min\{L_-(t),L_+(t)\}$, and set $r:=m(1)/m(0)$.  Then, for every
$\beta>0$,
\begin{equation}\label{JC:eq:min-affine-payment}
 \int_0^1(1-t)\left\{1-
 \left(\frac{m(t)}{m(0)}\right)^{2\beta}\right\}\,dt
 \le
 \begin{cases}
 I_\beta(r),&0<r\le1,\\
 0,&r\ge1.
 \end{cases}
\end{equation}
Consequently every \emph{upper} bound in
Lemma~\ref{JC:lem:support-time-log} remains valid for the positive part of the
loss when a single affine length is replaced by the shortest of the two
incident physical lengths.  No switch-point boundary term is present.
\end{lemma}

\begin{proof}
The pointwise minimum of affine functions is concave.  If $r\le1$,
concavity gives
\[
 \frac{m(t)}{m(0)}\ge(1-t)+tr.
\]
The function $x\mapsto1-x^{2\beta}$ is decreasing, so integration with
the nonnegative weight $1-t$ proves the first line of
\eqref{JC:eq:min-affine-payment}.  If $r\ge1$, the chord joining the endpoint
values lies above $m(0)$, and concavity gives $m(t)\ge m(0)$ throughout;
the integrand is therefore nonpositive.  The argument differentiates
neither $m$ nor $m^{2\beta}$, hence a change of the minimizing affine side
creates no distributional boundary contribution.
\end{proof}

\begin{remark}[Exact terminal normalization]
For an actual physical side length with the affine form
$L(t)=L_0\{(1-t)+tr\}$, the squared corner-flux power uses
\begin{equation}\label{JC:eq:physical-corner-exponent}
 \beta_i^{\rm ph}=\frac{\alpha_i}{\pi-\alpha_i},
\end{equation}
not $\alpha_i/\pi$.  Put
$\Delta_i=a_i^--a_i^+=2\sin(\alpha_i/2)\tau_i$.  The integrated terminal
amplitude loss obeys
\begin{equation}\label{JC:eq:terminal-log-payment}
 \alpha_i\tau_i^2 I_{\beta_i^{\rm ph}}(r)
 \le |\Delta_i|^2\min\left\{
 \frac{\alpha_i^2}{12(\pi-\alpha_i)\sin^2(\alpha_i/2)}
       \log\frac1r,
 \frac{\alpha_i^2}{8\pi\sin^2(\alpha_i/2)}\right\}.
\end{equation}
Thus the logarithmic coefficient is $1/(3\pi)+O(\alpha_i)$, while the
exact absolute cap is
\begin{equation}\label{JC:eq:terminal-absolute-payment}
 \alpha_i\tau_i^2 I_{\beta_i^{\rm ph}}(r)
 \le c_T(S)|\Delta_i|^2,
 \qquad
 c_T(S):=\sup_{0<\alpha\le S}
 \frac{\alpha^2}{8\pi\sin^2(\alpha/2)},
\end{equation}
and $c_T(S)\to1/(2\pi)$ as $S\downarrow0$.  This is the sharp
support-time cap for a terminal length tending to zero.  A growing
terminal length has the favorable sign and is separated before the
absolute estimate.
The exponent $\alpha_i/\pi$ belongs instead to a conformal-prevertex
power.  No retained node proves that a conformal prevertex interval is
affine in support time, so that substitution is not allowed.  If the
terminal cutoff is the smaller of two incident physical lengths,
Lemma~\ref{JC:lem:support-time-min} applies globally: concavity of the
minimum dominates the affine endpoint chord, so neither subdivision nor a
switch correction is needed.  (Restarting the Taylor weight after a
subdivision would in fact change the ledger and is not the argument used
here.)
The right side is the descendant-annulus/terminal-cap jump charge on the
same physical scale.  Thus taking a side-ratio supremum before the
support-time integral destroys a real cancellation and is inadmissible.
\end{remark}

Formula \eqref{JC:eq:integrated-form} is the reason the two old obligations
must not be estimated independently.  In the three-periodic obstruction
$e_\alpha$ has a non-vanishing normalized leading term; the surplus in the
integrated Hessian is available in \eqref{JC:eq:integrated-form} to complete
the corresponding Schur square.

\begin{proposition}[JC--RS: exact quadratic root certificate]
\label{JC:prop:JC-RS}
Put
\begin{equation}\label{JC:eq:root-form}
 Q_\alpha[\xi,\xi]
 :=H_{\alpha,0}[\xi,\xi]-g_\alpha(\xi,\xi).
\end{equation}
Define the frozen quadratic root model
\begin{equation}\label{JC:eq:quadratic-root-model}
 \mathscr A_\alpha^{(2)}(\xi)
 :=\frac12g_\alpha(d_\alpha,d_\alpha)
   +\frac\kappa2\Jfour(\alpha)
   +b_\alpha[\xi]+\frac12Q_\alpha[\xi,\xi].
\end{equation}
Then $\mathscr A_\alpha^{(2)}(\xi)\ge0$ for every $\xi$ if and only if
\begin{equation}\label{JC:eq:root-Schur-necessary}
 \boxed{
 Q_\alpha\ge0,
 \qquad
 \|b_\alpha\|_{Q_\alpha^*}^2
 \le g_\alpha(d_\alpha,d_\alpha)
       +\kappa\Jfour(\alpha).}
\end{equation}
For a semidefinite $Q_\alpha$, the second statement includes that
$b_\alpha$ annihilates its kernel and uses the quotient dual norm.
\end{proposition}

\begin{proof}
This is Hilbert completion of the quadratic form
\eqref{JC:eq:quadratic-root-model}.  If it is nonnegative on the whole vector
space, its quadratic part is semidefinite and its linear part annihilates
the kernel.  Its infimum is
\[
 \frac12g_\alpha(d_\alpha,d_\alpha)
 +\frac\kappa2\Jfour(\alpha)
 -\frac12\|b_\alpha\|_{Q_\alpha^*}^2,
\]
which is nonnegative exactly under~\eqref{JC:eq:root-Schur-necessary}.  The
reverse implication follows from the same completed square.
\end{proof}

Proposition~\ref{JC:prop:JC-RS} is an exact certificate for the frozen
quadratic root model, not a necessary consequence of local nonnegativity
of the full action: its positive constant term prevents that inference.
To use it for JC--PL one must additionally compare the full integrated
Hessian to this frozen model on the whole admissible amplitude range.  The
closed JP--DE estimates control $p_\alpha$ and hence $d_\alpha$, but do not
prove either that comparison or the displayed sharp dual bound for
$b_\alpha=\ell_\alpha-p_\alpha/2$.  The period-three computation passes
the quadratic certificate with a sizeable numerical margin; a uniform
no-ratio proof over all profiles is still required.

\begin{proposition}[JC--RC: exact root/path sufficient reduction]
\label{JC:prop:JC-RC}
Suppose a positive semidefinite form $L_\alpha$ satisfies, throughout the
radial/metric path chart,
\begin{equation}\label{JC:eq:path-loss-form}
 \int_0^1(1-t)
 \{H_{\alpha,tz}-H_{\alpha,0}\}[z,z]\,dt
 \ge-L_\alpha[z,z].
\end{equation}
Define
\begin{equation}\label{JC:eq:root-after-path-loss}
 \widehat Q_\alpha:=H_{\alpha,0}-g_\alpha-2L_\alpha.
\end{equation}
If
\begin{equation}\label{JC:eq:root-path-Schur}
 \widehat Q_\alpha\ge0,
 \qquad
 \|b_\alpha\|_{\widehat Q_\alpha^*}^2
 \le g_\alpha(d_\alpha,d_\alpha)+\kappa\Jfour(\alpha),
\end{equation}
then JC--PL holds.
\end{proposition}

\begin{proof}
Split the Hessian integral in~\eqref{JC:eq:pathwise-action} at
$H_{\alpha,0}$.  The Bregman weight has mass $1/2$, so
\eqref{JC:eq:path-loss-form} gives
\[
 \mathscr A_\alpha(z)
 \ge\frac12\{g_\alpha(d_\alpha,d_\alpha)
                 +\kappa\Jfour(\alpha)\}
     +b_\alpha[z]
     +\frac12\widehat Q_\alpha[z,z].
\]
The right side is nonnegative by the same Hilbert completion as
Proposition~\ref{JC:prop:JC-RS}.
\end{proof}

This reduction records the required slack correctly.  Terminal and
nonterminal estimates must first be assembled into a genuine form
$L_\alpha$; one may not subtract their scalar upper bounds twice or
replace~\eqref{JC:eq:path-loss-form} by a pointwise support-time supremum.

\begin{remark}[Why the tempting exact Gram shortcut does not close the root]
Let $\mathcal H_D$ be the disk boundary Hilbert space, let $K$ be the
fan-normal sine trace, $T$ the genuine material/kinematic trace at the
tangential endpoint, and $B$ the fan source.  Put
\[
 G=K^*K,\qquad p=2K^*B,\qquad
 d=-G^{-1}p/2,\qquad c=-Kd=P_{K\mathsf X_\alpha}B,
 \qquad E=T-K.
\]
For the pure quadratic disk principal block, where
$\ell=2T^*B$ and $H_0=2T^*T$, direct expansion gives, with $V=2T-K$,
\begin{align}
 b[z]&=\langle c,Vz\rangle_D
       +2\langle B-c,Ez\rangle_D,
 \label{JC:eq:Gram-linear-defect}\\
 H_0-G&=V^*V-2E^*E,
 \label{JC:eq:Gram-quadratic-defect}\\
 \frac12G(d)+b[z]+\frac12(H_0-G)[z,z]
 &=\frac12\|c+Kz+2Ez\|_D^2-
   \|Ez\|_D^2+2\langle B-c,Ez\rangle_D.
 \label{JC:eq:Gram-failed-square}
\end{align}
Thus choosing $U=c$ and $V=2T-K$ leaves the negative form $-2E^*E$
and an orthogonal-source cross term; they do not vanish merely because
$c$ is the projection of $B$ onto $K\mathsf X_\alpha$.

For the genuine polygonal block write
$\Gamma=H_{\alpha,0}-2T^*T$ and let $r$ be the difference between
$b_\alpha$ and the first pairing in~\eqref{JC:eq:Gram-linear-defect}.  A
valid Gram repair must, at minimum, pay the block
\[
 \begin{pmatrix}
  \kappa\Jfour(\alpha) & r^*\\ r & \Gamma-2E^*E
 \end{pmatrix}
 \ge0
\]
after adding any explicitly retained fan square and path-loss slack.
No closed TR or JP--DE row proves this block.  Equations
\eqref{JC:eq:Gram-linear-defect}--\eqref{JC:eq:Gram-failed-square} explain why
the numerically favorable root ratio cannot be promoted by a bare
Cauchy--Schwarz citation.
\end{remark}

\begin{remark}[Audited physical-trace branch: exact algebra, failed chart]
There is an algebraically cleaner inactive branch.  Let
$\mathbf q_D=2j_{0,1}^2q_D^0$ be the physical disk half-Hessian, put
\[
 \widetilde K_\alpha z=\Pi_{\rm ng}\widetilde T_\alpha z,
 \qquad
 \widetilde G_\alpha(z)=\mathbf q_D(\widetilde K_\alpha z),
 \qquad
 \widetilde p_\alpha[z]
 =2\langle B_\alpha,\widetilde K_\alpha z\rangle_D,
\]
and let
$\widetilde d_\alpha=-\tfrac12\widetilde{\mathcal R}_\alpha^{-1}
\widetilde p_\alpha$.  The normalized TR endpoint gives the exact square
\[
 \mathbf q_D(B_\alpha+\widetilde K_\alpha z)
 =\widetilde P_D(\alpha)
  +\widetilde G_\alpha(z-\widetilde d_\alpha).
\]
The corresponding physical residual and action satisfy the same exact
Schur identity as Proposition~\ref{JC:prop:JC-ALG}; in its pure disk block
the root is exactly
\[
 \frac12\widetilde G_\alpha(z-\widetilde d_\alpha)
 +\frac\kappa2\Jfour(\alpha)\ge0.
\]
Thus the $E=T-K$ defect in~\eqref{JC:eq:Gram-quadratic-defect} disappears
from the disk block.

This does not close the polygonal action.  The large component of
$\widetilde T_\alpha$ produced by a short angle is a tangential slide of a
straight side viewed in the disk radial frame; it can be nonzero even
when the true side-normal speed is zero.  It is cancelled at the same
order by the nonlinear reparametrization remainder $\mathcal N$:
\[
 D_z^2\widetilde{\mathscr R}_\alpha(0)[\xi,\xi]
 =H_{\alpha,0}[\xi,\xi]-2\widetilde G_\alpha(\xi).
\]
By the exact value split in TR, the difference on the right is precisely
the full second $z$-derivative of $\mathcal N(W_{\alpha,z})$ at $z=0$;
this formulation includes all chain-rule terms at the vector-field base
point $W_{\alpha,0}$.
The cancellation cannot be estimated from the radial image alone.  For
example, if
$X(\theta)=e_r(\theta+\delta\sin2\theta)$ and $W=X-e_r$, the boundary
image is the unit circle, so $\mathcal L(W)=\mathcal L(0)$, whereas
\[
 f=W\cdot e_r=\cos(\delta\sin2\theta)-1,
 \qquad
 \mathcal N(W)=-\mathbf q_D(\Pi_{\rm ng}f).
\]
Hence no image-small argument can yield
$\mathcal N(W)\ge-\vartheta\mathbf q_D(\Pi_{\rm ng}f)$ with a universal
$\vartheta<1$.  The radial-not-affine example in JP further shows that
the full material $W^{1,\infty}$ norm can diverge on the radial chart.
Consequently the old physical CF+CM$\to$NR support-tail argument cannot
be reactivated.  The physical branch replaces JC--PL by an equivalent
open full-radial-chart action; it is recorded as an inactive alternative,
not as a second proof of J0.
\end{remark}

\begin{remark}[Audited Minkowski shortcut: two concavities do not order the ratio]
For a fixed positive normal fan, support interpolation is Minkowski
interpolation.  Consequently
\[
 U(h)=|P(h)|^{1/2},
 \qquad V(h)=\lambda_1(P(h))^{-1/2}
\]
are positive, one-homogeneous concave functions, while
$\Phi=(U/V)^2$.  At a stationary point,
\[
 \frac{D^2\Phi}{\Phi}
 =2\left(\frac{D^2U}{U}-\frac{D^2V}{V}\right)
\]
on a nonradial direction.  Both displayed curvatures are nonpositive,
but their difference has no determined sign.

This is a genuine logical obstruction, not only a missing inequality.
On the Lorentz cone $\{(x,y):x>|y|\}$ put, with $c=11/20$,
\[
 U(x,y)=\sqrt{x^2-y^2},
 \qquad V(x,y)=x-c\frac{y^2}{x}.
\]
Both functions are positive, one-homogeneous and concave, since
\[
 D^2U=-\frac1{(x^2-y^2)^{3/2}}
 \begin{pmatrix}y^2&-xy\\-xy&x^2\end{pmatrix},
 \quad
 D^2V=-\frac{2c}{x^3}
 \begin{pmatrix}y^2&-xy\\-xy&x^2\end{pmatrix}.
\]
Moreover $U^2=x^2-y^2$ is exactly quadratic.  At $(1,0)$ the ratio
$(U/V)^2$ is stationary and has second derivative $1/5>0$ in the
$y$ direction.  Nevertheless the same-area point $(5/3,4/3)$ satisfies
\[
 U(5/3,4/3)=1,
 \qquad
 \left(\frac UV\right)^2(5/3,4/3)=\frac{625}{729}<1.
\]
Thus area and eigenvalue Brunn--Minkowski concavity, even with exact
quadratic area and strict local support minimality, do not give a global
support comparison.  The missing relative curvature comparison is again
part of JC--PL.
\end{remark}

\begin{remark}[Audited stationary-endpoint bypass: why it is insufficient]
The independent stationary O--H endpoint theorem cannot by itself replace
JC--PL.  Indeed, put the regular polygon and a hypothetical active global
minimizer in one fixed-area joint chart at $0$ and $x$.  Since both are
stationary, the exact identity is
\begin{equation}\label{JC:eq:stationary-gradient-chord}
 0=\langle\nabla\Phi_N(x)-\nabla\Phi_N(0),x\rangle
 =\int_0^1D^2\Phi_N(tx)[x,x]\,dt.
\end{equation}
The interior points $tx$ are not stationary, so a Hessian theorem whose
hypothesis is stationarity does not control this integral.

Endpoint coercivity alone is logically insufficient even at the relevant
super-near scale.  For $a>0$ and $0<c<1/3$ consider
\[
 F_{a,c}(t,y)=a t^2(a-t)^2
 -ca^5\{3(t/a)^2-2(t/a)^3\}+|y|^2,
 \qquad E_a(u,v)=a^3u^2+|v|^2.
\]
The points $(0,0)$ and $(a,0)$ are stationary, while
$F_{a,c}(a,0)=-ca^5<F_{a,c}(0,0)=0$.  Both endpoint Hessians are strictly
positive: their $t$-eigenvalues are respectively
\[
 (2-6c)a^3,\qquad (2+6c)a^3,
\]
and the $y$-block is $2I$.  Moreover, with $u=t/a$,
\[
 \frac{F_{a,c}(t,0)}{a^5}+c
 =(1-u)^2\{(u+c)^2+c-c^2\}\ge0,
\]
so $(a,0)$ is the global minimum.  Nevertheless
$D_{tt}^2F_{a,c}(a/2,0)=-a^3$ and, for $x=(a,0)$,
\[
 \int_0^1D^2F_{a,c}(sx)[x,x]\,ds=0,
 \qquad E_a(x)=a^5.
\]
Thus positive endpoint spectra, including uniform spectra in the naturally
scaled norm, do not prove the path monotonicity needed in
\eqref{JC:eq:stationary-gradient-chord}.

The minimal stationary-route replacement is a super-near joint statement,
not an endpoint citation.  After gauges, it would have to prove uniformly
\begin{equation}\label{JC:eq:stationary-super-near-cut}
 \int_0^1D^2\Phi_N(tx)[x,x]\,dt
 \ge\mu\left\{6a^2\sum_i d_i^2
 +\Gjp_{\alpha_*}(z-Z_Nd)\right\}
\end{equation}
for nonzero $x=(d,z)$ with
$\Jfour(\alpha_*+d)+\Gjp_{\alpha_*}(z)\lesssim a^5$.
At the regular polygon this requires the full cyclic fan/support Schur
symbol, including its nonzero period-three mixed block; it must then be
propagated diagonally along the super-near path.  No retained node proves
either the regular Schur complement or this relative path stability.
Consequently the stationary route only narrows JC--PL; it does not close
the main theorem.
\end{remark}

\subsection{The unique open pathwise payment}

The closed disk endpoint and value rows provide universal
$S_D,\kappa,C_p>0$ and a modulus $\varepsilon_V(S)\to0$ such that, for
$S\le S_D$,
\begin{align}
 P_D^{\rm ns}(\alpha)-P_D^{\rm ns}(\alpha_*)
 &\ge2\kappa\Jfour(\alpha),\label{JC:eq:D0}\\
 \|p_\alpha\|_{(\Gjp_\alpha)^*}
 &\le C_p\Jfour(\alpha)^{1/2},\label{JC:eq:P}\\
 \left|\Phi_\alpha(h_\alpha\mathbf1)
 -\Phi_{\alpha_*}(h_*\mathbf1)
 -\{q_D(B_\alpha)-q_D(B_{\alpha_*})\}\right|
 &\le\varepsilon_V(S)\Jfour(\alpha).
 \label{JC:eq:V}
\end{align}

\begin{proposition}[Regular disk-joint Schur symbol]
\label{JC:prop:regular-disk-joint}
Let $\alpha_*=a\mathbf1$ with $2a<S_D$.  Choose the canonical analytic
scale/translation trivialization constructed in
Proposition~\ref{JC:prop:super-near-bridge-active},
\[
 J_\alpha:\mathsf X_{\alpha_*}\longrightarrow\mathsf X_\alpha,
 \qquad J_{\alpha_*}=I,
\]
and put
\[
 F_D^\sharp(\alpha,z):=
 q_D(B_\alpha+K_\alpha J_\alpha z),\qquad
 \bar z_\alpha^D:=J_\alpha^{-1}z_\alpha^D,\qquad
 Z_N\delta:=
 D_\alpha\bar z_\alpha^D\big|_{\alpha_*}[\delta].
\]
For every $\sum_i\delta_i=0$ and every regular-quotient support vector
$\zeta$,
\begin{align}
 \mathcal D_N(\delta,\zeta)
 &:=
 \frac12D^2F_D^\sharp(\alpha_*,0)
 [(\delta,\zeta),(\delta,\zeta)]
 \label{JC:eq:regular-disk-joint-identity}\\
 &=
 \frac12D^2P_D^{\rm ns}(\alpha_*)[\delta,\delta]
 +\Gjp_{\alpha_*}(\zeta-Z_N\delta)
 \notag\\
 &\ge
 12\kappa a^2\sum_i\delta_i^2
 +\Gjp_{\alpha_*}(\zeta-Z_N\delta).
 \label{JC:eq:regular-disk-joint-gap}
\end{align}
Thus every cyclic Fourier block of the regular
\emph{normal-sine disk-joint} symbol has a uniform completed Schur gap.
This statement concerns the fixed disk scalar; it is not a lower bound
for the polygonal joint Hessian.
\end{proposition}

\begin{proof}
The endpoint data and $J_\alpha$ are analytic in the fixed quotient
chart.  If
$\bar\Gjp_\alpha(w):=\Gjp_\alpha(J_\alpha w)$, exact disk completion gives
\[
 F_D^\sharp(\alpha,z)
 =P_D^{\rm ns}(\alpha)
  +\bar\Gjp_\alpha(z-\bar z_\alpha^D),
\qquad \bar z_{\alpha_*}^D=0.
\]
Along $\alpha(t)=\alpha_*+t\delta$, $z(t)=t\zeta$, one has
$\bar z_{\alpha(t)}^D=tZ_N\delta+O(t^2)$, while variation of
$\bar\Gjp_{\alpha(t)}$ contributes only from order $t^3$ onward.  Taking
the coefficient of $t^2$ proves
\eqref{JC:eq:regular-disk-joint-identity}.  Moreover
\[
 \Jfour(\alpha_*+t\delta)
 =6a^2t^2\sum_i\delta_i^2
  +4at^3\sum_i\delta_i^3+t^4\sum_i\delta_i^4.
\]
Apply the closed estimate
$P_D^{\rm ns}(\alpha)-P_D^{\rm ns}(\alpha_*)
\ge2\kappa\Jfour(\alpha)$.  For fixed $\delta$, both signs of sufficiently
small $t$ keep $\alpha_*+t\delta$ in the positive-fan region
$S\le S_D$.  Dividing the two signed inequalities first by $t$ shows
$DP_D^{\rm ns}(\alpha_*)[\delta]=0$; subtracting that linear term and
then dividing by $t^2$ gives
$\frac12D^2P_D^{\rm ns}(\alpha_*)[\delta,\delta]
\ge12\kappa a^2\sum_i\delta_i^2$, proving
\eqref{JC:eq:regular-disk-joint-gap}.
\end{proof}

\begin{contract}[JC--AR: regular completed-action Schur block]
\label{JC:con:JC-AR}
In the trivialization of Proposition~\ref{JC:prop:regular-disk-joint}, put
\[
 F_P^\sharp(\alpha,z)
 :=\Phi_\alpha(h_\alpha\mathbf1+J_\alpha z),
 \qquad
 \mathcal E_N(d,\zeta)
 :=\frac12D^2(F_P^\sharp-F_D^\sharp)(\alpha_*,0)
 [(d,\zeta),(d,\zeta)].
\]
Write \(\bar z_\alpha^D=J_\alpha^{-1}z_\alpha^D\), so that
\(D\bar z_{\alpha_*}^D[d]=Z_Nd\), and use the exact completed coordinates
\[
 \mathbf A_N(d,w):=
 \mathscr A_{\alpha_*+d}
 \bigl(J_{\alpha_*+d}(\bar z_{\alpha_*+d}^D+w)\bigr).
\]
Then \(\mathbf A_N(0,0)=0\), \(D\mathbf A_N(0,0)=0\), and the exact
regular quadratic root is
\begin{equation}\label{JC:eq:regular-action-root}
 \boxed{
 \mathcal Q_N(d,w):=\frac12D^2\mathbf A_N(0,0)[(d,w),(d,w)]
 =\mathcal E_N(d,w+Z_Nd)-\mathcal E_N(d,0)
 +\frac12\Gjp_{\alpha_*}(w)
 +3\kappa a^2\sum_i d_i^2.}
\end{equation}
The subtraction \(\mathcal E_N(d,0)\) and the coefficients
\(1/2\) and \(3\kappa\) are forced respectively by the base-value
subtraction in \(\mathscr R_\alpha\), the completed disk square, and
\(\frac\kappa2\Jfour\).  They may not be replaced by the larger margins
of the full disk-joint Hessian.

Let
\[
 \mathcal P_N(d,\zeta):=
 \frac12D^2F_P^\sharp(\alpha_*,0)
 [(d,\zeta),(d,\zeta)]
\]
be the genuine polygonal joint half-Hessian.  Exact disk completion gives
the following sharper algebraic reduction:
\begin{equation}\label{JC:eq:regular-action-polygon-reduction}
 \boxed{
 \mathcal Q_N(d,w)=
 \mathcal P_N(d,w+Z_Nd)-\mathcal P_N(d,0)
 +\Gjp_{\alpha_*}(Z_Nd)-\frac12\Gjp_{\alpha_*}(w)
 +3\kappa a^2\sum_i d_i^2.}
\end{equation}
Indeed Proposition~\ref{JC:prop:regular-disk-joint} yields
\[
 \frac12D^2F_D^\sharp(\alpha_*,0)[(d,\zeta)]^2
 =\frac12D^2P_D^{\rm ns}(\alpha_*)[d,d]
  +\Gjp_{\alpha_*}(\zeta-Z_Nd).
\]
Subtracting its values at $\zeta=w+Z_Nd$ and $\zeta=0$ gives
$\Gjp(w)-\Gjp(Z_Nd)$, which proves
\eqref{JC:eq:regular-action-polygon-reduction}.  In particular the pure
polygonal fan block cancels exactly.  JC--AR therefore asks only for the
polygonal support block and its fan--support coupling after the corrector
is inserted; estimating the entire polygonal joint Hessian is unnecessary.

For $d=0$, equation~\eqref{JC:eq:regular-action-polygon-reduction} becomes
\begin{equation}\label{JC:eq:regular-action-pure-support-identity}
 \mathcal Q_N(0,w)=
 \frac12D_h^2\Phi_{\alpha_*}(h_*\mathbf1)[w,w]
 -\frac12\Gjp_{\alpha_*}(w).
\end{equation}
Consequently the archived regular specialization of JP--S0,
\[
 D_h^2\Phi_{\alpha_*}(h_*\mathbf1)[w,w]
 \ge(2-\omega_0(a))\Gjp_{\alpha_*}(w),
 \qquad \omega_0(a)\to0,
\]
would close the pure-support row of JC--AR with any fixed
$\mu_s<1/2$ for all sufficiently large $N$.  The remaining part is then
the modewise Schur payment of the fan--support coupling in
\eqref{JC:eq:regular-action-polygon-reduction}.

The open regular action target is to prove constants
\begin{equation}\label{JC:eq:regular-action-constants}
 \mu_f>0,\qquad \mu_s>0,
\end{equation}
independent of $N$, such that
\begin{equation}\label{JC:eq:regular-action-Schur}
 \boxed{
 \mathcal Q_N(d,w)
 \ge \mu_fa^2\sum_i d_i^2+
      \mu_s\Gjp_{\alpha_*}(w).}
\end{equation}
This is not a consequence of Proposition~\ref{JC:prop:regular-disk-joint},
which controls the full disk scalar rather than the completed action.

More explicitly, cyclic equivariance diagonalizes every surviving quotient
mode $k$.  With the normalization
\begin{align*}
 \mathcal Q_{N,k}(d_k,w_k)
 &=q^{ff}_{N,k}|d_k|^2
   +2\operatorname {Re}(\overline{q^{fs}_{N,k}}d_k\overline{w_k})
   +q^{ss}_{N,k}|w_k|^2,\\
 \Gjp_{\alpha_*}(w_k)&=g_{N,k}|w_k|^2,
 \qquad g_{N,k}>0,
\end{align*}
write the polygonal half-Hessian on the same character as
\[
 \mathcal P_{N,k}(d_k,\zeta_k)
 =p^{ff}_{N,k}|d_k|^2
  +2\operatorname {Re}
    (\overline{p^{fs}_{N,k}}d_k\overline{\zeta_k})
  +p^{ss}_{N,k}|\zeta_k|^2,
 \qquad (Z_Nd)_k=Z_{N,k}d_k.
\]
Expanding the already proved polygonal reduction
\eqref{JC:eq:regular-action-polygon-reduction} gives the exact coefficients
\begin{equation}\label{JC:eq:regular-action-polygon-mode-coefficients}
 \boxed{
 \begin{aligned}
 q^{ss}_{N,k}&=p^{ss}_{N,k}-\frac12g_{N,k},\\
 q^{fs}_{N,k}&=p^{fs}_{N,k}+p^{ss}_{N,k}\overline{Z_{N,k}},\\
 q^{ff}_{N,k}&=3\kappa a^2
 {}+2\operatorname {Re}(p^{fs}_{N,k}Z_{N,k})
 +(p^{ss}_{N,k}+g_{N,k})|Z_{N,k}|^2.
 \end{aligned}}
\end{equation}
Thus the pure polygonal fan coefficient $p^{ff}_{N,k}$ cancels.  Before
positive margins are inserted, the remaining Schur complement is
\[
 3\kappa a^2+\frac12g_{N,k}|Z_{N,k}|^2
 -\frac{|p^{fs}_{N,k}+\frac12g_{N,k}\overline{Z_{N,k}}|^2}
 {p^{ss}_{N,k}-\frac12g_{N,k}}.
\]
This is an algebraic simplification, not a positivity assertion: the
polygonal support coefficient and the fan--support covector remain to be
estimated.

The same cancellation remains exact after margins are inserted.  Put
\[
 A_{N,k}:=p^{fs}_{N,k}
 +\frac12g_{N,k}\overline{Z_{N,k}},
 \qquad
 r_{N,k}:=p^{ss}_{N,k}
 -\left(\frac12+\mu_s\right)g_{N,k}.
\]
Here $A_{N,k}$ is precisely the regular fan derivative of the completed
covector $b_\alpha=\ell_\alpha-p_\alpha/2$ in the same orthonormal
Fourier normalization.  When $r_{N,k}>0$, the full modal Schur row is
equivalent to
\begin{equation}\label{JC:eq:regular-action-reduced-Schur}
 \boxed{
 (3\kappa-\mu_f)a^2
 +\left(\frac12-\mu_s\right)g_{N,k}|Z_{N,k}|^2
 -\frac{|A_{N,k}+\mu_sg_{N,k}\overline{Z_{N,k}}|^2}
 {r_{N,k}}\ge0.}
\end{equation}
Thus the mixed analytic blocker is a sharp all-character estimate for
$D_\alpha(\ell_\alpha-p_\alpha/2)$, not for the discarded covector
$\ell_\alpha-p_\alpha$.

Equation~\eqref{JC:eq:regular-action-Schur} is equivalent, mode by mode, to
\begin{equation}\label{JC:eq:regular-action-mode-matrix}
 \boxed{
 \begin{pmatrix}
  q^{ff}_{N,k}-\mu_fa^2&q^{fs}_{N,k}\\
  \overline{q^{fs}_{N,k}}&q^{ss}_{N,k}-\mu_sg_{N,k}
 \end{pmatrix}\succeq0.}
\end{equation}
Equivalently, with the usual quotient interpretation when the lower-right
entry vanishes,
\begin{align}
 q^{ss}_{N,k}-\mu_sg_{N,k}&\ge0,
 \label{JC:eq:regular-action-support-row}\\
 q^{ff}_{N,k}-\mu_fa^2
 &\ge\frac{|q^{fs}_{N,k}|^2}
 {q^{ss}_{N,k}-\mu_sg_{N,k}}.
 \label{JC:eq:regular-action-fan-Schur-row}
\end{align}
In particular, the pure support specialization already asks for
\begin{equation}\label{JC:eq:regular-action-support-gap}
 D_h^2\Phi_{\alpha_*}(h_*\mathbf1)[z,z]
 \ge(1+2\mu_s)\Gjp_{\alpha_*}(z).
\end{equation}
The exact fixed-fan support Hessian identifies the left side with the
regular-polygon reduced Dirichlet-to-Neumann block plus its geometric
area row.  More precisely, for the orthonormal character
$z_i=N^{-1/2}e^{\mathrm i k a i}$ let
$f_{N,k}=\mathcal I_{\alpha_*}z$,
$g_{N,k}=q_D(f_{N,k})$, let $p_N=-\partial_\nu u_N$, and let
$\psi_{N,k}$ be the character-$k$ reduced Helmholtz solution on the
area-$\pi$ regular polygon with boundary value $z_ip_N$ on side $i$.
Writing
\[
 \mathscr D_{N,k}=\int_{R_N}|\nabla\psi_{N,k}|^2
 -\lambda_N\int_{R_N}|\psi_{N,k}|^2,
\]
the exact support Hessian and area symbol give
\begin{equation}\label{JC:eq:regular-support-DtN-sector}
 \boxed{
 p^{ss}_{N,k}=\mathscr D_{N,k}
 +\frac{\lambda_N}{\pi}
   \frac{\cos(ka)-\cos a}{\sin a},
 \qquad
 q^{ss}_{N,k}=p^{ss}_{N,k}-\frac12g_{N,k}.}
\end{equation}
Hence the pure-support row is exactly the uniform sector inequality
\[
 \mathscr D_{N,k}
 +\frac{\lambda_N}{\pi}
   \frac{\cos(ka)-\cos a}{\sin a}
 \ge\left(\frac12+\mu_s\right)g_{N,k},
 \qquad 2\le k\le N-2.
\]
All terms except the regular-polygon reduced Helmholtz DtN energy are
explicit.  No closed terminal/LCA theorem proves this inequality, even
before the mixed Schur row
\eqref{JC:eq:regular-action-fan-Schur-row} is considered.

The exact Schwarz--Christoffel representation also shows why the DtN and
area rows must be estimated together.  After a rotation,
\[
 F_N'(z)=c_N(1-z^N)^{-2/N},\qquad J_N=|F_N'|,
\]
and, for the unnormalized character $z_i=e^{ikai}$, let
$\sigma_{N,k}$ be the step function on $\mathbb T$ which equals
$e^{ikai}$ on the preimage arc of side $i$ (its endpoint values are
irrelevant).  The pulled-back boundary trace is
\[
 h_{N,k}=\sigma_{N,k}c_N^{-1}
 |1-e^{iN\theta}|^{2/N}q_N,
 \qquad q_N=-\partial_r(u_N\circ F_N)|_{\mathbb T}.
\]
Uniform disk $W^{2,p}$ estimates give $q_N\to d_\circ$ uniformly,
including at the prevertices.  Keeping only the cross-prevertex
rectangles in the Gagliardo norm gives the harmonic-extension lower bound
$H(h_{N,k})\ge cN^2\{1-\cos(ka)\}-C$.  To pass to the reduced Helmholtz
energy, use the uniform boundary Friedrichs inequality
\[
 \|\phi\|_{L^2(R_N)}^2
 \le\varepsilon\|\nabla\phi\|_{L^2(R_N)}^2
 +C_\varepsilon\|h_{N,k}\|_{L^2(\partial R_N)}^2.
\]
Choose $\varepsilon<(2\sup_N\lambda_N)^{-1}$.  Rellich's identity gives
$\|h_{N,k}\|_{L^2(\partial R_N)}^2
=\int_{\partial R_N}p_N^2=2\lambda_N/h_N\le C$ in the unnormalized
character.  Hence every admissible reduced competitor has energy at least
$\frac12H(h_{N,k})-C$; the ground-state orthogonality only decreases the
competitor set.  This proves the rigorous DtN-only lower bound
\begin{equation}\label{JC:eq:corner-layer-DtN-lower}
 \mathscr D^{\rm un}_{N,k}
 \ge cN^2\{1-\cos(ka)\}-C.
\end{equation}
This does \emph{not} close the support row.  Formula
\eqref{JC:eq:regular-support-DtN-sector} uses the orthonormal character;
in the unnormalized convention its area term must also be multiplied by
$N$ and is of order
$-N^2\{1-\cos(ka)\}$.  The two leading corner-scale terms cancel in the
disk limit.  More precisely, the uncancelled coefficient contains
\[
 \frac{2}{a^2}
 \{\pi c_N^{-2}q_{N,*}^2-\lambda_N\}
 \{1-\cos(ka)\},
\]
for the prevertex endpoint amplitude $q_{N,*}$.  Rate-free convergence
alone does not control this coefficient after multiplication by $a^{-2}$.
One must consequently estimate the jointly renormalized remainder,
rather than combine
\eqref{JC:eq:corner-layer-DtN-lower} with an area symbol written in the other
normalization.  The next proposition supplies precisely that expansion,
uniformly through the mesoscopic regime $ka\downarrow0$.
\end{contract}

\begin{lemma}[Fixed-character polygon--disk support limit]
\label{JC:lem:fixed-character-support-limit}
For every fixed integer $k\ge2$,
\begin{equation}\label{JC:eq:fixed-character-support-limit}
 \frac{p^{ss}_{N,k}}{g_{N,k}}\longrightarrow1
 \qquad(N\to\infty).
\end{equation}
\end{lemma}

\begin{proof}
Use the unnormalized support character $z_i=e^{\mathrm i kia}$; the
normalization cancels from the quotient.  The material vertex velocity
at the vertex with bisector angle $\varphi_i=(i+\tfrac12)a$ is exactly
\[
 e^{\mathrm i k\varphi_i}\left\{
 \frac{\cos(ka/2)}{\cos(a/2)}e_r(\varphi_i)
 +\mathrm i\frac{\sin(ka/2)}{\sin(a/2)}
 e_\theta(\varphi_i)\right\}.
\]
For fixed $k$ its physical side-affine interpolant converges in
$H^{1/2}(\mathbb T;\mathbb C^2)$ to the smooth disk field
\[
 V_k^\circ(\theta)=e^{\mathrm i k\theta}
 \{e_r(\theta)+\mathrm i k e_\theta(\theta)\},
 \qquad V_k^\circ\cdot e_r=e^{\mathrm i k\theta}.
\]

We record the fixed-slot Hessian passage rather than citing an
unassembled stage.  Pull the area-normalized regular polygons to the
disk by the normalized near-identity conformal charts.  Their stiffness
forms and mass densities converge to the disk forms in every fixed
$L^p$, their first eigenfunctions converge strongly in $H_0^1$, and
the uniform convex-domain gradient bound makes
$\nabla u_N\otimes\nabla u_N\to
\nabla u_\circ\otimes\nabla u_\circ$ in $L^2$.
Polarize the exact volume formula for the normalized eigenvalue
Hessian with the fixed smooth extension of $V_k^\circ$ in both slots.
Every direct stiffness, mass, area, and product-of-first-jets row then
converges by the preceding strong coefficient convergence.

For the Schur row, put the fixed first-variation functional into the
reduced resolvent.  The spectral gap above the ground state is uniform,
the functionals converge in $H^{-1}$, and the resolvent identity on the
moving ground-state complements gives strong $H_0^1$ convergence of
the reduced states.  Thus the Schur pairing converges as well.  This
proves convergence of the complete polygonal normalized Hessian to the
disk Hessian on $V_k^\circ$; no mesh-ratio estimate is used because the
slot is fixed and smooth.

At the disk, scale and translation modes have already been projected
out and the normalized Hessian is
\[
 D^2\Phi_\circ[V_k^\circ,V_k^\circ]
 =2q_D(e^{\mathrm i k\theta}).
\]
On the other hand, the exact normal-sine interpolant
$f_{N,k}=\mathcal I_{\alpha_*}z$ converges to
$e^{\mathrm i k\theta}$ in $H^{1/2}$, so
$g_{N,k}\to q_D(e^{\mathrm i k\theta})$.
Since $p^{ss}_{N,k}$ is the polygonal half-Hessian coefficient,
\eqref{JC:eq:fixed-character-support-limit} follows.
\end{proof}

\begin{proposition}[Jointly renormalized regular support sector]
\label{JC:prop:regular-support-renormalized}
There is $N_s$ such that, for $N\ge N_s$ and every quotient character
$2\le k\le N-2$,
\begin{equation}\label{JC:eq:regular-support-closed-gap}
 \boxed{
 p^{ss}_{N,k}\ge\frac{11}{20}g_{N,k}
 =\left(\frac12+\frac1{20}\right)g_{N,k}.}
\end{equation}
Consequently the pure-support row
\eqref{JC:eq:regular-action-support-row} is closed with
$\mu_s=1/20$.
\end{proposition}

\begin{proof}
We keep one normalization until the last step.  By complex conjugation
assume $2\le k\le N/2$, and put
\[
 \beta=\frac2N,\qquad a=\pi\beta,\qquad
 x=\frac{k}{N}\in(0,\tfrac12],\qquad \xi=2\pi x.
\]
Use first the unnormalized complex character
$z_i=e^{2\pi\mathrm i x i}$.
Every quadratic value is then $N$ times its coefficient in the
orthonormal complex character.  The orthonormal real vector
$\sqrt{2/N}\operatorname{Re}z$ has the same modal coefficient.

Let
\[
 w_\beta(t)=|2\sin(\pi t)|^\beta,\qquad0<t<1.
\]
For the constant disk flux $d_\circ$, the exact corner trace is
$d_\circ\sigma_{N,k}w_\beta$.  Its only circular frequencies are
$k+mN$, and the beta integral gives
\begin{equation}\label{JC:eq:corner-Gamma-coefficient}
 I_{\beta,m}(x):=\int_0^1w_\beta(t)e^{-2\pi\mathrm i(m+x)t}\,dt
 =
 e^{-\pi\mathrm i(m+x)}
 \frac{\Gamma(1+\beta)}
 {\Gamma(1+\beta/2+m+x)\Gamma(1+\beta/2-m-x)}.
\end{equation}
Since $\pi d_\circ^2=\lambda_\circ$ by Rellich's identity, harmonic
energy and the area row, divided by $Nd_\circ^2$, combine into
\begin{equation}\label{JC:eq:corner-joint-finite-N-symbol}
 P_\beta(x):=
 2\pi\sum_{m\in\mathbb Z}|m+x||I_{\beta,m}(x)|^2
 +\frac{\cos(2\pi x)-\cos(\pi\beta)}{\sin(\pi\beta)}.
\end{equation}
The two terms in this formula separately contain the same
$\beta^{-1}$ pole with opposite signs.

We extract the finite part before estimating.  For $0<y<1$, put
\[
 T_\beta(y):=\sum_{n\ge0}(n+y)
 \left\{\frac{\Gamma(n+y-\beta/2)}
 {\Gamma(n+y+1+\beta/2)}\right\}^2.
\]
Reflection in \eqref{JC:eq:corner-Gamma-coefficient} gives the positive
identity
\begin{align}
 &\sum_m|m+x||I_{\beta,m}(x)|^2\notag\\
 &\quad=
 \frac{\Gamma(1+\beta)^2}{\pi^2}
 \left\{\sin^2\!\pi(x-\beta/2)T_\beta(x)
 +\sin^2\!\pi(x+\beta/2)T_\beta(1-x)\right\}.
 \label{JC:eq:corner-reflected-positive-series}
\end{align}
Separate the $n=0$ term.  Taylor's formula for $\log\Gamma$, together
with $|\psi_{\rm dg}^{(r)}(u)|\le C_r(1+u^{-r-1})$, gives, uniformly
for $0<\beta\le x/2$ and $x\le1/2$,
\begin{align*}
 T_\beta(x)
 &=\zeta(1+2\beta,x)
 +O\left(\frac\beta{x^2}
 +\frac{\beta|\log x|}{x}
 +\frac{\beta^2}{x^3}+\beta\right),\\
 T_\beta(1-x)&=\zeta(1+2\beta,1-x)+O(\beta),\\
 \zeta(1+2\beta,x)
 &=\frac1{2\beta}-\psi_{\rm dg}(x)
 +O\left(\frac{\beta|\log x|}{x}
 +\frac{\beta^2|\log x|^2}{x}+\beta\right),
\end{align*}
where $\psi_{\rm dg}=\Gamma'/\Gamma$.  Indeed, after comparison with
$(n+y)^{-1-2\beta}$ the derivative of the logarithmic quotient is
\[
 -2\psi_{\rm dg}(n+y)-(n+y)^{-1}+2\log(n+y)
 =O((n+y)^{-2}),
\]
so all $n\ge1$ remainders are absolutely summable; the singular powers
come only from $n=0$.

Expand the sine squares in
\eqref{JC:eq:corner-reflected-positive-series}, retaining their opposite
linear terms until cancellation, and use
$\Gamma(1+\beta)^2=1-2\gamma\beta+O(\beta^2)$.  We obtain
\begin{align}
 P_\beta(x)&=P(x)+
 O\left(\beta+\beta x|\log x|+\frac{\beta^2}{x}\right),
 \label{JC:eq:corner-mesoscopic-remainder}\\
 P(x)&:=-\frac{2\sin^2(\pi x)}{\pi}
 \{\psi_{\rm dg}(x)+\psi_{\rm dg}(1-x)+2\gamma\}.
 \label{JC:eq:corner-renormalized-symbol}
\end{align}

The high-character normal-sine disk symbol is
\begin{equation}\label{JC:eq:corner-disk-symbol}
 Q(x):=\frac{2\sin^4(\pi x)}{\pi^3}
 \{\zeta(3,x)+\zeta(3,1-x)\}.
\end{equation}
It satisfies $Q(x)\asymp x$ on $(0,1/2]$.  We now prove a strict
analytic gap.  Set
\[
 D(x):=-\{\psi_{\rm dg}(x)+\psi_{\rm dg}(1-x)+2\gamma\},
 \qquad H_3(x):=\zeta(3,x)+\zeta(3,1-x).
\]
The digamma and paired Hurwitz series give
\[
 D(x)=\frac1x+2\sum_{r\ge1}\zeta(2r+1)x^{2r},
\]
and
\[
 H_3(x)=x^{-3}+\sum_{n\ge1}\{(n-x)^{-3}+(n+x)^{-3}\}
 \le x^{-3}+14\zeta(3)-8.
\]
The tail increases with $x$ and the last constant is its value at
$x=1/2$.  Using $\sin(\pi x)\le\pi x$ and the integral-tail bound
$\zeta(3)<121/100$, we get
\begin{align*}
 D(x)-\frac35\frac{\sin^2(\pi x)}{\pi^2}H_3(x)
 &\ge\frac{2}{5x}+\frac{24-32\zeta(3)}5x^2\\
 &>\frac{2}{5x}-\frac{368}{125}x^2>0
 \qquad(0<x\le\tfrac12).
\end{align*}
Therefore
\begin{equation}\label{JC:eq:corner-three-fifths}
 \boxed{P(x)>\frac35Q(x),\qquad0<x\le\frac12.}
\end{equation}
Dividing \eqref{JC:eq:corner-mesoscopic-remainder} by $Q(x)\asymp x$
also shows
\begin{equation}\label{JC:eq:corner-relative-mesoscopic}
 P_\beta(x)=P(x)+o(Q(x))
\end{equation}
along every sequence for which $k=2x/\beta\to\infty$.

It remains to replace the constant flux by the true polygonal flux.
The exact area normalization of
$F_N'(z)=c_N(1-z^N)^{-\beta}$ is
\[
 c_N^{-2}=\sum_{m\ge0}
 \frac{(\beta)_m^2}{(m!)^2(mN+1)}.
\]
Since $(\beta)_m/m!\le C\beta m^{\beta-1}$ for $m\ge1$,
$c_N=1+O(\beta^3)$.  Put
$W_N=|F_N'|^2=c_N^2|1-z^N|^{-2\beta}$.  For every fixed $p>4$,
the change of variables $t=r^N$ gives
\begin{equation}\label{JC:eq:corner-weight-Lp-rate}
 \|W_N-1\|_{L^p(\mathbb D)}
 \le C_p\beta^{1+1/p}.
\end{equation}
Indeed, on $t\le1/2$ one has
$|\log|1-te^{\mathrm i\phi}||\le Ct$, while on $t\ge1/2$ all fixed
logarithmic moments remain uniformly integrable against
$|1-te^{\mathrm i\phi}|^{-2p\beta}$.  Since
$r\,dr=(\beta/2)t^{\beta-1}dt$, the $p$-th power of the norm is
$O(\beta^{p+1})$.

The pulled eigenfunction satisfies
\[
 -\Delta v_N=\lambda_NW_Nv_N,\qquad v_N|_{\mathbb T}=0,
 \qquad\int_{\mathbb D}W_Nv_N^2=1.
\]
The disk spectral gap first gives an $H^1$ estimate with the rate in
\eqref{JC:eq:corner-weight-Lp-rate}.  Apply a $W^{2,q}$ estimate with
$q>1$ to obtain $L^\infty$ control, substitute it back into the
difference equation, and then apply the $W^{2,p}$ estimate.  Thus
\begin{equation}\label{JC:eq:corner-flux-holder-rate}
 |\lambda_N-\lambda_\circ|
 +\|v_N-u_\circ\|_{W^{2,p}(\mathbb D)}
 \le C_p\beta^{1+1/p},
 \qquad
 \|q_N-d_\circ\|_{C^\alpha(\mathbb T)}
 \le C_p\beta^{1+1/p},
\end{equation}
where $q_N=-\partial_rv_N|_{\mathbb T}$ and
$\alpha=1-2/p>1/2$.

For $r\in C^\alpha$ and $f\in H^{1/2}(\mathbb T)$, the Gagliardo
double integral gives
\[
 [rf]_{H^{1/2}}^2
 \le2\|r\|_\infty^2[f]_{H^{1/2}}^2
 +C[r]_{C^\alpha}^2\|f\|_2^2.
\]
The positive Gamma series
\eqref{JC:eq:corner-reflected-positive-series} and
$T_\beta(y)\le C(\beta^{-1}+y^{-1})$ imply
\begin{equation}\label{JC:eq:corner-trace-energy-upper}
 H(\sigma_{N,k}w_\beta)\le Ck^2,
 \qquad2\le k\le N/2.
\end{equation}
Consequently polarization and
\eqref{JC:eq:corner-flux-holder-rate} give, with
$\varepsilon_N=C_p\beta^{1+1/p}$,
\begin{equation}\label{JC:eq:corner-true-flux-error}
 \left|H(\sigma_{N,k}c_N^{-1}w_\beta q_N)
 -d_\circ^2H(\sigma_{N,k}w_\beta)\right|
 \le C\varepsilon_Nk^2=o(k)
\end{equation}
uniformly along every growing-character sequence, because
$\varepsilon_Nk\le C\beta^{1/p}$.  The eigenvalue error in the area
row has the same bound.

The sector compatibility
$\int_{\partial R_N}z_ip_N\partial_\nu u_N\,ds=0$
follows from cyclic character orthogonality.  The uniform gap on
$u_N^\perp$, the Poisson $L^2$ estimate, and
$h_N\int_{\partial R_N}p_N^2\,ds=2\lambda_N$ show that reduced
Helmholtz energy differs from harmonic extension by $O(1)$ in the
unnormalized character.  This is also $o(k)$ when $k\to\infty$.
It follows from
\eqref{JC:eq:corner-relative-mesoscopic} and
\eqref{JC:eq:corner-true-flux-error} that
\begin{equation}\label{JC:eq:corner-polygon-symbol-limit}
 \frac{p^{ss,\mathrm{un}}_{N,k}}{Nd_\circ^2}
 =P(k/N)+o(Q(k/N))
\end{equation}
on every growing-character sequence.

Finally, the exact Bessel alias formula for $g_{N,k}$, together with
\[
 \omega_n=|n|+1+O(|n|^{-1}),\qquad
 (1-n^2)^{-2}=n^{-4}\{1+O(n^{-2})\},
\]
gives
\begin{equation}\label{JC:eq:corner-disk-alias-limit}
 \frac{g^{\mathrm{un}}_{N,k}}{Nd_\circ^2}
 =Q(k/N)\{1+o(1)\}
\end{equation}
uniformly for $k\to\infty$.  Indeed every alias has
$|k+mN|\ge k$, and replacing
$\cos(ka)-\cos a$ by $\cos(ka)-1$ costs $O(k^{-2})$
relatively.

If \eqref{JC:eq:regular-support-closed-gap} failed along a sequence, choose
a violating character and conjugate it into $2\le k_N\le N/2$.
If $k_N$ is bounded, pass to a constant subsequence; the retained
one-fixed-slot disk limit gives
$p^{ss}_{N,k_N}/g_{N,k_N}\to1$.  If $k_N\to\infty$, then
\eqref{JC:eq:corner-three-fifths},
\eqref{JC:eq:corner-polygon-symbol-limit}, and
\eqref{JC:eq:corner-disk-alias-limit} give
\[
 \liminf_N\frac{p^{ss}_{N,k_N}}{g_{N,k_N}}\ge\frac35.
\]
Both alternatives contradict a ratio below $11/20$.  This proves
\eqref{JC:eq:regular-support-closed-gap}.
\end{proof}

\begin{lemma}[Polygonal mixed Schwarz--Christoffel finite part]
\label{JC:lem:regular-polygonal-mixed-finite-part}
Fix $0<x<1/2$, put
\[
 \theta=\pi x,\qquad \rho=e^{2\pi ix},\qquad
 D(x)=-\{\psi(x)+\psi(1-x)+2\gamma\},
 \qquad T(x)=\psi_1(x)-\psi_1(1-x)=-D'(x).
\]
For $\beta=2/N$, $a=\pi\beta$, and a regular fan character with
$k/N\to x$, the constant-flux polygonal fan--support coefficient has the
phase-realified limit
\begin{equation}\label{JC:eq:regular-polygonal-mixed-finite-part}
 \boxed{
 L_{\rm poly}(x):=
 \lim_{\beta\downarrow0}
 \left[-e^{i\theta}\frac{p^{fs}_\beta}{d_\circ^2a}\right]
 =\frac{\cos\theta}{2\pi}D(x)
 -\frac{\sin\theta}{2\pi^2}T(x)
 =\frac{\sin\theta}{2\pi^2}
 \{D'(x)+\pi\cot\theta D(x)\}.}
\end{equation}
Let
\[
 F_p(y):=\sum_{n\ge0}\frac{(-1)^n}{(n+y)^p}
 =2^{-p}\{\zeta(p,y/2)-\zeta(p,(y+1)/2)\}
\]
and define the exact disk-forcing symbol
\begin{equation}\label{JC:eq:regular-disk-forcing-finite-part}
 R_I(x)=-\frac{\sin\theta}{4\pi^3}
 \{F_3(x)-F_3(1-x)\}
 +\frac{\cos\theta}{4\pi^2}
 \{F_2(x)+F_2(1-x)\}.
\end{equation}
Then the genuine polygonal completed-covector finite part exists and is
\begin{equation}\label{JC:eq:regular-completed-covector-finite-part}
 \boxed{
 R_{b,{\rm poly}}(x)=
 \frac{D'(x)+\pi\cot\theta D(x)}{2\pi^2}-R_I(x),
 \qquad
 \frac{A_{N,k}}{d_\circ^2a}
 =-e^{-i\theta}\sin\theta R_{b,{\rm poly}}(x)+o_x(1).}
\end{equation}
The formulas extend continuously to $x=1/2$, where both
$R_I$ and $R_{b,{\rm poly}}$ vanish.
\end{lemma}

\begin{proof}
On one regular Schwarz--Christoffel side let
\[
 w_\beta(t)=|2\sin\pi t|^\beta,\qquad
 J_\beta(t)=\int_{1/2}^t w_\beta(u)^{-1}\,du,\qquad 0<t<1.
\]
For the complex fan character $d_i=\rho^i$, differentiating the turning
angles gives $\dot\theta_i=(\rho-1)^{-1}\rho^i$ and $\dot h=0$.
Centred physical arclength satisfies
$ds=c_Naw_\beta^{-1}dt$ and $s=c_NaJ_\beta(t)$, while the pulled-back
flux is $p_N=c_N^{-1}w_\beta q_N$.  Hence the exact fan--flux datum is
\[
 -a(\rho-1)^{-1}\rho^i w_\beta(t)J_\beta(t)q_N(t).
\]
For $y=m+x$ put
\[
 I_\beta(y)=\int_0^1w_\beta(t)e^{-2\pi iyt}\,dt,\qquad
 S_\beta(y)=\int_0^1w_\beta(t)J_\beta(t)e^{-2\pi iyt},dt,
\]
and $M_\beta(x)=\sum_m|m+x|S_\beta(m+x)\overline{I_\beta(m+x)}$.
The harmonic mixed row is
$-2\pi a(\rho-1)^{-1}M_\beta(x)$ after division by $d_\circ^2$; the area
row is
\[
 \frac14h_N\sec^2(a/2)(1+\rho^{-1}),\qquad
 h_N=\sqrt{\frac\pi{N\tan(a/2)}}.
\]

It remains to extract the joint finite part.  Set
\[
 C_\beta=\int_0^{1/2}w_\beta^{-1}
 =2^{-\beta}\frac{\Gamma((1-\beta)/2)}
 {2\sqrt\pi\,\Gamma(1-\beta/2)}.
\]
Then $C_0=1/2$, $C'_0=0$, and direct integration at the two endpoints
gives
\begin{align*}
 w_\beta(t)J_\beta(t)
 &=-C_\beta(2\pi t)^\beta+\frac{t}{1-\beta}
   +O(t^{\beta+2}+t^3),\\
 w_\beta(1-u)J_\beta(1-u)
 &= C_\beta(2\pi u)^\beta-\frac{u}{1-\beta}
   +O(u^{\beta+2}+u^3).
\end{align*}
The endpoint expansion can be used here with a summable differentiated
remainder, as follows.  Choose smooth cutoffs at $0$ and $1$, subtract
from $w_\beta$ its two endpoint models of type $t^\beta$, and subtract
from $w_\beta J_\beta$ the four endpoint models displayed above.  The
remainders and their first $\beta$--derivatives have two integrable
derivatives; differentiating in $\beta$ merely inserts a logarithm.

Here is the resulting coefficient ledger, which also fixes the two tail
signs.  Put
\[
 a_\beta=\frac{\Gamma(1+\beta)}{2\pi},\qquad
 \phi_\beta=\frac\pi2(1+\beta).
\]
For $\sigma=\operatorname{sgn}y$, the two endpoint models give
\begin{align*}
 A_+(\beta,x)&=a_\beta
  \{e^{-i\phi_\beta}+e^{-2\pi ix}e^{i\phi_\beta}\},&
 B_+(\beta,x)&=C_\beta a_\beta
  \{-e^{-i\phi_\beta}+e^{-2\pi ix}e^{i\phi_\beta}\},\\
 A_-(\beta,x)&=a_\beta
  \{e^{i\phi_\beta}+e^{-2\pi ix}e^{-i\phi_\beta}\},&
 B_-(\beta,x)&=C_\beta a_\beta
  \{-e^{i\phi_\beta}+e^{-2\pi ix}e^{-i\phi_\beta}\}.
\end{align*}
Thus, for $j=0,1$, fixed $x\in(0,1/2]$, and
$0\le\beta\le\beta_x<1$,
\begin{align}
 I_\beta(y)&=A_\sigma(\beta,x)|y|^{-1-\beta}
             +E^I_\beta(y),\notag\\
 S_\beta(y)&=B_\sigma(\beta,x)|y|^{-1-\beta}
             +E^S_\beta(y),\notag\\
 |\partial_\beta^jE^I_\beta(y)|
 +|\partial_\beta^jE^S_\beta(y)|
 &\le C_x\frac{\{1+\log(2+|y|)\}^{j}}{(1+|y|)^2}.
 \label{JC:eq:regular-mixed-Fourier-remainders}
\end{align}
The linear endpoint models are included in $E^S$; their coefficients are
$O(|y|^{-2})$.  The displayed estimate follows by two integrations by
parts on the cutoff remainders.  It is uniform down to $\beta=0$ because
$t^\beta\log t$ and the second derivatives of
$t^{\beta+2}\log t$ are integrable uniformly on
$0\le\beta\le\beta_x$.

Direct multiplication gives
\[
 B_+(\beta,x)\overline{A_+(\beta,x)}
 =2iC_\beta a_\beta^2\sin(2\pi x-\pi\beta),
 \qquad
 B_-(\beta,x)\overline{A_-(\beta,x)}
 =2iC_\beta a_\beta^2\sin(2\pi x+\pi\beta).
\]
If $c_\pm(\beta,x)$ denote these two coefficients and
\[
 R_m(\beta,x):=|m+x|S_\beta(m+x)\overline{I_\beta(m+x)}
 -c_{\operatorname{sgn}(m+x)}(\beta,x)|m+x|^{-1-2\beta},
\]
then~\eqref{JC:eq:regular-mixed-Fourier-remainders}, followed by the product
rule, gives
\begin{equation}\label{JC:eq:regular-mixed-endpoint-remainder}
 |R_m(\beta,x)|+|\partial_\beta R_m(\beta,x)|
 \le C_x\frac{1+\log(2+|m|)}{(1+|m|)^2}.
\end{equation}
Indeed the coefficients of every subtracted singular model follow from
\[
 \int_0^\infty t^\alpha e^{-2\pi iyt}\,dt
 =\Gamma(1+\alpha)e^{-i\pi(1+\alpha)/2}
  (2\pi y)^{-1-\alpha},\qquad y>0,
\]
first with an exponential cutoff and then by continuation.  This proves
the displayed bound rather than assuming a formal asymptotic series.
Its majorant is summable, so the subtracted series is $C^1$ at
$\beta=0$ and termwise passage to the finite part is justified.  The
only non-summable tail coefficients are
\[
 c_\pm(\beta,x)=2iC_\beta
 \left\{\frac{\Gamma(1+\beta)}{2\pi}\right\}^{\!2}
 \sin(2\pi x\mp\pi\beta),
\]
multiplying the corresponding $|m+x|^{-1-2\beta}$ tails.  The linear
endpoint terms produce summable order $|m|^{-2-\beta}$ terms, so no
additional $\beta\zeta(1+\beta)$ finite part is possible.

More explicitly,
\[
 M_\beta(x)=c_+(\beta,x)\zeta(1+2\beta,x)
 +c_-(\beta,x)\zeta(1+2\beta,1-x)
 +\sum_{m\in\mathbb Z}R_m(\beta,x).
\]
At $\beta=0$ the exact formulas below give
\[
 R_m(0,x)=-\frac{i\sin^2\theta}{2\pi^3}
 \frac{\operatorname{sgn}(m+x)}{|m+x|^2},
 \qquad
 \sum_mR_m(0,x)=-\frac{i\sin^2\theta}{2\pi^3}T(x).
\]
Also $C'_0=0$ and
$\partial_\beta\Gamma(1+\beta)^2|_{\beta=0}=-2\gamma$; the opposite
$\mp\pi\cos(2\theta)$ derivatives of the two sine coefficients cancel.
These identities account for every finite term in the next display.

At $\beta=0$ one has
\[
 I_0(y)=e^{-\pi iy}\frac{\sin\pi y}{\pi y},\qquad
 S_0(y)=\frac i2I_0(y)
 \left(\cot\theta-\frac1{\pi y}\right).
\]
Using
$\zeta(1+2\beta,q)=(2\beta)^{-1}-\psi(q)+O(\beta)$ in the two tails and
summing the $C^1$ remainder from
\eqref{JC:eq:regular-mixed-endpoint-remainder} yields, locally uniformly in
$x$,
\[
 M_\beta(x)=\frac{i\sin(2\theta)}{4\pi^2\beta}
 +i\left\{\frac{\sin(2\theta)}{4\pi^2}D(x)
 -\frac{\sin^2\theta}{2\pi^3}T(x)\right\}+O_x(\beta).
\]
Because $2i(\rho-1)^{-1}\sin(2\theta)=1+\rho^{-1}$ and
$h_N\sec^2(a/2)=1+O(a^2)$, the pole cancels the area row.  Multiplication
by $-e^{i\theta}$ proves~\eqref{JC:eq:regular-polygonal-mixed-finite-part}.
The exact disk alias calculation gives
\eqref{JC:eq:regular-disk-forcing-finite-part}; adding the corrector
contribution $-\sin\theta R_I$ proves
\eqref{JC:eq:regular-completed-covector-finite-part}.  Reflection about
$x=1/2$ gives the stated continuous endpoint values.
\end{proof}

\begin{proposition}[Strict support--double-disk gap]
\label{JC:prop:regular-support-double-disk-gap}
For every $0<x\le1/2$ one has
\begin{equation}\label{JC:eq:regular-support-double-disk-gap}
 \boxed{P(x)>2Q_{\mathrm{mat}}(x)},\qquad
 Q_{\mathrm{mat}}(x):=\frac{\sin^4(\pi x)}{\pi^3}
 \{\zeta(3,x)+\zeta(3,1-x)\}.
\end{equation}
\end{proposition}

\begin{proof}
Write $H_3(x)=\zeta(3,x)+\zeta(3,1-x)$.  The alternating Taylor
series for $\sin^2y$ on $0\le y\le\pi/2$ gives
\[
 \sin^2(\pi x)\le \pi^2x^2(1-2x^2).
\]
Indeed the alternating expansion gives
$\sin^2y\le y^2-y^4/3+2y^6/45$.  Since $y\le\pi/2$ and
$1/3-2/\pi^2-\pi^2/90>0$, its right side is at most
$y^2-(2/\pi^2)y^4$.  The convergent series for $D$ gives
\[
 D(x)=\frac1x+2x^2\sum_{n\ge1}\frac1{n(n^2-x^2)}
 \ge\frac1x+2\zeta(3)x^2.
\]
Moreover $H_3(x)-x^{-3}$ is increasing on $[0,1/2]$, term by term, and
therefore
\[
 H_3(x)\le x^{-3}+H_3(1/2)-8
 =x^{-3}+14\zeta(3)-8.
\]
Using the elementary bounds $6/5<\zeta(3)<121/100$ (sum the first ten
terms and bound the decreasing tail between
$\int_{11}^\infty t^{-3}dt$ and $\int_{10}^\infty t^{-3}dt$), it follows that
\begin{align*}
 D(x)-\frac{\sin^2(\pi x)}{\pi^2}H_3(x)
 &\ge 2x+(8-12\zeta(3))x^2+(28\zeta(3)-16)x^4\\
 &\ge x\left(2-\frac{163}{25}x+\frac{88}{5}x^3\right)>0.
\end{align*}
For the last inequality, the cubic bracket has its minimum at
$x_0=\sqrt{163/1320}<44/125$, where it equals
$2-(326/75)x_0>0$.  Since
$P(x)>2Q_{\mathrm{mat}}(x)$ is exactly the last displayed strict
inequality multiplied by $2\sin^2(\pi x)/\pi$, the proof is complete.
\end{proof}

\begin{remark}[Retraction of the old period-three completion]
\label{JC:rem:period-three-completion-retracted}
The revision-21 period-three completion used the material support constant
$E_3$.  The jointly renormalized Schwarz--Christoffel calculation above
shows instead that the true real-paired polygonal support finite part is
\[
 \frac12P(1/3)=\frac{9\log3}{4\pi}=0.786823093\ldots,
\]
whereas the material computation gives $E_3=0.77956467\ldots$.  The
missing difference is a same-order corner finite part.  Therefore the old
period-three proposition remains withdrawn.  The proofs below use only
the genuine polygonal diagonal and mixed finite parts computed above; they
do not reinstate the material transfer.
\end{remark}

\begin{proposition}[Certified completed-covector ratio]
\label{JC:prop:regular-completed-covector-ratio}
For $0<x<1/2$, put
\[
 J(x):=\frac{D'(x)+\pi\cot(\pi x)D(x)}{2\pi^2}.
\]
Then
\begin{equation}\label{JC:eq:regular-completed-covector-ratio}
 \boxed{2R_I(x)<J(x)<R_I(x)<0.}
\end{equation}
Consequently, for $0\le\mu\le1/4$ and
$Q_{\rm mat}=Q/2$,
put
\[
 D_{\rm poly}(x;\mu)
 :=P(x)-(1+2\mu)Q_{\rm mat}(x).
\]
Then $D_{\rm poly}(x;\mu)>0$ and
\begin{equation}\label{JC:eq:regular-completed-covector-determinant}
 \begin{aligned}
 \Delta_\mu(x):={}&(1-2\mu)R_I(x)^2
 D_{\rm poly}(x;\mu)\\
 &-Q_{\rm mat}(x)
 \{R_{b,{\rm poly}}(x)-2\mu R_I(x)\}^2>0
 \qquad(0<x<1/2).
 \end{aligned}
\end{equation}
At $x=1/2$ both covectors vanish and $\Delta_\mu(1/2)=0$.
\end{proposition}

\begin{proof}
Write $s=\sin(\pi x)$, $c=\cos(\pi x)$ and
$A=F_2(x)+F_2(1-x)$.  Since
$A'=-2\{F_3(x)-F_3(1-x)\}$, one has the exact identities
\begin{equation}\label{JC:eq:regular-ratio-derivative-identities}
 J=\frac{(sD)'}{2\pi^2s},\qquad
 R_I=\frac{(s^2A)'}{8\pi^3s}.
\end{equation}
Thus it suffices to prove the three strict signs
\[
 -R_I>0,\qquad U:=R_I-J>0,\qquad V:=J-2R_I>0.
\]
We give the finite rational certificate, including the endpoint
regularizations, so that no sampled floating-point value is used.

At the left endpoint the convergent expansions are
\begin{align*}
 D&=x^{-1}+2\sum_{r\ge1}\zeta(2r+1)x^{2r},\\
 \pi\cot(\pi x)&=x^{-1}-2\sum_{r\ge1}\zeta(2r)x^{2r-1},\\
 F_3(x)-F_3(1-x)
 &=x^{-3}-2\sum_{r\ge0}\binom{2r+2}{2}\eta(2r+3)x^{2r},\\
 F_2(x)+F_2(1-x)
 &=x^{-2}+4\sum_{r\ge0}(r+1)\eta(2r+3)x^{2r+1}.
\end{align*}
The singular terms cancel symbolically and give
\[
 J(0)=-\frac16,\qquad R_I(0)=-\frac1{12},\qquad
 \lim_{x\downarrow0}\frac{V(x)}x
 =\frac{3\zeta(3)}{4\pi^2}>0.
\]
Ten retained orders, $q=x^2$, and the elementary bounds
$\zeta(r)<121/100$ for $r\ge3$ and $\zeta(r)<5/3$ for $r\ge2$
bound the omitted $D,D',D/x,\cot$, and divided-cotangent tails by
\begin{gather*}
 \frac{121}{50}\frac{x^{22}}{1-q},\quad
 \frac{121}{25}x^{21}\frac{11-10q}{(1-q)^2},\quad
 \frac{121}{50}\frac{x^{21}}{1-q},\\
 \frac{10}{3}\frac{x^{21}}{1-q},\qquad
 \frac{10}{3}\frac{x^{20}}{1-q}.
\end{gather*}
The two alternating-$F$ tails are dominated by
\[
 \frac{8(12)^2q^{11}}{(1-q)^3},\qquad
 \frac{4x(12)q^{11}}{(1-q)^2},
\]
with the divided second bound obtained by deleting $x$.  Directed rational
evaluation on $0\le x\le1/8$ then gives
\begin{equation}\label{JC:eq:regular-ratio-left-certificate}
 R_I<-0.064086,\qquad U>0.035346,\qquad V/x>0.034494.
\end{equation}

For the right endpoint put $y=1/2-x$ and $C=\cos(\pi y)$.  Then
\begin{align*}
 D&=4\log2+2\sum_{n\ge1}(2^{2n+1}-1)\zeta(2n+1)y^{2n},\\
 A&=\sum_{n\ge0}(2n+1)2^{2n+3}\beta_D(2n+2)y^{2n}.
\end{align*}
If $\mathcal A=(CD)_y$ and $\mathcal B=(C^2A)_y$, the endpoint zeros
factor exactly as
\[
 -R_I/y=\frac{\mathcal B/y}{8\pi^3C},\quad
 U/y=\frac{4\pi\mathcal A/y-\mathcal B/y}{8\pi^3C},\quad
 V/y=\frac{\mathcal B/y-2\pi\mathcal A/y}{4\pi^3C}.
\]
With nine retained orders and $q=4y^2$, the positive tails for
$D,D_y/y,A,A_y/y$ are bounded respectively by
\[
 \frac{121}{25}\frac{q^{10}}{1-q},\quad
 \frac{968}{25}q^9\frac{10-9q}{(1-q)^2},\quad
 \frac{8(21)q^{10}}{(1-q)^2},\quad
 \frac{128(10)(21)q^9}{(1-q)^3}.
\]
Directed rational evaluation on $0\le y\le1/32$ gives
\begin{equation}\label{JC:eq:regular-ratio-right-certificate}
 -R_I/y>0.167528,\qquad U/y>0.105954,\qquad V/y>0.001213.
\end{equation}

It remains only the compact interval $[1/8,15/32]$.  The executable
certificate
\texttt{tests/check\_jc\_ar\_ratio\_symbol.py} partitions it into the
11,264 exact rational cells
\[
 [i/32768,(i+1)/32768],\qquad4096\le i\le15359.
\]
On each cell it evaluates twelve terms of the positive $D,D'$ series and
the alternating $F_2,F_3$ series.  Their omitted tails are bounded by
\[
 \frac{4x^2}{3\cdot12^2},\qquad
 \frac{32x}{9\cdot12^2},\qquad (12+z)^{-p},
\]
respectively.  The rational enclosure
\[
 \frac{103993}{33102}<\pi<\frac{104348}{33215}
\]
comes from Machin's identity; sine and cosine use alternating Taylor
series with their signed next terms.  Every arithmetic operation is
rounded outwards at 64 decimal digits.  In particular, sign reversal is
exact and integer powers use only directed interval multiplication.  The
finite minima are
\begin{equation}\label{JC:eq:regular-ratio-middle-certificate}
 -R_I>0.0056578993130883869,\quad
 U>0.0038837669640606282,\quad
 V>0.0009703006655867654.
\end{equation}
Equations~\eqref{JC:eq:regular-ratio-left-certificate}--
\eqref{JC:eq:regular-ratio-middle-certificate} prove
\eqref{JC:eq:regular-completed-covector-ratio}.  The script was independently
audited against exact rational arithmetic, including a hostile seven-digit
ambient decimal context; the same 64-digit enclosures result.

Finally Lemma~\ref{JC:lem:regular-polygonal-mixed-finite-part} gives
$R_{b,{\rm poly}}=J-R_I$.  Hence
$0<t:=R_{b,{\rm poly}}/R_I<1$.  For $0\le\mu\le1/4$,
$|t-2\mu|\le1-2\mu$.  Therefore
\[
 \Delta_\mu\ge(1-2\mu)R_I^2\{P-2Q_{\rm mat}\}>0
\]
by Proposition~\ref{JC:prop:regular-support-double-disk-gap}.  This proves
\eqref{JC:eq:regular-completed-covector-determinant}.
\end{proof}

\begin{lemma}[Uniform polygonal mixed bridge]
\label{JC:lem:regular-mixed-uniform-bridge}
Let $\beta=2/N$, $a=\pi\beta$, $x=k/N$, and
$2\le k\le N/2$.  The constant-flux harmonic-plus-area mixed coefficient
$p^{fs,0}_{\beta}$ satisfies
\begin{equation}\label{JC:eq:regular-mixed-CF2P}
 \left|-e^{i\pi x}\frac{p^{fs,0}_{\beta}}{d_\circ^2a}
       -L_{\rm poly}(x)\right|
 \le C\left(\beta+\frac{\beta^2}{x}\right).
\end{equation}
The actual reduced Helmholtz-minus-harmonic mixed modal row satisfies
\begin{equation}\label{JC:eq:regular-mixed-RR-ABS}
 \boxed{|E^{fs}_{N,k}|\le C_pa^2.}
\end{equation}
After replacing the constant flux by the true polygonal flux, the complete
mixed and diagonal ledgers have the following consequence.  If $k_N\to
\infty$, their Schur quotient differs from the limiting proxy determined
by~\eqref{JC:eq:regular-completed-covector-determinant} by $o(a^2)$; if
$k_N$ is bounded, the entire mixed Schur quotient is $o(a^2)$.
\end{lemma}

\begin{proof}
For~\eqref{JC:eq:regular-mixed-CF2P}, keep the $m=0$ term of
$M_\beta(x)$ exact.  Reflection gives $S_\beta(0)=0$, so that term is
$x^2O_{C^2}(1)$.  On $|m|\ge1$, subtract the $t^\beta$ endpoint models
from $I_{\beta,m}$ and the $t^\beta$ and linear endpoint models from
$S_{\beta,m}$.  Two integrations by parts give, for
$0\le\ell\le1$ and $0\le j\le2$,
\begin{equation}\label{JC:eq:regular-mixed-joint-remainder}
 |\partial_x^\ell\partial_\beta^jR_m(\beta,x)|
 \le C\frac{1+\log^2(2+|m|)}{(1+|m|)^2}.
\end{equation}
The removed $|m+x|^{-1-2\beta}$ families are first summed as paired
Hurwitz zeta functions.  After subtracting their common pole and finite
part, their Laurent remainders are jointly $C^{1,2}$ for
$0\le\beta\le x\le1/2$.  Reflection makes the first $\beta$ coefficient
vanish at $x=0$.  Two-variable Taylor's formula therefore gives
\[
 \left|M_\beta(x)-\frac{i\sin(2\pi x)}{4\pi^2\beta}
 -i\left\{\frac{\sin(2\pi x)}{4\pi^2}D(x)
 -\frac{\sin^2(\pi x)}{2\pi^3}T(x)\right\}\right|
 \le C(\beta x+\beta^2).
\]
The exact area pole cancellation proves~\eqref{JC:eq:regular-mixed-CF2P}.

We prove~\eqref{JC:eq:regular-mixed-RR-ABS}, which is the reduced row omitted
in the first bridge attempt.  Pull the regular polygon to the disk and
write $V_N=\lambda_N|F_N'|^2$.  In a nontrivial character sector let
$U_Nf$ solve $(-\Delta-V_N)U_Nf=0$ with boundary value $f$, and set
\[
 \mathcal K_N(f,g)=\langle(\Lambda_{V_N}-\Lambda_0)f,g\rangle.
\]
The ground state is invariant, so every retained sector is orthogonal to
it.  The first/second spectral gap is uniform; hence, for fixed $p>4$ and
$q=2p/(p-1)$,
\begin{equation}\label{JC:eq:regular-sector-solution-bound}
 \|U_Nf\|_{L^q}+\|Hf\|_{L^q}\le C_p\|f\|_\infty.
\end{equation}
The retained SC weight estimate gives
$\|V_N-\lambda_\circ\|_{L^p}\le C_p\beta^{1+1/p}$.
Alessandrini's identity and~\eqref{JC:eq:regular-sector-solution-bound} yield
\begin{equation}\label{JC:eq:regular-sector-potential-perturbation}
 |\mathcal K_N(f,g)-\mathcal K_\circ(f,g)|
 \le C_p\beta^{1+1/p}\|f\|_\infty\|g\|_\infty.
\end{equation}

For the unnormalized constant-flux traces,
\[
 f^s=d_\circ\sigma_{N,k}w_\beta,\qquad
 f^f=-d_\circ a(\rho-1)^{-1}\sigma_{N,k}w_\beta J_\beta.
\]
Their Fourier coefficients at $n=k+mN$ are, apart from the displayed
fixed factors, $I_{\beta,m}$ and $S_{\beta,m}$.  Uniformly for
$0<\beta\le x\le1/2$,
\[
 |I_{\beta,0}|\le C,\quad |S_{\beta,0}|\le Cx,\qquad
 |I_{\beta,m}|\le\frac{Cx}{1+|m|},\quad
 |S_{\beta,m}|\le\frac C{1+|m|}\quad(m\ne0).
\]
For the constant disk potential the reduced-minus-harmonic multiplier is
\[
 \varkappa_n=j_{0,1}\frac{J_{|n|}'(j_{0,1})}{J_{|n|}(j_{0,1})}-|n|
 =-j_{0,1}\frac{J_{|n|+1}(j_{0,1})}{J_{|n|}(j_{0,1})},
 \qquad |\varkappa_n|\le\frac C{1+|n|}.
\]
Here $|n|\ge2$, so there is no ground-state resonance.  It follows that
\[
 |\mathcal K_\circ(f^f,f^s)|
 \le C\frac ax\left\{\frac{x}{1+Nx}
 +\sum_{m\ne0}\frac{x}{N(1+|m|)^3}\right\}\le Ca.
\]
The true flux obeys
$\|q_N-d_\circ\|_\infty\le C_pa^{1+1/p}$; the Green identity
$\mathcal K_N(f,g)=-\int V_NU_Nf\,\overline{Hg}$ and
\eqref{JC:eq:regular-sector-solution-bound} show that replacing either slot
changes the preceding unnormalized form by $o(a)$.  Finally an
unnormalized character is $\sqrt N$ times an orthonormal one in each
mixed slot.  Division by $N=2\pi/a$ proves
\eqref{JC:eq:regular-mixed-RR-ABS}.

For completeness, define the positive side-mass normalization by
\[
 \bar q_N^2:=
 \frac{\int_0^1w_\beta q_N^2}{\int_0^1w_\beta}.
\]
The regular-side stationarity identity
$\int_{S_i}p_N^2\,ds=(\lambda_N/\pi)L_i$, together with
$p_N=c_N^{-1}w_\beta q_N$ and
$ds=c_Na w_\beta^{-1}dt$, gives exactly
\[
 \bar q_N^2=\frac{\lambda_N}{\pi}c_N^2
 \frac{\int_0^1w_\beta^{-1}}{\int_0^1w_\beta}
 =\frac{\lambda_N}{\pi}\{1+O(\beta^3)\}.
\]
Here $c_N=1+O(\beta^3)$ and the expansions of
$\int w_\beta^{\pm1}$ have the same quadratic coefficient because
$\int_0^1\log|2\sin\pi t|\,dt=0$.  Moreover $q_N$ is cell-periodic and
$\bar q_N$ lies between its minimum and maximum on a cell.  Rescaling a
cell of length $O(\beta)$ to $[0,1]$ in
\eqref{JC:eq:corner-flux-holder-rate} therefore proves
\[
 \|q_N-\bar q_N\|_{C^\alpha_{\rm cell}}
 \le C_p\beta^{2-1/p},\quad \alpha=1-2/p>1/2.
\]
The endpoint coefficient bounds used above also give
\[
 \sum_m|m+x||I_{\beta,m}|^2\le C\frac{x^2}{\beta},\qquad
 \sum_m|m+x||S_{\beta,m}|^2\le\frac C\beta.
\]
Multiplication by a $C^\alpha_{\rm cell}$ function is bounded on this
Bloch $H^{1/2}$ space; polarization and
$a|\rho-1|^{-1}\le C\beta/x$ therefore show that replacing the constant
flux by its cell oscillation changes the harmonic mixed row by at most
$C_p\beta^{2-1/p}$.  The disk aliases give
uniformly
\begin{equation}\label{JC:eq:regular-mixed-alias-ledger}
 cx\le g_{N,k}\le Cx,\qquad |Z_{N,k}|\le Ca,\qquad
 |g_{N,k}Z_{N,k}|\le Cax.
\end{equation}
On every growing-character sequence,~\eqref{JC:eq:regular-mixed-CF2P},
\eqref{JC:eq:regular-mixed-RR-ABS}, the support asymptotic, and the preceding
true-flux estimate give, for the limiting proxies denoted by bars,
\begin{align*}
 |\bar A|&\le Cax,& r,\bar r&\ge cx,\\
 |A-\bar A|&\le o(ax)+C_pa^{2-1/p}+Ca^2,&
 |r-\bar r|&=o(x).
\end{align*}
Elementary perturbation of $|A+\mu_sg\overline Z|^2/r$ then costs
\[
 o(a^2x)+C_pa^{3-1/p}+C_pa^{4-2/p}/x+O(a^3+a^4/x)=o(a^2),
\]
because $x\ge\beta\asymp a$ and $p>4$.  If $k$ stays bounded, the
fixed-slot limit and~\eqref{JC:eq:regular-mixed-alias-ledger} give
$r\ge c_ka$, while the numerator is
$O_k(a^2)+O_p(a^{2-1/p})$; its quotient is
$O_p(a^{3-2/p})=o(a^2)$.  This proves the last assertion.
\end{proof}

\begin{proposition}[Closure of the regular completed-action Schur block]
\label{JC:prop:regular-action-Schur-closed}
For all sufficiently large $N$, Contract~\ref{JC:con:JC-AR} holds with
\begin{equation}\label{JC:eq:regular-action-closed-constants}
 \boxed{\mu_s=\frac1{40},\qquad \mu_f=\kappa.}
\end{equation}
In particular, the all-character matrix
\eqref{JC:eq:regular-action-mode-matrix} is positive semidefinite.
\end{proposition}

\begin{proof}
The connection between the certified determinant and the Schur row is
division-free.  The limiting ledgers give
\[
 \bar g=2d_\circ^2Q_{\rm mat},\qquad
 \bar g\,\overline{\bar Z}
 =2d_\circ^2a e^{-i\pi x}\sin(\pi x)R_I,
\]
\[
 \bar A=-d_\circ^2a e^{-i\pi x}\sin(\pi x)R_{b,{\rm poly}},
 \qquad
 \bar r=d_\circ^2D_{\rm poly}(x;\mu_s).
\]
Consequently the limiting part of
\eqref{JC:eq:regular-action-reduced-Schur} after deletion of the fan reserve
is exactly
\begin{equation}\label{JC:eq:regular-proxy-Schur-determinant}
 \left(\frac12-\mu_s\right)\bar g|\bar Z|^2
 -\frac{|\bar A+\mu_s\bar g\,\overline{\bar Z}|^2}{\bar r}
 =\frac{d_\circ^2a^2\sin^2(\pi x)}
 {Q_{\rm mat}D_{\rm poly}(x;\mu_s)}\,\Delta_{\mu_s}(x).
\end{equation}
Both denominators are positive by
Proposition~\ref{JC:prop:regular-support-double-disk-gap}.

The support theorem gives
$p^{ss}_{N,k}\ge(1/2+1/20)g_{N,k}$, hence the denominator in
\eqref{JC:eq:regular-action-reduced-Schur}, with $\mu_s=1/40$, satisfies
$r_{N,k}\ge g_{N,k}/40>0$.  Suppose the claimed Schur row failed along a
sequence of characters.  After conjugation take $2\le k_N\le N/2$.
If $k_N$ is bounded, Lemma~\ref{JC:lem:regular-mixed-uniform-bridge} makes
the negative quotient and the corrector term $o(a^2)$; the untouched fan
term is $(3\kappa-\mu_f)a^2=2\kappa a^2$, a contradiction.

Otherwise pass to a subsequence with $k_N\to\infty$.  The same lemma
compares the quotient, including its denominator and corrector, with the
limiting proxy up to $o(a^2)$.  Proposition~\ref{JC:prop:regular-completed-covector-ratio}
makes that proxy nonnegative for $\mu_s=1/40$, including sequences with
$x_N\downarrow0$; at $x_N\to1/2$ its possible zero is again paid by the
same $2\kappa a^2$ fan reserve.  Thus the actual row is
$2\kappa a^2-o(a^2)>0$, another contradiction.  Every character sequence
has one of these two alternatives, proving the assertion uniformly.
\end{proof}

\begin{remark}[Why stationary O--H does not prove JC--AR]
On an even regular fan take the alternating support vector
$z_i=(-1)^i$.  The exact polygonal reconstruction in the CT node gives
\[
 \|V_z\|_{H^{1/2}(\mathbb T;\mathbb R^2)}^2\asymp Na^{-2},
 \qquad
 \Gjp_{\alpha_*}(z)\asymp N.
\]
Thus the full-vector norm in stationary O--H is larger than the stable
normal-sine metric by the sharp factor $a^{-2}$.  Its error can therefore
be as large as
\[
 C\omega_{OH}(a)a^{-2}\Gjp_{\alpha_*}(z).
\]
The proved statement $\omega_{OH}(a)\to0$ supplies neither the rate
$O(a^2)$ nor a sufficiently small leading constant.  Moreover its radial
disk trace is not literally the normal-sine completed trace, and it
controls the full Hessian rather than the base-subtracted action root
\eqref{JC:eq:regular-action-root}.  Hence O--H cannot be used to assert either
\eqref{JC:eq:regular-action-support-row} or
\eqref{JC:eq:regular-action-mode-matrix}.
\end{remark}

\begin{contract}[JC--SP: remaining super-near completed-action modulus and output]
\label{JC:con:JC-SP}
Suppose the regular completed-action contract JC--AR holds.  Let
\begin{equation}\label{JC:eq:super-near-energy}
 \mathcal N_N(d,w):=
 a^2\sum_i d_i^2+\Gjp_{\alpha_*}(w),
 \qquad
 \mathcal N_N(d,w)\le Ma^5.
\end{equation}
For the exact completed chart of JC--AR also put
\begin{equation}\label{JC:eq:finite-completed-energy}
 \mathfrak E_N(d,w):=
 \Jfour(\alpha_*+d)+
 \Gjp_{\alpha_*+d}(J_{\alpha_*+d}w).
\end{equation}
The finite-coordinate row needed here is closed by
Proposition~\ref{JC:prop:super-near-bridge-active}: uniformly for fixed $M$,
\begin{equation}\label{JC:eq:super-near-energy-bridge}
 c_M\mathcal N_N(d,w)\le\mathfrak E_N(d,w)
 \le C_M\mathcal N_N(d,w)
\end{equation}
whenever either side is at most a fixed multiple of $a^5$.  This row
includes the moving metric, the $J_\alpha$ transport, the quotient gauges,
and their uniform constants.  The exact finite corrector cancels here
because $z-z_\alpha^D=J_\alpha w$; its nonlinear variation instead belongs
to the completed chart and Hessian modulus below.

The exact fourth-defect identity gives the fan part of this bridge and the
radius
\begin{equation}\label{JC:eq:super-near-radius}
 \frac{\|d\|_{\ell^\infty}}a+
 a^{-2}\Gjp_{\alpha_*}(w)^{1/2}
 \le C_Ma^{1/2}.
\end{equation}
Indeed,
\[
 \Jfour(\alpha_*+d)
 =\sum_i d_i^2(6a^2+4ad_i+d_i^2)
 \ge3a^2\sum_i d_i^2
\]
for $d_i\ge-a$ and $\sum_i d_i=0$; the last coefficient is sharp at
$d_i=-a$.  The support part of
\eqref{JC:eq:super-near-energy-bridge} is the moving-low-mode transversality
content of Proposition~\ref{JC:prop:super-near-bridge-active}; it is not
inferred from a bare interchange of the two meshes.
Moreover, applying the closed corrector estimate to
$\alpha_*+td$, dividing by $t^2$, and letting $t\to0$ gives
\begin{equation}\label{JC:eq:corrector-derivative-bound}
 \Gjp_{\alpha_*}(Z_Nd)\le Ca^2\sum_i d_i^2.
\end{equation}
Thus the initial physical support velocity
$\zeta=w+Z_Nd$ also satisfies
$\Gjp_{\alpha_*}(\zeta)\le C_Ma^5$.

Use the affine path in the exact completed coordinates
\begin{equation}\label{JC:eq:completed-coordinate-path}
 q_t=(td,tw),\qquad
 \alpha_t=\alpha_*+td,\qquad
 z_t=J_{\alpha_t}(\bar z_{\alpha_t}^D+tw),
 \qquad0\le t\le1.
\end{equation}
Proposition~\ref{JC:prop:super-near-bridge-active} proves that this path
remains in the active convex chart and that all chart maps above are
defined uniformly.  Its Hessian is taken after
composition with the exact scalar \(\mathbf A_N\); hence the moving fan,
support, corrector, gauge, and coordinate-acceleration terms are all part
of the Hessian.  The earlier proposed affine chord of physical vertices
has zero vertex acceleration, but its initial chart tangent need not equal
the endpoint difference $(d,w)$.  It is therefore not used here without a
separate tangent-energy comparison.

For later estimates the exact completed scalar has the following
corrector-free normal form.  Put $c_\alpha=\bar z_\alpha^D$, so that
$z_\alpha^D=J_\alpha c_\alpha$.  Then
\begin{equation}\label{JC:eq:completed-action-corrector-normal-form}
 \boxed{
 \begin{aligned}
 \mathbf A_N(d,w)={}&
 \Phi_\alpha\!\left(h_\alpha\mathbf1
       +J_\alpha(c_\alpha+w)\right)
 -\Phi_\alpha(h_\alpha\mathbf1)\\
 &+\Gjp_\alpha(J_\alpha c_\alpha)
 -\frac12\Gjp_\alpha(J_\alpha w)
 +\frac\kappa2\Jfour(\alpha),
 \qquad \alpha=\alpha_*+d.
 \end{aligned}}
\end{equation}
Indeed $p_\alpha[\xi]=-2g_\alpha(z_\alpha^D,\xi)$, so the forcing,
metric, and completed square combine exactly into
$\Gjp_\alpha(z_\alpha^D)-\frac12\Gjp_\alpha(J_\alpha w)$.
Thus~\eqref{JC:eq:completed-action-corrector-normal-form} retains the moving
fan, support chart, quotient gauge, corrector, base-value subtraction,
moving metric, $1/2$ normalization, and fourth-defect row in one scalar.

The only remaining analytic row of JC--SP is a relative modulus for the
Hessian of the completed action itself.  For every fixed $M$ there must be
$\varepsilon_M(a)\to0$ such that
\begin{equation}\label{JC:eq:super-near-path-modulus}
 \boxed{
 \left|D^2\mathbf A_N(q_t)[(d,w),(d,w)]
       -D^2\mathbf A_N(0)[(d,w),(d,w)]\right|
 \le\varepsilon_M(a)\mathcal N_N(d,w)
 \quad(0\le t\le1).}
\end{equation}
A bound $\varepsilon_M(a)=O_M(a^{1/2})$ would be sufficient.  The finite
rank scale/translation gauge change, moving corrector, and base-value
subtraction are part of the left side and must not be dropped.

One stronger uniform operator formulation, sufficient for
\eqref{JC:eq:super-near-path-modulus}, is the following.  Let $R_N$ be the Riesz operator of
$\mathcal N_N$ on the fixed regular quotient and put
$H_N(q)=D^2\mathbf A_N(q)$.  With
\[
 r_N(q)=\frac{\|d(q)\|_{\ell^\infty}}a
 +a^{-2}\Gjp_{\alpha_*}(w(q))^{1/2},
\]
the natural operator statement is
\begin{equation}\label{JC:eq:super-near-operator-modulus}
 \boxed{
 \Omega_N(r):=
 \sup_{r_N(q)\le r}
 \left\|R_N^{-1/2}
 \{H_N(q)-H_N(0)\}R_N^{-1/2}\right\|_{\rm op}
 \longrightarrow0}
\end{equation}
as $r\downarrow0$, uniformly in $N$.  A Lipschitz estimate
$\Omega_N(r)\le Cr$, together with
$r_N(q_t)\le C_Ma^{1/2}$ from the chart/energy bridge, would imply the
stated $O_M(a^{1/2})$ rate.  This operator statement is stronger than the
diagonal path estimate and is recorded only as a sufficient formulation.
In block
form it must simultaneously contain
\begin{align*}
 |\Delta H_{ff}[d_1,d_2]|
 &\le Cr\,a^2\|d_1\|_2\|d_2\|_2,\\
 |\Delta H_{fs}[d,u]|
 &\le Cr\,a\|d\|_2\Gjp_{\alpha_*}(u)^{1/2},\\
 |\Delta H_{ss}[u_1,u_2]|
 &\le Cr\,
 \Gjp_{\alpha_*}(u_1)^{1/2}
 \Gjp_{\alpha_*}(u_2)^{1/2}.
\end{align*}
The nonzero period-three mixed row is therefore retained rather than
estimated as a separate remainder.

The absolute analytic CT estimate is not
\eqref{JC:eq:super-near-path-modulus}.  On the alternating support mode it
again loses $a^{-2}$, so a generic full-vector modulus gives only
$a^{-2}\omega(C_Ma^{1/2})\mathcal N_N$.  Analytic Lipschitz continuity or
$\omega(r)\to0$ does not make this quantity small.  The remaining part of
JC--SP therefore is an open natural-operator statement.

There is a scalar reduction which is weaker than proving the full operator
modulus.  For $q=(d,w)$ define
\[
 f_q(\zeta):=\mathbf A_N(\zeta d,\zeta w),\qquad
 r_N(q):=\frac{\|d\|_{\ell^\infty}}a
 +a^{-2}\Gjp_{\alpha_*}(w)^{1/2}.
\]
Suppose there are absolute $c_0,C_0>0$ such that every exact completed
coordinate in~\eqref{JC:eq:completed-action-corrector-normal-form} extends
holomorphically to $|\zeta|r_N(q)\le c_0$ and
\begin{equation}\label{JC:eq:super-near-complex-ray-bound}
 \boxed{|f_q''(\zeta)|\le C_0\mathcal N_N(d,w)}
\end{equation}
there.  If $r_N(q)\le c_0/2$, Cauchy's estimate on the disk of radius
$c_0/(2r_N(q))$ about $t\in[0,1]$ gives
\[
 |f_q'''(t)|\le\frac{2C_0}{c_0}r_N(q)\mathcal N_N(d,w).
\]
Since the completed-coordinate path is affine,
$f_q''(t)=D^2\mathbf A_N(q_t)[q,q]$.  Integration therefore proves
\begin{equation}\label{JC:eq:super-near-Cauchy-output}
 \boxed{
 |D^2\mathbf A_N(q_t)[q,q]-D^2\mathbf A_N(0)[q,q]|
 \le\frac{2C_0}{c_0}r_N(q)\mathcal N_N(d,w).}
\end{equation}
The bridge gives $r_N(q)=O_M(a^{1/2})$, so
\eqref{JC:eq:super-near-complex-ray-bound} would close the required path
modulus with $\varepsilon_M(a)=O_M(a^{1/2})$.

The precise analytic blocker behind
\eqref{JC:eq:super-near-complex-ray-bound} is a nonstationary quasiuniform
resummed completed-jet identity, abbreviated QRCI.  It must be obtained by
differentiating the exact scalar
\eqref{JC:eq:completed-action-corrector-normal-form}; the stage-130
polygon-minus-disk longitudinal ledger is retracted and is not such a
derivative.  Terminal two-ray reduced-DtN, its full-cell complement, the
negative cotangent/area pole, and every endpoint-localized derivative of
$h_\alpha,J_\alpha,c_\alpha$, the moving metric, and the base subtraction
must be combined before absolute values are taken.  The existing
full-vector Hessian estimate loses $a^{-2}$ on alternating support data,
while the available mesh-free completed estimate is stationary and uses a
degenerate radial comparison form.  Neither proves QRCI.

\begin{remark}[The scalar radial-chart shortcut fails at the critical index]
\label{JC:rem:super-near-radial-Hhalf-obstruction}
One cannot obtain the missing bound by differentiating only the scalar
radial graph.  If the fixed-normal support numbers move as
$h_i(t)=h_i+tz_i$, then on side $i$
$R_t(\varphi)=(h_i+tz_i)/\cos(\varphi-\theta_i)$.  At a regular vertex
the two one-sided traces of the formal derivative differ by
\begin{equation}\label{JC:eq:super-near-radial-Hhalf-jump}
 J_i=\frac{z_{i+1}-z_i}{\cos(a/2)}.
\end{equation}
Every nonconstant cyclic character has a nonzero such jump.  A jump is
not in $H^{1/2}(\mathbb T)$, because its Slobodetskii integral contains
\[
 \frac{|J_i|^2}{4}
 \int_0^{\varepsilon/2}\!\int_0^{\varepsilon/2}
 \frac{du\,dv}{(u+v)^2}=+\infty.
\]
Thus $t\mapsto R_t$ is not differentiable as an $H^{1/2}$-valued curve,
so the scalar CT analytic-Hessian theorem cannot be invoked.  Freezing
the vertices by a continuous physical affine-side field removes the jump
but restores the already recorded $a^{-2}$ loss.  This leaves the
endpoint-resummed NQCH cancellation as a genuine analytic input.
\end{remark}

\begin{remark}[Exact radius resummation, retraction of DAR, and the live completed third jet]
\label{JC:rem:super-near-completed-third-jet}
The artificial endpoint radius is not itself an obstruction.  If a
homogeneous full-cell corner power has amplitude $A$, exponent $\beta$, and
the terminal window has relative radius $\vartheta$, then in either the
common or jump channel, with $U_\beta$ and $D_\beta$ denoting the exact
two-ray common and jump coefficients,
\[
 A^2\vartheta^{2\beta}K_\beta
 +A^2(1-\vartheta^{2\beta})K_\beta=A^2K_\beta,
 \qquad K_\beta\in\{U_\beta,D_\beta\}.
\]
Hence the apparent $\beta\log\beta$ loss created by stopping at
$\vartheta=\beta$ cancels exactly, including after differentiation along
the completed ray.

The former amplitude row DAR is retracted as an independent contract.  It
was written, with
\[
 K_i(t)=\frac{D_{\beta_i(t)}}\pi j_i(q;t)^2,
\]
as $\sum_i(A_i^2)'K_i+(\sum_iC_i^{\rm end})'$.  But the exact endpoint
jet is unchanged by the analytic reallocation
\[
 \widetilde A_i^2=A_i^2+\gamma_i,
 \qquad
 \widetilde C_i^{\rm end}=C_i^{\rm end}-\gamma_iK_i,
\]
whereas the displayed DAR expression changes by
$-\sum_i\gamma_iK_i'$.  Choosing a sufficiently small constant nonzero
$\gamma_i$ and a nonconstant $K_i$ makes this change nonzero while
preserving $\widetilde A_i^2>0$ locally whenever $A_i^2>0$.  Thus the old
row depends on how the full-cell amplitude is
separated from its endpoint descendants; it is not a gauge-invariant
statement about the scalar~\eqref{JC:eq:completed-action-corrector-normal-form}.

The replacement live row is the third derivative of that scalar itself.
For $q=(d,w)$ and $q_t=tq$, put
$f_q(t)=\mathbf A_N(q_t)$.  The required real estimate is
\begin{equation}\label{JC:eq:super-near-completed-third-jet}
 \boxed{
 |f_q'''(t)|
 =|D^3\mathbf A_N(q_t)[q,q,q]|
 \le C r_N(q)\mathcal N_N(q),\qquad0\le t\le1.}
\end{equation}
It is gauge invariant and includes the full-cell flux, both terminal
channels, their exact complements, every endpoint and nonendpoint
chain-rule descendant, the moving metric, the corrector, the quotient
transport, and the base subtraction.

For clarity, it is obtained directly from the normal form.  If
\[
 Y_t=h_{\alpha_t}\mathbf1+J_{\alpha_t}(c_{\alpha_t}+tw),\quad
 B_t=h_{\alpha_t}\mathbf1,\quad
 C_t=J_{\alpha_t}c_{\alpha_t},\quad
 W_t=tJ_{\alpha_t}w,
\]
and $X_t=(\alpha_t,Y_t)$, then for every scalar family $F$
\[
 \frac{d^3}{dt^3}F(X_t)
 =D^3F(X_t)[X_t',X_t',X_t']
  +3D^2F(X_t)[X_t',X_t'']+DF(X_t)[X_t'''].
\]
Apply this identity to the two $\Phi$ terms and the two $\Gjp$ terms in
\eqref{JC:eq:completed-action-corrector-normal-form}, with $Y_t,B_t,C_t,W_t$
respectively, and add the exact fourth-defect derivative.  In particular
\[
 \frac\kappa2\frac{d^3}{dt^3}\Jfour(\alpha_*+td)
 =12\kappa\left(a\sum_i d_i^3+t\sum_i d_i^4\right),
\]
whose absolute value is at most
$Cr_N(q)a^2\sum_i d_i^2$.  No stage-130 ledger enters this chain rule.

The complex-ray bound~\eqref{JC:eq:super-near-complex-ray-bound} implies
\eqref{JC:eq:super-near-completed-third-jet} by the Cauchy estimate already
displayed above.  Conversely,
\eqref{JC:eq:super-near-completed-third-jet} directly implies the real path
modulus~\eqref{JC:eq:super-near-path-modulus} by integration, but does not by
itself assert a holomorphic-tube bound.  Either route would suffice for
JC--SP.  Neither is proved here; renaming the missing estimate does not
close QRCI or the main theorem.
\end{remark}

Its required output is nevertheless exact.  Put
$\mu_0=\min\{\mu_f,\mu_s\}$.  Taylor's formula in the completed chart gives
\[
 \mathbf A_N(d,w)=
 \int_0^1(1-t)D^2\mathbf A_N(q_t)[(d,w),(d,w)]\,dt.
\]
Thus JC--AR and~\eqref{JC:eq:super-near-path-modulus} imply, for sufficiently
small $a$ depending only on $M$,
\begin{equation}\label{JC:eq:JC-SP-output}
 \boxed{
 \mathbf A_N(d,w)\ge
 \left(\mu_0-\frac12\varepsilon_M(a)\right)\mathcal N_N(d,w)
 \ge0.}
\end{equation}
Together with the energy bridge~\eqref{JC:eq:super-near-energy-bridge}, this
is precisely the super-near action statement
\eqref{JC:eq:super-near-action-bound}; it is an output of JC--SP, not a
fourth assumed row.
\end{contract}

\begin{proposition}[Canonical finite bridge and completed-path activity]
\label{JC:prop:super-near-bridge-active}
There is a canonical analytic choice of the scale/translation
trivialization $J_\alpha$ near the regular fan for which
\eqref{JC:eq:super-near-energy-bridge} holds, with constants independent of
$N$, whenever either side is at most $Ma^5$.  For this choice the whole
path~\eqref{JC:eq:completed-coordinate-path} remains in the active convex
chart, and $J_{\alpha_t}$, its inverse, and the exact corrector are defined
uniformly for $0\le t\le1$.
\end{proposition}

\begin{proof}
Write $\alpha_i=a+d_i$, $\sum_i d_i=0$, and
\(\varrho=\|d\|_{\ell^\infty}/a\).  Choose normal angles with
$\theta_0=0$ and define the material map, for $0\le s\le a$, by
\[
 F_d(ia+s)=\theta_i+\frac{\alpha_i}{a}s .
\]
Thus
\begin{equation}\label{JC:eq:material-map-small}
 \|F_d'-1\|_{L^\infty}\le\varrho,
 \qquad
 \|F_d-\operatorname{id}\|_{L^\infty}\le2\pi\varrho .
\end{equation}
Composition by $F_d$ and $F_d^{-1}$ is consequently bounded on
$H^{1/2}(\mathbb T)$ uniformly when $\varrho\le1/2$, as follows directly
from the Gagliardo double-integral norm.

Direct integration of the two sine basis functions gives
\begin{equation}\label{JC:eq:exact-scale-functional}
 L_\alpha z:=\sum_i\left\{
 \tan\frac{\alpha_{i-1}}2+\tan\frac{\alpha_i}2\right\}z_i
 =\int_{\mathbb T}\mathcal I_\alpha z .
\end{equation}
The side with normal $\theta_i$ in the tangential reference polygon has
length $h_\alpha\{\tan(\alpha_{i-1}/2)+\tan(\alpha_i/2)\}$, so
$L_\alpha z=0$ is exactly the affine scale row.  For $b\in\mathbb R^2$
put $\tau_\alpha(b)_i=b\cdot e_r(\theta_i)$.  Since first harmonics solve
$f''+f=0$ on each cell,
\begin{equation}\label{JC:eq:exact-translation-interpolation}
 \mathcal I_\alpha\tau_\alpha(b)=b\cdot e_r
\end{equation}
exactly.

Let $\mathsf X_*$ be the regular nodal complement of characters
$0,1,N-1$.  For $w\in\mathsf X_*$ define
\[
 c_\alpha(w)=\frac{L_\alpha w}{L_\alpha\mathbf1},
 \qquad
 f_\alpha(w)=\mathcal I_\alpha(w-c_\alpha(w)\mathbf1),
\]
let $\beta_\alpha(w)$ be the $L^2$ projection coefficient of $f_\alpha(w)$
onto $\{\cos\theta,\sin\theta\}$, and set
\begin{equation}\label{JC:eq:canonical-Jalpha}
 J_\alpha w=w-c_\alpha(w)\mathbf1-\tau_\alpha(\beta_\alpha(w)).
\end{equation}
Equations~\eqref{JC:eq:exact-scale-functional} and
\eqref{JC:eq:exact-translation-interpolation} show that $J_\alpha w$ lies
in the moving scale/translation quotient and that $J_{\alpha_*}=I$.

We prove uniform transversality.  On a material cell, with $r=s/a$, the
two pulled-back basis functions are
\[
 B_0(\alpha_i,r)=\frac{\sin((1-r)\alpha_i)}{\sin\alpha_i},
 \qquad
 B_1(\alpha_i,r)=\frac{\sin(r\alpha_i)}{\sin\alpha_i}.
\]
Taylor expansion, cellwise integration, and real interpolation give, for
$\|v\|_{\ell_a^2}^2=a\sum_i|v_i|^2$,
\begin{align}
 \|\mathcal I_\alpha v\circ F_d-\mathcal I_*v\|_{L^2}
 &\le C\varrho a^2\|v\|_{\ell_a^2},\label{JC:eq:material-L2-error}\\
 \|\mathcal I_\alpha v\circ F_d-\mathcal I_*v\|_{H^{1/2}}
 &\le C\varrho a^{3/2}\|v\|_{\ell_a^2}.
 \label{JC:eq:material-Hhalf-error}
\end{align}
Indeed the pointwise basis error is $O(\varrho a^2)$ and its physical
derivative is $O(\varrho a)$.

Cyclic equivariance implies that the continuous Fourier support of the
regular sine interpolant of a nodal character $k$ lies in
$k+N\mathbb Z$.  Hence $\mathcal I_*w$ is exactly orthogonal to
$1,\cos\theta,\sin\theta$, not merely asymptotically so.  The regular
mass matrix and the disk multiplier give
\begin{equation}\label{JC:eq:regular-mass-G-control}
 \|w\|_{\ell_a^2}\le C\Gjp_{\alpha_*}(w)^{1/2}.
\end{equation}
Since
$|\tan(\alpha_i/2)-\tan(a/2)|\le C\varrho a$ and
$L_\alpha\mathbf1\asymp1$, weighted Cauchy--Schwarz yields
\begin{equation}\label{JC:eq:scale-coefficient-control}
 |c_\alpha(w)|\le C\varrho\Gjp_{\alpha_*}(w)^{1/2}.
\end{equation}
Changing variables $\theta=F_d(s)$ in the first-harmonic coefficients,
subtracting the exactly zero regular coefficient, and using
\eqref{JC:eq:material-map-small}--\eqref{JC:eq:scale-coefficient-control} gives
\begin{equation}\label{JC:eq:translation-coefficient-control}
 |\beta_\alpha(w)|\le C\varrho\Gjp_{\alpha_*}(w)^{1/2}.
\end{equation}
No inverse estimate enters this step: the multiplier
$e_r(F_d)F_d'-e_r$ is $O(\varrho)$ in $L^\infty$.

Put $T_\alpha w=\mathcal I_\alpha J_\alpha w$.  The functions
$(\mathcal I_\alpha\mathbf1)\circ F_d$ and $e_r\circ F_d$ have uniformly
bounded $H^{1/2}$ norms.  Therefore
\eqref{JC:eq:material-Hhalf-error},
\eqref{JC:eq:scale-coefficient-control}, and
\eqref{JC:eq:translation-coefficient-control} imply
\begin{equation}\label{JC:eq:moving-low-mode-transversality}
 \|T_\alpha w\circ F_d-\mathcal I_*w\|_{H^{1/2}}
 \le C\varrho\Gjp_{\alpha_*}(w)^{1/2}.
\end{equation}
Now $T_\alpha w$ is exactly orthogonal to the three disk null modes, so
$q_D(T_\alpha w)\asymp\|T_\alpha w\|_{H^{1/2}}^2$.  Absorbing the right
side of~\eqref{JC:eq:moving-low-mode-transversality} for a universal
$\varrho\le\varrho_0$ proves
\begin{equation}\label{JC:eq:moving-metric-equivalence}
 c\Gjp_{\alpha_*}(w)\le
 \Gjp_\alpha(J_\alpha w)\le C\Gjp_{\alpha_*}(w).
\end{equation}
It also proves that $J_\alpha$ is injective.  The two quotient spaces
have dimension $N-3$, hence it is onto and has a uniform inverse.

The exact identity
\[
 \Jfour(\alpha_*+d)
 =\sum_i d_i^2(6a^2+4ad_i+d_i^2)
\]
has coefficient at least $3a^2$ for every $d_i\ge-a$.  If either side of
\eqref{JC:eq:super-near-energy-bridge} is at most $Ma^5$, it follows first
that $\varrho\le C_Ma^{1/2}$, hence $\varrho\le\varrho_0$ for large $N$.
The same coefficient is then bounded above by a universal multiple of
$a^2$.  Combining this fact with
\eqref{JC:eq:moving-metric-equivalence} proves
\eqref{JC:eq:super-near-energy-bridge}.  The corrector does not occur in
this comparison because $z-z_\alpha^D=J_\alpha w$ identically.

It remains to prove activity.  If $\Lambda_{\alpha,i}(z)$ denotes the
change of side $i$ under a support increment $z$, the exact support-line
formula is
\begin{equation}\label{JC:eq:exact-side-length-jump}
 \Lambda_{\alpha,i}(z)=
 \frac{z_{i+1}-z_i\cos\alpha_i}{\sin\alpha_i}
 +\frac{z_{i-1}-z_i\cos\alpha_{i-1}}{\sin\alpha_{i-1}}
 =f'(\theta_i+)-f'(\theta_i-),
 \quad f=\mathcal I_\alpha z.
\end{equation}
After subtracting an exact translation, which changes no side length,
the quasiuniform one-dimensional finite-element inverse estimates give
\[
 \|f\|_{H^1}\le Ca^{-1/2}\|f\|_{H^{1/2}},
 \qquad
 \|f'\|_{L^\infty}\le Ca^{-1/2}\|f'\|_{L^2}.
\]
Because $L_\alpha z=0$ removes the constant mode, the disk multiplier is
equivalent to the remaining $H^{1/2}$ norm.  Thus
\begin{equation}\label{JC:eq:side-length-inverse}
 \max_i|\Lambda_{\alpha,i}(z)|
 \le Ca^{-1}\Gjp_\alpha(z)^{1/2}.
\end{equation}

Along~\eqref{JC:eq:completed-coordinate-path}, exact linearity gives
$z_t=z_{\alpha_t}^D+J_{\alpha_t}(tw)$.  The closed corrector estimate,
the fourth-defect identity, and
\eqref{JC:eq:moving-metric-equivalence} give
\[
 \Gjp_{\alpha_t}(z_{\alpha_t}^D)\le C_Ma^5,
 \qquad
 \Gjp_{\alpha_t}(J_{\alpha_t}(tw))\le C_Ma^5.
\]
Consequently $\Gjp_{\alpha_t}(z_t)^{1/2}\le C_Ma^{5/2}$.
The reference side lengths are at least $ca$, whereas
\eqref{JC:eq:side-length-inverse} bounds every change by
$C_Ma^{3/2}=o(a)$.  Thus every side remains positive.  The companion
nodal and derivative inverse bounds give normal and tangential vertex
displacements $O_M(a^2)$ and $O_M(a^{3/2})$, respectively, so the whole
path remains in a fixed smaller radial $W^{1,\infty}$ chart.  This proves
all assertions.
\end{proof}

\begin{remark}[Scope of the regular/super-near split]
JC--AR and JC--SP isolate the regular completed-action root, the finite
energy bridge, and its diagonal propagation.  Their exact output is only
the super-near action bound~\eqref{JC:eq:JC-SP-output}.  The full no-ratio
Contract~\ref{JC:con:JC-PL} additionally requires JC--CL to absorb every
endpoint outside~\eqref{JC:eq:super-near-energy}.  No far implication is
assumed in JC--AR or JC--SP.
\end{remark}

\begin{contract}[JC--CL: coarse/far action localization]
\label{JC:con:JC-CL}
Write $d=d_\alpha=z_\alpha^D$ and $z=d+w$.  Exact disk completion first
removes every disk covector from the coarse problem:
\begin{equation}\label{JC:eq:coarse-action-normal-form}
 \boxed{
 \mathscr A_\alpha(d+w)=
 \Phi_\alpha(h_\alpha\mathbf1+d+w)
 -\Phi_\alpha(h_\alpha\mathbf1)
 -\frac12\Gjp_\alpha(w)+\Gjp_\alpha(d)
 +\frac\kappa2\Jfour(\alpha).}
\end{equation}
Indeed
$q_D(B_\alpha+K_\alpha z)-q_D(B_\alpha)
=\Gjp_\alpha(z-d)-\Gjp_\alpha(d)$, and
Proposition~\ref{JC:prop:JC-ALG} gives the formula.  This identity also fixes
the admissible path.  The active chart supplies only
$h_\alpha\mathbf1+[0,1](d+w)$; it does not imply that
$h_\alpha\mathbf1+d$ or
$h_\alpha\mathbf1+d+[0,1]w$ is active.  In the presence of an arbitrarily
short side, the energy bound $\Gjp_\alpha(d)\lesssim\Jfour(\alpha)$ has no
ratio-free inverse side-length consequence.  Thus a coarse proof may not
Taylor-expand about the disk corrector without a new activity lemma.

Put
\begin{equation}\label{JC:eq:coarse-action-energy}
 E_\alpha(z):=\Jfour(\alpha)
 +\Gjp_\alpha(z-d_\alpha).
\end{equation}
A sufficient coarse estimate is the existence of constants
$c_F,C_F>0$ such that, on the whole active no-ratio chart,
\begin{equation}\label{JC:eq:coarse-action-bound}
 \boxed{
 \mathscr A_\alpha(z)
 \ge c_FE_\alpha(z)-C_Fa^5.}
\end{equation}
For $E_\alpha(z)\le Ma^5$, set
$d=\alpha-\alpha_*$ and
$w=J_\alpha^{-1}(z-d_\alpha)$.  The energy bridge and output required in
JC--SP then give
\begin{equation}\label{JC:eq:super-near-action-bound}
 \mathscr A_\alpha(z)\ge0
 \quad\hbox{whenever}\quad
 E_\alpha(z)\le Ma^5
\end{equation}
for every fixed $M$.  Taking $M\ge C_F/c_F$, Contract~\ref{JC:con:JC-PL}
follows.  Indeed,
\eqref{JC:eq:super-near-action-bound} handles the first branch, while on
$E_\alpha(z)>Ma^5$ equation~\eqref{JC:eq:coarse-action-bound} gives
\[
 \mathscr A_\alpha(z)
 >(c_FM-C_F)a^5\ge0.
\]

No retained estimate proves~\eqref{JC:eq:coarse-action-bound}.  The existing
two-regime estimates for $F_4$, the normal-sine/Bessel/secant transfers,
and JE--V concern pure fan disk endpoints or scalar tails.  They do not
control arbitrary support amplitudes.  In particular, at
$\alpha=\alpha_*$ one may have
$\Jfour=0$ and $\Gjp_{\alpha_*}(z)\gg a^5$ while remaining in the radial
chart.  The fixed-fan Hessian identity JP--EH supplies no lower bound for
that branch.  The nonzero physical-curvature mixed row also cannot be
inserted as a vanishing error, by the CT mixed counterexample.

This fixed-fan obstruction has an exact scalar form.  Cyclic symmetry at
$\alpha_*$ gives $b_*=d_*=0$, so JC--ALG reduces to
\[
 \mathscr A_*(z)=
 \Phi_*(h_*\mathbf1+z)-\Phi_*(h_*\mathbf1)
 -\frac12\Gjp_*(z).
\]
Consequently the restriction of~\eqref{JC:eq:coarse-action-bound} to the
regular fan requires the following regular-fan Bregman coercivity row,
with $c_s=c_F$ and $C=C_F$:
\begin{equation}\label{JC:eq:regular-fan-Bregman-coercivity}
 \boxed{
 \Phi_*(h_*\mathbf1+z)-\Phi_*(h_*\mathbf1)
 \ge\left(\frac12+c_s\right)\Gjp_*(z)-Ca^5.}
\end{equation}
Equivalently, JP--EH rewrites its left side as the support-time integral
of the exact regular-fan reduced Helmholtz DtN, area, and lower rows.  The
two classical Minkowski concavities for $|P|^{1/2}$ and
$\lambda_1(P)^{-1/2}$ do not order their ratio and hence do not prove this
inequality.  Thus~\eqref{JC:eq:regular-fan-Bregman-coercivity}, followed by
the genuinely nonregular coarse localization, is the smallest presently
isolated content of JC--CL.

At the regular fan cyclic symmetry gives $\ell_*=d_*=0$.  Define
\begin{equation}\label{JC:eq:regular-fan-integrated-modulus}
 \mathcal L_*(z):=\int_0^1(1-t)
 \{H_{*,tz}-H_{*,0}\}[z,z],dt.
\end{equation}

\begin{proposition}[The fixed-reserve RF--IM target is false]
\label{JC:prop:RF-IM-period-three-counterexample}
Let \(N=3M\to\infty\), \(a=2\pi/N\), and, on the regular fan, put
\[
 d_i=\cos\frac{2\pi i}{3},\qquad
 s_*=\frac{4h_*\sin^2(a/2)}{1+2\cos a},\qquad
 z_N=(1-a^2)s_*d.
\]
The segment \(h_*\mathbf1+[0,1]z_N\) is strictly active and remains in
the radial small chart.  Moreover
\begin{align}
 \Gjp_*(z_N)
 &=\frac{13d_\circ^2\zeta(3)}{4\pi^2}a^3\{1+o(1)\},
 \label{JC:eq:RF-counterexample-G}\\
 \frac{\Phi_*(h_*\mathbf1+z_N)-\Phi_*(h_*\mathbf1)}
 {\Gjp_*(z_N)}
 &\longrightarrow\frac{10}{13},
 \label{JC:eq:RF-counterexample-value}\\
 \frac{\frac12H_{*,0}[z_N,z_N]}{\Gjp_*(z_N)}
 &\longrightarrow\frac{2\pi^2\log3}{13\zeta(3)}.
 \label{JC:eq:RF-counterexample-root}
\end{align}
Consequently
\begin{equation}\label{JC:eq:regular-fan-integrated-modulus-disproved}
 \frac{\mathcal L_*(z_N)}{\Gjp_*(z_N)}
 \longrightarrow
 \frac{10}{13}-\frac{2\pi^2\log3}{13\zeta(3)}
 =-0.6185019068\ldots<-\frac1{20}.
\end{equation}
Thus no \(N\)-independent \(C\) and \(\eta<1/20\) can make the former
RF--IM bound
\(\mathcal L_*(z)\ge-\eta\Gjp_*(z)-Ca^5\) valid on every regular active
radial path.  This does not disprove RF--BC: on the same family
\begin{equation}\label{JC:eq:RF-counterexample-action}
 \frac{\mathscr A_*(z_N)}{\Gjp_*(z_N)}
 \longrightarrow\frac7{26}>0.
\end{equation}
\end{proposition}

\begin{proof}
For a fixed normal fan,
\[
 L_i(h)=\frac{h_{i-1}+h_{i+1}-2h_i\cos a}{\sin a}.
\]
Since \(d_{3j}=1\) and \(d_{3j+1}=d_{3j+2}=-1/2\), exactly the
\(3j\)-sides first vanish at \(s=s_*\), while all other sides grow.
Taking \(s_N=(1-a^2)s_*\) makes the entire segment strictly active.
Also \(s_*=a^2/3+O(a^4)\), so the support displacement is \(O(a^2)\)
and its adjacent difference quotient is \(O(a)\).

At \(s=s_*\), deletion of the zero sides leaves equal support numbers and
alternating normal gaps \(a,2a\).  Denote the area-normalized disk sources
of the regular and alternating fans by \(B_{\rm reg}\) and \(B_{\rm alt}\).
In the Fourier normalization used here, their leading material profiles
satisfy
\[
 F_{\rm reg}''=1-\sum_{j\in\mathbb Z}\delta_{j+1/2},
\qquad
 F_{\rm alt}''=1-\delta_{1/2}-2\delta_2
\]
with periods \(1\) and \(3\), respectively.  Hence
\[
 \|B_{\rm reg}\|_{\dot H^{1/2}}^2
 =\frac{\zeta(3)}{4\pi^3}a^3+O(a^4).
\]
For \(\varkappa=2\pi/3\),
\[
 \widehat F_{{\rm alt},\ell}
 =\frac{e^{-i\varkappa\ell/2}+2e^{-2i\varkappa\ell}}
 {3(\varkappa\ell)^2}.
\]
The numerator has squared modulus \(9\) for even \(\ell\) and \(1\) for
odd \(\ell\).  The even and odd \(\ell^{-3}\) sums are
\(\zeta(3)/8\) and \(7\zeta(3)/8\), so
\[
 \|B_{\rm alt}\|_{\dot H^{1/2}}^2
 =\frac{3\zeta(3)}{2\pi^3}a^3+O(a^4).
\]
The source-profile Taylor remainder is \(O(a^{7/2})\) in \(H^{1/2}\).
Since the first nonconstant frequencies are of order \(a^{-1}\) and the
disk multiplier is \(\omega_n=|n|+1+O(|n|^{-1})\),
\begin{equation}\label{JC:eq:RF-counterexample-source}
 q_D(B_{\rm alt})-q_D(B_{\rm reg})
 =\frac{5d_\circ^2\zeta(3)}{2\pi^2}a^3+O(a^4).
\end{equation}

The closed disk-chart scalar expansion applies separately to both
quasiuniform fans:
\[
 \Phi_\alpha(h_\alpha\mathbf1)
 =\lambda_\circ+q_D(B_\alpha)+O(a^4).
\]
Indeed the complete cubic term is bounded by
\(\|B\|_{W^{1,\infty}}\|B\|_{H^{1/2}}^2=O(a^4)\), and the tame
fourth-and-higher remainder is \(O(a^5)\).  The support distance from
\(s_N\) to the collapsed endpoint is \(O(a^4)\); two-sided parallel-body
inclusions and Dirichlet monotonicity change \(\Phi\) by \(O(a^4)\).
Thus~\eqref{JC:eq:RF-counterexample-source} gives the numerator in
\eqref{JC:eq:RF-counterexample-value}.

For the regular normal-sine metric,
\[
 Q(x)=\frac{2\sin^4(\pi x)}{\pi^3}
 \{\zeta(3,x)+\zeta(3,1-x)\}.
\]
The real character \(d\) has half the energy of the unnormalized complex
character.  Since
\[
 \zeta(3,1/3)+\zeta(3,2/3)=26\zeta(3),\qquad
 Q(1/3)=\frac{117\zeta(3)}{4\pi^3},
\]
and \(s_N^2=a^4/9+O(a^6)\), the factor
\((N/2)d_\circ^2Q(1/3)\) proves~\eqref{JC:eq:RF-counterexample-G}, and
then~\eqref{JC:eq:RF-counterexample-value}.

Finally the true polygonal support symbol gives
\[
 \frac{\frac12H_{*,0}[z_N,z_N]}{\Gjp_*(z_N)}
 \longrightarrow\frac{P(1/3)}{Q(1/3)}.
\]
The digamma third-point formula yields
\[
 P(1/3)=\frac{9\log3}{2\pi},\qquad
 \frac{P(1/3)}{Q(1/3)}
 =\frac{2\pi^2\log3}{13\zeta(3)},
\]
which proves~\eqref{JC:eq:RF-counterexample-root}.  The exact Taylor identity
\[
 \mathcal L_*(z_N)
 =\Phi_*(h_*\mathbf1+z_N)-\Phi_*(h_*\mathbf1)
 -\frac12H_{*,0}[z_N,z_N]
\]
now proves~\eqref{JC:eq:regular-fan-integrated-modulus-disproved}.  The strict
comparison with \(-1/20\) follows already from
\(\pi^2>9\), \(\log3>1\), and \(\zeta(3)<5/4\).  Since
\(\Gjp_*(z_N)\asymp a^3\), an \(Ca^5\) remainder cannot absorb it.
Subtracting \(1/2\) from the value ratio gives~\eqref{JC:eq:RF-counterexample-action}.
\end{proof}

\begin{contract}[RF--REL: relative preservation of the full root reserve]
\label{JC:con:RF-relative-root-reserve}
There exist \(\delta>0\) and \(C\), independent of \(N\), such that every
regular active radial path satisfies
\begin{equation}\label{JC:eq:RF-relative-root-reserve}
 \boxed{\mathcal L_*(z)\ge-\frac{1-\delta}{2}
 \{H_{*,0}[z,z]-\Gjp_*(z)\}-Ca^5.}
\end{equation}
\end{contract}

\begin{lemma}[RF--REL implies RF--BC]
\label{JC:lem:RF-relative-implies-Bregman}
Contract~\ref{JC:con:RF-relative-root-reserve} implies
\[
 \mathscr A_*(z)\ge\frac{\delta}{20}\Gjp_*(z)-Ca^5.
\]
\end{lemma}

\begin{proof}
Taylor's identity and the regular action give
\[
 \mathscr A_*(z)
 =\frac12\{H_{*,0}[z,z]-\Gjp_*(z)\}+\mathcal L_*(z).
\]
RF--REL and
\(H_{*,0}\ge(11/10)\Gjp_*\) imply
\[
 \mathscr A_*(z)
 \ge\frac{\delta}{2}\{H_{*,0}[z,z]-\Gjp_*(z)\}-Ca^5
 \ge\frac{\delta}{20}\Gjp_*(z)-Ca^5.
\]
\end{proof}

Thus the regular coarse row remains open, but its current sufficient
interface is RF--REL rather than the disproved fixed \(1/20\)-reserve
target.

\begin{remark}[Collapse-word audit and the exact JC--CL cut]
\label{JC:rem:RF-collapse-word-audit}
Every fixed periodic pattern of deleted regular-fan sides can be encoded
by a positive integer gap word $(m_1,\ldots,m_r)$.  At its collapse
endpoint the support profile on a gap of length $m$ is
$u_j=j(m-j)/2$, and the disk-source, metric, and root-reserve
coefficients reduce to finite Fourier/Hurwitz-zeta sums.  These formulas
recover exactly the period-three ratios $10/13$ and $7/26$ above.
The reproducible integer-word precheck
\texttt{diagnostics/explore\_regular\_collapse\_words.py} enumerates all
words of total period at most $16$ and several families through period
$4096$.  It finds no negative action and no vanishing relative reserve;
the smallest sampled ratios are approximately
\[
 \mathscr A_*/\Gjp_*=0.0912513,
 \qquad \mathscr A_*/\mathcal R_0=0.181689,
 \qquad \mathcal R_0=(H_{*,0}-\Gjp_*)/2.
\]
These numbers are not lower bounds on the full active microscopic class.
For $z_i=a^2u_i$, the exact leading envelope is
\[
 B_u(x)=\min_i\left\{u_i+\frac{(x-i)^2}{2}\right\},
 \quad s_i=i+\frac12+u_{i+1}-u_i,
 \quad g_i=s_i-s_{i-1}\ge0.
\]
At fixed microscopic period, the complete \emph{leading} active class is
the gap simplex
$g_i\ge0$, $\sum_i g_i=p$, and along support time
$g_i(t)=(1-t)+tg_i$.  If $S_p(t)$ is the resulting leading disk-source
energy, its exact pair ledger is
\[
 S_p''(t)=\frac1{\pi p}\sum_{i,j}(v_i-v_j)^2
 \log\left|2\sin\frac{\pi\{(i-j)+t(v_i-v_j)\}}p\right|,
 \qquad v_i=u_{i+1}-u_i.
\]
Here the diagonal terms are defined by \(0^2\log0:=0\).  If distinct
corner atoms collide at \(t=1\), this formula and the support-time Taylor
integral are understood as improper identities; the resulting logarithmic
singularity is integrable against \(1-t\).
Put
\[
 \widehat u_k=\frac1p\sum_{j=0}^{p-1}u_j e^{-2\pi i k j/p},
 \qquad
 \mathcal Q_p[u]=2\pi\sum_{k=1}^{p-1}
       |\widehat u_k|^2Q(k/p).
\]
Consequently the smallest leading regular target is the universal
monotone-transport inequality
\begin{equation}\label{JC:eq:RF-obstacle-BC}
 2\pi\{S_p(1)-S_p(0)\}
 \ge\left(\frac12+c_0\right)\mathcal Q_p[u]
\end{equation}
for every period and every point of that simplex.  Continuous-simplex
searches lower the sampled ratios to about $0.082293$ relative to
$\mathcal Q_p$ and $0.163711$ relative to the root reserve at $p=384$;
they find no negative value, but do not prove a uniform constant.
Equation~\eqref{JC:eq:RF-obstacle-BC}, its finite-$N$ secant/cubic/tail
transfer, and the nonregular original-path row all remain open.
This is a floating-point necessary-family precheck, not a proof of
RF--REL.  It also shows that the regular path loss is order one rather
than an $o(1)\Gjp_*$ error.  Even a proof of RF--REL would leave the
genuinely nonregular original-path inequality below as a second,
independent input.
\end{remark}

For a genuinely nonregular fan the exact original-path form is, with
$z=d+w$,
\begin{equation}\label{JC:eq:nonregular-original-path-action}
 \boxed{
 \begin{aligned}
 \mathscr A_\alpha(d+w)={}&\ell_\alpha[d+w]
 +\int_0^1(1-t)
 H_{\alpha,t(d+w)}[d+w,d+w],dt\\
 &-\frac12\Gjp_\alpha(w)+\Gjp_\alpha(d)
 +\frac\kappa2\Jfour(\alpha).
 \end{aligned}}
\end{equation}
The nonregular part of JC--CL is precisely a lower bound for this one
budget by $c_F\{\Jfour(\alpha)+\Gjp_\alpha(w)\}-C_Fa^5$, uniformly in
all gap ratios.  Separating its covector from its integrated Hessian would
recreate the unavailable JE--FM step.

For a stationary-minimizer bypass one may replace
\eqref{JC:eq:coarse-action-bound} by the weaker localization statement:
every stationary competitor $x$ with
$\Phi_N(x)\le\Phi_N(0)$ satisfies
\begin{equation}\label{JC:eq:stationary-coarse-localization}
 6a^2\sum_i d_i^2
 +\Gjp_{\alpha_*}(z-Z_Nd)\le Ma^5.
\end{equation}
Together with JC--AR and JC--SP, this would put the competitor in the
super-near action-positive branch and exclude a lower value through the
exact JP--JC/NR assembly.  The present global wrapper
gives only qualitative near-disk entry, not
\eqref{JC:eq:stationary-coarse-localization}.  Thus the stationary route has
three open analytic rows, not two.
\end{contract}

\begin{contract}[JC--PL: pathwise LCA payment]
\label{JC:con:JC-PL}
After decreasing the no-ratio fan and radial endpoint-chart thresholds so that
$\varepsilon_V(S)\le\kappa/2$, prove, for every active normalized support
vector $z\in\mathsf X_\alpha$ in that chart,
\begin{equation}\label{JC:eq:JC-PL}
 \boxed{\mathscr A_\alpha(z)\ge0.}
\end{equation}
The proof must use a disjoint, complete pair ledger.  The terminal two-ray
power and the negative cotangent/area jump row must be joined before any
estimate; every unused terminal annulus must remain in the nonterminal
ledger; same-cell, adjacent, and first-separation LCA pairs must be counted
exactly once.  The support-time integral must be taken before a side-ratio
supremum, using Lemma~\ref{JC:lem:support-time-log}.  The true ground-state
multiplier, reduced-energy, mixed-flux, area, and fixed-rank rows must
retain the same completed covector $b_\alpha$ of
\eqref{JC:eq:completed-covector}.  All constants must be independent of the
smallest angle and adjacent ratios.  In particular the contract includes
endpoints with an arbitrarily short positive side: smallness of an
affine-side parametrization relative to the tangential fan polygon is not
an assumption.  Its admissible chart is exactly the closed JP--NS output
\[
 h_\alpha\mathbf1+[0,1]z\ \hbox{active},\qquad
 \sup_{0\le t\le1}\|\rho_t\|_{W^{1,\infty}}
 +\Gjp_\alpha(z)^{1/2}\le C(S+\rho_0),
\]
not the false labelled affine-side chart of
Proposition JP--NS~\textup{(radial-not-affine)}.

Equivalently, by Proposition~\ref{JC:prop:JC-ALG}, prove
\begin{equation}\label{JC:eq:JC}
\boxed{
 \mathscr R_\alpha(z)
 \ge-\frac\kappa2\Jfour(\alpha)
      -\frac12\Gjp_\alpha(z-z_\alpha^D).}
\end{equation}
The stage-130 moving-disk longitudinal calculation omits material
chain-rule terms and is not an admissible proof of either display.
\end{contract}

\begin{theorem}[JC--PL closes the active-chart theorem]\label{JC:thm:JC-to-J0}
Assume Contract~\ref{JC:con:JC-PL}.  Then every normalized polygon in the active
near-disk chart satisfies
\begin{equation}\label{JC:eq:J0}
\boxed{
 \Phi_\alpha(h_\alpha\mathbf1+z)
 -\Phi_{\alpha_*}(h_*\mathbf1)
 \ge\kappa\Jfour(\alpha)
      +\frac12\Gjp_\alpha(z-z_\alpha^D).}
\end{equation}
Moreover
\begin{equation}\label{JC:eq:corrector-bound}
 \Gjp_\alpha(z_\alpha^D)
 \le\frac{C_p^2}{4}\Jfour(\alpha).
\end{equation}
\end{theorem}

\begin{proof}
Add and subtract the two disk endpoint values.  Equations
\eqref{JC:eq:joint-residual}, \eqref{JC:eq:disk-completion}, and \eqref{JC:eq:V}
give
\begin{align*}
 \Phi_\alpha(h_\alpha\mathbf1+z)
 -\Phi_{\alpha_*}(h_*\mathbf1)
 \ge{}&P_D^{\rm ns}(\alpha)-P_D^{\rm ns}(\alpha_*)
       +\Gjp_\alpha(z-z_\alpha^D)\\
 &-\varepsilon_V(S)\Jfour(\alpha)+\mathscr R_\alpha(z).
\end{align*}
Insert \eqref{JC:eq:D0}, \eqref{JC:eq:JC}, and
$\varepsilon_V(S)\le\kappa/2$.  The fan coefficient is
$2\kappa-\kappa/2-\kappa/2=\kappa$, and the support coefficient is
$1-1/2=1/2$.  This proves \eqref{JC:eq:J0}.  Equations \eqref{JC:eq:P} and
\eqref{JC:eq:corrector-size} prove \eqref{JC:eq:corrector-bound}.
\end{proof}

\begin{remark}[Precise mathematical cut]
Contract~\ref{JC:con:JC-PL} is not a restatement of either discarded separated
estimate.  At $\alpha=\alpha_*$ it asks only for the coarse one-sided
support bound encoded by \eqref{JC:eq:JC}; away from the regular fan it gives
the forcing access to the complete integrated Hessian surplus before a
Schur complement is taken.  The exact algebra and the support-time
logarithm cancellation are now proved.  Revision 24 further resolves the
former vague ``regular plus path'' description into one closed and two
open exact subcontracts.  The all-character regular completed-action Schur matrix
JC--AR is now closed.  The two remaining open contracts are the
completed-action Hessian modulus JC--SP and the coarse/far localization
JC--CL.  The former period-three completed-covector claim remains
retracted, and the fixed-reserve RF--IM target is disproved.  What is
still missing is the ratio-free
pathwise LCA payment for the true nonstationary ground-state multiplier and
all remaining lower rows.  No retained node implies either open analytic
estimate.
\end{remark}

\subsection{DAG disposition}

\begin{center}
\begin{longtable}{>{\raggedright\arraybackslash}p{.34\linewidth}
                  >{\raggedright\arraybackslash}p{.54\linewidth}}
\toprule
item & revision-24 disposition\\
\midrule
$D0_{\rm ns}$, JE--P, JE--V & closed inputs, equations
\eqref{JC:eq:D0}--\eqref{JC:eq:V};\\
JE--FM & unavailable; the old material falsification is retracted and the
row is absent from the active DAG;\\
JP--SC & retained only as a possible stronger auxiliary estimate; no
longer an active prerequisite;\\
JC--ALG and JC--TL & closed by Proposition~\ref{JC:prop:JC-ALG} and
Lemmas~\ref{JC:lem:support-time-log}--\ref{JC:lem:support-time-min};\\
pure-support JC--AR row & closed for all characters by
Lemma~\ref{JC:lem:fixed-character-support-limit} and
Proposition~\ref{JC:prop:regular-support-renormalized};\\
period-three completed-covector row & retracted by
Remark~\ref{JC:rem:period-three-completion-retracted}; its true polygonal
mixed finite part is now closed at every fixed Bloch ratio, as is the
strict support--double-disk gap; the ratio certificate and uniform
all-character bridge are closed by
Proposition~\ref{JC:prop:regular-action-Schur-closed};\\
RF--IM & disproved by
Proposition~\ref{JC:prop:RF-IM-period-three-counterexample}; RF--REL is the
replacement open sufficient interface for RF--BC;\\
super-near bridge/activity & closed by
Proposition~\ref{JC:prop:super-near-bridge-active};\\
terminal-radius pure-power resummation & closed by
Remark~\ref{JC:rem:super-near-completed-third-jet}; the old DAR row is
retracted as noninvariant, while the completed third-jet/QRCI row stays open;\\
JC--AR & closed all-character regular action Schur contract;\\
JC--SP, JC--CL & remaining natural Hessian-modulus and coarse/far
subcontracts inside JC--PL;\\
JC--PL & unique open conceptual node, Contract~\ref{JC:con:JC-PL};\\
JC--PL $\to$ JP--JC $\to$ J0 & closed algebraically by
Theorem~\ref{JC:thm:JC-to-J0}.\\
\bottomrule
\end{longtable}
\end{center}
\endgroup

\begingroup
\section{Node NR: joint endpoint assembly and J0 handoff}
\label{module:NR}
\noindent\textit{Audited source: \nolinkurl{NR_assembly.tex}.}
\subsection{Incoming joint comparison}

Use the notation of the JP and JC nodes.  In particular,
\begin{equation}\label{NR:eq:corrector}
 z_\alpha^D=-\frac12\mathcal R_\alpha^{-1}p_\alpha,
 \qquad
 \Gjp_\alpha(z_\alpha^D)
 \le\frac{C_p^2}{4}\Jfour(\alpha),
\end{equation}
and the exact disk completion is
\begin{equation}\label{NR:eq:disk-completion}
 q_D(B_\alpha+K_\alpha z)
 =P_D^{\rm ns}(\alpha)+\Gjp_\alpha(z-z_\alpha^D).
\end{equation}
The closed disk endpoint has margin $2\kappa\Jfour$, and closed value
matching has error at most $\varepsilon_V(S)\Jfour$.  JC--ALG and JC--TL
are closed.  The sole unavailable input is JC--PL; its exact algebraic
output is the fixed-fan scalar comparison JP--JC,
\begin{equation}\label{NR:eq:JC}
 \mathscr R_\alpha(z)
 \ge-\frac\kappa2\Jfour(\alpha)
      -\frac12\Gjp_\alpha(z-z_\alpha^D),
 \qquad \varepsilon_V(S)\le\frac\kappa2.
\end{equation}

\subsection{Conditional NR/J0 theorem}

\begin{theorem}[JP--JC implies J0]\label{NR:thm:NR}
Assume JC--PL, equivalently its JP--JC output \eqref{NR:eq:JC}.  Then every normalized polygon in
the active near-disk chart satisfies
\begin{equation}\label{NR:eq:J0}
\boxed{
 \Phi_\alpha(h_\alpha\mathbf1+z)
 -\Phi_{\alpha_*}(h_*\mathbf1)
 \ge\kappa\Jfour(\alpha)
      +\frac12\Gjp_\alpha(z-z_\alpha^D).}
\end{equation}
The corrector obeys the first estimate in \eqref{NR:eq:corrector}, uniformly
in the smallest angle and all adjacent ratios.
\end{theorem}

\begin{proof}
By the definition of the honest joint residual,
\begin{align*}
 \Phi_\alpha(h_\alpha\mathbf1+z)
 -\Phi_{\alpha_*}(h_*\mathbf1)
 ={}&q_D(B_\alpha+K_\alpha z)-q_D(B_{\alpha_*})\\
 &+\mathscr R_\alpha(z)+E_{V,\alpha},
 \qquad |E_{V,\alpha}|\le\varepsilon_V(S)\Jfour(\alpha).
\end{align*}
Insert \eqref{NR:eq:disk-completion}, the closed estimate
$P_D^{\rm ns}(\alpha)-P_D^{\rm ns}(\alpha_*)
\ge2\kappa\Jfour(\alpha)$, and \eqref{NR:eq:JC}.  Since
$\varepsilon_V(S)\le\kappa/2$, the remaining coefficients are exactly
$\kappa$ and $1/2$, which proves \eqref{NR:eq:J0}.  The corrector estimate is
the closed JE--P row followed by exact Hilbert completion.
\end{proof}

\begin{corollary}[Exact downstream handoff]\label{NR:cor:J0}
The implication
\[
 \boxed{\mathrm{JP\text{-}DE+JP\text{-}JC}
 \Longrightarrow\mathrm{NR}\Longrightarrow\mathrm{J0}}
\]
is proved.  JP--DE abbreviates the closed $D0_{\rm ns}$, JE--P, and JE--V
rows.  JC--PL is the only unavailable premise; JP--JC is dependent on it.
\end{corollary}

\subsection{DAG certificate}

\begin{center}
\begin{longtable}{>{\raggedright\arraybackslash}p{.34\linewidth}
                  >{\raggedright\arraybackslash}p{.54\linewidth}}
\toprule
edge & disposition\\
\midrule
JP--EH and JP--NS & closed exact scalar/Hessian chart and disk metric;\\
JP--DE & closed disk Schur gap, corrector bound, and value matching;\\
JP--SC, JE--FM, JP--FS, JP--BR & removed from the active chain; the old
period-three material argument for JE--FM is retracted, and all four rows
are unnecessary on the joint route;\\
JC--ALG and JC--TL & closed exact action and support-time cancellation;\\
JC--PL & unique open pathwise LCA payment;\\
JP--JC & dependent exact output of JC--PL;\\
JP--DE$+$JP--JC $\to$ NR $\to$ J0 & proved by
Theorem~\ref{NR:thm:NR}.\\
\bottomrule
\end{longtable}
\end{center}
\endgroup


\part{Side activation and the global theorem}

\begingroup
\section{Foundation F57: stationary side mass and first moment}
\label{module:F57}
\noindent\textit{Audited source: \nolinkurl{convex_largeN_polya_szego_stage57_stationary_flux_reduction.tex}.}
\subsection{Active stationary polygons}

Let
\begin{equation}
 P=\bigcap_{i=1}^M\{x\in\mathbb R^2:x\cdot\nu_i<h_i\}
\end{equation}
be a bounded convex active polygon.  Write
\begin{equation}
 \nu_i=(\cos\theta_i,\sin\theta_i),\qquad
 \tau_i=(-\sin\theta_i,\cos\theta_i),
\end{equation}
and let $S_i$ be the side on the line $x\cdot\nu_i=h_i$.
Every point of $S_i$ has the unique representation
\begin{equation}\label{F57:eq:side-coordinate}
 x=h_i\nu_i+s\tau_i,\qquad s_i^-<s<s_i^+.
\end{equation}
Thus its length is
$\ell_i=s_i^+-s_i^-$.  Let $u>0$ be the $L^2(P)$-normalized first
eigenfunction,
\begin{equation}\label{F57:eq:eigen}
 -\Delta u=\lambda u\quad\hbox{in }P,\qquad
 u=0\quad\hbox{on }\partial P,\qquad
 \int_Pu^2=1,
\end{equation}
and put
\begin{equation}
 p_i(s)=-\partial_\nu u(h_i\nu_i+s\tau_i)\ge0.
\end{equation}
Convex polygonal regularity gives $u\in H^2(P)$ and
$p_i\in L^2(S_i)$.  The distributed Lipschitz-domain derivative, or
equivalently approximation inside a fixed active chart, gives the Hadamard
formula
\begin{equation}\label{F57:eq:Hadamard}
 D\lambda(P)[V]=-
 \sum_i\int_{S_i}p_i^2(V\cdot\nu_i)\,ds
\end{equation}
for every vertex-compatible polygonal material field $V$.

We use the scale-invariant functional
\begin{equation}
 \Lambda(P)=\frac{|P|}{\pi}\lambda_1(P).
\end{equation}
Stationarity of $\Lambda$ is equivalent to stationarity of $\lambda_1$
under the constraint $|P|=A$.

\subsection{The two exact sidewise moment identities}

\begin{theorem}[Stationary flux equidistribution]\label{F57:thm:flux}
Suppose that $P$ is stationary for $\Lambda$ with respect to all active
vertex perturbations.  Then, on every side $S_i$,
\begin{align}
 \frac1{\ell_i}\int_{s_i^-}^{s_i^+}p_i(s)^2\,ds
 &=\frac{\lambda}{|P|},\label{F57:eq:mean}\\
 \int_{s_i^-}^{s_i^+}
 \left(s-\frac{s_i^-+s_i^+}{2}\right)p_i(s)^2\,ds
 &=0.\label{F57:eq:first-moment}
\end{align}
Both identities are independent of the smallest side and of all mesh
ratios.
\end{theorem}

\begin{proof}
Work first with the Lagrangian
\begin{equation}
 \mathcal F(P)=\lambda_1(P)+\mu |P|.
\end{equation}

Fix the normal fan and replace $h_i$ by $h_i+t z_i$.  On the open side
$S_i$ the induced normal velocity is $z_i$; on every other open side it is
zero.  The endpoint motions are tangential to the neighboring sides and do
not contribute to the boundary integrals.  Consequently
\begin{align}
 D\lambda(P)[z]&=-\sum_i z_i\int_{S_i}p_i^2\,ds,
 &D|P|[z]&=\sum_i z_i\ell_i.
\end{align}
The $z_i$ are independent in an active chart.  Stationarity gives
\begin{equation}\label{F57:eq:multiplier-row}
 \int_{S_i}p_i^2\,ds=\mu\ell_i
 \qquad(1\le i\le M).
\end{equation}

To identify the multiplier, use the dilation field $V(x)=x$.  Scaling gives
\begin{equation}
 D\lambda(P)[x]=-2\lambda,\qquad D|P|[x]=2|P|.
\end{equation}
Hence $-2\lambda+2\mu|P|=0$, or
\begin{equation}\label{F57:eq:mu}
 \mu=\lambda/|P|.
\end{equation}
Equations \eqref{F57:eq:multiplier-row}--\eqref{F57:eq:mu} prove
\eqref{F57:eq:mean}.

It remains to rotate one supporting line.  Replace
$\theta_i$ by $\theta_i+tq_i$ while keeping $h_i$ fixed.  Differentiating
\begin{equation}
 x_t\cdot\nu_i(t)=h_i
\end{equation}
at a point $x=h_i\nu_i+s\tau_i$ gives
\begin{equation}\label{F57:eq:rot-velocity}
 V\cdot\nu_i=-s q_i.
\end{equation}
The normal velocity is zero on every other open side.  Therefore
\begin{align}
 D\lambda(P)[q_i]
 &=q_i\int_{s_i^-}^{s_i^+}s p_i(s)^2\,ds,\notag\\
 D|P|[q_i]
 &=-q_i\int_{s_i^-}^{s_i^+}s\,ds.
\end{align}
Stationarity of $\lambda+\mu|P|$ and the sign convention in
\eqref{F57:eq:multiplier-row} give
\begin{equation}\label{F57:eq:raw-first}
 \int_{s_i^-}^{s_i^+}s p_i(s)^2\,ds
 =\mu\int_{s_i^-}^{s_i^+}s\,ds.
\end{equation}
Subtract the midpoint times \eqref{F57:eq:multiplier-row} from
\eqref{F57:eq:raw-first}.  This is \eqref{F57:eq:first-moment}.
\end{proof}

\begin{corollary}[Sidewise probability law]\label{F57:cor:probability}
For every side define
\begin{equation}\label{F57:eq:probability}
 d\mathbb P_i(s)=
 \frac{p_i(s)^2}{(\lambda/|P|)\ell_i}\,ds.
\end{equation}
Then $\mathbb P_i$ is a probability measure and
\begin{equation}
 \int s\,d\mathbb P_i(s)=\frac{s_i^-+s_i^+}{2}.
\end{equation}
In particular, a stationary short side cannot be analyzed using only its
total flux: its first flux moment is fixed as well.
\end{corollary}

\begin{proof}
This is exactly \eqref{F57:eq:mean} and \eqref{F57:eq:first-moment}.
\end{proof}

\begin{remark}[Why the moment row is not redundant]
For a shallow corner model
$p(s)^2\simeq C s^{2\varepsilon_-}
(\ell-s)^{2\varepsilon_+}$, the normalized first moment is the mean of a
beta distribution.  Its displacement from $\ell/2$ has the sign of
$\varepsilon_--\varepsilon_+$.  Thus \eqref{F57:eq:first-moment} detects an
imbalance of the two endpoint corner exponents which the average identity
\eqref{F57:eq:mean} cannot see.
\end{remark}
\endgroup

\begingroup
\section{Node A0: strict transverse side activation}
\label{module:A0}
\noindent\textit{Audited source: \nolinkurl{A0_side_activation.tex}.}
\subsection{A sharp one-dimensional tent inequality}

For \(L>0\) and \(x\in(0,L)\), let
\begin{equation}\label{A0:eq:green}
 K_x(s)=
 \begin{cases}
  \displaystyle\frac{s(L-x)}{L},&0\le s\le x,\\[2mm]
  \displaystyle\frac{x(L-s)}{L},&x\le s\le L.
 \end{cases}
\end{equation}
Thus \(K_x\) is the Dirichlet Green kernel of \(-d^2/ds^2\), viewed as a
function of \(s\), and the unit-height triangular tent is
\begin{equation}\label{A0:eq:tent}
 h_x(s)=\frac{L}{x(L-x)}K_x(s).
\end{equation}

\begin{lemma}[Log-concave density dominates its constant rearrangement]
\label{A0:lem:tent}
Let \(q:(0,L)\to(0,\infty)\) be log-concave and integrable, with
\begin{equation}\label{A0:eq:endpoint}
 \lim_{s\downarrow0}q(s)=\lim_{s\uparrow L}q(s)=0.
\end{equation}
Put
\begin{equation}
 \mu=\frac1L\int_0^Lq(s)\,ds.
\end{equation}
If
\begin{equation}\label{A0:eq:barycenter}
 \int_0^L\left(s-\frac L2\right)q(s)\,ds=0,
\end{equation}
then, for every \(x\in(0,L)\),
\begin{equation}\label{A0:eq:tent-positive}
 \boxed{\int_0^Lh_x(s)\{q(s)-\mu\}\,ds>0.}
\end{equation}
\end{lemma}

\begin{proof}
Set \(g=q-\mu\) and
\begin{equation}
 G(y)=\int_0^yg(s)\,ds,\qquad
 H(y)=\int_0^yG(r)\,dr.
\end{equation}
Mass and barycenter matching give
\begin{equation}\label{A0:eq:moments}
 \int_0^Lg=0,\qquad \int_0^Lsg(s)\,ds=0,
\end{equation}
and therefore
\begin{equation}\label{A0:eq:Hend}
 H(L)=\int_0^L(L-s)g(s)\,ds=0.
\end{equation}

Log-concavity says that every superlevel set of \(q\) is an interval.
In view of \eqref{A0:eq:endpoint} and \(\int q=\mu L\), the set
\(\{q>\mu\}\) is a nonempty interval \((a,b)\Subset(0,L)\), up to
irrelevant closed plateau endpoints.  Thus \(g\) has the sign pattern
\[
 -,\quad +,\quad -.
\]
Consequently \(G\) first decreases, then increases, and finally decreases
to \(G(L)=0\).  On the last interval it must be nonnegative, because it is
decreasing to zero.  Hence \(G\) changes sign exactly once, from negative
to positive.  It follows that \(H\) first decreases and then increases.
Together with \(H(0)=H(L)=0\), this gives
\begin{equation}\label{A0:eq:Hnegative}
 H(y)<0,\qquad 0<y<L.
\end{equation}
Strictness follows because \(q\) cannot equal the positive constant
\(\mu\) while satisfying \eqref{A0:eq:endpoint}.

Now set
\[
 F(x)=\int_0^LK_x(s)g(s)\,ds.
\]
The Green-kernel identity gives
\[
 -F''=g,\qquad F(0)=F(L)=0.
\]
Moreover \(F'(0)=L^{-1}\int_0^L(L-s)g(s)\,ds=0\) by
\eqref{A0:eq:moments}.  Hence \(F'=-G\) and \(F=-H\).  Equation
\eqref{A0:eq:Hnegative} yields \(F(x)>0\).  Multiplication by the positive
factor in \eqref{A0:eq:tent} proves \eqref{A0:eq:tent-positive}.
\end{proof}

\begin{remark}
The proof is a convex-order statement, but no external convex-order
theorem is needed.  The only special input is that a log-concave density
crosses its average on one interval; the two moment equalities then force
the sign of every tent pairing.
\end{remark}

\subsection{Log-concavity of the side flux}

Let \(P\) be a bounded convex polygon and let \(u>0\) be its first
Dirichlet eigenfunction:
\begin{equation}
 -\Delta u=\lambda u\quad\hbox{in }P,\qquad
 u=0\quad\hbox{on }\partial P.
\end{equation}
Fix an open side \(S\), identify it with \(0<s<L\), and let
\(\nu_{\rm in}\) be its inward unit normal.  Define
\begin{equation}\label{A0:eq:p}
 p(s)=\partial_{\nu_{\rm in}}u(s)>0.
\end{equation}

\begin{lemma}[The boundary flux is log-concave on a side]
\label{A0:lem:flux-logconcave}
The function \(p\) in \eqref{A0:eq:p} is log-concave on \((0,L)\).
Moreover,
\begin{equation}\label{A0:eq:p-endpoints}
 p(s)\longrightarrow0
 \quad\hbox{as }s\downarrow0\hbox{ or }s\uparrow L.
\end{equation}
Consequently \(p^2\) satisfies the hypotheses of
Lemma~\ref{A0:lem:tent}.
\end{lemma}

\begin{proof}
The Brascamp--Lieb theorem~\cite{BL76} gives log-concavity of \(u\) in every bounded
convex domain.  Fix \(s_0,s_1\in(0,L)\) and \(0<\theta<1\).  For all
sufficiently small \(r>0\), the three points
\[
 s_j+r\nu_{\rm in}\quad(j=0,1),\qquad
 ((1-\theta)s_0+\theta s_1)+r\nu_{\rm in}
\]
belong to \(P\).  Log-concavity of \(u\) gives
\begin{multline}
 \frac{u(((1-\theta)s_0+\theta s_1)+r\nu_{\rm in})}{r}\\
 \ge
 \left(\frac{u(s_0+r\nu_{\rm in})}{r}\right)^{1-\theta}
 \left(\frac{u(s_1+r\nu_{\rm in})}{r}\right)^\theta.
\end{multline}
Letting \(r\downarrow0\) proves log-concavity of \(p\).

At a convex vertex of interior angle \(\omega<\pi\), the standard corner
expansion of a Dirichlet solution starts with
\(r^{\pi/\omega}\sin(\pi\vartheta/\omega)\).  Hence its gradient is
\(O(r^{\pi/\omega-1})\), which tends to zero.  This proves
\eqref{A0:eq:p-endpoints}.
\end{proof}

\subsection{Activating one new side}

Use the scale-invariant functional
\begin{equation}\label{A0:eq:Lambda}
 \Lambda(P)=\frac{|P|}{\pi}\lambda_1(P).
\end{equation}
Suppose \(P\) is stationary under all active vertex variations.  On each
side \(S\) of length \(L\), stages 70--72 give the exact identities
\begin{align}
 \frac1L\int_0^Lp(s)^2\,ds&=\frac{\lambda}{|P|}=:\mu,
 \label{A0:eq:mean}\\
 \int_0^L\left(s-\frac L2\right)p(s)^2\,ds&=0.
 \label{A0:eq:first-moment}
\end{align}

Choose \(x\in(0,L)\), insert at that point a collinear marked vertex, and
move it a distance \(\varepsilon>0\) in the outward normal direction,
keeping the two endpoints of \(S\) fixed.  For small \(\varepsilon\) this
replaces \(S\) by two sides of a convex polygon and adds a triangular cap.
At \(\varepsilon=0\), the normal velocity on \(S\) is exactly the tent
\(h_x\) of \eqref{A0:eq:tent}.

\begin{theorem}[Strict side-activation descent]\label{A0:thm:descent}
For every side and every \(x\in(0,L)\), the outward splitting path
\(P_\varepsilon\) satisfies
\begin{equation}\label{A0:eq:strict}
 \boxed{
 \frac d{d\varepsilon}\Lambda(P_\varepsilon)
 \bigg|_{\varepsilon=0+}<0.}
\end{equation}
Thus a stationary convex \(M\)-gon is not stationary, and is not a local
minimum, in the closed class of convex polygons with at most \(M+1\)
active sides.
\end{theorem}

\begin{proof}
Hadamard's formula and the area derivative give, for an outward normal
velocity \(h\),
\begin{align}
 D\Lambda(P)[h]
 &=\frac{\lambda}{\pi}\int_{\partial P}h\,ds
   -\frac{|P|}{\pi}\int_{\partial P}p^2h\,ds\notag\\
 &=\frac{|P|}{\pi}\int_{\partial P}
   \left(\frac{\lambda}{|P|}-p^2\right)h\,ds.
 \label{A0:eq:Hadamard}
\end{align}
For the splitting path, \(h=h_x\) on \(S\) and is zero on the other
sides.  Lemma~\ref{A0:lem:flux-logconcave} applies to \(q=p^2\), while
\eqref{A0:eq:mean}--\eqref{A0:eq:first-moment} are precisely its mass and
barycenter hypotheses.  Lemma~\ref{A0:lem:tent} therefore gives
\[
 \int_0^Lh_x(p^2-\mu)\,ds>0.
\]
Insert this inequality and \(\mu=\lambda/|P|\) into
\eqref{A0:eq:Hadamard} to obtain \eqref{A0:eq:strict}.

The path has one additional active vertex for every sufficiently small
\(\varepsilon>0\).  Since its scale-invariant value decreases to first
order, the final assertion follows.
\end{proof}

\begin{corollary}[No minimizing side loss]\label{A0:cor:no-loss}
Let \(P\) minimize \(\Lambda\) in the class of convex polygons with at
most \(N\) sides.  If \(P\) has at least three sides and is stationary in
its active stratum, then it has exactly \(N\) active sides.
\end{corollary}

\begin{proof}
If \(P\) had \(M<N\) active sides, Theorem~\ref{A0:thm:descent} would produce
a convex \((M+1)\)-gon with strictly smaller \(\Lambda\), contradicting
minimality.
\end{proof}
\endgroup

\begingroup
\section{Node G0: compact global assembly and equality rigidity}
\label{module:G0}
\noindent\textit{Audited source: \nolinkurl{G0_global_wrapper.tex}.}
\subsection{Statement and normalization}

For $A>0$, let $\Pclass_N^{\rm conv}(A)$ be the class of bounded convex
planar polygons of area $A$ with exactly $N$ genuine sides, after
consecutive collinear sides are merged.  Let $R_N(A)$ be the regular
$N$-gon of area $A$.

\begin{theorem}[Global theorem conditional on J0]
\label{G0:thm:main}
Assume the active near-disk coercivity input in
\ref{G0:thm:J0}.  Then
For every $A>0$ there is $N_0$ such that, for all $N\ge N_0$ and all
$P\in\Pclass_N^{\rm conv}(A)$,
\begin{equation}\label{G0:eq:main}
 \lambda_1(P)\ge\lambda_1(R_N(A)).
\end{equation}
Equality holds if and only if $P$ is congruent to $R_N(A)$.
\end{theorem}

Put
\begin{equation}\label{G0:eq:Lambda}
 \LambdaF(\Omega)=\frac{|\Omega|}{\pi}\lambda_1(\Omega).
\end{equation}
It is invariant under dilations.  Hence it is enough to work at area
$\pi$; the general statement follows by the dilation
$P\mapsto(\pi/A)^{1/2}P$.

Let $\Kclass_{\le N}^{\rm conv}(\pi)$ be the Hausdorff-closed class of
area-$\pi$ convex polygons with at most $N$ active sides.  This class is
used only to take a global minimum.  The theorem itself concerns exactly
$N$ sides.

\subsection{The two retained local inputs}

We isolate the only two nonclassical inputs used in the global assembly.

\begin{theorem}[Strict transverse side activation; stage 132]
\label{G0:thm:activation}
Let $Q$ be a stationary convex $M$-gon for $\LambdaF$ in its active
stratum.  Insert a marked point in the relative interior of any side and
move it outward while keeping the two endpoints fixed.  For all
sufficiently small positive displacements this gives a convex
$(M+1)$-gon $Q_\varepsilon$ and
\begin{equation}\label{G0:eq:activation}
 \frac d{d\varepsilon}\LambdaF(Q_\varepsilon)
 \bigg|_{\varepsilon=0+}<0.
\end{equation}
The assertion is independent of all side lengths and mesh ratios.
\end{theorem}

The proof in stage 132 uses the stationary mass and barycenter identities
for the squared ground-state flux on a side, log-concavity of that flux,
and Hadamard's formula.  In particular, it is a first-variation theorem at
the lower-side point; it is not a mountain-pass deformation theorem.

\begin{obligation}[Active near-disk value coercivity; current J0 input]
\label{G0:thm:J0}
There are $\eta_0,S_0,c_0>0$ such that the following holds.  Let $P$ be an
area-$\pi$ active convex $N$-gon with positive normal fan
\begin{equation*}
 \alpha_i>0,\qquad \sum_i\alpha_i=2\pi,\qquad
 a=\frac{2\pi}{N},\qquad S=a+\max_i\alpha_i\le S_0.
\end{equation*}
Assume, after translation, that its radial graph is within $\eta_0$ of
the disk in $W^{1,\infty}$.  Let $h$ be its support vector in the fixed
translation complement of JP.  Put $h_0=h_\alpha\mathbf1$,
$A_0=A_\alpha(h_0)$, and
\(
s_\alpha(h)=A_\alpha'[h_0;h]/(2A_0).
\)
Normalize scale by
\[
 \widetilde h=\frac{h}{s_\alpha(h)},
 \qquad s_\alpha(\widetilde h)=1,
 \qquad \widetilde h=h_\alpha\mathbf1+z,
 \qquad z\in\mathsf X_\alpha.
\]
Shrinking the physical chart if necessary, $s_\alpha(h)>0$ and this
normalization stays in the same uniform near-disk chart.  Then
\begin{equation}\label{G0:eq:J0}
 \boxed{
 \lambda_1(P)-\lambda_1(R_N(\pi))
 \ge c_0\left\{
 \Jfour(\alpha)
 +\Gjp_\alpha(z-z_\alpha^\sharp)\right\},}
\end{equation}
where
\begin{equation}\label{G0:eq:J4}
 \Jfour(\alpha)=\sum_i\alpha_i^4-Na^4.
\end{equation}
The support metric is positive on the fixed affine scale section
$\mathsf X_\alpha$ after the fixed translation normalization.  Since
$\Phi_\alpha$ is scale invariant and the original $P$ has area $\pi$,
the JP/NR value inequality for $\widetilde h$ is exactly
\eqref{G0:eq:J0} for $\lambda_1(P)$.
Every constant is independent of the smallest physical, polar, and
conformal cells and of adjacent ratios.
\end{obligation}

For clarity, the conditional analytic chain behind input~\ref{G0:thm:J0} is
\begin{equation}\label{G0:eq:analytic-chain}
 \begin{gathered}
 \text{TR: exact physical trace, disk endpoint gap, value split}\\
 \downarrow\\
 \text{CT core}\longrightarrow\{\text{JP--EH},\text{JP--NS}\}\\
 \text{JP--NS+TR+CT core+QB+DF+DM}\longrightarrow\text{JP--DE}\\
 \text{JP--EH+JP--NS+JP--DE}\longrightarrow
 \{\text{JC--ALG},\text{JC--TL}\}\\
 \text{JC--ALG+JC--TL}\longrightarrow\text{JC--PL}
 \longrightarrow\text{JP--JC}\\
 \downarrow\\
 \text{JP--DE+JP--JC}\longrightarrow\text{NR}\longrightarrow\text{J0}.
 \end{gathered}
\end{equation}
No stationary Hessian is integrated on the nonstationary chart.

\subsection{Compact minimization in the closed side class}

\begin{lemma}[Compact spectral sublevels]\label{G0:lem:compact}
For fixed $N$ and $L<\infty$, the family
\begin{equation}\label{G0:eq:sublevel}
 \{P\in\Kclass_{\le N}^{\rm conv}(\pi):
       \lambda_1(P)\le L\}
\end{equation}
is Hausdorff precompact up to translations.  Its Hausdorff limits remain
in $\Kclass_{\le N}^{\rm conv}(\pi)$, and $\lambda_1$ is continuous along
such nondegenerate convex-domain convergence.  Consequently the infimum
of $\lambda_1$ over $\Kclass_{\le N}^{\rm conv}(\pi)$ is attained.
\end{lemma}

\begin{proof}
If $w(P)$ is the minimum width, containment in a strip and the
one-dimensional Poincar\'e inequality give
\[
 \lambda_1(P)\ge\frac{\pi^2}{w(P)^2}.
\]
Thus a spectral sublevel has a positive width bound.  For convex bodies,
fixed area and a positive minimum-width bound give a diameter bound.  To
see this without an implicit eccentricity lemma, let $[A,B]\subset P$ be
a diameter, $D=|A-B|$, and use it as the horizontal axis.  If $y_+$ and
$y_-$ are the extreme signed heights of $P$ above and below that axis,
the two triangles with base $[A,B]$ and apices at the corresponding
support points have disjoint interiors and lie in $P$.  Hence
\[
 |P|\ge \frac D2(y_+-y_-)
      \ge \frac D2 w(P),
 \qquad
 D\le \frac{2\pi}{w(P)}.
\]
After fixing, for example, the Steiner point at the origin, Blaschke
compactness then gives a Hausdorff convergent
subsequence.  Area is continuous for convex bodies, and having at most
$N$ genuine sides is closed (parametrize by at most $N$ cyclic vertices
and allow adjacent vertices to coalesce).  Dirichlet spaces of bounded convex
domains converge in the Mosco sense under this Hausdorff convergence, so
the first eigenvalue converges.  A minimizing sequence may be restricted
to the sublevel determined by any one competitor, proving attainment.
\end{proof}

\begin{proposition}[A closed-class minimizer loses no side]
\label{G0:prop:no-loss}
Every minimizer $Q_N$ of $\lambda_1$ over
$\Kclass_{\le N}^{\rm conv}(\pi)$ has exactly $N$ active sides.
\end{proposition}

\begin{proof}
After collinear sides are merged, suppose that $Q_N$ has $M<N$ genuine
sides.  The active $M$-gon stratum is a smooth finite-dimensional
manifold near $Q_N$ after the area and rigid gauges are fixed.  Since
$Q_N$ is a global minimizer, it is stationary on that stratum.
Theorem~\ref{G0:thm:activation} supplies a convex $(M+1)$-gon with strictly
smaller $\LambdaF$, hence with strictly smaller $\lambda_1$ at fixed
area.  It still belongs to $\Kclass_{\le N}^{\rm conv}(\pi)$, a
contradiction.
\end{proof}

\begin{remark}[Why SB disappears]
Stage 133 correctly observes that different one-side activation branches
need not be connected below the value of their common lower stratum.
Proposition~\ref{G0:prop:no-loss} does not connect them.  It selects a global
minimizer and needs only one strict feasible descent to contradict
minimality.  Therefore the side-branch bypass SB is not an open edge of
this proof.
\end{remark}

\subsection{Qualitative global-to-local entry}

Let $B$ be the area-$\pi$ disk and
$\lambda_B=\lambda_1(B)$.

\begin{lemma}[Regular polygons converge spectrally to the disk]
\label{G0:lem:regular-limit}
The area-$\pi$ regular polygons satisfy
\begin{equation}\label{G0:eq:regular-limit}
 R_N(\pi)\longrightarrow B\quad\text{in Hausdorff distance},
 \qquad
 \lambda_1(R_N(\pi))\longrightarrow\lambda_B.
\end{equation}
\end{lemma}

\begin{proof}
The inradius and circumradius of the area-normalized regular $N$-gon both
tend to one.  This proves Hausdorff convergence.  Spectral convergence is
the convex-domain stability used in Lemma~\ref{G0:lem:compact}.
\end{proof}

\begin{lemma}[Near-disk entry from a competing value]
\label{G0:lem:entry}
Let $N_j\to\infty$ and let
$P_j\in\Kclass_{\le N_j}^{\rm conv}(\pi)$ satisfy
\begin{equation}\label{G0:eq:competing}
 \lambda_1(P_j)\le\lambda_1(R_{N_j}(\pi)).
\end{equation}
Then, after translations,
\begin{equation}\label{G0:eq:support-limit}
 P_j\longrightarrow B\quad\text{in Hausdorff distance},\qquad
 \|h_{P_j}-1\|_{L^\infty(\mathbb S^1)}\longrightarrow0.
\end{equation}
If every $P_j$ is active and $\alpha_{\max}(P_j)$ denotes its largest
gap between consecutive outer normals, then
\begin{equation}\label{G0:eq:angle-limit}
 \alpha_{\max}(P_j)\longrightarrow0.
\end{equation}
Finally its radial function $r_j=1+\rho_j$ satisfies
\begin{equation}\label{G0:eq:radial-limit}
 \|\rho_j\|_{W^{1,\infty}(\mathbb S^1)}\longrightarrow0.
\end{equation}
\end{lemma}

\begin{proof}
Faber--Krahn and \eqref{G0:eq:competing}--\eqref{G0:eq:regular-limit} give
\[
 \lambda_B\le\lambda_1(P_j)
 \le\lambda_1(R_{N_j}(\pi))\longrightarrow\lambda_B.
\]
The sequence lies in one fixed spectral sublevel.  The proof of
Lemma~\ref{G0:lem:compact}, with no fixed side bound needed for the convex
compactness part, gives Hausdorff precompactness after translation.
Every limit has area $\pi$ and first eigenvalue $\lambda_B$.
Faber--Krahn rigidity makes the limit the disk.  Since every subsequential
limit is the disk, the whole sequence converges, and convergence of convex
support functions gives \eqref{G0:eq:support-limit}.

Suppose \eqref{G0:eq:angle-limit} failed.  Along a subsequence two adjacent
supporting normals would have a gap at least $\varepsilon>0$.  Their
support numbers tend uniformly to one by \eqref{G0:eq:support-limit}; their
intersection vertex in the bisector direction would therefore have
distance at least $\sec(\varepsilon/2)+o(1)>1$ from the origin.  This
contradicts Hausdorff convergence to $B$.

For a ray meeting the side with normal angle $\theta_i$ and support
number $h_i$,
\begin{equation}\label{G0:eq:radial-side}
 r_j(\theta)=\frac{h_i}{\cos(\theta-\theta_i)}.
\end{equation}
Hausdorff convergence and the inclusions
$(1-\varepsilon_j)B\subset P_j\subset(1+\varepsilon_j)B$ give
$r_j\to1$ uniformly, while support convergence gives $h_i\to1$
uniformly in $i$.  Thus, whenever the ray at $\theta$ meets side $i$,
\[
 \cos(\theta-\theta_i)=\frac{h_i}{r_j(\theta)}\longrightarrow1
\]
uniformly.  In particular $|\theta-\theta_i|\to0$; this also proves the
needed statement at one-sided sector endpoints.  Equation
\eqref{G0:eq:radial-side} and its one-sided derivative therefore give
\eqref{G0:eq:radial-limit}.
\end{proof}

\subsection{Contradiction and equality rigidity}

\begin{proof}[Proof of Theorem~\ref{G0:thm:main}]
By scale invariance, take $A=\pi$.

Assume first that the inequality fails for arbitrarily large side counts.
Then there are $N_j\to\infty$ and exact $N_j$-gons $P_j$ such that
\begin{equation}\label{G0:eq:strict-counter}
 \lambda_1(P_j)<\lambda_1(R_{N_j}(\pi)).
\end{equation}
Let $Q_j$ minimize $\lambda_1$ over
$\Kclass_{\le N_j}^{\rm conv}(\pi)$.  Lemma~\ref{G0:lem:compact} and
Proposition~\ref{G0:prop:no-loss} show that $Q_j$ exists and has exactly
$N_j$ active sides.  Moreover
\begin{equation}\label{G0:eq:min-counter}
 \lambda_1(Q_j)\le\lambda_1(P_j)
 <\lambda_1(R_{N_j}(\pi)).
\end{equation}
Lemma~\ref{G0:lem:entry} gives
\[
 \|\rho_{Q_j}\|_{W^{1,\infty}}\to0,\qquad
 S_j=\frac{2\pi}{N_j}+\alpha_{\max}(Q_j)\to0.
\]
Thus input~\ref{G0:thm:J0} applies for all sufficiently large $j$ and
gives
\[
 \lambda_1(Q_j)-\lambda_1(R_{N_j}(\pi))
 \ge c_0\{\Jfour(\alpha^{(j)})
 +\Gjp_{\alpha^{(j)}}(
 z_j-z_{\alpha^{(j)}}^\sharp)\}\ge0,
\]
contradicting \eqref{G0:eq:min-counter}.  Hence the inequality in
\eqref{G0:eq:main} holds for every sufficiently large $N$.

Suppose next that nonregular equality cases exist for arbitrarily large
$N$.  Choose $N_j\to\infty$ and nonregular exact $N_j$-gons $P_j$ with
\[
 \lambda_1(P_j)=\lambda_1(R_{N_j}(\pi)).
\]
The inequality just proved shows that this common value is the minimum
over the exact $N_j$-side class.  If the minimum over the closed class
with at most $N_j$ sides were smaller, its minimizer would have exactly
$N_j$ sides by Proposition~\ref{G0:prop:no-loss}, contradicting the proved
inequality.  Hence each $P_j$ is also a closed-class minimizer.
Lemma~\ref{G0:lem:entry} and input~\ref{G0:thm:J0} apply and force
\begin{equation}\label{G0:eq:defect-zero}
 \Jfour(\alpha^{(j)})=0,\qquad
 \Gjp_{\alpha^{(j)}}(
 z_j-z_{\alpha^{(j)}}^\sharp)=0.
\end{equation}
Strict convexity of $x^4$ under $\sum_i\alpha_i=2\pi$ makes the first
identity equivalent to $\alpha_i=2\pi/N_j$ for every $i$.  At the uniform
fan cyclic symmetry gives
$z_{\alpha^{(j)}}^\sharp=0$ on the exact quotient.  Positivity
of the support metric then makes the support vector constant modulo
translations and scale.  Thus $P_j$ is congruent to the regular
$N_j$-gon, a contradiction.

Both a strict failure and a nonregular equality can therefore occur only
for finitely many side counts.  Increasing $N_0$ proves the theorem.
\end{proof}
\endgroup


\section{Final interface verification}
\label{sec:final-interface-verification}

\subsection{Discharge ledger}

The QB targets \ref{QB:obl:DF} and \ref{QB:obl:DM} are proved by
Corollary~\ref{ZG:cor:DF} and Theorem~\ref{PT:thm:DM}.  The adaptive
sampling target \ref{DF0A:obl:AVS} is Theorem~\ref{ZG:thm:AVS}; the
provisional ELCA target is bypassed by the zero-gap source decomposition.
The strong ZG--$T1$ square, full-vector reconstruction, separate-curvature
estimate, and CT--JM estimate are retracted rather than assumed.
The old two-sided vanishing-modulus TR--NR remainder is likewise archived
and excluded from this assembly; the current NR node uses only the
one-sided fixed-loss output of JC--PL through JP--JC.

The old CT--JP/CT--M4 support Taylor branch is retired.  The closed theorem
JP--EH, Theorem~\ref{JP:thm:EH}, gives the exact nonstationary support
Hessian, including reduced energy, mixed flux, and area second variation.
Proposition~\ref{JP:prop:normal-speed-chart-entry} closes JP--NS: its sine
interpolant reproduces translations exactly, is controlled by the physical
vertex displacement without inverse-angle loss, and keeps the whole
support segment active.
Proposition~\ref{JP:prop:cubic-counterexample} supplies a certified
alternating-mesh negative direction for the old cubic candidate.  Thus the
proof neither assumes nor silently reuses that candidate.

Theorem~\ref{JP:thm:normal-disk-endpoint} proves the $D0_{\rm ns}$ endpoint
and JE--P principal forcing estimate; Theorem
\ref{JP:thm:normal-value-matching} proves JE--V.  These closed rows form
JP--DE.  The near-regular period-three material proposition in the JP node
is retracted after the polygonal corner finite-part audit and proves
neither JE--FM nor its negation.  The JC node
therefore defines the honest residual \ref{JC:eq:joint-residual}, keeps its
forcing and integrated Hessian together in \ref{JC:eq:integrated-form}, and
proves the exact completed-action identity and support-time logarithm
cancellation, and isolates the sole open pathwise payment JC--PL,
Contract~\ref{JC:con:JC-PL}.  Its algebraic output is JP--JC,
equation~\ref{JC:eq:JC}.
Theorem~\ref{JC:thm:JC-to-J0} and its synchronized NR handoff
\ref{NR:thm:NR} prove the exact implication to J0.  Hence NR, J0, G0, and
the main theorem are dependent only on JC--PL through JP--JC.

\subsection{Conditional composition}

The correct-source gap and no-ratio interpolant give the angular disk
endpoint; F129 and TR transfer it to physical traces.  The CT core and QB
dictionary identify the exact disk cubic rows while retaining curvature.
ZG proves DF and PT proves DM.  Exact disk Hilbert completion gives
$P_D^{\rm ns}(\alpha)+\Gjp_\alpha(z-z_\alpha^D)$.  JP--JC permits at most
half of the closed fan and support margins to be lost by the genuine
polygonal-to-disk residual.  The assembly in Theorem~\ref{NR:thm:NR}
therefore gives J0 without a separate flux estimate.  Strict activation
and compact minimization then yield Theorem~\ref{G0:thm:main}, which is the
conditional conclusion of Theorem~\ref{thm:assembled-main} after scale
normalization.

Since JC--PL remains open, this verifies the conditional implication but
does not certify the unconditional target theorem.

\begin{thebibliography}{9}
\bibitem{BL76}
H.~J.~Brascamp and E.~H.~Lieb,
\emph{On extensions of the Brunn--Minkowski and Pr\'ekopa--Leindler
theorems, including inequalities for log concave functions, and with an
application to the diffusion equation},
J. Functional Analysis \textbf{22} (1976), 366--389.
\end{thebibliography}

\end{document}
